% !TeX program = xelatex
% ================================================================
%  《收获与播种》(Récoltes et Semailles) 中文译本
%  卷二 · 第一部 —— 愚妄与更新 (Fatuité et Renouvellement)
%  原书 p.123-254, 第 5-12 章:
%    5. 劳动与发现  6. 梦与做梦者  7. 双重面孔  8. 师与生
%    9. 收获       10. 孩子嬉戏  11. 孤独的历险  12. 第一部分的注
%  --------------------------------------------------------------
%  五卷本之一。共用导言区 common.tex, 封面 cover-art.tex。
%  构建: xelatex ×4 (目录页码需 4 遍才收敛)
% ================================================================
\PassOptionsToPackage{AutoFallBack}{xeCJK}
\documentclass[
    PaperSize=a4paper,
    HeadingFont=hei,
    Biblatex=false,
    BookTitleCN={收获与播种},
    Translator={数译},
    LogoFile=logo.pdf
]{mathtranslation}

% ================================================================
%  《收获与播种》中文译本 —— 各卷共用导言区
%  由 vol-prelude / vol-part1 / vol-part2 / vol-part3 / vol-part4
%  以及全书合订 main-all 共同 \input。基类 mathtranslation.cls v3.2 不改。
%  --------------------------------------------------------------
%  样式均在此改写, 以贴近原书:
%    · 章标题仿原书 "5. Travail et découverte"(块状左对齐黑体)
%    · 节/小节: 黑体左对齐, 编号后带句点
%    · 脚注: 人工编号照录原书——前奏与第一部分页内重新编号,
%      第二至四部各自连续编号并带(*)角标; 第一部分另有 \fnoteP
%      处理 "(n)" 式连续注号
%    · 边码: 原打字稿页码 \origpage{N}(前奏作 A1、P1、L1、罗马数字等)
%    · ⋄ 记号: 段首补写标记 \pmark 与上标形式 \dmark
% ================================================================

% 题名含空格, 不能走类选项(会被 kvoptions 吃掉), 在此赋值
\renewcommand{\BookTitleEN}{Récoltes et Semailles}
\renewcommand{\OriginalAuthor}{Alexandre Grothendieck}
\renewcommand{\TranslationDate}{\today}

% --------------------------------------------------------------
%  罕用字后备(如"祐"): Fandol 不载时退到 SimSun
% --------------------------------------------------------------
\setCJKfallbackfamilyfont{\CJKrmdefault}{{SimSun}}

% --------------------------------------------------------------
%  hyperref 元信息(各卷在 % ================================================================
%  《收获与播种》中文译本 —— 各卷共用导言区
%  由 vol-prelude / vol-part1 / vol-part2 / vol-part3 / vol-part4
%  以及全书合订 main-all 共同 \input。基类 mathtranslation.cls v3.2 不改。
%  --------------------------------------------------------------
%  样式均在此改写, 以贴近原书:
%    · 章标题仿原书 "5. Travail et découverte"(块状左对齐黑体)
%    · 节/小节: 黑体左对齐, 编号后带句点
%    · 脚注: 人工编号照录原书——前奏与第一部分页内重新编号,
%      第二至四部各自连续编号并带(*)角标; 第一部分另有 \fnoteP
%      处理 "(n)" 式连续注号
%    · 边码: 原打字稿页码 \origpage{N}(前奏作 A1、P1、L1、罗马数字等)
%    · ⋄ 记号: 段首补写标记 \pmark 与上标形式 \dmark
% ================================================================

% 题名含空格, 不能走类选项(会被 kvoptions 吃掉), 在此赋值
\renewcommand{\BookTitleEN}{Récoltes et Semailles}
\renewcommand{\OriginalAuthor}{Alexandre Grothendieck}
\renewcommand{\TranslationDate}{\today}

% --------------------------------------------------------------
%  罕用字后备(如"祐"): Fandol 不载时退到 SimSun
% --------------------------------------------------------------
\setCJKfallbackfamilyfont{\CJKrmdefault}{{SimSun}}

% --------------------------------------------------------------
%  hyperref 元信息(各卷在 % ================================================================
%  《收获与播种》中文译本 —— 各卷共用导言区
%  由 vol-prelude / vol-part1 / vol-part2 / vol-part3 / vol-part4
%  以及全书合订 main-all 共同 \input。基类 mathtranslation.cls v3.2 不改。
%  --------------------------------------------------------------
%  样式均在此改写, 以贴近原书:
%    · 章标题仿原书 "5. Travail et découverte"(块状左对齐黑体)
%    · 节/小节: 黑体左对齐, 编号后带句点
%    · 脚注: 人工编号照录原书——前奏与第一部分页内重新编号,
%      第二至四部各自连续编号并带(*)角标; 第一部分另有 \fnoteP
%      处理 "(n)" 式连续注号
%    · 边码: 原打字稿页码 \origpage{N}(前奏作 A1、P1、L1、罗马数字等)
%    · ⋄ 记号: 段首补写标记 \pmark 与上标形式 \dmark
% ================================================================

% 题名含空格, 不能走类选项(会被 kvoptions 吃掉), 在此赋值
\renewcommand{\BookTitleEN}{Récoltes et Semailles}
\renewcommand{\OriginalAuthor}{Alexandre Grothendieck}
\renewcommand{\TranslationDate}{\today}

% --------------------------------------------------------------
%  罕用字后备(如"祐"): Fandol 不载时退到 SimSun
% --------------------------------------------------------------
\setCJKfallbackfamilyfont{\CJKrmdefault}{{SimSun}}

% --------------------------------------------------------------
%  hyperref 元信息(各卷在 \input{common} 之后可用 \hypersetup 覆写)
% --------------------------------------------------------------
\hypersetup{
  pdftitle={收获与播种},
  pdfauthor={Alexandre Grothendieck / 数译 译},
  pdfsubject={Récoltes et Semailles — Réflexions et témoignage sur un passé de mathématicien}
}

% --------------------------------------------------------------
%  标题样式: 仿原书 —— 黑体左对齐, 编号后带句点
% --------------------------------------------------------------
\titleformat{\section}[block]
  {\normalfont\Large\headfont}{\thesection.}{0.5em}{}
\titlespacing*{\section}{0pt}{3.2ex plus 1.2ex minus .2ex}{1.6ex plus .2ex}
\titleformat{\subsection}[block]
  {\normalfont\large\headfont}{\thesubsection.}{0.5em}{}
\titlespacing*{\subsection}{0pt}{2.6ex plus 1ex minus .2ex}{1.2ex plus .2ex}
\titleformat{\subsubsection}[block]
  {\normalfont\normalsize\headfont}{\thesubsubsection.}{0.5em}{}
\titlespacing*{\subsubsection}{0pt}{2.2ex plus .8ex minus .2ex}{1ex plus .2ex}
\titleformat{name=\subsubsection,numberless}[block]
  {\normalfont\normalsize\headfont}{}{0em}{}
\setcounter{secnumdepth}{3}
\setcounter{tocdepth}{3}
% 为含长外文书名与行内公式的段落提供必要的断行余量
\setlength{\emergencystretch}{3em}

% 章标题样式: 原书章标题与节同级排印("13. A) HERITAGE..."),
% 故章标题取块状左对齐黑体, 字号介于节与部之间
\titleformat{\chapter}[block]
  {\normalfont\LARGE\headfont}{\thechapter.}{0.5em}{}
\titlespacing*{\chapter}{0pt}{3.6ex plus 1.4ex minus .2ex}{2ex plus .2ex}
% 目录中章号照原书用阿拉伯数字("13. A) ..." 而非 "第十三章")
\renewcommand{\CTEXthechapter}{\thechapter\hspace{.3em}}

% --------------------------------------------------------------
%  目录观感: 条目与引导点之间、引导点与页码之间各留 0.5em
%  内核 \@dottedtocline 让 \leaders 从标题末字起排, 第一个点几乎
%  贴着末字, 读起来像多标了一个句点; 页码前也没有余量。两处各补
%  0.5em。
% --------------------------------------------------------------
\makeatletter
\patchcmd{\@dottedtocline}{\nobreak\leaders}{\nobreak\hspace{0.5em}\leaders}%
  {\typeout{MT: TOC 引导点前已留 0.5em}}%
  {\typeout{MT-WARN: TOC 引导点前间距补丁未生效}}
\patchcmd{\@dottedtocline}{\hfill\nobreak\hb@xt@\@pnumwidth}%
  {\hfill\nobreak\hspace{0.5em}\hb@xt@\@pnumwidth}%
  {\typeout{MT: TOC 页码前已留 0.5em}}%
  {\typeout{MT-WARN: TOC 页码前间距补丁未生效}}
\makeatother

% --------------------------------------------------------------
%  正文半角标点后的间距
%  本书体例用半角 , ; : ? ! 与 ASCII 括号, 由 xeCJK 按 HalfRight
%  (右半留白) 处理。放宽到 0.45em, 与全角标点右半留白大致齐平。
% --------------------------------------------------------------
\xeCJKsetup{CJKecglue={\hskip 0.45em plus 0.06em minus 0.09em}}

% --------------------------------------------------------------
%  页眉页脚: 偶数页眉右对齐斜体节题; 页码页脚外侧
% --------------------------------------------------------------
\fancyhf{}
\fancyhead[RE]{\small\itshape\rightmark}
\fancyfoot[LE,RO]{\small\thepage}
\renewcommand{\headrulewidth}{0pt}
\renewcommand{\sectionmark}[1]{\markright{\thesection.\ #1}}
\fancypagestyle{plain}{%
  \fancyhf{}%
  \fancyfoot[LE,RO]{\small \thepage}%
  \renewcommand{\headrulewidth}{0pt}%
}

% --------------------------------------------------------------
%  原书版式记号宏
% --------------------------------------------------------------
% 边码: 原打字稿页码(G. 手稿系统里的自引页码), 排在页边
\usepackage{marginnote}
\newcommand{\origpage}[1]{\marginnote{\normalfont\scriptsize p.~#1}}

% ⋄ 记号: 段首/行内正常大小; 上标形式为作者后补文字的位置标记
\DeclareRobustCommand{\pmark}{\ensuremath{\diamond}\hspace{0.15em}}
\DeclareRobustCommand{\dmark}{\textsuperscript{\ensuremath{\diamond}}}

% 脚注人工编号: \fnote{795}{(*)}{文本} —— 编号与(*)角标均照录原书
\newcommand{\fnote}[3]{%
  \begingroup
  \renewcommand{\thefootnote}{#1#2}%
  \footnote{#3}%
  \endgroup
}

% 第一部分的脚注: 页内序号作上标(#1), 注文行首印 "序号(连续注号)",
% 正文在序号之后另行括注同一连续注号 —— \fnoteP{1}{(3)}{注文}
% 原书有三条注无连续注号而改标日期(ch07 注3 = 30 Septembre,
% ch09 注11/12 = 8 août), 第二参数留空, 日期写进注文之首。
\newcommand{\fnoteP}[3]{%
  \begingroup
  \renewcommand{\thefootnote}{#1}\footnotemark
  \renewcommand{\thefootnote}{#1#2}\footnotetext{#3}%
  \endgroup
}

% 正文呼号与注文行首标号不同(原书照录)
\newcommand{\fnotevariant}[4]{%
  \begingroup
  \renewcommand{\thefootnote}{#1#2}\footnotemark
  \renewcommand{\thefootnote}{#1#3}\footnotetext{#4}%
  \endgroup
}

% 大注块标题: \notehead{171³}{斜体评论} —— 正文随后另起
\newcommand{\notehead}[2]{%
  \par\bigskip
  \noindent{\headfont 注~#1}\ifblank{#2}{}{\hspace{0.5em}{\itshape #2}}%
  \par\nobreak\smallskip\nobreak
}

% ∗ ∗ ∗ 分隔
\newcommand{\asterisms}{%
  \par\vspace{1.2ex}%
  \noindent\hspace*{2em}$\ast$\hspace{1.2em}$\ast$\hspace{1.2em}$\ast$%
  \par\vspace{1.2ex}%
}

% 无编号子目的目录条目
\newcommand{\tocsubsub}[1]{%
  \phantomsection
  \addcontentsline{toc}{subsubsection}{#1}%
}

% 部扉: 仿原书部题页
\newcommand{\partdivider}[3]{%
  \clearpage
  \thispagestyle{empty}%
  \phantomsection
  \addcontentsline{toc}{part}{#1\quad #3}%
  \vspace*{6.5cm}
  \begin{center}
    {\fontsize{16pt}{26pt}\selectfont\headfont #1\par}
    \vspace{0.9cm}
    {\fontsize{22pt}{34pt}\selectfont\headfont #2\par}
  \end{center}
  \clearpage
}

% 献辞页: 仿原书部题后的题献, 独占一页且不入目录
\newcommand{\dedication}[1]{%
  \thispagestyle{empty}%
  \vspace*{4cm}%
  \begin{center}#1\end{center}%
  \clearpage
}

% --------------------------------------------------------------
%  各卷"译者说明"的共用体例段(各卷的范围段自写, 后接本段)
% --------------------------------------------------------------
\newcommand{\CommonNoteBody}{%
为便于与原书及学界通行的打字稿页码对勘, 本册保留:

\begin{itemize}
  \item 原书全部章节编号(章~1--18, 节如~5.2、18.2, 子目如~18.1.1)与原书小节标题;
  \item 原书脚注的编号与角标——
        \begin{itemize}
          \item 前奏(第~1--4 章)与第一部的注在各页内重新编号, 正文中的上标即该
                页序号; 第一部紧随其后的括注(如“(3)”)是该部分的连续注号, 与第~12
                章“第一部分的注”里的“注~3”一一对应, 注文行首照原书印作“1(3)”式。
                其中三条注原书未标连续注号而改标日期(ch7 注~3 = 30 Septembre、
                ch9 注~11/12 = 8 août), 本册照录其日期;
          \item 第二部第~13--16~章各自重新编号, 分别为~1--41、1--86、1--132、1--29;
                第三部第~17~章自~1~起至~12, 第~18~章自~1~重新起算并延续至~794;
                末段自~\textbf{795(*)}~续至~1053(*)。各段与原书一一对应;
        \end{itemize}
  \item 页边码“p.~1”“p.~A3”“p.~1066”式小注——这是作者引用自己打字稿原稿时
        所用的页码, 与本排印本页码不同。前奏四章的边码各有前缀: 第~1 章作
        \texttt{A1…A6}, 第~2 章作 \texttt{P1…P65}, 第~3 章作 \texttt{L1…L56},
        第~4 章作罗马数字 \texttt{I…XXII}, 自第一部起作纯数字 \texttt{1…171};
  \item 作者用以标记段落的~$\diamond$~记号, 以及标示后来补写文字的上标
        ~\textsuperscript{$\diamond$}~记号;
  \item 原书大小注(“注~44'”“注~46$^1$”“注~171³”等)的编号与层级。
\end{itemize}

第~12 章“第一部分的注”把第~5 至第~11 章各页脚注汇为一编, 并给每条注配上节题
(如“12.19. 耶稣与十二使徒”), 因此部分注文与正文脚注重复——这是原书的体例, 非
重出之误。这一部分属随笔式的自省文字, 同一句法语在正文脚注与章末汇编注中
各有措辞, 本册有意保留各自的译笔, 不强求统一。

除上述标记外, 译文力求贴近作者驳杂、跳跃而辛辣的文风; 数学名词的译法见各卷
术语表(如~\emph{topos}~不译, \emph{faisceaux pervers}~作“反常层”)。文中的日期
括注(如“(5~月~30~日)”“(1984年3月补写)”)为作者历次补写时的纪日, 一律照录。文中
另有一些对注号、页码的引用(如“n°~1695”“打字稿第~227~页”)沿袭作者打字稿的编号
系统, 与本排印本不尽一致, 亦无法在本书中落实, 悉照原文录出。

全套五卷(另有全书合订一册)的分工是:

\begin{itemize}
  \item \textbf{卷一·前奏}——主题的呈现: 或四个乐章的前奏(原书 p.23--122);
  \item \textbf{卷二·第一部}——愚妄与更新(原书 p.123--254);
  \item \textbf{卷三·第二部}——葬礼(I): 或中国皇帝的新衣(原书 p.255--412);
  \item \textbf{卷四·第三部}——葬礼(II): 或阴阳之钥(原书 p.413--632);
  \item \textbf{卷五·第四部}——葬礼(III): 或四种运算(原书 p.633--929)。
\end{itemize}
}
 之后可用 \hypersetup 覆写)
% --------------------------------------------------------------
\hypersetup{
  pdftitle={收获与播种},
  pdfauthor={Alexandre Grothendieck / 数译 译},
  pdfsubject={Récoltes et Semailles — Réflexions et témoignage sur un passé de mathématicien}
}

% --------------------------------------------------------------
%  标题样式: 仿原书 —— 黑体左对齐, 编号后带句点
% --------------------------------------------------------------
\titleformat{\section}[block]
  {\normalfont\Large\headfont}{\thesection.}{0.5em}{}
\titlespacing*{\section}{0pt}{3.2ex plus 1.2ex minus .2ex}{1.6ex plus .2ex}
\titleformat{\subsection}[block]
  {\normalfont\large\headfont}{\thesubsection.}{0.5em}{}
\titlespacing*{\subsection}{0pt}{2.6ex plus 1ex minus .2ex}{1.2ex plus .2ex}
\titleformat{\subsubsection}[block]
  {\normalfont\normalsize\headfont}{\thesubsubsection.}{0.5em}{}
\titlespacing*{\subsubsection}{0pt}{2.2ex plus .8ex minus .2ex}{1ex plus .2ex}
\titleformat{name=\subsubsection,numberless}[block]
  {\normalfont\normalsize\headfont}{}{0em}{}
\setcounter{secnumdepth}{3}
\setcounter{tocdepth}{3}
% 为含长外文书名与行内公式的段落提供必要的断行余量
\setlength{\emergencystretch}{3em}

% 章标题样式: 原书章标题与节同级排印("13. A) HERITAGE..."),
% 故章标题取块状左对齐黑体, 字号介于节与部之间
\titleformat{\chapter}[block]
  {\normalfont\LARGE\headfont}{\thechapter.}{0.5em}{}
\titlespacing*{\chapter}{0pt}{3.6ex plus 1.4ex minus .2ex}{2ex plus .2ex}
% 目录中章号照原书用阿拉伯数字("13. A) ..." 而非 "第十三章")
\renewcommand{\CTEXthechapter}{\thechapter\hspace{.3em}}

% --------------------------------------------------------------
%  目录观感: 条目与引导点之间、引导点与页码之间各留 0.5em
%  内核 \@dottedtocline 让 \leaders 从标题末字起排, 第一个点几乎
%  贴着末字, 读起来像多标了一个句点; 页码前也没有余量。两处各补
%  0.5em。
% --------------------------------------------------------------
\makeatletter
\patchcmd{\@dottedtocline}{\nobreak\leaders}{\nobreak\hspace{0.5em}\leaders}%
  {\typeout{MT: TOC 引导点前已留 0.5em}}%
  {\typeout{MT-WARN: TOC 引导点前间距补丁未生效}}
\patchcmd{\@dottedtocline}{\hfill\nobreak\hb@xt@\@pnumwidth}%
  {\hfill\nobreak\hspace{0.5em}\hb@xt@\@pnumwidth}%
  {\typeout{MT: TOC 页码前已留 0.5em}}%
  {\typeout{MT-WARN: TOC 页码前间距补丁未生效}}
\makeatother

% --------------------------------------------------------------
%  正文半角标点后的间距
%  本书体例用半角 , ; : ? ! 与 ASCII 括号, 由 xeCJK 按 HalfRight
%  (右半留白) 处理。放宽到 0.45em, 与全角标点右半留白大致齐平。
% --------------------------------------------------------------
\xeCJKsetup{CJKecglue={\hskip 0.45em plus 0.06em minus 0.09em}}

% --------------------------------------------------------------
%  页眉页脚: 偶数页眉右对齐斜体节题; 页码页脚外侧
% --------------------------------------------------------------
\fancyhf{}
\fancyhead[RE]{\small\itshape\rightmark}
\fancyfoot[LE,RO]{\small\thepage}
\renewcommand{\headrulewidth}{0pt}
\renewcommand{\sectionmark}[1]{\markright{\thesection.\ #1}}
\fancypagestyle{plain}{%
  \fancyhf{}%
  \fancyfoot[LE,RO]{\small \thepage}%
  \renewcommand{\headrulewidth}{0pt}%
}

% --------------------------------------------------------------
%  原书版式记号宏
% --------------------------------------------------------------
% 边码: 原打字稿页码(G. 手稿系统里的自引页码), 排在页边
\usepackage{marginnote}
\newcommand{\origpage}[1]{\marginnote{\normalfont\scriptsize p.~#1}}

% ⋄ 记号: 段首/行内正常大小; 上标形式为作者后补文字的位置标记
\DeclareRobustCommand{\pmark}{\ensuremath{\diamond}\hspace{0.15em}}
\DeclareRobustCommand{\dmark}{\textsuperscript{\ensuremath{\diamond}}}

% 脚注人工编号: \fnote{795}{(*)}{文本} —— 编号与(*)角标均照录原书
\newcommand{\fnote}[3]{%
  \begingroup
  \renewcommand{\thefootnote}{#1#2}%
  \footnote{#3}%
  \endgroup
}

% 第一部分的脚注: 页内序号作上标(#1), 注文行首印 "序号(连续注号)",
% 正文在序号之后另行括注同一连续注号 —— \fnoteP{1}{(3)}{注文}
% 原书有三条注无连续注号而改标日期(ch07 注3 = 30 Septembre,
% ch09 注11/12 = 8 août), 第二参数留空, 日期写进注文之首。
\newcommand{\fnoteP}[3]{%
  \begingroup
  \renewcommand{\thefootnote}{#1}\footnotemark
  \renewcommand{\thefootnote}{#1#2}\footnotetext{#3}%
  \endgroup
}

% 正文呼号与注文行首标号不同(原书照录)
\newcommand{\fnotevariant}[4]{%
  \begingroup
  \renewcommand{\thefootnote}{#1#2}\footnotemark
  \renewcommand{\thefootnote}{#1#3}\footnotetext{#4}%
  \endgroup
}

% 大注块标题: \notehead{171³}{斜体评论} —— 正文随后另起
\newcommand{\notehead}[2]{%
  \par\bigskip
  \noindent{\headfont 注~#1}\ifblank{#2}{}{\hspace{0.5em}{\itshape #2}}%
  \par\nobreak\smallskip\nobreak
}

% ∗ ∗ ∗ 分隔
\newcommand{\asterisms}{%
  \par\vspace{1.2ex}%
  \noindent\hspace*{2em}$\ast$\hspace{1.2em}$\ast$\hspace{1.2em}$\ast$%
  \par\vspace{1.2ex}%
}

% 无编号子目的目录条目
\newcommand{\tocsubsub}[1]{%
  \phantomsection
  \addcontentsline{toc}{subsubsection}{#1}%
}

% 部扉: 仿原书部题页
\newcommand{\partdivider}[3]{%
  \clearpage
  \thispagestyle{empty}%
  \phantomsection
  \addcontentsline{toc}{part}{#1\quad #3}%
  \vspace*{6.5cm}
  \begin{center}
    {\fontsize{16pt}{26pt}\selectfont\headfont #1\par}
    \vspace{0.9cm}
    {\fontsize{22pt}{34pt}\selectfont\headfont #2\par}
  \end{center}
  \clearpage
}

% 献辞页: 仿原书部题后的题献, 独占一页且不入目录
\newcommand{\dedication}[1]{%
  \thispagestyle{empty}%
  \vspace*{4cm}%
  \begin{center}#1\end{center}%
  \clearpage
}

% --------------------------------------------------------------
%  各卷"译者说明"的共用体例段(各卷的范围段自写, 后接本段)
% --------------------------------------------------------------
\newcommand{\CommonNoteBody}{%
为便于与原书及学界通行的打字稿页码对勘, 本册保留:

\begin{itemize}
  \item 原书全部章节编号(章~1--18, 节如~5.2、18.2, 子目如~18.1.1)与原书小节标题;
  \item 原书脚注的编号与角标——
        \begin{itemize}
          \item 前奏(第~1--4 章)与第一部的注在各页内重新编号, 正文中的上标即该
                页序号; 第一部紧随其后的括注(如“(3)”)是该部分的连续注号, 与第~12
                章“第一部分的注”里的“注~3”一一对应, 注文行首照原书印作“1(3)”式。
                其中三条注原书未标连续注号而改标日期(ch7 注~3 = 30 Septembre、
                ch9 注~11/12 = 8 août), 本册照录其日期;
          \item 第二部第~13--16~章各自重新编号, 分别为~1--41、1--86、1--132、1--29;
                第三部第~17~章自~1~起至~12, 第~18~章自~1~重新起算并延续至~794;
                末段自~\textbf{795(*)}~续至~1053(*)。各段与原书一一对应;
        \end{itemize}
  \item 页边码“p.~1”“p.~A3”“p.~1066”式小注——这是作者引用自己打字稿原稿时
        所用的页码, 与本排印本页码不同。前奏四章的边码各有前缀: 第~1 章作
        \texttt{A1…A6}, 第~2 章作 \texttt{P1…P65}, 第~3 章作 \texttt{L1…L56},
        第~4 章作罗马数字 \texttt{I…XXII}, 自第一部起作纯数字 \texttt{1…171};
  \item 作者用以标记段落的~$\diamond$~记号, 以及标示后来补写文字的上标
        ~\textsuperscript{$\diamond$}~记号;
  \item 原书大小注(“注~44'”“注~46$^1$”“注~171³”等)的编号与层级。
\end{itemize}

第~12 章“第一部分的注”把第~5 至第~11 章各页脚注汇为一编, 并给每条注配上节题
(如“12.19. 耶稣与十二使徒”), 因此部分注文与正文脚注重复——这是原书的体例, 非
重出之误。这一部分属随笔式的自省文字, 同一句法语在正文脚注与章末汇编注中
各有措辞, 本册有意保留各自的译笔, 不强求统一。

除上述标记外, 译文力求贴近作者驳杂、跳跃而辛辣的文风; 数学名词的译法见各卷
术语表(如~\emph{topos}~不译, \emph{faisceaux pervers}~作“反常层”)。文中的日期
括注(如“(5~月~30~日)”“(1984年3月补写)”)为作者历次补写时的纪日, 一律照录。文中
另有一些对注号、页码的引用(如“n°~1695”“打字稿第~227~页”)沿袭作者打字稿的编号
系统, 与本排印本不尽一致, 亦无法在本书中落实, 悉照原文录出。

全套五卷(另有全书合订一册)的分工是:

\begin{itemize}
  \item \textbf{卷一·前奏}——主题的呈现: 或四个乐章的前奏(原书 p.23--122);
  \item \textbf{卷二·第一部}——愚妄与更新(原书 p.123--254);
  \item \textbf{卷三·第二部}——葬礼(I): 或中国皇帝的新衣(原书 p.255--412);
  \item \textbf{卷四·第三部}——葬礼(II): 或阴阳之钥(原书 p.413--632);
  \item \textbf{卷五·第四部}——葬礼(III): 或四种运算(原书 p.633--929)。
\end{itemize}
}
 之后可用 \hypersetup 覆写)
% --------------------------------------------------------------
\hypersetup{
  pdftitle={收获与播种},
  pdfauthor={Alexandre Grothendieck / 数译 译},
  pdfsubject={Récoltes et Semailles — Réflexions et témoignage sur un passé de mathématicien}
}

% --------------------------------------------------------------
%  标题样式: 仿原书 —— 黑体左对齐, 编号后带句点
% --------------------------------------------------------------
\titleformat{\section}[block]
  {\normalfont\Large\headfont}{\thesection.}{0.5em}{}
\titlespacing*{\section}{0pt}{3.2ex plus 1.2ex minus .2ex}{1.6ex plus .2ex}
\titleformat{\subsection}[block]
  {\normalfont\large\headfont}{\thesubsection.}{0.5em}{}
\titlespacing*{\subsection}{0pt}{2.6ex plus 1ex minus .2ex}{1.2ex plus .2ex}
\titleformat{\subsubsection}[block]
  {\normalfont\normalsize\headfont}{\thesubsubsection.}{0.5em}{}
\titlespacing*{\subsubsection}{0pt}{2.2ex plus .8ex minus .2ex}{1ex plus .2ex}
\titleformat{name=\subsubsection,numberless}[block]
  {\normalfont\normalsize\headfont}{}{0em}{}
\setcounter{secnumdepth}{3}
\setcounter{tocdepth}{3}
% 为含长外文书名与行内公式的段落提供必要的断行余量
\setlength{\emergencystretch}{3em}

% 章标题样式: 原书章标题与节同级排印("13. A) HERITAGE..."),
% 故章标题取块状左对齐黑体, 字号介于节与部之间
\titleformat{\chapter}[block]
  {\normalfont\LARGE\headfont}{\thechapter.}{0.5em}{}
\titlespacing*{\chapter}{0pt}{3.6ex plus 1.4ex minus .2ex}{2ex plus .2ex}
% 目录中章号照原书用阿拉伯数字("13. A) ..." 而非 "第十三章")
\renewcommand{\CTEXthechapter}{\thechapter\hspace{.3em}}

% --------------------------------------------------------------
%  目录观感: 条目与引导点之间、引导点与页码之间各留 0.5em
%  内核 \@dottedtocline 让 \leaders 从标题末字起排, 第一个点几乎
%  贴着末字, 读起来像多标了一个句点; 页码前也没有余量。两处各补
%  0.5em。
% --------------------------------------------------------------
\makeatletter
\patchcmd{\@dottedtocline}{\nobreak\leaders}{\nobreak\hspace{0.5em}\leaders}%
  {\typeout{MT: TOC 引导点前已留 0.5em}}%
  {\typeout{MT-WARN: TOC 引导点前间距补丁未生效}}
\patchcmd{\@dottedtocline}{\hfill\nobreak\hb@xt@\@pnumwidth}%
  {\hfill\nobreak\hspace{0.5em}\hb@xt@\@pnumwidth}%
  {\typeout{MT: TOC 页码前已留 0.5em}}%
  {\typeout{MT-WARN: TOC 页码前间距补丁未生效}}
\makeatother

% --------------------------------------------------------------
%  正文半角标点后的间距
%  本书体例用半角 , ; : ? ! 与 ASCII 括号, 由 xeCJK 按 HalfRight
%  (右半留白) 处理。放宽到 0.45em, 与全角标点右半留白大致齐平。
% --------------------------------------------------------------
\xeCJKsetup{CJKecglue={\hskip 0.45em plus 0.06em minus 0.09em}}

% --------------------------------------------------------------
%  页眉页脚: 偶数页眉右对齐斜体节题; 页码页脚外侧
% --------------------------------------------------------------
\fancyhf{}
\fancyhead[RE]{\small\itshape\rightmark}
\fancyfoot[LE,RO]{\small\thepage}
\renewcommand{\headrulewidth}{0pt}
\renewcommand{\sectionmark}[1]{\markright{\thesection.\ #1}}
\fancypagestyle{plain}{%
  \fancyhf{}%
  \fancyfoot[LE,RO]{\small \thepage}%
  \renewcommand{\headrulewidth}{0pt}%
}

% --------------------------------------------------------------
%  原书版式记号宏
% --------------------------------------------------------------
% 边码: 原打字稿页码(G. 手稿系统里的自引页码), 排在页边
\usepackage{marginnote}
\newcommand{\origpage}[1]{\marginnote{\normalfont\scriptsize p.~#1}}

% ⋄ 记号: 段首/行内正常大小; 上标形式为作者后补文字的位置标记
\DeclareRobustCommand{\pmark}{\ensuremath{\diamond}\hspace{0.15em}}
\DeclareRobustCommand{\dmark}{\textsuperscript{\ensuremath{\diamond}}}

% 脚注人工编号: \fnote{795}{(*)}{文本} —— 编号与(*)角标均照录原书
\newcommand{\fnote}[3]{%
  \begingroup
  \renewcommand{\thefootnote}{#1#2}%
  \footnote{#3}%
  \endgroup
}

% 第一部分的脚注: 页内序号作上标(#1), 注文行首印 "序号(连续注号)",
% 正文在序号之后另行括注同一连续注号 —— \fnoteP{1}{(3)}{注文}
% 原书有三条注无连续注号而改标日期(ch07 注3 = 30 Septembre,
% ch09 注11/12 = 8 août), 第二参数留空, 日期写进注文之首。
\newcommand{\fnoteP}[3]{%
  \begingroup
  \renewcommand{\thefootnote}{#1}\footnotemark
  \renewcommand{\thefootnote}{#1#2}\footnotetext{#3}%
  \endgroup
}

% 正文呼号与注文行首标号不同(原书照录)
\newcommand{\fnotevariant}[4]{%
  \begingroup
  \renewcommand{\thefootnote}{#1#2}\footnotemark
  \renewcommand{\thefootnote}{#1#3}\footnotetext{#4}%
  \endgroup
}

% 大注块标题: \notehead{171³}{斜体评论} —— 正文随后另起
\newcommand{\notehead}[2]{%
  \par\bigskip
  \noindent{\headfont 注~#1}\ifblank{#2}{}{\hspace{0.5em}{\itshape #2}}%
  \par\nobreak\smallskip\nobreak
}

% ∗ ∗ ∗ 分隔
\newcommand{\asterisms}{%
  \par\vspace{1.2ex}%
  \noindent\hspace*{2em}$\ast$\hspace{1.2em}$\ast$\hspace{1.2em}$\ast$%
  \par\vspace{1.2ex}%
}

% 无编号子目的目录条目
\newcommand{\tocsubsub}[1]{%
  \phantomsection
  \addcontentsline{toc}{subsubsection}{#1}%
}

% 部扉: 仿原书部题页
\newcommand{\partdivider}[3]{%
  \clearpage
  \thispagestyle{empty}%
  \phantomsection
  \addcontentsline{toc}{part}{#1\quad #3}%
  \vspace*{6.5cm}
  \begin{center}
    {\fontsize{16pt}{26pt}\selectfont\headfont #1\par}
    \vspace{0.9cm}
    {\fontsize{22pt}{34pt}\selectfont\headfont #2\par}
  \end{center}
  \clearpage
}

% 献辞页: 仿原书部题后的题献, 独占一页且不入目录
\newcommand{\dedication}[1]{%
  \thispagestyle{empty}%
  \vspace*{4cm}%
  \begin{center}#1\end{center}%
  \clearpage
}

% --------------------------------------------------------------
%  各卷"译者说明"的共用体例段(各卷的范围段自写, 后接本段)
% --------------------------------------------------------------
\newcommand{\CommonNoteBody}{%
为便于与原书及学界通行的打字稿页码对勘, 本册保留:

\begin{itemize}
  \item 原书全部章节编号(章~1--18, 节如~5.2、18.2, 子目如~18.1.1)与原书小节标题;
  \item 原书脚注的编号与角标——
        \begin{itemize}
          \item 前奏(第~1--4 章)与第一部的注在各页内重新编号, 正文中的上标即该
                页序号; 第一部紧随其后的括注(如“(3)”)是该部分的连续注号, 与第~12
                章“第一部分的注”里的“注~3”一一对应, 注文行首照原书印作“1(3)”式。
                其中三条注原书未标连续注号而改标日期(ch7 注~3 = 30 Septembre、
                ch9 注~11/12 = 8 août), 本册照录其日期;
          \item 第二部第~13--16~章各自重新编号, 分别为~1--41、1--86、1--132、1--29;
                第三部第~17~章自~1~起至~12, 第~18~章自~1~重新起算并延续至~794;
                末段自~\textbf{795(*)}~续至~1053(*)。各段与原书一一对应;
        \end{itemize}
  \item 页边码“p.~1”“p.~A3”“p.~1066”式小注——这是作者引用自己打字稿原稿时
        所用的页码, 与本排印本页码不同。前奏四章的边码各有前缀: 第~1 章作
        \texttt{A1…A6}, 第~2 章作 \texttt{P1…P65}, 第~3 章作 \texttt{L1…L56},
        第~4 章作罗马数字 \texttt{I…XXII}, 自第一部起作纯数字 \texttt{1…171};
  \item 作者用以标记段落的~$\diamond$~记号, 以及标示后来补写文字的上标
        ~\textsuperscript{$\diamond$}~记号;
  \item 原书大小注(“注~44'”“注~46$^1$”“注~171³”等)的编号与层级。
\end{itemize}

第~12 章“第一部分的注”把第~5 至第~11 章各页脚注汇为一编, 并给每条注配上节题
(如“12.19. 耶稣与十二使徒”), 因此部分注文与正文脚注重复——这是原书的体例, 非
重出之误。这一部分属随笔式的自省文字, 同一句法语在正文脚注与章末汇编注中
各有措辞, 本册有意保留各自的译笔, 不强求统一。

除上述标记外, 译文力求贴近作者驳杂、跳跃而辛辣的文风; 数学名词的译法见各卷
术语表(如~\emph{topos}~不译, \emph{faisceaux pervers}~作“反常层”)。文中的日期
括注(如“(5~月~30~日)”“(1984年3月补写)”)为作者历次补写时的纪日, 一律照录。文中
另有一些对注号、页码的引用(如“n°~1695”“打字稿第~227~页”)沿袭作者打字稿的编号
系统, 与本排印本不尽一致, 亦无法在本书中落实, 悉照原文录出。

全套五卷(另有全书合订一册)的分工是:

\begin{itemize}
  \item \textbf{卷一·前奏}——主题的呈现: 或四个乐章的前奏(原书 p.23--122);
  \item \textbf{卷二·第一部}——愚妄与更新(原书 p.123--254);
  \item \textbf{卷三·第二部}——葬礼(I): 或中国皇帝的新衣(原书 p.255--412);
  \item \textbf{卷四·第三部}——葬礼(II): 或阴阳之钥(原书 p.413--632);
  \item \textbf{卷五·第四部}——葬礼(III): 或四种运算(原书 p.633--929)。
\end{itemize}
}


% --------------------------------------------------------------
%  封面参数
% --------------------------------------------------------------
% ================================================================
%  《收获与播种》统一封面 —— TikZ 原生矢量绘制(不依赖任何外部图片)
%  母题: 暮色中的麦田 —— 犁沟向地平线的灭点收束, 前景麦穗剪影,
%        一位播种者把一把种子撒向天空, 种子沿弧线落入田里。
%  用法(在 \begin{document} 之前, 依次):
%      % ================================================================
%  《收获与播种》中文译本 —— 各卷共用导言区
%  由 vol-prelude / vol-part1 / vol-part2 / vol-part3 / vol-part4
%  以及全书合订 main-all 共同 \input。基类 mathtranslation.cls v3.2 不改。
%  --------------------------------------------------------------
%  样式均在此改写, 以贴近原书:
%    · 章标题仿原书 "5. Travail et découverte"(块状左对齐黑体)
%    · 节/小节: 黑体左对齐, 编号后带句点
%    · 脚注: 人工编号照录原书——前奏与第一部分页内重新编号,
%      第二至四部各自连续编号并带(*)角标; 第一部分另有 \fnoteP
%      处理 "(n)" 式连续注号
%    · 边码: 原打字稿页码 \origpage{N}(前奏作 A1、P1、L1、罗马数字等)
%    · ⋄ 记号: 段首补写标记 \pmark 与上标形式 \dmark
% ================================================================

% 题名含空格, 不能走类选项(会被 kvoptions 吃掉), 在此赋值
\renewcommand{\BookTitleEN}{Récoltes et Semailles}
\renewcommand{\OriginalAuthor}{Alexandre Grothendieck}
\renewcommand{\TranslationDate}{\today}

% --------------------------------------------------------------
%  罕用字后备(如"祐"): Fandol 不载时退到 SimSun
% --------------------------------------------------------------
\setCJKfallbackfamilyfont{\CJKrmdefault}{{SimSun}}

% --------------------------------------------------------------
%  hyperref 元信息(各卷在 % ================================================================
%  《收获与播种》中文译本 —— 各卷共用导言区
%  由 vol-prelude / vol-part1 / vol-part2 / vol-part3 / vol-part4
%  以及全书合订 main-all 共同 \input。基类 mathtranslation.cls v3.2 不改。
%  --------------------------------------------------------------
%  样式均在此改写, 以贴近原书:
%    · 章标题仿原书 "5. Travail et découverte"(块状左对齐黑体)
%    · 节/小节: 黑体左对齐, 编号后带句点
%    · 脚注: 人工编号照录原书——前奏与第一部分页内重新编号,
%      第二至四部各自连续编号并带(*)角标; 第一部分另有 \fnoteP
%      处理 "(n)" 式连续注号
%    · 边码: 原打字稿页码 \origpage{N}(前奏作 A1、P1、L1、罗马数字等)
%    · ⋄ 记号: 段首补写标记 \pmark 与上标形式 \dmark
% ================================================================

% 题名含空格, 不能走类选项(会被 kvoptions 吃掉), 在此赋值
\renewcommand{\BookTitleEN}{Récoltes et Semailles}
\renewcommand{\OriginalAuthor}{Alexandre Grothendieck}
\renewcommand{\TranslationDate}{\today}

% --------------------------------------------------------------
%  罕用字后备(如"祐"): Fandol 不载时退到 SimSun
% --------------------------------------------------------------
\setCJKfallbackfamilyfont{\CJKrmdefault}{{SimSun}}

% --------------------------------------------------------------
%  hyperref 元信息(各卷在 \input{common} 之后可用 \hypersetup 覆写)
% --------------------------------------------------------------
\hypersetup{
  pdftitle={收获与播种},
  pdfauthor={Alexandre Grothendieck / 数译 译},
  pdfsubject={Récoltes et Semailles — Réflexions et témoignage sur un passé de mathématicien}
}

% --------------------------------------------------------------
%  标题样式: 仿原书 —— 黑体左对齐, 编号后带句点
% --------------------------------------------------------------
\titleformat{\section}[block]
  {\normalfont\Large\headfont}{\thesection.}{0.5em}{}
\titlespacing*{\section}{0pt}{3.2ex plus 1.2ex minus .2ex}{1.6ex plus .2ex}
\titleformat{\subsection}[block]
  {\normalfont\large\headfont}{\thesubsection.}{0.5em}{}
\titlespacing*{\subsection}{0pt}{2.6ex plus 1ex minus .2ex}{1.2ex plus .2ex}
\titleformat{\subsubsection}[block]
  {\normalfont\normalsize\headfont}{\thesubsubsection.}{0.5em}{}
\titlespacing*{\subsubsection}{0pt}{2.2ex plus .8ex minus .2ex}{1ex plus .2ex}
\titleformat{name=\subsubsection,numberless}[block]
  {\normalfont\normalsize\headfont}{}{0em}{}
\setcounter{secnumdepth}{3}
\setcounter{tocdepth}{3}
% 为含长外文书名与行内公式的段落提供必要的断行余量
\setlength{\emergencystretch}{3em}

% 章标题样式: 原书章标题与节同级排印("13. A) HERITAGE..."),
% 故章标题取块状左对齐黑体, 字号介于节与部之间
\titleformat{\chapter}[block]
  {\normalfont\LARGE\headfont}{\thechapter.}{0.5em}{}
\titlespacing*{\chapter}{0pt}{3.6ex plus 1.4ex minus .2ex}{2ex plus .2ex}
% 目录中章号照原书用阿拉伯数字("13. A) ..." 而非 "第十三章")
\renewcommand{\CTEXthechapter}{\thechapter\hspace{.3em}}

% --------------------------------------------------------------
%  目录观感: 条目与引导点之间、引导点与页码之间各留 0.5em
%  内核 \@dottedtocline 让 \leaders 从标题末字起排, 第一个点几乎
%  贴着末字, 读起来像多标了一个句点; 页码前也没有余量。两处各补
%  0.5em。
% --------------------------------------------------------------
\makeatletter
\patchcmd{\@dottedtocline}{\nobreak\leaders}{\nobreak\hspace{0.5em}\leaders}%
  {\typeout{MT: TOC 引导点前已留 0.5em}}%
  {\typeout{MT-WARN: TOC 引导点前间距补丁未生效}}
\patchcmd{\@dottedtocline}{\hfill\nobreak\hb@xt@\@pnumwidth}%
  {\hfill\nobreak\hspace{0.5em}\hb@xt@\@pnumwidth}%
  {\typeout{MT: TOC 页码前已留 0.5em}}%
  {\typeout{MT-WARN: TOC 页码前间距补丁未生效}}
\makeatother

% --------------------------------------------------------------
%  正文半角标点后的间距
%  本书体例用半角 , ; : ? ! 与 ASCII 括号, 由 xeCJK 按 HalfRight
%  (右半留白) 处理。放宽到 0.45em, 与全角标点右半留白大致齐平。
% --------------------------------------------------------------
\xeCJKsetup{CJKecglue={\hskip 0.45em plus 0.06em minus 0.09em}}

% --------------------------------------------------------------
%  页眉页脚: 偶数页眉右对齐斜体节题; 页码页脚外侧
% --------------------------------------------------------------
\fancyhf{}
\fancyhead[RE]{\small\itshape\rightmark}
\fancyfoot[LE,RO]{\small\thepage}
\renewcommand{\headrulewidth}{0pt}
\renewcommand{\sectionmark}[1]{\markright{\thesection.\ #1}}
\fancypagestyle{plain}{%
  \fancyhf{}%
  \fancyfoot[LE,RO]{\small \thepage}%
  \renewcommand{\headrulewidth}{0pt}%
}

% --------------------------------------------------------------
%  原书版式记号宏
% --------------------------------------------------------------
% 边码: 原打字稿页码(G. 手稿系统里的自引页码), 排在页边
\usepackage{marginnote}
\newcommand{\origpage}[1]{\marginnote{\normalfont\scriptsize p.~#1}}

% ⋄ 记号: 段首/行内正常大小; 上标形式为作者后补文字的位置标记
\DeclareRobustCommand{\pmark}{\ensuremath{\diamond}\hspace{0.15em}}
\DeclareRobustCommand{\dmark}{\textsuperscript{\ensuremath{\diamond}}}

% 脚注人工编号: \fnote{795}{(*)}{文本} —— 编号与(*)角标均照录原书
\newcommand{\fnote}[3]{%
  \begingroup
  \renewcommand{\thefootnote}{#1#2}%
  \footnote{#3}%
  \endgroup
}

% 第一部分的脚注: 页内序号作上标(#1), 注文行首印 "序号(连续注号)",
% 正文在序号之后另行括注同一连续注号 —— \fnoteP{1}{(3)}{注文}
% 原书有三条注无连续注号而改标日期(ch07 注3 = 30 Septembre,
% ch09 注11/12 = 8 août), 第二参数留空, 日期写进注文之首。
\newcommand{\fnoteP}[3]{%
  \begingroup
  \renewcommand{\thefootnote}{#1}\footnotemark
  \renewcommand{\thefootnote}{#1#2}\footnotetext{#3}%
  \endgroup
}

% 正文呼号与注文行首标号不同(原书照录)
\newcommand{\fnotevariant}[4]{%
  \begingroup
  \renewcommand{\thefootnote}{#1#2}\footnotemark
  \renewcommand{\thefootnote}{#1#3}\footnotetext{#4}%
  \endgroup
}

% 大注块标题: \notehead{171³}{斜体评论} —— 正文随后另起
\newcommand{\notehead}[2]{%
  \par\bigskip
  \noindent{\headfont 注~#1}\ifblank{#2}{}{\hspace{0.5em}{\itshape #2}}%
  \par\nobreak\smallskip\nobreak
}

% ∗ ∗ ∗ 分隔
\newcommand{\asterisms}{%
  \par\vspace{1.2ex}%
  \noindent\hspace*{2em}$\ast$\hspace{1.2em}$\ast$\hspace{1.2em}$\ast$%
  \par\vspace{1.2ex}%
}

% 无编号子目的目录条目
\newcommand{\tocsubsub}[1]{%
  \phantomsection
  \addcontentsline{toc}{subsubsection}{#1}%
}

% 部扉: 仿原书部题页
\newcommand{\partdivider}[3]{%
  \clearpage
  \thispagestyle{empty}%
  \phantomsection
  \addcontentsline{toc}{part}{#1\quad #3}%
  \vspace*{6.5cm}
  \begin{center}
    {\fontsize{16pt}{26pt}\selectfont\headfont #1\par}
    \vspace{0.9cm}
    {\fontsize{22pt}{34pt}\selectfont\headfont #2\par}
  \end{center}
  \clearpage
}

% 献辞页: 仿原书部题后的题献, 独占一页且不入目录
\newcommand{\dedication}[1]{%
  \thispagestyle{empty}%
  \vspace*{4cm}%
  \begin{center}#1\end{center}%
  \clearpage
}

% --------------------------------------------------------------
%  各卷"译者说明"的共用体例段(各卷的范围段自写, 后接本段)
% --------------------------------------------------------------
\newcommand{\CommonNoteBody}{%
为便于与原书及学界通行的打字稿页码对勘, 本册保留:

\begin{itemize}
  \item 原书全部章节编号(章~1--18, 节如~5.2、18.2, 子目如~18.1.1)与原书小节标题;
  \item 原书脚注的编号与角标——
        \begin{itemize}
          \item 前奏(第~1--4 章)与第一部的注在各页内重新编号, 正文中的上标即该
                页序号; 第一部紧随其后的括注(如“(3)”)是该部分的连续注号, 与第~12
                章“第一部分的注”里的“注~3”一一对应, 注文行首照原书印作“1(3)”式。
                其中三条注原书未标连续注号而改标日期(ch7 注~3 = 30 Septembre、
                ch9 注~11/12 = 8 août), 本册照录其日期;
          \item 第二部第~13--16~章各自重新编号, 分别为~1--41、1--86、1--132、1--29;
                第三部第~17~章自~1~起至~12, 第~18~章自~1~重新起算并延续至~794;
                末段自~\textbf{795(*)}~续至~1053(*)。各段与原书一一对应;
        \end{itemize}
  \item 页边码“p.~1”“p.~A3”“p.~1066”式小注——这是作者引用自己打字稿原稿时
        所用的页码, 与本排印本页码不同。前奏四章的边码各有前缀: 第~1 章作
        \texttt{A1…A6}, 第~2 章作 \texttt{P1…P65}, 第~3 章作 \texttt{L1…L56},
        第~4 章作罗马数字 \texttt{I…XXII}, 自第一部起作纯数字 \texttt{1…171};
  \item 作者用以标记段落的~$\diamond$~记号, 以及标示后来补写文字的上标
        ~\textsuperscript{$\diamond$}~记号;
  \item 原书大小注(“注~44'”“注~46$^1$”“注~171³”等)的编号与层级。
\end{itemize}

第~12 章“第一部分的注”把第~5 至第~11 章各页脚注汇为一编, 并给每条注配上节题
(如“12.19. 耶稣与十二使徒”), 因此部分注文与正文脚注重复——这是原书的体例, 非
重出之误。这一部分属随笔式的自省文字, 同一句法语在正文脚注与章末汇编注中
各有措辞, 本册有意保留各自的译笔, 不强求统一。

除上述标记外, 译文力求贴近作者驳杂、跳跃而辛辣的文风; 数学名词的译法见各卷
术语表(如~\emph{topos}~不译, \emph{faisceaux pervers}~作“反常层”)。文中的日期
括注(如“(5~月~30~日)”“(1984年3月补写)”)为作者历次补写时的纪日, 一律照录。文中
另有一些对注号、页码的引用(如“n°~1695”“打字稿第~227~页”)沿袭作者打字稿的编号
系统, 与本排印本不尽一致, 亦无法在本书中落实, 悉照原文录出。

全套五卷(另有全书合订一册)的分工是:

\begin{itemize}
  \item \textbf{卷一·前奏}——主题的呈现: 或四个乐章的前奏(原书 p.23--122);
  \item \textbf{卷二·第一部}——愚妄与更新(原书 p.123--254);
  \item \textbf{卷三·第二部}——葬礼(I): 或中国皇帝的新衣(原书 p.255--412);
  \item \textbf{卷四·第三部}——葬礼(II): 或阴阳之钥(原书 p.413--632);
  \item \textbf{卷五·第四部}——葬礼(III): 或四种运算(原书 p.633--929)。
\end{itemize}
}
 之后可用 \hypersetup 覆写)
% --------------------------------------------------------------
\hypersetup{
  pdftitle={收获与播种},
  pdfauthor={Alexandre Grothendieck / 数译 译},
  pdfsubject={Récoltes et Semailles — Réflexions et témoignage sur un passé de mathématicien}
}

% --------------------------------------------------------------
%  标题样式: 仿原书 —— 黑体左对齐, 编号后带句点
% --------------------------------------------------------------
\titleformat{\section}[block]
  {\normalfont\Large\headfont}{\thesection.}{0.5em}{}
\titlespacing*{\section}{0pt}{3.2ex plus 1.2ex minus .2ex}{1.6ex plus .2ex}
\titleformat{\subsection}[block]
  {\normalfont\large\headfont}{\thesubsection.}{0.5em}{}
\titlespacing*{\subsection}{0pt}{2.6ex plus 1ex minus .2ex}{1.2ex plus .2ex}
\titleformat{\subsubsection}[block]
  {\normalfont\normalsize\headfont}{\thesubsubsection.}{0.5em}{}
\titlespacing*{\subsubsection}{0pt}{2.2ex plus .8ex minus .2ex}{1ex plus .2ex}
\titleformat{name=\subsubsection,numberless}[block]
  {\normalfont\normalsize\headfont}{}{0em}{}
\setcounter{secnumdepth}{3}
\setcounter{tocdepth}{3}
% 为含长外文书名与行内公式的段落提供必要的断行余量
\setlength{\emergencystretch}{3em}

% 章标题样式: 原书章标题与节同级排印("13. A) HERITAGE..."),
% 故章标题取块状左对齐黑体, 字号介于节与部之间
\titleformat{\chapter}[block]
  {\normalfont\LARGE\headfont}{\thechapter.}{0.5em}{}
\titlespacing*{\chapter}{0pt}{3.6ex plus 1.4ex minus .2ex}{2ex plus .2ex}
% 目录中章号照原书用阿拉伯数字("13. A) ..." 而非 "第十三章")
\renewcommand{\CTEXthechapter}{\thechapter\hspace{.3em}}

% --------------------------------------------------------------
%  目录观感: 条目与引导点之间、引导点与页码之间各留 0.5em
%  内核 \@dottedtocline 让 \leaders 从标题末字起排, 第一个点几乎
%  贴着末字, 读起来像多标了一个句点; 页码前也没有余量。两处各补
%  0.5em。
% --------------------------------------------------------------
\makeatletter
\patchcmd{\@dottedtocline}{\nobreak\leaders}{\nobreak\hspace{0.5em}\leaders}%
  {\typeout{MT: TOC 引导点前已留 0.5em}}%
  {\typeout{MT-WARN: TOC 引导点前间距补丁未生效}}
\patchcmd{\@dottedtocline}{\hfill\nobreak\hb@xt@\@pnumwidth}%
  {\hfill\nobreak\hspace{0.5em}\hb@xt@\@pnumwidth}%
  {\typeout{MT: TOC 页码前已留 0.5em}}%
  {\typeout{MT-WARN: TOC 页码前间距补丁未生效}}
\makeatother

% --------------------------------------------------------------
%  正文半角标点后的间距
%  本书体例用半角 , ; : ? ! 与 ASCII 括号, 由 xeCJK 按 HalfRight
%  (右半留白) 处理。放宽到 0.45em, 与全角标点右半留白大致齐平。
% --------------------------------------------------------------
\xeCJKsetup{CJKecglue={\hskip 0.45em plus 0.06em minus 0.09em}}

% --------------------------------------------------------------
%  页眉页脚: 偶数页眉右对齐斜体节题; 页码页脚外侧
% --------------------------------------------------------------
\fancyhf{}
\fancyhead[RE]{\small\itshape\rightmark}
\fancyfoot[LE,RO]{\small\thepage}
\renewcommand{\headrulewidth}{0pt}
\renewcommand{\sectionmark}[1]{\markright{\thesection.\ #1}}
\fancypagestyle{plain}{%
  \fancyhf{}%
  \fancyfoot[LE,RO]{\small \thepage}%
  \renewcommand{\headrulewidth}{0pt}%
}

% --------------------------------------------------------------
%  原书版式记号宏
% --------------------------------------------------------------
% 边码: 原打字稿页码(G. 手稿系统里的自引页码), 排在页边
\usepackage{marginnote}
\newcommand{\origpage}[1]{\marginnote{\normalfont\scriptsize p.~#1}}

% ⋄ 记号: 段首/行内正常大小; 上标形式为作者后补文字的位置标记
\DeclareRobustCommand{\pmark}{\ensuremath{\diamond}\hspace{0.15em}}
\DeclareRobustCommand{\dmark}{\textsuperscript{\ensuremath{\diamond}}}

% 脚注人工编号: \fnote{795}{(*)}{文本} —— 编号与(*)角标均照录原书
\newcommand{\fnote}[3]{%
  \begingroup
  \renewcommand{\thefootnote}{#1#2}%
  \footnote{#3}%
  \endgroup
}

% 第一部分的脚注: 页内序号作上标(#1), 注文行首印 "序号(连续注号)",
% 正文在序号之后另行括注同一连续注号 —— \fnoteP{1}{(3)}{注文}
% 原书有三条注无连续注号而改标日期(ch07 注3 = 30 Septembre,
% ch09 注11/12 = 8 août), 第二参数留空, 日期写进注文之首。
\newcommand{\fnoteP}[3]{%
  \begingroup
  \renewcommand{\thefootnote}{#1}\footnotemark
  \renewcommand{\thefootnote}{#1#2}\footnotetext{#3}%
  \endgroup
}

% 正文呼号与注文行首标号不同(原书照录)
\newcommand{\fnotevariant}[4]{%
  \begingroup
  \renewcommand{\thefootnote}{#1#2}\footnotemark
  \renewcommand{\thefootnote}{#1#3}\footnotetext{#4}%
  \endgroup
}

% 大注块标题: \notehead{171³}{斜体评论} —— 正文随后另起
\newcommand{\notehead}[2]{%
  \par\bigskip
  \noindent{\headfont 注~#1}\ifblank{#2}{}{\hspace{0.5em}{\itshape #2}}%
  \par\nobreak\smallskip\nobreak
}

% ∗ ∗ ∗ 分隔
\newcommand{\asterisms}{%
  \par\vspace{1.2ex}%
  \noindent\hspace*{2em}$\ast$\hspace{1.2em}$\ast$\hspace{1.2em}$\ast$%
  \par\vspace{1.2ex}%
}

% 无编号子目的目录条目
\newcommand{\tocsubsub}[1]{%
  \phantomsection
  \addcontentsline{toc}{subsubsection}{#1}%
}

% 部扉: 仿原书部题页
\newcommand{\partdivider}[3]{%
  \clearpage
  \thispagestyle{empty}%
  \phantomsection
  \addcontentsline{toc}{part}{#1\quad #3}%
  \vspace*{6.5cm}
  \begin{center}
    {\fontsize{16pt}{26pt}\selectfont\headfont #1\par}
    \vspace{0.9cm}
    {\fontsize{22pt}{34pt}\selectfont\headfont #2\par}
  \end{center}
  \clearpage
}

% 献辞页: 仿原书部题后的题献, 独占一页且不入目录
\newcommand{\dedication}[1]{%
  \thispagestyle{empty}%
  \vspace*{4cm}%
  \begin{center}#1\end{center}%
  \clearpage
}

% --------------------------------------------------------------
%  各卷"译者说明"的共用体例段(各卷的范围段自写, 后接本段)
% --------------------------------------------------------------
\newcommand{\CommonNoteBody}{%
为便于与原书及学界通行的打字稿页码对勘, 本册保留:

\begin{itemize}
  \item 原书全部章节编号(章~1--18, 节如~5.2、18.2, 子目如~18.1.1)与原书小节标题;
  \item 原书脚注的编号与角标——
        \begin{itemize}
          \item 前奏(第~1--4 章)与第一部的注在各页内重新编号, 正文中的上标即该
                页序号; 第一部紧随其后的括注(如“(3)”)是该部分的连续注号, 与第~12
                章“第一部分的注”里的“注~3”一一对应, 注文行首照原书印作“1(3)”式。
                其中三条注原书未标连续注号而改标日期(ch7 注~3 = 30 Septembre、
                ch9 注~11/12 = 8 août), 本册照录其日期;
          \item 第二部第~13--16~章各自重新编号, 分别为~1--41、1--86、1--132、1--29;
                第三部第~17~章自~1~起至~12, 第~18~章自~1~重新起算并延续至~794;
                末段自~\textbf{795(*)}~续至~1053(*)。各段与原书一一对应;
        \end{itemize}
  \item 页边码“p.~1”“p.~A3”“p.~1066”式小注——这是作者引用自己打字稿原稿时
        所用的页码, 与本排印本页码不同。前奏四章的边码各有前缀: 第~1 章作
        \texttt{A1…A6}, 第~2 章作 \texttt{P1…P65}, 第~3 章作 \texttt{L1…L56},
        第~4 章作罗马数字 \texttt{I…XXII}, 自第一部起作纯数字 \texttt{1…171};
  \item 作者用以标记段落的~$\diamond$~记号, 以及标示后来补写文字的上标
        ~\textsuperscript{$\diamond$}~记号;
  \item 原书大小注(“注~44'”“注~46$^1$”“注~171³”等)的编号与层级。
\end{itemize}

第~12 章“第一部分的注”把第~5 至第~11 章各页脚注汇为一编, 并给每条注配上节题
(如“12.19. 耶稣与十二使徒”), 因此部分注文与正文脚注重复——这是原书的体例, 非
重出之误。这一部分属随笔式的自省文字, 同一句法语在正文脚注与章末汇编注中
各有措辞, 本册有意保留各自的译笔, 不强求统一。

除上述标记外, 译文力求贴近作者驳杂、跳跃而辛辣的文风; 数学名词的译法见各卷
术语表(如~\emph{topos}~不译, \emph{faisceaux pervers}~作“反常层”)。文中的日期
括注(如“(5~月~30~日)”“(1984年3月补写)”)为作者历次补写时的纪日, 一律照录。文中
另有一些对注号、页码的引用(如“n°~1695”“打字稿第~227~页”)沿袭作者打字稿的编号
系统, 与本排印本不尽一致, 亦无法在本书中落实, 悉照原文录出。

全套五卷(另有全书合订一册)的分工是:

\begin{itemize}
  \item \textbf{卷一·前奏}——主题的呈现: 或四个乐章的前奏(原书 p.23--122);
  \item \textbf{卷二·第一部}——愚妄与更新(原书 p.123--254);
  \item \textbf{卷三·第二部}——葬礼(I): 或中国皇帝的新衣(原书 p.255--412);
  \item \textbf{卷四·第三部}——葬礼(II): 或阴阳之钥(原书 p.413--632);
  \item \textbf{卷五·第四部}——葬礼(III): 或四种运算(原书 p.633--929)。
\end{itemize}
}
                       % 共用导言区
%      % ================================================================
%  《收获与播种》统一封面 —— TikZ 原生矢量绘制(不依赖任何外部图片)
%  母题: 暮色中的麦田 —— 犁沟向地平线的灭点收束, 前景麦穗剪影,
%        一位播种者把一把种子撒向天空, 种子沿弧线落入田里。
%  用法(在 \begin{document} 之前, 依次):
%      % ================================================================
%  《收获与播种》中文译本 —— 各卷共用导言区
%  由 vol-prelude / vol-part1 / vol-part2 / vol-part3 / vol-part4
%  以及全书合订 main-all 共同 \input。基类 mathtranslation.cls v3.2 不改。
%  --------------------------------------------------------------
%  样式均在此改写, 以贴近原书:
%    · 章标题仿原书 "5. Travail et découverte"(块状左对齐黑体)
%    · 节/小节: 黑体左对齐, 编号后带句点
%    · 脚注: 人工编号照录原书——前奏与第一部分页内重新编号,
%      第二至四部各自连续编号并带(*)角标; 第一部分另有 \fnoteP
%      处理 "(n)" 式连续注号
%    · 边码: 原打字稿页码 \origpage{N}(前奏作 A1、P1、L1、罗马数字等)
%    · ⋄ 记号: 段首补写标记 \pmark 与上标形式 \dmark
% ================================================================

% 题名含空格, 不能走类选项(会被 kvoptions 吃掉), 在此赋值
\renewcommand{\BookTitleEN}{Récoltes et Semailles}
\renewcommand{\OriginalAuthor}{Alexandre Grothendieck}
\renewcommand{\TranslationDate}{\today}

% --------------------------------------------------------------
%  罕用字后备(如"祐"): Fandol 不载时退到 SimSun
% --------------------------------------------------------------
\setCJKfallbackfamilyfont{\CJKrmdefault}{{SimSun}}

% --------------------------------------------------------------
%  hyperref 元信息(各卷在 \input{common} 之后可用 \hypersetup 覆写)
% --------------------------------------------------------------
\hypersetup{
  pdftitle={收获与播种},
  pdfauthor={Alexandre Grothendieck / 数译 译},
  pdfsubject={Récoltes et Semailles — Réflexions et témoignage sur un passé de mathématicien}
}

% --------------------------------------------------------------
%  标题样式: 仿原书 —— 黑体左对齐, 编号后带句点
% --------------------------------------------------------------
\titleformat{\section}[block]
  {\normalfont\Large\headfont}{\thesection.}{0.5em}{}
\titlespacing*{\section}{0pt}{3.2ex plus 1.2ex minus .2ex}{1.6ex plus .2ex}
\titleformat{\subsection}[block]
  {\normalfont\large\headfont}{\thesubsection.}{0.5em}{}
\titlespacing*{\subsection}{0pt}{2.6ex plus 1ex minus .2ex}{1.2ex plus .2ex}
\titleformat{\subsubsection}[block]
  {\normalfont\normalsize\headfont}{\thesubsubsection.}{0.5em}{}
\titlespacing*{\subsubsection}{0pt}{2.2ex plus .8ex minus .2ex}{1ex plus .2ex}
\titleformat{name=\subsubsection,numberless}[block]
  {\normalfont\normalsize\headfont}{}{0em}{}
\setcounter{secnumdepth}{3}
\setcounter{tocdepth}{3}
% 为含长外文书名与行内公式的段落提供必要的断行余量
\setlength{\emergencystretch}{3em}

% 章标题样式: 原书章标题与节同级排印("13. A) HERITAGE..."),
% 故章标题取块状左对齐黑体, 字号介于节与部之间
\titleformat{\chapter}[block]
  {\normalfont\LARGE\headfont}{\thechapter.}{0.5em}{}
\titlespacing*{\chapter}{0pt}{3.6ex plus 1.4ex minus .2ex}{2ex plus .2ex}
% 目录中章号照原书用阿拉伯数字("13. A) ..." 而非 "第十三章")
\renewcommand{\CTEXthechapter}{\thechapter\hspace{.3em}}

% --------------------------------------------------------------
%  目录观感: 条目与引导点之间、引导点与页码之间各留 0.5em
%  内核 \@dottedtocline 让 \leaders 从标题末字起排, 第一个点几乎
%  贴着末字, 读起来像多标了一个句点; 页码前也没有余量。两处各补
%  0.5em。
% --------------------------------------------------------------
\makeatletter
\patchcmd{\@dottedtocline}{\nobreak\leaders}{\nobreak\hspace{0.5em}\leaders}%
  {\typeout{MT: TOC 引导点前已留 0.5em}}%
  {\typeout{MT-WARN: TOC 引导点前间距补丁未生效}}
\patchcmd{\@dottedtocline}{\hfill\nobreak\hb@xt@\@pnumwidth}%
  {\hfill\nobreak\hspace{0.5em}\hb@xt@\@pnumwidth}%
  {\typeout{MT: TOC 页码前已留 0.5em}}%
  {\typeout{MT-WARN: TOC 页码前间距补丁未生效}}
\makeatother

% --------------------------------------------------------------
%  正文半角标点后的间距
%  本书体例用半角 , ; : ? ! 与 ASCII 括号, 由 xeCJK 按 HalfRight
%  (右半留白) 处理。放宽到 0.45em, 与全角标点右半留白大致齐平。
% --------------------------------------------------------------
\xeCJKsetup{CJKecglue={\hskip 0.45em plus 0.06em minus 0.09em}}

% --------------------------------------------------------------
%  页眉页脚: 偶数页眉右对齐斜体节题; 页码页脚外侧
% --------------------------------------------------------------
\fancyhf{}
\fancyhead[RE]{\small\itshape\rightmark}
\fancyfoot[LE,RO]{\small\thepage}
\renewcommand{\headrulewidth}{0pt}
\renewcommand{\sectionmark}[1]{\markright{\thesection.\ #1}}
\fancypagestyle{plain}{%
  \fancyhf{}%
  \fancyfoot[LE,RO]{\small \thepage}%
  \renewcommand{\headrulewidth}{0pt}%
}

% --------------------------------------------------------------
%  原书版式记号宏
% --------------------------------------------------------------
% 边码: 原打字稿页码(G. 手稿系统里的自引页码), 排在页边
\usepackage{marginnote}
\newcommand{\origpage}[1]{\marginnote{\normalfont\scriptsize p.~#1}}

% ⋄ 记号: 段首/行内正常大小; 上标形式为作者后补文字的位置标记
\DeclareRobustCommand{\pmark}{\ensuremath{\diamond}\hspace{0.15em}}
\DeclareRobustCommand{\dmark}{\textsuperscript{\ensuremath{\diamond}}}

% 脚注人工编号: \fnote{795}{(*)}{文本} —— 编号与(*)角标均照录原书
\newcommand{\fnote}[3]{%
  \begingroup
  \renewcommand{\thefootnote}{#1#2}%
  \footnote{#3}%
  \endgroup
}

% 第一部分的脚注: 页内序号作上标(#1), 注文行首印 "序号(连续注号)",
% 正文在序号之后另行括注同一连续注号 —— \fnoteP{1}{(3)}{注文}
% 原书有三条注无连续注号而改标日期(ch07 注3 = 30 Septembre,
% ch09 注11/12 = 8 août), 第二参数留空, 日期写进注文之首。
\newcommand{\fnoteP}[3]{%
  \begingroup
  \renewcommand{\thefootnote}{#1}\footnotemark
  \renewcommand{\thefootnote}{#1#2}\footnotetext{#3}%
  \endgroup
}

% 正文呼号与注文行首标号不同(原书照录)
\newcommand{\fnotevariant}[4]{%
  \begingroup
  \renewcommand{\thefootnote}{#1#2}\footnotemark
  \renewcommand{\thefootnote}{#1#3}\footnotetext{#4}%
  \endgroup
}

% 大注块标题: \notehead{171³}{斜体评论} —— 正文随后另起
\newcommand{\notehead}[2]{%
  \par\bigskip
  \noindent{\headfont 注~#1}\ifblank{#2}{}{\hspace{0.5em}{\itshape #2}}%
  \par\nobreak\smallskip\nobreak
}

% ∗ ∗ ∗ 分隔
\newcommand{\asterisms}{%
  \par\vspace{1.2ex}%
  \noindent\hspace*{2em}$\ast$\hspace{1.2em}$\ast$\hspace{1.2em}$\ast$%
  \par\vspace{1.2ex}%
}

% 无编号子目的目录条目
\newcommand{\tocsubsub}[1]{%
  \phantomsection
  \addcontentsline{toc}{subsubsection}{#1}%
}

% 部扉: 仿原书部题页
\newcommand{\partdivider}[3]{%
  \clearpage
  \thispagestyle{empty}%
  \phantomsection
  \addcontentsline{toc}{part}{#1\quad #3}%
  \vspace*{6.5cm}
  \begin{center}
    {\fontsize{16pt}{26pt}\selectfont\headfont #1\par}
    \vspace{0.9cm}
    {\fontsize{22pt}{34pt}\selectfont\headfont #2\par}
  \end{center}
  \clearpage
}

% 献辞页: 仿原书部题后的题献, 独占一页且不入目录
\newcommand{\dedication}[1]{%
  \thispagestyle{empty}%
  \vspace*{4cm}%
  \begin{center}#1\end{center}%
  \clearpage
}

% --------------------------------------------------------------
%  各卷"译者说明"的共用体例段(各卷的范围段自写, 后接本段)
% --------------------------------------------------------------
\newcommand{\CommonNoteBody}{%
为便于与原书及学界通行的打字稿页码对勘, 本册保留:

\begin{itemize}
  \item 原书全部章节编号(章~1--18, 节如~5.2、18.2, 子目如~18.1.1)与原书小节标题;
  \item 原书脚注的编号与角标——
        \begin{itemize}
          \item 前奏(第~1--4 章)与第一部的注在各页内重新编号, 正文中的上标即该
                页序号; 第一部紧随其后的括注(如“(3)”)是该部分的连续注号, 与第~12
                章“第一部分的注”里的“注~3”一一对应, 注文行首照原书印作“1(3)”式。
                其中三条注原书未标连续注号而改标日期(ch7 注~3 = 30 Septembre、
                ch9 注~11/12 = 8 août), 本册照录其日期;
          \item 第二部第~13--16~章各自重新编号, 分别为~1--41、1--86、1--132、1--29;
                第三部第~17~章自~1~起至~12, 第~18~章自~1~重新起算并延续至~794;
                末段自~\textbf{795(*)}~续至~1053(*)。各段与原书一一对应;
        \end{itemize}
  \item 页边码“p.~1”“p.~A3”“p.~1066”式小注——这是作者引用自己打字稿原稿时
        所用的页码, 与本排印本页码不同。前奏四章的边码各有前缀: 第~1 章作
        \texttt{A1…A6}, 第~2 章作 \texttt{P1…P65}, 第~3 章作 \texttt{L1…L56},
        第~4 章作罗马数字 \texttt{I…XXII}, 自第一部起作纯数字 \texttt{1…171};
  \item 作者用以标记段落的~$\diamond$~记号, 以及标示后来补写文字的上标
        ~\textsuperscript{$\diamond$}~记号;
  \item 原书大小注(“注~44'”“注~46$^1$”“注~171³”等)的编号与层级。
\end{itemize}

第~12 章“第一部分的注”把第~5 至第~11 章各页脚注汇为一编, 并给每条注配上节题
(如“12.19. 耶稣与十二使徒”), 因此部分注文与正文脚注重复——这是原书的体例, 非
重出之误。这一部分属随笔式的自省文字, 同一句法语在正文脚注与章末汇编注中
各有措辞, 本册有意保留各自的译笔, 不强求统一。

除上述标记外, 译文力求贴近作者驳杂、跳跃而辛辣的文风; 数学名词的译法见各卷
术语表(如~\emph{topos}~不译, \emph{faisceaux pervers}~作“反常层”)。文中的日期
括注(如“(5~月~30~日)”“(1984年3月补写)”)为作者历次补写时的纪日, 一律照录。文中
另有一些对注号、页码的引用(如“n°~1695”“打字稿第~227~页”)沿袭作者打字稿的编号
系统, 与本排印本不尽一致, 亦无法在本书中落实, 悉照原文录出。

全套五卷(另有全书合订一册)的分工是:

\begin{itemize}
  \item \textbf{卷一·前奏}——主题的呈现: 或四个乐章的前奏(原书 p.23--122);
  \item \textbf{卷二·第一部}——愚妄与更新(原书 p.123--254);
  \item \textbf{卷三·第二部}——葬礼(I): 或中国皇帝的新衣(原书 p.255--412);
  \item \textbf{卷四·第三部}——葬礼(II): 或阴阳之钥(原书 p.413--632);
  \item \textbf{卷五·第四部}——葬礼(III): 或四种运算(原书 p.633--929)。
\end{itemize}
}
                       % 共用导言区
%      % ================================================================
%  《收获与播种》统一封面 —— TikZ 原生矢量绘制(不依赖任何外部图片)
%  母题: 暮色中的麦田 —— 犁沟向地平线的灭点收束, 前景麦穗剪影,
%        一位播种者把一把种子撒向天空, 种子沿弧线落入田里。
%  用法(在 \begin{document} 之前, 依次):
%      \input{common}                       % 共用导言区
%      \input{cover-art}                    % 本文件
%      \CoverSetup{prelude}{卷一}{前奏}%
%                 {主题的呈现——或四个乐章的前奏}{p.23–122}
%  正文中再 \makecover 即可。配色关键字:
%      prelude / part1 / part2 / part3 / part4 / all
%  五卷共用同一幅画, 只换配色与卷次文字, 以保证"一套书"的观感。
% ================================================================

\RequirePackage{tikz}

% --------------------------------------------------------------
%  配色: 六套(前奏 / 四部 / 全书), 各为同一暮色母题的变调
% --------------------------------------------------------------
\newcommand{\CoverPalette}[1]{%
  \ifstrequal{#1}{prelude}{%
    \definecolor{cvSkyTop}{RGB}{10,26,48}%
    \definecolor{cvSkyMid}{RGB}{46,84,116}%
    \definecolor{cvSkyBot}{RGB}{232,164,92}%
    \definecolor{cvSun}{RGB}{255,224,168}%
    \definecolor{cvFieldTop}{RGB}{150,110,52}%
    \definecolor{cvFieldBot}{RGB}{58,36,18}%
    \definecolor{cvFurrow}{RGB}{40,22,10}%
    \definecolor{cvSil}{RGB}{30,18,10}%
    \definecolor{cvCream}{RGB}{247,240,228}%
  }{%
  \ifstrequal{#1}{part1}{%
    \definecolor{cvSkyTop}{RGB}{22,70,100}%
    \definecolor{cvSkyMid}{RGB}{92,150,168}%
    \definecolor{cvSkyBot}{RGB}{240,200,104}%
    \definecolor{cvSun}{RGB}{255,244,200}%
    \definecolor{cvFieldTop}{RGB}{176,140,58}%
    \definecolor{cvFieldBot}{RGB}{84,54,22}%
    \definecolor{cvFurrow}{RGB}{45,28,12}%
    \definecolor{cvSil}{RGB}{34,22,10}%
    \definecolor{cvCream}{RGB}{250,245,232}%
  }{%
  \ifstrequal{#1}{part2}{%
    \definecolor{cvSkyTop}{RGB}{46,20,30}%
    \definecolor{cvSkyMid}{RGB}{128,52,54}%
    \definecolor{cvSkyBot}{RGB}{226,120,64}%
    \definecolor{cvSun}{RGB}{255,206,150}%
    \definecolor{cvFieldTop}{RGB}{132,84,44}%
    \definecolor{cvFieldBot}{RGB}{48,26,18}%
    \definecolor{cvFurrow}{RGB}{36,16,14}%
    \definecolor{cvSil}{RGB}{26,12,12}%
    \definecolor{cvCream}{RGB}{248,236,224}%
  }{%
  \ifstrequal{#1}{part3}{%
    \definecolor{cvSkyTop}{RGB}{30,20,54}%
    \definecolor{cvSkyMid}{RGB}{92,62,124}%
    \definecolor{cvSkyBot}{RGB}{216,132,140}%
    \definecolor{cvSun}{RGB}{255,214,196}%
    \definecolor{cvFieldTop}{RGB}{128,90,74}%
    \definecolor{cvFieldBot}{RGB}{44,26,30}%
    \definecolor{cvFurrow}{RGB}{34,20,30}%
    \definecolor{cvSil}{RGB}{22,14,22}%
    \definecolor{cvCream}{RGB}{246,238,238}%
  }{%
  \ifstrequal{#1}{part4}{%
    \definecolor{cvSkyTop}{RGB}{12,28,28}%
    \definecolor{cvSkyMid}{RGB}{44,86,72}%
    \definecolor{cvSkyBot}{RGB}{186,168,86}%
    \definecolor{cvSun}{RGB}{240,232,168}%
    \definecolor{cvFieldTop}{RGB}{110,104,50}%
    \definecolor{cvFieldBot}{RGB}{30,34,20}%
    \definecolor{cvFurrow}{RGB}{26,32,20}%
    \definecolor{cvSil}{RGB}{14,20,14}%
    \definecolor{cvCream}{RGB}{240,238,226}%
  }{%
    % 默认: 全书(麦浪金褐)
    \definecolor{cvSkyTop}{RGB}{28,40,32}%
    \definecolor{cvSkyMid}{RGB}{128,110,64}%
    \definecolor{cvSkyBot}{RGB}{232,178,88}%
    \definecolor{cvSun}{RGB}{255,236,180}%
    \definecolor{cvFieldTop}{RGB}{168,124,52}%
    \definecolor{cvFieldBot}{RGB}{62,40,20}%
    \definecolor{cvFurrow}{RGB}{42,26,12}%
    \definecolor{cvSil}{RGB}{28,18,10}%
    \definecolor{cvCream}{RGB}{250,242,226}%
  }}}}}%
}

% --------------------------------------------------------------
%  一键设定:  \CoverSetup{配色}{卷次}{卷名}{副题}{原书页码范围}
% --------------------------------------------------------------
\newcommand{\CoverSetup}[5]{%
  \CoverPalette{#1}%
  \newcommand{\VolNum}{#2}%
  \newcommand{\VolName}{#3}%
  \newcommand{\VolSub}{#4}%
  \newcommand{\VolPages}{#5}%
}

% --------------------------------------------------------------
%  一支麦穗:  \cvWheatEar{x}{底 y}{高}{倾角}
%  茎略弯, 麦粒成对左右斜出, 顶端一芒三针
% --------------------------------------------------------------
\newcommand{\cvWheatEar}[4]{%
  \begin{scope}[shift={(#1,#2)}, rotate=#4]
    \draw[cvSil, line width=0.85mm, line cap=round]
      (0,0) .. controls (0.5,{0.35*#3}) and (0.8,{0.72*#3}) .. ({0.55*#3},#3);
    \foreach \cvK in {0,...,6}{%
      \pgfmathsetmacro{\cvT}{0.50+0.50*\cvK/6}
      \pgfmathsetmacro{\cvGY}{\cvT*#3}
      \pgfmathsetmacro{\cvGX}{0.55*\cvT*\cvT*#3}
      \pgfmathsetmacro{\cvGS}{1.0-0.34*\cvT}
      \begin{scope}[shift={(\cvGX,\cvGY)}, rotate=28]
        \fill[cvSil] (0,{1.9*\cvGS}) ellipse [x radius={0.92*\cvGS}, y radius={2.3*\cvGS}];
      \end{scope}
      \begin{scope}[shift={(\cvGX,\cvGY)}, rotate=-28]
        \fill[cvSil] (0,{-1.9*\cvGS}) ellipse [x radius={0.92*\cvGS}, y radius={2.3*\cvGS}];
      \end{scope}
    }
    \begin{scope}[shift={({0.55*#3},#3)}]
      \fill[cvSil] (0,1.5) ellipse [x radius=0.78, y radius=1.95];
    \end{scope}
    \draw[cvSil, line width=0.26mm, line cap=round]
      ({0.55*#3},{#3+2.3}) -- ({0.55*#3-1.7},{#3+9.5});
    \draw[cvSil, line width=0.26mm, line cap=round]
      ({0.55*#3},{#3+2.3}) -- ({0.55*#3+2.1},{#3+10.5});
    \draw[cvSil, line width=0.26mm, line cap=round]
      ({0.55*#3},{#3+2.3}) -- ({0.55*#3},{#3+11.5});
  \end{scope}
}

% --------------------------------------------------------------
%  封面
% --------------------------------------------------------------
\makeatletter
\renewcommand{\makecover}{%
  \begin{titlepage}
    \thispagestyle{empty}
    \begin{tikzpicture}[remember picture, overlay, x=1mm, y=1mm,
                        shift={(current page.south west)}]
      % ===== 天空 =====
      \shade[top color=cvSkyTop, middle color=cvSkyMid, bottom color=cvSkyBot]
        (0,105) rectangle (210,297);
      % ===== 田野 =====
      \shade[top color=cvFieldTop, bottom color=cvFieldBot]
        (0,0) rectangle (210,105);
      % ===== 落日与光晕 =====
      \begin{scope}
        \clip (0,105) rectangle (210,297);
        \fill[cvSun, opacity=0.14] (152,132) circle (46);
        \fill[cvSun, opacity=0.24] (152,132) circle (30);
      \end{scope}
      \fill[cvSun, opacity=0.95] (152,132) circle (16.5);
      % ===== 犁沟: 向地平线灭点收束 =====
      \begin{scope}
        \clip (0,0) rectangle (210,105);
        \foreach \cvI in {1,...,20}{%
          \pgfmathsetmacro{\cvDX}{3.1*pow(\cvI,1.55)}
          \draw[cvFurrow, opacity=0.30, line width=0.32mm] ({105+\cvDX},0) -- (105,105);
          \draw[cvFurrow, opacity=0.30, line width=0.32mm] ({105-\cvDX},0) -- (105,105);
        }
        \foreach \cvY in {8,18,31,47,68,92}{%
          \draw[cvFurrow, opacity=0.15, line width=0.3mm]
            (0,\cvY) .. controls (60,{\cvY+3.5}) and (150,{\cvY-3}) .. (210,{\cvY+1.5});
        }
      \end{scope}
      % 地平线
      \draw[cvSil, opacity=0.45, line width=0.45mm] (0,105) -- (210,105);
      % ===== 播种者剪影(立于地平线) =====
      \begin{scope}[cvSil]
        \draw[line width=1.9mm, line cap=round, line join=round]
          (55.5,116) -- (50,105)  (55.5,116) -- (60.5,105.5);
        \draw[line width=2.3mm, line cap=round, line join=round]
          (55.5,116) -- (55,123.5);
        \draw[line width=1.6mm, line cap=round, line join=round]
          (55.2,123.2) -- (61.5,126.6) -- (68,129.4);
        \draw[line width=1.5mm, line cap=round, line join=round]
          (55.2,123.2) -- (51.4,119.8) -- (48.6,116.2);
        \fill (55.6,127.3) circle (2.35);
      \end{scope}
      % ===== 种子: 沿贝塞尔弧撒出 =====
      \draw[cvCream, opacity=0.13, line width=0.5mm]
        (68,129.4) .. controls (100,168) and (160,162) .. (196,104);
      % 掌心那一把(尚未撒出)
      \foreach \cvI in {1,...,5}{%
        \pgfmathsetmacro{\cvJX}{1.1*sin(\cvI*53)}
        \pgfmathsetmacro{\cvJY}{1.1*cos(\cvI*83)}
        \fill[cvCream, opacity=0.9] ({68+\cvJX},{129.4+\cvJY}) circle (1.0);
      }
      \foreach \cvI in {1,...,26}{%
        \pgfmathsetmacro{\cvT}{0.035+0.92*\cvI/26}
        \pgfmathsetmacro{\cvU}{1-\cvT}
        \pgfmathsetmacro{\cvSX}{\cvU*\cvU*\cvU*68 + 3*\cvU*\cvU*\cvT*100
                                + 3*\cvU*\cvT*\cvT*160 + \cvT*\cvT*\cvT*196}
        \pgfmathsetmacro{\cvSY}{\cvU*\cvU*\cvU*129.4 + 3*\cvU*\cvU*\cvT*168
                                + 3*\cvU*\cvT*\cvT*162 + \cvT*\cvT*\cvT*104}
        \pgfmathsetmacro{\cvR}{1.55-0.85*\cvT}
        \pgfmathsetmacro{\cvJX}{1.7*sin(\cvI*97)}
        \pgfmathsetmacro{\cvJY}{1.7*cos(\cvI*71)}
        \fill[cvCream, opacity={0.88-0.28*\cvT}] ({\cvSX+\cvJX},{\cvSY+\cvJY}) circle (\cvR);
      }
      % ===== 前景麦穗(两排) =====
      \foreach \cvI in {1,...,11}{%
        \pgfmathsetmacro{\cvWX}{4+19*(\cvI-1)+1.6*sin(\cvI*2.1)}
        \pgfmathsetmacro{\cvWH}{23+7*sin(\cvI*1.7)}
        \pgfmathsetmacro{\cvWL}{3.2*sin(\cvI*2.9)}
        \begin{scope}[opacity=0.45]
          \cvWheatEar{\cvWX}{30}{\cvWH}{\cvWL}
        \end{scope}
      }
      \foreach \cvI in {1,...,10}{%
        \pgfmathsetmacro{\cvWX}{9+20.4*(\cvI-1)+1.5*sin(\cvI*1.3)}
        \pgfmathsetmacro{\cvWH}{41+10*sin(\cvI*2.3)}
        \pgfmathsetmacro{\cvWL}{4.5*sin(\cvI*3.7)}
        \cvWheatEar{\cvWX}{30}{\cvWH}{\cvWL}
      }
      % ===== 压暗一层, 统一色调并保证文字可读 =====
      \fill[black, opacity=0.10] (0,105) rectangle (210,297);
      % ===== 底部信息带 =====
      \fill[cvSil, opacity=0.72] (0,0) rectangle (210,30);
      \draw[cvCream, opacity=0.40, line width=0.4mm] (0,30) -- (210,30);
      % ===== 文字 =====
      \node[anchor=north west, inner sep=0, text=cvCream] at (22,268)
        {\fontsize{38pt}{45pt}\selectfont\headfont \ziju{0.08}\BookTitleCN};
      \node[anchor=north west, inner sep=0, text=cvCream] at (22.6,242)
        {\fontsize{15pt}{20pt}\selectfont\itshape \BookTitleEN};
      \draw[cvCream, opacity=0.55, line width=0.5mm] (22,233) -- (120,233);
      \node[anchor=north west, inner sep=0, text=cvCream] at (22,224)
        {\fontsize{11pt}{15pt}\selectfont 一位数学家对往昔的反思与见证};
      \node[anchor=north west, inner sep=0, text=cvCream, opacity=0.78] at (22.4,208)
        {\fontsize{10pt}{14pt}\selectfont \ziju{0.14}\VolNum};
      \node[anchor=north west, inner sep=0, text=cvCream] at (22,199)
        {\fontsize{30pt}{37pt}\selectfont\headfont \ziju{0.1}\VolName};
      \node[anchor=north west, inner sep=0, text=cvCream, opacity=0.92] at (22,180)
        {\fontsize{13pt}{18pt}\selectfont \VolSub};
      \node[anchor=north west, inner sep=0, text=cvCream] at (22,19)
        {\fontsize{11.5pt}{15pt}\selectfont \OriginalAuthor\quad 著};
      \node[anchor=north west, inner sep=0, text=cvCream, opacity=0.85] at (22,7)
        {\fontsize{9.5pt}{13pt}\selectfont \Translator\ 译\quad ·\quad
         \ifdefempty{\TranslationDate}{\today}{\TranslationDate}};
      \node[anchor=north east, inner sep=0, text=cvCream, opacity=0.85] at (188,13)
        {\fontsize{10pt}{14pt}\selectfont 原书 \VolPages};
    \end{tikzpicture}
  \end{titlepage}%
}
\makeatother
                    % 本文件
%      \CoverSetup{prelude}{卷一}{前奏}%
%                 {主题的呈现——或四个乐章的前奏}{p.23–122}
%  正文中再 \makecover 即可。配色关键字:
%      prelude / part1 / part2 / part3 / part4 / all
%  五卷共用同一幅画, 只换配色与卷次文字, 以保证"一套书"的观感。
% ================================================================

\RequirePackage{tikz}

% --------------------------------------------------------------
%  配色: 六套(前奏 / 四部 / 全书), 各为同一暮色母题的变调
% --------------------------------------------------------------
\newcommand{\CoverPalette}[1]{%
  \ifstrequal{#1}{prelude}{%
    \definecolor{cvSkyTop}{RGB}{10,26,48}%
    \definecolor{cvSkyMid}{RGB}{46,84,116}%
    \definecolor{cvSkyBot}{RGB}{232,164,92}%
    \definecolor{cvSun}{RGB}{255,224,168}%
    \definecolor{cvFieldTop}{RGB}{150,110,52}%
    \definecolor{cvFieldBot}{RGB}{58,36,18}%
    \definecolor{cvFurrow}{RGB}{40,22,10}%
    \definecolor{cvSil}{RGB}{30,18,10}%
    \definecolor{cvCream}{RGB}{247,240,228}%
  }{%
  \ifstrequal{#1}{part1}{%
    \definecolor{cvSkyTop}{RGB}{22,70,100}%
    \definecolor{cvSkyMid}{RGB}{92,150,168}%
    \definecolor{cvSkyBot}{RGB}{240,200,104}%
    \definecolor{cvSun}{RGB}{255,244,200}%
    \definecolor{cvFieldTop}{RGB}{176,140,58}%
    \definecolor{cvFieldBot}{RGB}{84,54,22}%
    \definecolor{cvFurrow}{RGB}{45,28,12}%
    \definecolor{cvSil}{RGB}{34,22,10}%
    \definecolor{cvCream}{RGB}{250,245,232}%
  }{%
  \ifstrequal{#1}{part2}{%
    \definecolor{cvSkyTop}{RGB}{46,20,30}%
    \definecolor{cvSkyMid}{RGB}{128,52,54}%
    \definecolor{cvSkyBot}{RGB}{226,120,64}%
    \definecolor{cvSun}{RGB}{255,206,150}%
    \definecolor{cvFieldTop}{RGB}{132,84,44}%
    \definecolor{cvFieldBot}{RGB}{48,26,18}%
    \definecolor{cvFurrow}{RGB}{36,16,14}%
    \definecolor{cvSil}{RGB}{26,12,12}%
    \definecolor{cvCream}{RGB}{248,236,224}%
  }{%
  \ifstrequal{#1}{part3}{%
    \definecolor{cvSkyTop}{RGB}{30,20,54}%
    \definecolor{cvSkyMid}{RGB}{92,62,124}%
    \definecolor{cvSkyBot}{RGB}{216,132,140}%
    \definecolor{cvSun}{RGB}{255,214,196}%
    \definecolor{cvFieldTop}{RGB}{128,90,74}%
    \definecolor{cvFieldBot}{RGB}{44,26,30}%
    \definecolor{cvFurrow}{RGB}{34,20,30}%
    \definecolor{cvSil}{RGB}{22,14,22}%
    \definecolor{cvCream}{RGB}{246,238,238}%
  }{%
  \ifstrequal{#1}{part4}{%
    \definecolor{cvSkyTop}{RGB}{12,28,28}%
    \definecolor{cvSkyMid}{RGB}{44,86,72}%
    \definecolor{cvSkyBot}{RGB}{186,168,86}%
    \definecolor{cvSun}{RGB}{240,232,168}%
    \definecolor{cvFieldTop}{RGB}{110,104,50}%
    \definecolor{cvFieldBot}{RGB}{30,34,20}%
    \definecolor{cvFurrow}{RGB}{26,32,20}%
    \definecolor{cvSil}{RGB}{14,20,14}%
    \definecolor{cvCream}{RGB}{240,238,226}%
  }{%
    % 默认: 全书(麦浪金褐)
    \definecolor{cvSkyTop}{RGB}{28,40,32}%
    \definecolor{cvSkyMid}{RGB}{128,110,64}%
    \definecolor{cvSkyBot}{RGB}{232,178,88}%
    \definecolor{cvSun}{RGB}{255,236,180}%
    \definecolor{cvFieldTop}{RGB}{168,124,52}%
    \definecolor{cvFieldBot}{RGB}{62,40,20}%
    \definecolor{cvFurrow}{RGB}{42,26,12}%
    \definecolor{cvSil}{RGB}{28,18,10}%
    \definecolor{cvCream}{RGB}{250,242,226}%
  }}}}}%
}

% --------------------------------------------------------------
%  一键设定:  \CoverSetup{配色}{卷次}{卷名}{副题}{原书页码范围}
% --------------------------------------------------------------
\newcommand{\CoverSetup}[5]{%
  \CoverPalette{#1}%
  \newcommand{\VolNum}{#2}%
  \newcommand{\VolName}{#3}%
  \newcommand{\VolSub}{#4}%
  \newcommand{\VolPages}{#5}%
}

% --------------------------------------------------------------
%  一支麦穗:  \cvWheatEar{x}{底 y}{高}{倾角}
%  茎略弯, 麦粒成对左右斜出, 顶端一芒三针
% --------------------------------------------------------------
\newcommand{\cvWheatEar}[4]{%
  \begin{scope}[shift={(#1,#2)}, rotate=#4]
    \draw[cvSil, line width=0.85mm, line cap=round]
      (0,0) .. controls (0.5,{0.35*#3}) and (0.8,{0.72*#3}) .. ({0.55*#3},#3);
    \foreach \cvK in {0,...,6}{%
      \pgfmathsetmacro{\cvT}{0.50+0.50*\cvK/6}
      \pgfmathsetmacro{\cvGY}{\cvT*#3}
      \pgfmathsetmacro{\cvGX}{0.55*\cvT*\cvT*#3}
      \pgfmathsetmacro{\cvGS}{1.0-0.34*\cvT}
      \begin{scope}[shift={(\cvGX,\cvGY)}, rotate=28]
        \fill[cvSil] (0,{1.9*\cvGS}) ellipse [x radius={0.92*\cvGS}, y radius={2.3*\cvGS}];
      \end{scope}
      \begin{scope}[shift={(\cvGX,\cvGY)}, rotate=-28]
        \fill[cvSil] (0,{-1.9*\cvGS}) ellipse [x radius={0.92*\cvGS}, y radius={2.3*\cvGS}];
      \end{scope}
    }
    \begin{scope}[shift={({0.55*#3},#3)}]
      \fill[cvSil] (0,1.5) ellipse [x radius=0.78, y radius=1.95];
    \end{scope}
    \draw[cvSil, line width=0.26mm, line cap=round]
      ({0.55*#3},{#3+2.3}) -- ({0.55*#3-1.7},{#3+9.5});
    \draw[cvSil, line width=0.26mm, line cap=round]
      ({0.55*#3},{#3+2.3}) -- ({0.55*#3+2.1},{#3+10.5});
    \draw[cvSil, line width=0.26mm, line cap=round]
      ({0.55*#3},{#3+2.3}) -- ({0.55*#3},{#3+11.5});
  \end{scope}
}

% --------------------------------------------------------------
%  封面
% --------------------------------------------------------------
\makeatletter
\renewcommand{\makecover}{%
  \begin{titlepage}
    \thispagestyle{empty}
    \begin{tikzpicture}[remember picture, overlay, x=1mm, y=1mm,
                        shift={(current page.south west)}]
      % ===== 天空 =====
      \shade[top color=cvSkyTop, middle color=cvSkyMid, bottom color=cvSkyBot]
        (0,105) rectangle (210,297);
      % ===== 田野 =====
      \shade[top color=cvFieldTop, bottom color=cvFieldBot]
        (0,0) rectangle (210,105);
      % ===== 落日与光晕 =====
      \begin{scope}
        \clip (0,105) rectangle (210,297);
        \fill[cvSun, opacity=0.14] (152,132) circle (46);
        \fill[cvSun, opacity=0.24] (152,132) circle (30);
      \end{scope}
      \fill[cvSun, opacity=0.95] (152,132) circle (16.5);
      % ===== 犁沟: 向地平线灭点收束 =====
      \begin{scope}
        \clip (0,0) rectangle (210,105);
        \foreach \cvI in {1,...,20}{%
          \pgfmathsetmacro{\cvDX}{3.1*pow(\cvI,1.55)}
          \draw[cvFurrow, opacity=0.30, line width=0.32mm] ({105+\cvDX},0) -- (105,105);
          \draw[cvFurrow, opacity=0.30, line width=0.32mm] ({105-\cvDX},0) -- (105,105);
        }
        \foreach \cvY in {8,18,31,47,68,92}{%
          \draw[cvFurrow, opacity=0.15, line width=0.3mm]
            (0,\cvY) .. controls (60,{\cvY+3.5}) and (150,{\cvY-3}) .. (210,{\cvY+1.5});
        }
      \end{scope}
      % 地平线
      \draw[cvSil, opacity=0.45, line width=0.45mm] (0,105) -- (210,105);
      % ===== 播种者剪影(立于地平线) =====
      \begin{scope}[cvSil]
        \draw[line width=1.9mm, line cap=round, line join=round]
          (55.5,116) -- (50,105)  (55.5,116) -- (60.5,105.5);
        \draw[line width=2.3mm, line cap=round, line join=round]
          (55.5,116) -- (55,123.5);
        \draw[line width=1.6mm, line cap=round, line join=round]
          (55.2,123.2) -- (61.5,126.6) -- (68,129.4);
        \draw[line width=1.5mm, line cap=round, line join=round]
          (55.2,123.2) -- (51.4,119.8) -- (48.6,116.2);
        \fill (55.6,127.3) circle (2.35);
      \end{scope}
      % ===== 种子: 沿贝塞尔弧撒出 =====
      \draw[cvCream, opacity=0.13, line width=0.5mm]
        (68,129.4) .. controls (100,168) and (160,162) .. (196,104);
      % 掌心那一把(尚未撒出)
      \foreach \cvI in {1,...,5}{%
        \pgfmathsetmacro{\cvJX}{1.1*sin(\cvI*53)}
        \pgfmathsetmacro{\cvJY}{1.1*cos(\cvI*83)}
        \fill[cvCream, opacity=0.9] ({68+\cvJX},{129.4+\cvJY}) circle (1.0);
      }
      \foreach \cvI in {1,...,26}{%
        \pgfmathsetmacro{\cvT}{0.035+0.92*\cvI/26}
        \pgfmathsetmacro{\cvU}{1-\cvT}
        \pgfmathsetmacro{\cvSX}{\cvU*\cvU*\cvU*68 + 3*\cvU*\cvU*\cvT*100
                                + 3*\cvU*\cvT*\cvT*160 + \cvT*\cvT*\cvT*196}
        \pgfmathsetmacro{\cvSY}{\cvU*\cvU*\cvU*129.4 + 3*\cvU*\cvU*\cvT*168
                                + 3*\cvU*\cvT*\cvT*162 + \cvT*\cvT*\cvT*104}
        \pgfmathsetmacro{\cvR}{1.55-0.85*\cvT}
        \pgfmathsetmacro{\cvJX}{1.7*sin(\cvI*97)}
        \pgfmathsetmacro{\cvJY}{1.7*cos(\cvI*71)}
        \fill[cvCream, opacity={0.88-0.28*\cvT}] ({\cvSX+\cvJX},{\cvSY+\cvJY}) circle (\cvR);
      }
      % ===== 前景麦穗(两排) =====
      \foreach \cvI in {1,...,11}{%
        \pgfmathsetmacro{\cvWX}{4+19*(\cvI-1)+1.6*sin(\cvI*2.1)}
        \pgfmathsetmacro{\cvWH}{23+7*sin(\cvI*1.7)}
        \pgfmathsetmacro{\cvWL}{3.2*sin(\cvI*2.9)}
        \begin{scope}[opacity=0.45]
          \cvWheatEar{\cvWX}{30}{\cvWH}{\cvWL}
        \end{scope}
      }
      \foreach \cvI in {1,...,10}{%
        \pgfmathsetmacro{\cvWX}{9+20.4*(\cvI-1)+1.5*sin(\cvI*1.3)}
        \pgfmathsetmacro{\cvWH}{41+10*sin(\cvI*2.3)}
        \pgfmathsetmacro{\cvWL}{4.5*sin(\cvI*3.7)}
        \cvWheatEar{\cvWX}{30}{\cvWH}{\cvWL}
      }
      % ===== 压暗一层, 统一色调并保证文字可读 =====
      \fill[black, opacity=0.10] (0,105) rectangle (210,297);
      % ===== 底部信息带 =====
      \fill[cvSil, opacity=0.72] (0,0) rectangle (210,30);
      \draw[cvCream, opacity=0.40, line width=0.4mm] (0,30) -- (210,30);
      % ===== 文字 =====
      \node[anchor=north west, inner sep=0, text=cvCream] at (22,268)
        {\fontsize{38pt}{45pt}\selectfont\headfont \ziju{0.08}\BookTitleCN};
      \node[anchor=north west, inner sep=0, text=cvCream] at (22.6,242)
        {\fontsize{15pt}{20pt}\selectfont\itshape \BookTitleEN};
      \draw[cvCream, opacity=0.55, line width=0.5mm] (22,233) -- (120,233);
      \node[anchor=north west, inner sep=0, text=cvCream] at (22,224)
        {\fontsize{11pt}{15pt}\selectfont 一位数学家对往昔的反思与见证};
      \node[anchor=north west, inner sep=0, text=cvCream, opacity=0.78] at (22.4,208)
        {\fontsize{10pt}{14pt}\selectfont \ziju{0.14}\VolNum};
      \node[anchor=north west, inner sep=0, text=cvCream] at (22,199)
        {\fontsize{30pt}{37pt}\selectfont\headfont \ziju{0.1}\VolName};
      \node[anchor=north west, inner sep=0, text=cvCream, opacity=0.92] at (22,180)
        {\fontsize{13pt}{18pt}\selectfont \VolSub};
      \node[anchor=north west, inner sep=0, text=cvCream] at (22,19)
        {\fontsize{11.5pt}{15pt}\selectfont \OriginalAuthor\quad 著};
      \node[anchor=north west, inner sep=0, text=cvCream, opacity=0.85] at (22,7)
        {\fontsize{9.5pt}{13pt}\selectfont \Translator\ 译\quad ·\quad
         \ifdefempty{\TranslationDate}{\today}{\TranslationDate}};
      \node[anchor=north east, inner sep=0, text=cvCream, opacity=0.85] at (188,13)
        {\fontsize{10pt}{14pt}\selectfont 原书 \VolPages};
    \end{tikzpicture}
  \end{titlepage}%
}
\makeatother
                    % 本文件
%      \CoverSetup{prelude}{卷一}{前奏}%
%                 {主题的呈现——或四个乐章的前奏}{p.23–122}
%  正文中再 \makecover 即可。配色关键字:
%      prelude / part1 / part2 / part3 / part4 / all
%  五卷共用同一幅画, 只换配色与卷次文字, 以保证"一套书"的观感。
% ================================================================

\RequirePackage{tikz}

% --------------------------------------------------------------
%  配色: 六套(前奏 / 四部 / 全书), 各为同一暮色母题的变调
% --------------------------------------------------------------
\newcommand{\CoverPalette}[1]{%
  \ifstrequal{#1}{prelude}{%
    \definecolor{cvSkyTop}{RGB}{10,26,48}%
    \definecolor{cvSkyMid}{RGB}{46,84,116}%
    \definecolor{cvSkyBot}{RGB}{232,164,92}%
    \definecolor{cvSun}{RGB}{255,224,168}%
    \definecolor{cvFieldTop}{RGB}{150,110,52}%
    \definecolor{cvFieldBot}{RGB}{58,36,18}%
    \definecolor{cvFurrow}{RGB}{40,22,10}%
    \definecolor{cvSil}{RGB}{30,18,10}%
    \definecolor{cvCream}{RGB}{247,240,228}%
  }{%
  \ifstrequal{#1}{part1}{%
    \definecolor{cvSkyTop}{RGB}{22,70,100}%
    \definecolor{cvSkyMid}{RGB}{92,150,168}%
    \definecolor{cvSkyBot}{RGB}{240,200,104}%
    \definecolor{cvSun}{RGB}{255,244,200}%
    \definecolor{cvFieldTop}{RGB}{176,140,58}%
    \definecolor{cvFieldBot}{RGB}{84,54,22}%
    \definecolor{cvFurrow}{RGB}{45,28,12}%
    \definecolor{cvSil}{RGB}{34,22,10}%
    \definecolor{cvCream}{RGB}{250,245,232}%
  }{%
  \ifstrequal{#1}{part2}{%
    \definecolor{cvSkyTop}{RGB}{46,20,30}%
    \definecolor{cvSkyMid}{RGB}{128,52,54}%
    \definecolor{cvSkyBot}{RGB}{226,120,64}%
    \definecolor{cvSun}{RGB}{255,206,150}%
    \definecolor{cvFieldTop}{RGB}{132,84,44}%
    \definecolor{cvFieldBot}{RGB}{48,26,18}%
    \definecolor{cvFurrow}{RGB}{36,16,14}%
    \definecolor{cvSil}{RGB}{26,12,12}%
    \definecolor{cvCream}{RGB}{248,236,224}%
  }{%
  \ifstrequal{#1}{part3}{%
    \definecolor{cvSkyTop}{RGB}{30,20,54}%
    \definecolor{cvSkyMid}{RGB}{92,62,124}%
    \definecolor{cvSkyBot}{RGB}{216,132,140}%
    \definecolor{cvSun}{RGB}{255,214,196}%
    \definecolor{cvFieldTop}{RGB}{128,90,74}%
    \definecolor{cvFieldBot}{RGB}{44,26,30}%
    \definecolor{cvFurrow}{RGB}{34,20,30}%
    \definecolor{cvSil}{RGB}{22,14,22}%
    \definecolor{cvCream}{RGB}{246,238,238}%
  }{%
  \ifstrequal{#1}{part4}{%
    \definecolor{cvSkyTop}{RGB}{12,28,28}%
    \definecolor{cvSkyMid}{RGB}{44,86,72}%
    \definecolor{cvSkyBot}{RGB}{186,168,86}%
    \definecolor{cvSun}{RGB}{240,232,168}%
    \definecolor{cvFieldTop}{RGB}{110,104,50}%
    \definecolor{cvFieldBot}{RGB}{30,34,20}%
    \definecolor{cvFurrow}{RGB}{26,32,20}%
    \definecolor{cvSil}{RGB}{14,20,14}%
    \definecolor{cvCream}{RGB}{240,238,226}%
  }{%
    % 默认: 全书(麦浪金褐)
    \definecolor{cvSkyTop}{RGB}{28,40,32}%
    \definecolor{cvSkyMid}{RGB}{128,110,64}%
    \definecolor{cvSkyBot}{RGB}{232,178,88}%
    \definecolor{cvSun}{RGB}{255,236,180}%
    \definecolor{cvFieldTop}{RGB}{168,124,52}%
    \definecolor{cvFieldBot}{RGB}{62,40,20}%
    \definecolor{cvFurrow}{RGB}{42,26,12}%
    \definecolor{cvSil}{RGB}{28,18,10}%
    \definecolor{cvCream}{RGB}{250,242,226}%
  }}}}}%
}

% --------------------------------------------------------------
%  一键设定:  \CoverSetup{配色}{卷次}{卷名}{副题}{原书页码范围}
% --------------------------------------------------------------
\newcommand{\CoverSetup}[5]{%
  \CoverPalette{#1}%
  \newcommand{\VolNum}{#2}%
  \newcommand{\VolName}{#3}%
  \newcommand{\VolSub}{#4}%
  \newcommand{\VolPages}{#5}%
}

% --------------------------------------------------------------
%  一支麦穗:  \cvWheatEar{x}{底 y}{高}{倾角}
%  茎略弯, 麦粒成对左右斜出, 顶端一芒三针
% --------------------------------------------------------------
\newcommand{\cvWheatEar}[4]{%
  \begin{scope}[shift={(#1,#2)}, rotate=#4]
    \draw[cvSil, line width=0.85mm, line cap=round]
      (0,0) .. controls (0.5,{0.35*#3}) and (0.8,{0.72*#3}) .. ({0.55*#3},#3);
    \foreach \cvK in {0,...,6}{%
      \pgfmathsetmacro{\cvT}{0.50+0.50*\cvK/6}
      \pgfmathsetmacro{\cvGY}{\cvT*#3}
      \pgfmathsetmacro{\cvGX}{0.55*\cvT*\cvT*#3}
      \pgfmathsetmacro{\cvGS}{1.0-0.34*\cvT}
      \begin{scope}[shift={(\cvGX,\cvGY)}, rotate=28]
        \fill[cvSil] (0,{1.9*\cvGS}) ellipse [x radius={0.92*\cvGS}, y radius={2.3*\cvGS}];
      \end{scope}
      \begin{scope}[shift={(\cvGX,\cvGY)}, rotate=-28]
        \fill[cvSil] (0,{-1.9*\cvGS}) ellipse [x radius={0.92*\cvGS}, y radius={2.3*\cvGS}];
      \end{scope}
    }
    \begin{scope}[shift={({0.55*#3},#3)}]
      \fill[cvSil] (0,1.5) ellipse [x radius=0.78, y radius=1.95];
    \end{scope}
    \draw[cvSil, line width=0.26mm, line cap=round]
      ({0.55*#3},{#3+2.3}) -- ({0.55*#3-1.7},{#3+9.5});
    \draw[cvSil, line width=0.26mm, line cap=round]
      ({0.55*#3},{#3+2.3}) -- ({0.55*#3+2.1},{#3+10.5});
    \draw[cvSil, line width=0.26mm, line cap=round]
      ({0.55*#3},{#3+2.3}) -- ({0.55*#3},{#3+11.5});
  \end{scope}
}

% --------------------------------------------------------------
%  封面
% --------------------------------------------------------------
\makeatletter
\renewcommand{\makecover}{%
  \begin{titlepage}
    \thispagestyle{empty}
    \begin{tikzpicture}[remember picture, overlay, x=1mm, y=1mm,
                        shift={(current page.south west)}]
      % ===== 天空 =====
      \shade[top color=cvSkyTop, middle color=cvSkyMid, bottom color=cvSkyBot]
        (0,105) rectangle (210,297);
      % ===== 田野 =====
      \shade[top color=cvFieldTop, bottom color=cvFieldBot]
        (0,0) rectangle (210,105);
      % ===== 落日与光晕 =====
      \begin{scope}
        \clip (0,105) rectangle (210,297);
        \fill[cvSun, opacity=0.14] (152,132) circle (46);
        \fill[cvSun, opacity=0.24] (152,132) circle (30);
      \end{scope}
      \fill[cvSun, opacity=0.95] (152,132) circle (16.5);
      % ===== 犁沟: 向地平线灭点收束 =====
      \begin{scope}
        \clip (0,0) rectangle (210,105);
        \foreach \cvI in {1,...,20}{%
          \pgfmathsetmacro{\cvDX}{3.1*pow(\cvI,1.55)}
          \draw[cvFurrow, opacity=0.30, line width=0.32mm] ({105+\cvDX},0) -- (105,105);
          \draw[cvFurrow, opacity=0.30, line width=0.32mm] ({105-\cvDX},0) -- (105,105);
        }
        \foreach \cvY in {8,18,31,47,68,92}{%
          \draw[cvFurrow, opacity=0.15, line width=0.3mm]
            (0,\cvY) .. controls (60,{\cvY+3.5}) and (150,{\cvY-3}) .. (210,{\cvY+1.5});
        }
      \end{scope}
      % 地平线
      \draw[cvSil, opacity=0.45, line width=0.45mm] (0,105) -- (210,105);
      % ===== 播种者剪影(立于地平线) =====
      \begin{scope}[cvSil]
        \draw[line width=1.9mm, line cap=round, line join=round]
          (55.5,116) -- (50,105)  (55.5,116) -- (60.5,105.5);
        \draw[line width=2.3mm, line cap=round, line join=round]
          (55.5,116) -- (55,123.5);
        \draw[line width=1.6mm, line cap=round, line join=round]
          (55.2,123.2) -- (61.5,126.6) -- (68,129.4);
        \draw[line width=1.5mm, line cap=round, line join=round]
          (55.2,123.2) -- (51.4,119.8) -- (48.6,116.2);
        \fill (55.6,127.3) circle (2.35);
      \end{scope}
      % ===== 种子: 沿贝塞尔弧撒出 =====
      \draw[cvCream, opacity=0.13, line width=0.5mm]
        (68,129.4) .. controls (100,168) and (160,162) .. (196,104);
      % 掌心那一把(尚未撒出)
      \foreach \cvI in {1,...,5}{%
        \pgfmathsetmacro{\cvJX}{1.1*sin(\cvI*53)}
        \pgfmathsetmacro{\cvJY}{1.1*cos(\cvI*83)}
        \fill[cvCream, opacity=0.9] ({68+\cvJX},{129.4+\cvJY}) circle (1.0);
      }
      \foreach \cvI in {1,...,26}{%
        \pgfmathsetmacro{\cvT}{0.035+0.92*\cvI/26}
        \pgfmathsetmacro{\cvU}{1-\cvT}
        \pgfmathsetmacro{\cvSX}{\cvU*\cvU*\cvU*68 + 3*\cvU*\cvU*\cvT*100
                                + 3*\cvU*\cvT*\cvT*160 + \cvT*\cvT*\cvT*196}
        \pgfmathsetmacro{\cvSY}{\cvU*\cvU*\cvU*129.4 + 3*\cvU*\cvU*\cvT*168
                                + 3*\cvU*\cvT*\cvT*162 + \cvT*\cvT*\cvT*104}
        \pgfmathsetmacro{\cvR}{1.55-0.85*\cvT}
        \pgfmathsetmacro{\cvJX}{1.7*sin(\cvI*97)}
        \pgfmathsetmacro{\cvJY}{1.7*cos(\cvI*71)}
        \fill[cvCream, opacity={0.88-0.28*\cvT}] ({\cvSX+\cvJX},{\cvSY+\cvJY}) circle (\cvR);
      }
      % ===== 前景麦穗(两排) =====
      \foreach \cvI in {1,...,11}{%
        \pgfmathsetmacro{\cvWX}{4+19*(\cvI-1)+1.6*sin(\cvI*2.1)}
        \pgfmathsetmacro{\cvWH}{23+7*sin(\cvI*1.7)}
        \pgfmathsetmacro{\cvWL}{3.2*sin(\cvI*2.9)}
        \begin{scope}[opacity=0.45]
          \cvWheatEar{\cvWX}{30}{\cvWH}{\cvWL}
        \end{scope}
      }
      \foreach \cvI in {1,...,10}{%
        \pgfmathsetmacro{\cvWX}{9+20.4*(\cvI-1)+1.5*sin(\cvI*1.3)}
        \pgfmathsetmacro{\cvWH}{41+10*sin(\cvI*2.3)}
        \pgfmathsetmacro{\cvWL}{4.5*sin(\cvI*3.7)}
        \cvWheatEar{\cvWX}{30}{\cvWH}{\cvWL}
      }
      % ===== 压暗一层, 统一色调并保证文字可读 =====
      \fill[black, opacity=0.10] (0,105) rectangle (210,297);
      % ===== 底部信息带 =====
      \fill[cvSil, opacity=0.72] (0,0) rectangle (210,30);
      \draw[cvCream, opacity=0.40, line width=0.4mm] (0,30) -- (210,30);
      % ===== 文字 =====
      \node[anchor=north west, inner sep=0, text=cvCream] at (22,268)
        {\fontsize{38pt}{45pt}\selectfont\headfont \ziju{0.08}\BookTitleCN};
      \node[anchor=north west, inner sep=0, text=cvCream] at (22.6,242)
        {\fontsize{15pt}{20pt}\selectfont\itshape \BookTitleEN};
      \draw[cvCream, opacity=0.55, line width=0.5mm] (22,233) -- (120,233);
      \node[anchor=north west, inner sep=0, text=cvCream] at (22,224)
        {\fontsize{11pt}{15pt}\selectfont 一位数学家对往昔的反思与见证};
      \node[anchor=north west, inner sep=0, text=cvCream, opacity=0.78] at (22.4,208)
        {\fontsize{10pt}{14pt}\selectfont \ziju{0.14}\VolNum};
      \node[anchor=north west, inner sep=0, text=cvCream] at (22,199)
        {\fontsize{30pt}{37pt}\selectfont\headfont \ziju{0.1}\VolName};
      \node[anchor=north west, inner sep=0, text=cvCream, opacity=0.92] at (22,180)
        {\fontsize{13pt}{18pt}\selectfont \VolSub};
      \node[anchor=north west, inner sep=0, text=cvCream] at (22,19)
        {\fontsize{11.5pt}{15pt}\selectfont \OriginalAuthor\quad 著};
      \node[anchor=north west, inner sep=0, text=cvCream, opacity=0.85] at (22,7)
        {\fontsize{9.5pt}{13pt}\selectfont \Translator\ 译\quad ·\quad
         \ifdefempty{\TranslationDate}{\today}{\TranslationDate}};
      \node[anchor=north east, inner sep=0, text=cvCream, opacity=0.85] at (188,13)
        {\fontsize{10pt}{14pt}\selectfont 原书 \VolPages};
    \end{tikzpicture}
  \end{titlepage}%
}
\makeatother

\CoverSetup{part1}{卷二}{第一部}{愚妄与更新}{p.123–254}

\hypersetup{
  pdftitle={收获与播种·卷二·第一部(愚妄与更新)},
  pdfauthor={Alexandre Grothendieck / 数译 译}
}

% --------------------------------------------------------------
\begin{document}

\frontmatter
\makecover

% ---------------------------------------------------------------
%  译者说明
% ---------------------------------------------------------------
\chapter*{译者说明}
\markboth{译者说明}{译者说明}
\addcontentsline{toc}{chapter}{译者说明}

本册是《收获与播种》中文译本的第二卷, 收原书的第一部“愚妄与更新”(Première
partie: Fatuité et Renouvellement), 即原书第~123 至第~254 页, 含第~5 至第~12~章:

\begin{itemize}
  \item 第~5~章“劳动与发现”;  第~6~章“梦与做梦者”;
  \item 第~7~章“双重面孔”;    第~8~章“师与生”;
  \item 第~9~章“收获”;        第~10~章“孩子嬉戏”;
  \item 第~11~章“孤独的历险”; 第~12~章“第一部分的注”。
\end{itemize}

本部的脚注照原书作双重标记: 正文中的上标是该页的页内序号, 其后括注的是本部
的连续注号(如“(3)”), 注文行首则印作“1(3)”式; 连续注号与第~12 章汇编注的
“注~3”等一一对应。其中三条注原书改标日期而不标连续注号(第~7~章注~3 =
30 Septembre, 第~9~章注~11、12 = 8 août), 本册照录。

\CommonNoteBody

\makecontents

\mainmatter

% ================================================================
%  第一部 (第 5-12 章, 原书 p.123-254)
% ================================================================
\partdivider{第一部}{愚妄与更新}{愚妄与更新}

% 原书第125页献辞
\dedication{%
献给曾是我的长辈的人,\\[0.5em]
他们手足般地接纳我\\[0.5em]
进入那个曾属于他们、\\[0.5em]
而后也成为我的世界。\\[2em]
献给曾是我的学生的人,\\[0.5em]
我给了他们我身上最好的\\[0.5em]
——也给了最坏的……%
}

\setcounter{chapter}{4}
\markboth{收获与播种·愚妄与更新}{}

% ================================================================
%  第 5 章 劳动与发现 (原书 p.127-132, 边码 p.1-p.8)
%  源文: _work/src1/task-ch05.md
% ================================================================
% ---------------- p.127 ----------------
\chapter{劳动与发现}

\noindent 1983年6月\par

\section{(1) 孩子与好上帝}

\origpage{1} 我眼下正在做的这些数学笔记,是十三年来头一批我打算拿去发表的文字。读者不会奇怪:沉默了这么久之后,我的表达风格变了。然而这种表达上的变化,并不是工作风格或工作方法有什么变化的标志\fnoteP{1}{(1)}{(1984年3月补写)说我的工作“风格”与“方法”没有变,恐怕失于武断——毕竟我在数学上的表达风格已经深刻改变了。这一年来用于《追寻栈》的大部分时间,是坐在打字机前敲下一些思考;它们就是要几乎原样发表的(只除了后来补加的若干相当简短的注,借互见、勘误之类方便阅读)。没有剪刀,没有浆糊,不必费劲去炮制一份“定稿”(那份定稿尤其不该让任何通向它的过程透出影来)——这毕竟也是“风格”与“方法”的变化!除非把数学工作本身与写作、与陈述结果的工作割裂开来,而那是人为的,因为那不合事情的实情:数学工作与写作是拆不开的。} (1),更不是说我的数学工作本身的性质发生了转变。不但这性质一如既往——而且我确信:发现的工作,其性质在一个人与另一个人身上是相同的,它超越种种制约与气质所造成的差异,那些差异千变万化,无穷无尽。

发现是孩子的特权。我要说的是小孩子——那个还不怕弄错的孩子,不怕显得像个傻瓜,不怕不正经,不怕不跟大家一样。他也不怕他瞧着的东西不识趣:与他对它们的期待不一样、与它们本该是的模样不一样,或者更确切地说:与人们理所当然认定它们是的模样不一样。他不懂那些无声而毫无缺口的共识——它们是我们呼吸的空气的一部分,是所有明白人、所有公认的明白人的空气。天知道,自时间的黑夜深处以来,出过多少这样的人——明白人,公认的明白人!

我们的头脑里塞满了一种杂七杂八的“知识”——恐惧与懒惰、渴念与禁忌纠缠成一团;还有随处涌来的信息和按一按就出来的解释——一个封闭的空间,信息、渴念与恐惧在里面堆积,而旷野的风从来吹不进去。除了那点例行公事的“会做”,这“知识”的主要作用,似乎就是排掉活的知觉,排掉对世间事物的认识。它的效果,主要是一种巨大的惰性,一种往往压得人喘不过气来的重负。

% ---------------- p.128 ----------------
小孩子发现世界,就像他呼吸——呼吸的涨落让他把世界接进自己娇嫩的存在里,也让他投身于接纳他的世界。成年人也在发现,在那些稀少的时刻:他忘了自己的恐惧和知识,用睁得大大的、渴于认知的眼睛看事物、看自己,一双崭新的眼睛——一双孩子的眼睛。

\asterisms

\origpage{2} 上帝一边发现世界,一边创造世界;更确切地说,他永恒地\textbf{创造}着世界,一边创造一边发现——也一边发现一边创造。他创造了世界,又日复一日地创造着它,反反复复,千百万个千百万次,不歇息,摸索着,错了千百万个千百万次又拨正航向,不知疲倦……每一次,在这场用探针敲击万物、听万物回应的游戏里(它们回答:“这一下敲得不赖”,或者:“这一下你可是彻底犯浑”,或者:“顺顺当当,就这么来”),在又一次探针中——它修正或重提前一次探针,回应前一个回答……在造物主与万物之间这场无限对话的每一次往复里(这对话在创世的每一刻、每一处发生),上帝学习着、发现着,他越来越亲切地认识万物,而万物也在他手中渐渐有了生命、有了形状、不断变形。

这就是发现与创造的路数;亘古以来,似乎一直如此(就我们所能认识的而言)。一直如此——并不需要人类姗姗来迟地登场(不过是一两百万年前的事),撸起袖子插手进来——再加上最近那些众所周知的糟糕后果。

我们当中,不时有人发现这样或那样的事物。有时他会在自己的生命里重新发现——怀着惊叹——\textbf{发现}究竟是怎么一回事。每个人身上都有他所需要的一切,去发现这广袤世界上吸引他的一切,包括他身上那份奇妙的能力——世上最简单、最显然不过的东西!(然而许多人已经忘了的东西,就像我们忘了唱歌,忘了像孩子那样呼吸……)

每个人都能重新发现发现与创造是怎么一回事,而没有人能把它发明出来。它们先于我们就在那里,它们就是它们之所是。

\section{(2) 错误与发现}

回到我真正意义上的数学工作的“风格”,或它的“性质”、它的“路数”:它们如今与从前一样,都是好上帝亲自教给我们的——不用言语,教给每一个人,天知道在什么时候,也许远在我们出生之前很久。我照他的样子做。这也是每个人出于本能所做的,一旦好奇心驱使他去认识某一样东西——从那一刻起,这东西就被这份欲望、这份渴念占据了……

\origpage{3} 当我对一样东西感到好奇,数学的也好,别的也好,我就去问它。我去问它,不在意我的问题也许很蠢、或者会显得很蠢,也不要求它非得先经过反复掂量。问题常常采取断言的形式——一个其实只是探针的断言。我对自己的断言信几分?这当然取决于我在理解眼前之物上走到了哪一步。往往,尤其在一项研究的开头,这断言干脆就是错的——可还是得把它做出来,才能让自己信服它是错的。往往,只消把它写下来,
% ---------------- p.129 ----------------
错就一目了然;而在写下来之前,只有一片模糊,一种不适,而不是这份明白。这样就能少一分无知,再次进攻,带着一个也许不那么“离谱”的问题—断言。更常见的是:断言照字面理解是错的,但那个仍旧笨拙地试图借它表达出来的直觉是对的,虽然还模糊。这直觉会渐渐从一层同样不成形的脉石中析出——起初是些错误或不恰当的观念;它会渐渐走出“未被理解”的幽冥(那只是等着被理解),走出“未知”的幽冥(那只是等着被认识),取得一个只属于它自己的形式,让轮廓变得精细、鲜明——随着我向眼前这些事物提出的问题越来越精确、越来越切题,把它圈得越来越近。

但也会这样:靠这条路子,反复的探针会聚成局势的某一幅图像,从雾中显出相当清晰的线条,足以让人开始相信这幅图像很好地表达了实情——其实满不是那么回事,只要这图像沾上一个分量十足的错,一个足以叫它深深走样的错。从最初发现的“对不上号”出发——所得图像与某些确凿事实对不上,或与另一些同样为我们所信任的图像对不上——逐步查出这样一个错误观念的工作,有时很艰苦,而且往往伴随着越来越紧的张力,随着人逼近矛盾的结,那矛盾先还模糊,渐渐越来越刺耳——直到它终于爆开:错误的发现,某一种世界观的坍塌,像一次解放。\origpage{4}\textbf{错误的发现,是一切发现工作中最关键的时刻之一,是其中最富创造性的时刻之一},无论那是数学工作,还是发现自我的工作。正是在那一刻,我们对所探究之物的认识突然更新了。

\textbf{怕错与怕真,是同一回事。}害怕弄错的人,无力发现。正是在我们害怕弄错的时候,我们身上的错误才变得岩石一般不可动摇。因为在恐惧中,我们死死抓住某一天被我们判定为“真”的东西,或者抓住那些历来就被当作“真”呈现给我们的东西。当我们不是被害怕失去一种虚幻安全的恐惧所驱动,而是被求知的渴念所驱动,那时错误就像痛苦或悲伤一样,从我们身上穿过,永不定形,而它经过的痕迹,是一种更新了的认识。

\section{(3) 不可告人的劳作}

一切真正研究的自发路数,可以说从来不出现在那些本应传达、传递“被找到”之物的实质的文本或话语里——这绝不是偶然。文本与话语,多半只满足于记下一些“结果”,用的形式,在凡夫俗子眼里,必须显得像一条条严峻而不变的定律,亘古以来就刻在某种巨型图书馆的花岗岩法版上,由某位全知的神口授给那些入门者—书记—学者之流;口授给那些写学术书和同样学术的论文的人,那些站在讲坛高处、或在讨论班更小的圈子里传递知识的人。有没有哪怕一本教科书、一本供小学生、中学生、大学生乃至“我们的研究者”用的手册,能给不幸的读者一点点关于研究究竟是什么的概念——除了那个普遍接受的概念:研究嘛,就是当你很有学问、考过一大堆考试甚至竞考的时候,就是那些大脑袋,巴斯德和居里,还有诺贝尔奖得主什么的……而我们这些读者或听众,勉勉强强把大人物们为了人类的利益而惠予记下的“知识”吞下去,我们顶多(如果肯下苦功)够格通过年终考试,而且
% ---------------- p.130 ----------------
还不一定……

\origpage{5} 这样的人有多少呢——包括那些不幸的“研究者”自己,那些为论文或文章发愁的人,甚至包括我们当中最有“学问”、最有声望的人——有谁天真到看得出:“寻找”不是别的,正是怀着激情\textbf{去问}万物——就像一个孩子\textbf{想知道}他自己或他的小妹妹是怎么来到这世上的。寻找与找到,也就是:提问与倾听,是世上最简单、最自然不过的事,世上没有人拥有它的特权。这是我们一出生就全都收到的一份“馈赠”——它注定要以无穷无尽的面孔表达自己、舒展自己,此一时与彼一时不同,一个人与另一个人也不同。

当你冒险把这类话说出口,你在这一位和那一位那里收到的,从最确信自己是笨学生的笨学生,到最确信自己有学问、远在凡夫之上的大学者,都是一样的笑:一半尴尬,一半心照不宣,好像你刚讲了一个有点粗俗的笑话,好像你在显摆一种一眼就能看穿的天真;“这话都不错,谁也不该看不起,这我同意——可也别太过分啊——笨学生就是笨学生,他不是爱因斯坦,也不是毕加索!”

面对这样一致的赞同,我再坚持就不识趣了。真是不可救药,我又一次失去了闭嘴的机会……

不,这绝不是偶然:各类教育书、劝世书和各种手册,异口同声地把“知识”呈现得像是一出娘胎就穿得整整齐齐,从那些天才的头脑里出来——那些头脑为了我们的益处把它记了下来。也不能说这是恶意,即便在少数场合,作者相当“门儿清”,知道那幅图像(他的文本免不了要暗示的图像)与实情毫不相符。在这种场合,陈述有时不止是一堆结果和诀窍:有一股气息穿过它,有一种活的视野赋予它生命,而这视野有时便从作者传到专心的读者那里。但一种默契——力量看来相当可观——使得文本不留下半点它由之产生的那些工作的痕迹,即使它以镌刻般的笔力表达出对事物有时相当深刻的视野,而那视野正是这工作真正的果实之一。

\origpage{6} 说实话,有些时候,我自己也隐隐感到这股力量、这份无声共识的分量——就在我筹划写作、出版《数学反思》的时候。如果我试着去探一探这份共识所采取的默不作声的形式,或者更确切地说,去探一探我内心对我这计划的抵抗(它由这份共识引发)所采取的形式,立刻冒出来的词是“不体面”。这份共识,不知从何时起已内化在我心里,它对我说(而这是我头一回费心把它这几周、也许是这几个月来带几分固执地对我嘟囔的东西,拉到日光下,拉进我的视野里):“把一项发现工作中的起起落落、边角上那些弄砸了的摸索、总之那些‘脏衣服’,摊在别人面前,甚至当众摊开——这不体面。这只是浪费读者的时间,而读者的时间是宝贵的。再说,这还要多出好多页,得排版、得印刷——多浪费啊,科学印刷品价钱这么贵!非得虚荣到什么地步,才会这样摊开那些对谁都没意思的东西,好像连我那些胡乱的折腾也成了了不起的事——总而言之,一个摆谱的机会。”还有更隐秘的一层:“发表这样一份反思的笔记——照它真正进行的样子——是不体面的,就像在公共广场上做爱是不体面的,或者把分娩的劳作留下的血污床单拿出来展示、哪怕只是让它晾在那里,也是不体面的……”

这里的禁忌,采取的是性禁忌那副既阴险又专横的面孔。正是在写这篇导论的时候,我才开始隐约看见它非凡的力量,以及这本身即非凡的事实所含的意义——这事实正是那股力量的证明:发现的真正路数,如此
% ---------------- p.131 ----------------
令人不知所措地简单,一种孩子气的简单,却几乎无处透出;它被无声地抹去、忽视、否认。甚至在科学发现这个相对无害的领域里也是如此——谢天谢地,那不是发现自己的“小鸡鸡”之类的事——总之是一项可以放心交到所有人手里的“发现”,而且(人们会以为)没什么可隐瞒的……

倘若我愿意顺着这里出现的这根“线”走下去——这根线一点也不细,反而粗壮结实得很——它肯定会把我带到远得多的地方,远过那几百页同调—同伦代数;那些页面我终究会写完,交给印刷厂的。

\section{(4) 他人之无谬与对己之轻蔑}

\origpage{7} 刚才我谨慎地说“我的表达风格”变了,还让人以为这没什么可奇怪的——说到底,那是个委婉说法。您明白:十三年没写过东西,当然跟从前不一样了,“表达风格”总得变,必然的……区别在于,从前我“表达”(原文如此)得像所有人一样:我先做工作,然后倒过来重做一遍,把所有的涂改都仔细擦掉。擦着擦着又出现新的涂改,把整个工作搅得有时比初稿还糟。于是又得重做——有时三遍,甚至四遍,直到一切都无可挑剔。不但没有可疑的角落,也没有悄悄扫到合宜的家具底下的垃圾(我从来就不喜欢角落里的垃圾,既然都费心扫地了);尤其是,读定稿的时候,它给人的印象固然体面(像任何别的科学文本一样),那就是:作者(此处即区区在下)是化为肉身的无谬。他万无一失地正好落在“那些”正确的概念上,然后是“那些”正确的陈述上,一环扣一环,像上了油的马达一样低吟着,那些证明“咚”的一声落下来,每一个都恰好在它该来的时候!

\origpage{8} 效果如何,请设想一个毫无戒心的读者:比如说一个正在学毕达哥拉斯定理或二次方程的中学生,甚至我的一位在所谓“高等”研究或教学机构里的同行(明白人一听就懂,此致!),(不妨说)正费劲地读某位显赫同行的某篇文章!这类经验在一个小学生——甚至大学生或研究者——的一生中重复成百上千次,再加上家庭里以及全世界所有国家的所有媒体里那场合宜的合唱,其效果可想而知。只要肯留心,人在自己身上、在别人身上都能看到:那是一种对自身毫无价值的内心确信,与那些“懂行的人”、“干事的人”的能力和重要性恰成对照。

这份内心的确信,有时会被一种能力所补偿——记住不理解之物的能力,甚至某种操作上的灵巧:矩阵相乘,用“正题”和“反题”拼出一篇法文作文……总之是鹦鹉或耍把戏的猴子的本领,如今比以往任何时候都更受看重,由人人艳羡的文凭认可,由安稳的职业生涯报偿。但即便那位文凭缝满一身、位置坐得安稳、也许还满身荣誉的人,在心底里也并不上当,并不把这些重要、这些“价值”的假标志当真。甚至那位更少见的人——他把全部赌注押在某样真正天赋的发展上,在职业生涯中懂得施展身手、做出创造性作品——在心底里也并不信服自己名声的光彩,尽管他常常想借这光彩让自己、也让别人看走眼。一种从未被审视过的同样的怀疑,住在这一位和那一位身上,也住在随便哪个笨学生身上;一种同样的确信,也许他们永远不敢去认识它。

正是这份怀疑、这份未曾说出的内心确信,推动着这一位和那一位不断地超越自己——
% ---------------- p.132 ----------------
在荣誉或作品的积累中;并且把他们私下里啃噬着自己的那份对己之轻蔑,投射到别人身上(尤其是那些他们对其多少有些权力的人身上……)——这是一场注定不可能的逃脱尝试,靠的是积累他们优于别人的“证据”\fnoteP{2}{(2)}{(1984年3月补写)重读最后这两段,我有些不适之感,因为写它们的时候,我牵连的是别人,而不是我自己。显然,写的时候,“我自己也可能与此有关”这个念头并没有碰到我。我肯定什么也没学到:我就这样满足于把多年来在别人身上察知、并以许多方式看着被证实的东西,白纸黑字写下来(无疑还带着某种满意)。在这份反思的后续部分,我不得不回想起:我的一生中,对别人的轻蔑姿态并不缺少。我所把握到的他人之轻蔑与对己之轻蔑之间的关联,若在我这个人身上竟然缺席,那才奇怪;健全的理性(以及类似情形下我对自身视而不见的经验——我终究还是察觉到了)告诉我:肯定不该是这样!然而眼下这只是一条纯粹的推论,它唯一可能的用处,是促使我亲眼看一看究竟发生了什么,去看、去审视(如果它确实存在过、或还存在着)这份仍属假设的对己之轻蔑——它埋得如此之深,至今完全逃过了我的目光。真的,可看的东西并不缺!这一件忽然在我看来是最要紧的几件事之一,恰恰因为它藏得如此之深。……\emph{[(1984年8月)不过关于这一点,参见“屠杀”一注(n\textsuperscript{o}~87)末两段的反思。]}} (2)。

% ================================================================
%  第 6 章 梦与做梦者 (原书 p.133-146, 前半: 6.1-6.8)
%  源文: _work/src1/task-ch06a.md
%  说明: 本章后半 (p.147-152, 6.9-6.11) 见 part1-ch06b.tex
% ================================================================
\chapter{梦与做梦者}

\noindent 1984年2月\par

% ---------------- p.133 ----------------
\section{(5) 被禁的梦}

\origpage{9} 借着《追寻栈》(À la Poursuite des Champs)写作中断三个月的机会,我要把《导论》接着去年六月停下的地方写下去。相隔半年多,我刚把它仔细重读了一遍,并添了几个小标题。

写这篇《导论》时,我很清楚:这类思考免不了招来一大堆“误解”——而想事先堵住它们也是白费力气,那无非是在原有的误解之上再堆上一些!关于这一点,我只想补充一句:我丝毫无意向那种由千年惯例确立起来的科学写作风格开战——我自己就曾勤勤恳恳地奉行它二十多年,还把它当作数学家这一行当的本质部分教给我的学生们。不管对错,今天我也仍然这样看待它,并且继续这样教。甚至可以说,我倒更显得老派:我坚持工作要做到底,从头到尾一针一线地缝,不放过任何一个稍有阴影的角落。近十年来我不得不往酒里掺水,那实在是形势所迫!“规范的写作”在我看来仍然是数学工作的重要一环,既是一种发现的工具——用来检验并深化对事物的理解,没有它,这理解始终是粗略而零碎的——也是一种传达这种理解的手段。从教学的角度看,严格的表述方式,也就是演绎的方式(它绝不排除铺开宏阔画卷的可能),有明显的好处:简洁,查考方便。这些好处是实实在在、有分量的——当讲稿面向
% ---------------- p.134 ----------------
\origpage{10} 数学家(比如说)、而且是已经对所述主题的某些来龙去脉、或对与之十分相近的主题足够熟悉的数学家时。

然而,若一篇讲稿的对象是孩子、是年轻人,或是那些事先完全不“在行”的成年人——他们的兴趣已经觉醒,而且多半(并将继续如此,而且事出有因……)对一项发现性工作的真正进程全然无知——那么这些好处就完全成了幻影。更确切地说,是这样的读者:他们连这种工作的存在都不知道——而这工作其实人人做得到,只要他有好奇心和常识;我们对宇宙万物的智性认识,正是从这工作中不断生出、又不断再生,包括那些以宏伟的编排体系表达出来的认识,比如欧几里得的《几何原本》,或达尔文的《物种起源》。对这种工作的存在与性质全然无知,几乎是普遍现象,各级教师也不例外,从小学教师到大学教授。这是一个非同寻常的事实;去年写这篇《导论》的第一部分、开始这番反思时,它才清清楚楚地显现在我眼前——同时也让我隐约看到了这件令人困惑之事深深的根……

即便讲稿面对的读者在各方面都确实“在行”,仍有一件重要的东西,是“严格的”表述方式不许自己传达的。而这东西,在严肃人士的圈子里——尤其在我们这些科学家当中——也完全不受待见!我说的就是梦。梦,以及梦向我们耳语的那些景象——起初像梦一样不可触摸,而且往往迟迟不肯成形。长年累月、甚至整整一生的紧张工作,也许都不足以让某个梦的景象完全显现出来,不足以看它凝聚、打磨,直到获得钻石的硬度与光华。这就是我们的工作,无论我们是用手还是用精神做工的工匠。当工作完成、或工作的某一部分完成时,我们把那有形的成果呈现在所能找到的最亮的光下,我们为它欢喜,还常常以它为荣。然而,我们长久打磨的这颗钻石里,并没有那个在打磨中激励我们的东西。也许我们造出了一件极精确的工具、一件好用的工具——但工具本身就是有限的,就像人手所做的一切,哪怕它在我们看来很了不起。一个景象,起初没有名字、没有轮廓,淡薄得像一片雾霭,它引导着我们的手,让我们伏在活儿上,感觉不到时辰流逝,也许也感觉不到年月。一片从雾与幽暗的无底之海中无声地剥落下来的雾霭……我们身上那无界限的东西,就是她,就是这片海——只要我们的渴望使她受孕,她便随时准备孕育、随时准备分娩。梦正是从这场婚配中涌出,像蜷伏在养育它的子宫里的胚胎,等待着幽暗的劳作把它带向第二次诞生,带向天光。

一个蔑视梦的世界有祸了——那也是我们身上深刻的东西遭蔑视的世界。我不知道在我们这种文化(电视、电脑和洲际火箭的文化)之前,是否还有别的文化宣扬过这种蔑视。这大概正是我们区别于前人的许多地方之一——我们把前人彻底取代、几乎可以说从地球表面清除掉了。我不曾知道有哪一种别的文化是不尊重梦的,不曾知道有哪一种文化里,梦的深根不被所有人感受到、不被承认。在一个人或一个民族的一生中,可曾有哪件宏大的作品,不是生于梦、不是先由梦哺育,而后才在大白天里绽放的?然而在我们这里(是不是已经该说:到处都一样?),对梦的敬重被叫作“迷信”;而众所周知,我们的心理学家和精神病学家早已把梦翻来覆去、量了个遍——那点分量,想必勉强够塞满一台小型电脑的存储器。也确实,“在我们这里”,再没有人会生火,也没有人敢在自己家里看着孩子出生,或看着母亲、父亲死去——
% ---------------- p.135 ----------------
\origpage{11}有诊所和医院就是干这个的。谢天谢地……我们这个世界,如此自豪于它的力量——以原子百万吨级计,以图书馆和电脑里储存的信息量计——恐怕也正是这样一个世界:在这里,每个人的无能,面对生活中那些简单而本质之物时的恐惧与蔑视,达到了顶点。

\origpage{12} 幸好,梦就像最压抑的社会里那原初的性冲动一样,命硬得很!不管是不是迷信,它仍旧偷偷地、执拗地向我们耳语一种认识——那是我们醒着的精神太笨重、或太怯懦,把握不住的认识——并且给它启发我们形成的种种计划注入生命、插上翅膀。

我刚才似乎暗示过,梦往往不肯成形;那只是表面现象,并未真正触及事情的底。那“不肯”多半来自我们醒着的精神,来自它平常的“心态”——而且“不肯”这个词还是轻描淡写!那更像是一种深深的戒心,它掩盖着一种古老的恐惧——认识的恐惧。就本来意义上的梦而言,梦的消息越是贴近我们,越是因为一旦被听见就会带来我们人格深处蜕变的威胁而显得沉重,这恐惧就越是起作用,它立起的屏障就越是有效。但看来,即便在数学的“梦”这种相对无害的场合,这份戒心也依然在场、依然有效;以至于一切梦似乎不仅被逐出了文本(反正我连一部记下梦的文本也没见过),而且也被逐出了同事间的谈话——不论是小圈子里,甚至一对一的时候。

事情若果然如此,那当然不是因为数学的梦不存在、或不再存在了——若真如此,我们的科学早就贫瘠了,而绝非如此;这种表面上的缺席、这种沉默的合谋,其原因想必与另一种共识紧密相连——那就是要小心抹去一切痕迹、一切提及:抹去那借以完成发现、借以更新我们对世界之认识的工作。或者不如说,围绕着梦、也围绕着梦所唤起、所启发、所滋养的工作的,是同一个沉默。以至于“数学的梦”这个说法本身,在许多人看来都会是无意义的胡话——我们如此经常地被一些按电钮式的陈词滥调驱动,而不是被我们对某种极其简单、日常而又重要的现实的直接经验驱动。

\section{(6) 做梦者}

事实上,经验告诉我:当精神渴望认识梦,而不是逃避它(或手里拿着一副现成的框框去套它——那是一样的),梦就丝毫不会“不肯成形”——它肯让人细致地描述,肯交出它的消息:那消息永远简单,从不愚蠢,有时还叫人震动。恰恰相反,我们心中的“做梦者”是一位无与伦比的大师:他每一次都能找到、或凭空造出最合用的语言,来绕过我们的恐惧、摇醒我们的麻木;他的舞台手段变化无穷,从完全没有一丝视觉或感官要素,到最令人瞠目的场面调度。当他显现时,绝不是为了躲开,而是为了鼓励我们(几乎总是白费力气,而他的好意从不疲倦……)走出我们自己,走出他看见我们深陷其中的那种沉重——有时他还会不动声色地拿这种沉重寻开心,用滑稽的色调把它戏仿一番。倾听我们心中的“做梦者”,就是与我们自己交流——逆着那些强大的堤坝,它们非要禁止我们这样做不可。

但能做大的人,也能做小。既然我们能借梦与自己交流、向自己显现自己,那么,把数学的梦——比如说——那绝不私密的消息传达给别人,想必也同样简单易行;这种梦并不牵动
% ---------------- p.136 ----------------
\origpage{13} 力量相当的抵抗。

而说实话,在我做数学家的过去里,除了追随一片片从浓重致密的雾的织物上一一剥落下来的梦,把它们“梦”到底——直到它们最显明、最坚实的显现,无可辩驳的显现——我还做过别的什么?又有多少次,我为自己那份固执而急得跺脚:我总要把每块凝结着我的梦的宝石、或半宝石,一丝不苟地打磨到最后一个刻面——而不是听从一股更深的冲动:去追随那母体织物千变万化的幽径,直到梦与它显在的、可以“发表”的化身(总之,按通行的规范)之间那片朦胧的边界!况且,我当时正要听从那股冲动,投入一项“数学科幻”的工作——一部关于“Motif”理论的“某种白日梦”,那理论当时还纯属假设,而且直到今天仍然如此,而且理所当然:因为再没有另一个“白日做梦者”投身这场冒险。那是六十年代末;那时我的生活(我丝毫没有想到)正要拐一个大弯,此后十来年间,这转折把我的数学热情挤到了边缘的位置,甚至遭到否认。

但细究起来,《追寻栈》——沉默十四年之后的第一部出版物——恰恰属于那场从未写出的“白日梦”的精神,而它似乎成了那场梦临时的后续。当然,至少初看之下,这两个梦的主题之间的差异,是两个数学主题之间所能有的最大差异;何况头一个(即 Motif 那个)似乎位于“用现有手段可做之事”的地平线之外,而第二个——即著名的“栈”及其同类——看起来完全伸手可及。这些差异,或许可以称为偶然的、意外的差异,它们也许会比人们预料的早得多就烟消云散\fnoteP{1}{(3)}{我在这里特别想到已故的莫德尔(Mordell)猜想、泰特猜想、沙法列维奇(Chafarévitch)猜想——去年,法尔廷斯(Faltings)一篇四十页的手稿把这三个猜想一并证明了;而就在此前,“在行”的人们早已确立的共识是:这些猜想“够不着”!恰好,我那心爱的“远阿贝尔代数几何”纲领当作拱心石的那条基本猜想,正与莫德尔猜想相近。(甚至据说后者可以由前者推出——这就更说明,这个纲领不是给严肃人士看的故事……)} (3)。在我看来,这些差异对于两个主题各自可能引发的这类工作,影响相对很小——因为这里说的正是“白日梦”,或者用不那么挑衅的说法:把概念上的粗加工一直做到一个足够融贯、足够精确的整体景象,足以或多或少完整地使人确信:这景象在根本上确实符合事物的实际。就本书所展开的这个主题而言,这大致意味着:对这一景象之有效性的逐项验证,变成了一件纯属本行手艺的事。这当然可能要求巨大的工作量,其中也需要机巧和想象,而且想必还有意外的转折与前景,幸而会使它不至于成为纯粹的例行公事(用安德烈·韦伊的话说,一场“长长的练习”)。

总之,这正是我过去做腻了又重做的那类工作,它就在我指尖上,因此在我余下的岁月里不必再做。只要我再度投入数学工作,我的精力肯定最该用在“白日梦”的边缘地带。在这个选择里,启发我的并不是什么效益考虑(假如这种考虑能启发得了谁的话),而恰恰是一个梦,或一些梦。我身上这股新的冲劲若能显出力量,那力量必定是从梦里汲取的!

% ---------------- p.137 ----------------
\section{(7) 伽罗瓦的遗产}

\origpage{14} 看起来,在所有的自然科学中,只有数学把我所谓的“梦”、或“白日梦”判以一道看来绝对的禁令,而且这禁令已施行两千多年。在别的科学里,包括物理学这样公认“精确”的科学,梦至少是被容忍的,甚至(视时代而定)是被鼓励的,只不过用的是些更“拿得出手”的名字,比如:“思辨”、“假设”(如著名的“原子假设”,它出自一个梦——对不起,出自德谟克利特的一种思辨)、“理论”……从那个“不敢道出自己名字的梦”的身份,过渡到“科学真理”的身份,是经由几乎察觉不到的层层台阶、经由一个逐渐扩大的共识完成的。相反,在数学里,这几乎总是(至少如今)一次突变,靠的是一篇证明的魔杖一挥\fnoteP{2}{(4)}{况且就在今天,我们也还遇得到身份可疑的“证明”。格劳尔特(Grauert)那个以他命名的有限性定理,多年来就是这样:他那个证明谁也不曾读完(而愿意读的人并不缺!)。这一困窘后来由另一些更透明的证明解决了,其中有些还走得更远,它们取代了最初那个证明。类似而更极端的情形,是所谓的“四色问题”的“解决”:其中计算的部分是用电脑(以及几百万美元)了结的。所以那是一个“证明”,它的依据不再是理解一个数学情境而来的内心确信,而是对一台没有理解能力的机器的信任——而使用它的数学家既不知道这台机器的结构,也不知道它如何运作。即便计算再由别的电脑、按别的计算程序核实一遍,我也不认为四色问题就此了结。它只是换了一副面孔:如今要做的几乎不再是找一个反例,而只是要一个证明(当然,得是读得懂的证明!)。} (4)。在那样一些年代——数学定义与数学证明的概念还不像今天这样清楚,还不成其为(或多或少)普遍的共识——却也有些显然重要的概念,身份暧昧不明,比如“负”数(帕斯卡拒斥过它),或“虚”数。这种暧昧,至今仍反映在通行的语言里。

\origpage{15} 定义、陈述、证明、数学理论这些概念的逐步澄清,在这方面极为有益。它让我们意识到我们手头的工具——尽管简单得像小孩玩意——有多么强大:靠它们,我们能把那看似无法表述的东西表述得完美精确——几乎只要足够严格地使用日常语言就行。如果说有什么东西从我童年起就让我着迷于数学,那恰恰是这种力量:用词语去圈住、并完美地表达某些数学事物的本质——这些事物乍看之下形态如此飘忽、如此神秘,仿佛在言语之外……

然而这种力量、这种完美精确与证明所提供的资源,有一个糟糕的心理上的反作用:它们更加强化了对“数学的梦”的传统禁忌;也就是说,强化了对一切不肯以常规的精确面貌出现的东西的禁忌(哪怕牺牲的是更宏阔的景象)——那种精确要有形式完整的证明来担保它“货真价实”,否则(而且如今越来越如此……)就得有据信可以整理成形的证明概要来担保。偶尔的猜想,勉强可以容忍,条件是它得符合问卷式的精确要求,而问卷上只允许“是”或“否”两种回答。(而且,还用说吗,还得是那个敢提猜想的人在数学界里有门有面。)据我所知,还没有过一个例子:把一门在其本质部分明明白白是猜想性的数学理论,当作“实验”来发展。确实,按现代规范,十七世纪以来发展起来的整套“无穷小”演算——后来成了微分学与积分学——会显得像一场白日梦,而它最终要等到两个世纪之后,才靠柯西的魔杖一挥,变成
% ---------------- p.138 ----------------
严肃的数学\origpage{16}。这必然让我想起埃瓦里斯特·伽罗瓦(Evariste Galois)的白日梦——他不巧也碰上了同一位柯西,运气不佳;但这回不到一百年就够了:又一下魔杖(这回是若尔当,如果我没记错),就让这场梦获得了公民权,并为此改名为“伽罗瓦理论”。

从这一切得出的结论——这结论对“1984年的数学”并不有利——是:幸好牛顿、莱布尼茨、伽罗瓦这样的人(肯定还漏掉很多,我对历史并不内行……)当年没有被我们今天的规范绊住;那时候,他们只管发现,没工夫立规范!

伽罗瓦这个例子不请自来,触动了我的一根敏感的弦。我似乎记得,第一次听人说起他和他的奇崛命运时(那时我大约还在念中学或大学),一种兄弟般的亲切感就在我心里醒来了。和他一样,我感到自己心中有一份对数学的激情——也和他一样,我觉得自己是个边缘人,是那个(在我看来)把他抛弃了的“上流社会”里的异乡人。可我自己最终也成了这上流社会的一员,又在某一天离开它,没有遗憾……这份有点被遗忘的亲近感,最近又以全新的面貌重新出现在我身上——当时我正在写《纲领概要》(Esquisse d'un Programme,为着我申请作为研究员进入国家科学研究中心[CNRS]的事)。这份报告主要概述我近十年来一些主要的思考主题。在所有这些主题中,最让我着迷、也是我打算在今后几年着重发展的那个,正是数学的梦的典型——它还接上了“Motif 之梦”,为后者提供了一条新进路。写这份《草稿》时,我想起了这十四年里我连续做过的、为期最长的一次数学思考。它从1981年1月一直持续到6月,我给它取名叫《穿越伽罗瓦理论的长征》。一针引一线,我意识到:我这几年来时断时续地追着的那场白日梦,后来得了“远阿贝尔代数几何”之名,其实无非是伽罗瓦理论的一个延续,“伽罗瓦理论的终极归宿,而且想必是合乎伽罗瓦本人的精神的”。

当我写到被引的那句话所在的那一段、这种连续性在我眼前出现时,一阵喜悦穿过我,至今未散。那是在完全的孤独中做下来的工作所得到的报偿之一。它的出现,与不久前我从两三位同事兼旧友那里受到的、颇为冷淡的接待同样出乎意料——他们明明是“在行”的人,其中一位还曾是我的学生;我曾有机会跟他们谈起我正在发现的东西,那时我还“热着”,心里满是欢喜……

这让我想到:今天接续伽罗瓦的遗产,必定也要接受他当年那份孤独的风险。也许时代的变化比我们想的要小;不过,这份“风险”在我眼里往往并不成为威胁。我若有时因所爱之人的故作冷淡或轻蔑而难过、失落,那么相反,多年来孤独——数学的或别的孤独——从不曾让我觉得沉重。若有一位忠实的女友,我每次离开她后都不断渴望再回到她身边,那就是她!

\section{(8) 梦与证明}

但让我们回到梦,回到几千年来压在它身上的那道数学禁令。这也许是所有先验之见中最顽固的一条——它往往不明说,扎根在习惯里,裁决着什么东西“算数学”、什么不算。要过几千年,那些像孩童玩意一样简单
% ---------------- p.139 ----------------
而又无所不在的东西——某些几何图形的对称群、另一些图形的拓扑形态、数字零、集合——才获准进入圣殿\origpage{17}!

\origpage{18} 当我跟学生们谈球面的拓扑,谈由球面添上环柄所得的那些形态——这些东西并不叫小孩子吃惊,却叫他们不知所措,因为他们自以为知道“数学”是什么——我得到的第一个自发反应总是:这哪算数学!数学嘛,当然是勾股定理、三角形的高、二次多项式……这些学生并不比你我更蠢;他们的反应,和古往今来、直到今天的世上所有数学家的反应一样,除了毕达哥拉斯或黎曼这样的人,也许还有五六位别的。就连庞加莱——那可不是等闲之辈——也能用一套深有感触的“甲加乙”式哲学,“证明”无穷集合算不上数学!想必也曾有过一个时候,三角形和正方形算不上数学——那只是孩子或陶匠在沙上、在瓶罐的泥坯上画的图,别搞混了……

精神被自己的“知识”憋住,这种根子上的惰性当然不是数学家独有的。我有点偏离正题了:我要谈的是压在数学的梦身上的禁令,以及借由它、压在一切不以制成品的惯常面貌出现(不是开箱即用)的东西身上的禁令。我对别的自然科学所知的那一点,已足以让我估量出:同样严厉的禁令早就会把它们判给贫瘠,或判给乌龟般的进展——有点像中世纪,那时容不得对《圣经》的字句刨根问底。但我也很清楚:发现的深源,以及发现之进程在其一切本质方面,在数学里与在宇宙中我们的身体和精神所能认识的任何别的领域、别的事物里都是一样的。放逐梦,就是放逐源头——判它只能在暗处存在。

我还很清楚(这经验自我少年时代与数学初恋爱慕以来从未被推翻):在展开一种广阔或深刻的数学事物之景象时,正是这景象与理解的展开、这渐进的深入,始终先于证明,使证明成为可能,并赋予它意义。当一个情境——从最卑微到最宏阔——在其本质方面被理解了,对所理解之物的证明(以及对其他东西的证明)就像熟透的果子一样落下来。而像青果一样从认识之树上硬拽下来的证明,会留下一丝不满足的余味,一种渴望未得平息的失落。在我做数学家的一生中,有两三次,我不得不在别无办法时下定决心去拽果子,而不是摘果子。我不是说我做得不好,也不是说我后悔。但我所做得最好的、我所最爱的东西,我都是心甘情愿取来的,不是强夺来的。如果说数学给了我丰盈的快乐,在我成熟之年仍让我着迷,那靠的不是我从它那里硬拽来的证明,而是我在它身上感到的无穷奥秘与完美和谐——它随时准备向一只爱它的手、一道爱它的目光显露自身。

\section{(9) 受欢迎的异乡人}

我觉得时候到了,该谈谈我与数学家世界的关系。这与我与数学的关系完全是两回事。后者从我年少时起就存在,而且很强,远在我想到世上竟有一个数学家的世界、一个数学家的圈子之前。那是一整个复杂的世界,有它的学会、它的期刊、它的聚会、讨论会、大会,有它的头牌女伶和苦力,有它的权力结构,有它的灰衣主教,还有同样灰扑扑的一大群任人使唤、任人摊派的人,为论文或文章发愁;也有另一些更少见的人,他们有的是手段和想法,却
% ---------------- p.140 ----------------
处处碰壁\origpage{19},绝望地找不到一位有权有势者撑腰——那些忙忙碌碌、叫人畏惧的强人握有一种魔法:让一篇文章得以发表……

\origpage{20} 我是在1948年、二十岁上到巴黎时,才发现有一个数学世界的;我那只瘪瘪的手提箱里,装着蒙彼利埃大学的理学学士学位,还有一部行距排得极密、双面书写、不留页边的手稿(纸很贵!),那是我三年孤独思考的记录,思考的东西(后来我才知道)当时早已以“测度论”或“勒贝格积分”之名广为人知。因为从没遇见过别人,我直到进京城那天都还以为,世上只有我一个人在“搞数学”,因此也就是唯一的数学家。(在我眼里这两者是一回事,而且直到今天也还多少是。)我曾玩弄那些我称之为可测的集合(其实也没遇到过一个不可测的集合……),玩弄几乎处处收敛,却不知道拓扑空间是什么。有十来个互不等价的“抽象空间”与紧性的概念,我一直没搞清,那是从一本小册子里捞来的(作者好像叫阿佩尔,收在《科学与工业时事》丛书里),这本小册子是我碰上的——天知道怎么碰上的。那时我还没听人说过——至少没在数学的语境里听过——群、域、环、模、复形、同调(还有别的,就不数了)这类古怪或野蛮的词;忽然之间,它们不打招呼,一齐朝我压下来。这一击不轻!

我之所以在这一击之下“活了下来”,还继续搞数学、甚至以此为业,是因为在那遥远的年头,数学世界还不太像它后来变成的样子。也可能是我运气好,落进了这个未曾料想的世界里一个格外好客的角落。我带着蒙彼利埃学院一位教授的一封含糊的推荐信,那是苏拉先生(跟他的同事们一样,他也没怎么在课堂上见过我!),他当年是嘉当的学生(父亲还是儿子,我已说不太准了)。因为埃利·嘉当(Elie Cartan)当时已经“退出赛场”,他儿子亨利·嘉当(Henri Cartan)便成了我有幸遇见的头一个“同类”。我当时并不知道这是多么吉利的兆头!他接待我时的那份礼貌中透着善意,那是他的标志,几代高师学生都熟悉——他们有幸在他手下初试身手。从他当时给我指点学业方向的那些建议看,他大概没有意识到我无知到了何等地步。无论如何,他的善意显然是冲着人来的,不是冲着行李或可能有的天分,也不是(后来)冲着名声或名气。

接下来那一年,我成了“高师”嘉当一门课上的常客(讲流形上的微分形式体系),我拼命跟;也常去“嘉当讨论班”,目瞪口呆地看着他和塞尔辩论,满嘴“谱序列”(呜!),满黑板的图(叫做“图表”),画满了箭头。那是“层”、“甲壳”以及一整套军火库的英勇时代,那些东西的意义我完全摸不着边,可我还是勉勉强强逼着自己把定义和陈述吞下去,把证明一一验过。在嘉当讨论班上,谢瓦莱、韦伊也定期露面;而到了布尔巴基讨论班的日子(参加者和听众加起来,撑死也就二三十人),就看见这个有名的布尔巴基帮的其余成员——迪厄多内、施瓦茨、戈德芒、德尔萨特——像一伙有点吵闹的哥们儿那样闯进来。他们彼此以“你”相称,说着一种我几乎完全听不懂的共同语言,烟抽得凶,也爱笑;只差几箱啤酒,气氛就齐全了——而啤酒是由粉笔和海绵代替的。那气氛与法兰西公学院勒雷的课截然不同(讲肖德尔关于无限维空间中拓扑度的理论,可怜的我!),我是听嘉当的建议去听的。我曾到法兰西公学院找过勒雷先生,问他(如果我没记错)这门课要讲什么。他可能给我的解释,我一点不记得了,
% ---------------- p.141 ----------------
\origpage{21} 也不记得自己有没有听懂什么——只记得:在那里我也感到一种善意的接待,那是给随便哪个初来乍到的外乡人的。肯定正是这一点,而非别的什么,让我去上那门课,并且像在嘉当讨论班那样硬着头皮跟下去——尽管勒雷在那里讲的东西,我当时几乎完全不懂。

奇怪的是:在这个世界里我是新来的,它的语言我几乎不懂,说得更差,可我并不觉得自己是个外乡人。虽然我几乎没机会(而且当然没本事)跟韦伊、迪厄多内这样快活的家伙说话,或跟嘉当、勒雷、谢瓦莱这样气度更雍容的先生们说话,我却觉得自己被接纳了,几乎可以说:是他们中的一个。我不记得有哪一次被这些人中的某一位居高临下地对待,也不记得有哪一次,我的求知欲——以及后来,又一次,我发现的喜悦——被一份自得或一份轻蔑挡了回去\fnoteP{3}{(5)}{这个事实尤其值得注意:直到1957年前后,布尔巴基小组的成员中不止一位对我还颇有保留,而他们最终把我增选进去时,我相信也带着几分不情愿。一句善意的玩笑话把我归入了“危险专家”(指泛函分析)之列。我有时感到嘉当身上有一种更深的、未说出的保留——有那么几年,我一定给了他这样的印象:一个热衷于空洞浮泛的推广的人。他曾惊讶地发现,我为布尔巴基写的那一篇(也是唯一一篇)稍长的稿子(论流形上的微分形式体系)里,竟有一点相当实在的想法——而我提议由我来写时,他并不怎么热心。(多年以后,在从凝聚对偶的观点发展留数形式主义时,这一想法又派上了用场。)况且在布尔巴基的会议上,我多半是跟不上的,尤其是逐篇共同宣读稿子的时候,我实在没法跟上宣读与讨论的速度。可能我确实不太适合集体工作。不管怎么说,我在融入共同工作上的这份困难,或我因别的缘故在嘉当和别一些人那里可能引起的保留,都不曾在任何时候给我招来讥讽或碰钉子,或哪怕一丝居高临下——至多在韦伊那里有过一两次(此人确实是个特例!)。嘉当对我的态度,从未有一刻离开过一贯的和善,其中带着亲切,也带着他特有的那点幽默——在我,那点幽默始终与他这个人分不开。} (5)。若不是这样,我也不会像人们说的那样“成为数学家”——我会另选一行,在那里我能施展自己,而不必去面对轻蔑……

尽管“客观上”我是这个世界的异乡人,一如我在法国是个异乡人,却有一条纽带把我和这些出身别的阶层、别的文化、别的命运的人连在一起:一份共同的激情。我怀疑,在我发现数学家世界的那个关键年头,他们中没有一个人——连嘉当也没有,我勉强算他的学生,可他还有许许多多别的学生(而且是不那么跟不上的学生!)——在我身上察觉出那份同样盘踞在他们心中的激情。在他们眼里,我大概只是课堂和讨论班一大群听众中的一个,记着笔记,显然不太跟得上。如果说我多少有点与众不同,那就是我不怕提问——而那些问题多半只是暴露出我在语言和数学两方面都无知得惊人。回答可能简短,甚至带着惊讶,但当时那个目瞪口呆的冒失鬼,从未碰过钉子,从未被“放回原位”——无论在布尔巴基那伙人不拘礼数的圈子里,还是在法兰西公学院勒雷课堂上更严峻的氛围里。那些年里,自从我口袋里揣着一封给埃利·嘉当的信来到巴黎,我从未觉得自己面对着一个宗派、一个封闭乃至敌意的世界。如果说我领教过、深谙面对轻蔑时那种内心的紧缩,那不是在那个世界里;至少不是在那时。对人的尊重,是我在那里呼吸的空气的一部分。不必去挣得尊重,不必先证明自己才能被接受、被和气地对待。说来也许奇怪:只要是一个人,只要有一张人的面孔,就够了。

% ---------------- p.142 ----------------
\section{(10) “数学共同体”:虚构与现实}

\origpage{22} 因此毫不奇怪:也许就在那一年,在我内心深处,我已感到自己是这个世界的一员;而在随后的年月里,这种感觉越来越清楚。我喜欢用“数学共同体”这个名称来指它,这名字对我满载着意义。在写下这些字句之前,我从未有机会细想过自己给这名字赋予了什么意思,尽管我在很大程度上认同于这个“共同体”。现在清楚了:它对我而言,不外是那个善意接待我、把我当作他们中一员的世界的理想延伸——在空间上和时间上的延伸;而且是一个与我生命中最主要的激情之一相连的世界。

\origpage{23} 我逐步认同的这个“共同体”,并不是对最初接待我的那个数学圈子所作的完全虚构的外推。最初的圈子一点点扩大了;我的意思是:在最初接触之后的十年、二十年里,由于共同感兴趣的主题和人与人之间的投缘,我经常来往的数学家圈子不断扩大。具体说来,就是同事与朋友的圈子,或者不如说那个同心结构:从与我关系最近的同事(先是迪厄多内、施瓦茨、戈德芒,后来尤其是塞尔,再后来是安德烈奥蒂、朗、泰特、扎里斯基、广中平祐、芒福德、博特、迈克·阿廷这样的人,还不算布尔巴基小组的人——那个小组也在一点点扩大——以及六十年代起投到我这里来的学生们……),直到我在此处彼处有机会遇到、由或强或弱的投缘而或紧或松地联系着的其他同行——总之,正是这个由相遇与投缘偶然凑成的微观世界,构成了那个于我充满温度与回响的名字的具体内容:数学共同体。当我把它认同为一个活生生的、温暖的实体时,实际上我认同的就是这个微观世界。

直到1970年那个“大转折”之后——我该说,那第一次觉醒之后——我才意识到:这个舒适可亲的微观世界只占“数学世界”极小的一部分;而我乐于赋予那个世界、那个我一直不去了解、从没想到要去关心的世界的那些特征,都是虚构的特征。

况且在这二十二年间,这个微观世界本身也换了面孔,而它周围的世界也在变。我自己当然也变了——年复一年,不知不觉,我和我周围的世界一样变了。我不知道我的朋友和同事们是否比我更觉察到这变化——周遭世界里的、他们自己的微观世界里的、以及他们自身之中的变化。我也说不清这奇怪的变化是何时、如何发生的——它想必是悄悄来的,蹑手蹑脚:有名望的人是叫人怕的。我自己就叫人生畏——不是叫我的学生或朋友怕,也不是叫那些亲自认识我的人怕,至少是叫那些只通过名声认识我、而自己并未感到有相当的名声庇护的人怕。

我是在将近十五年前那次“觉醒”之后,才意识到数学世界里(在其他科学圈子里同样如此,如果不是更甚)盛行的恐惧的。在此前的十五年里,我逐步地、不知不觉地进入了“大老板”的角色,进入了数学界《名人录》的世界。也同样是不知不觉,我成了这个角色的囚徒,它把我与所有人隔开,只除了少数几位“同侪”和少数几个学生(而且还……)——那些确实“有心气”的人。只有当我走出这个角色,环绕着它的恐惧至少有一部分才落了下来。那些在我面前多年沉默的舌头,松开了。

它们带给我的见证,不只是恐惧的见证,也是轻
% ---------------- p.143 ----------------
\origpage{24} 蔑的见证。尤其是身居其位的人对别人的轻蔑——一种唤起恐惧、也滋养恐惧的轻蔑。

我几乎不曾体验过恐惧,却实实在在体验过轻蔑——那是在一个人的存在与生活分量很轻的年月。我本来乐于忘掉轻蔑的时代,可它偏偏来向我问好了!也许它从未停止过,只是我满足于换了个世界(在我看来如此),满足于把目光移开,或者干脆:除了那些激动人心、没完没了的数学讨论之外,装作什么也看不见、什么也听不见?那些日子里,我终于肯去了解:轻蔑就在我周围到处盛行,在我选为自己的世界、我所认同、我所背书、并且曾娇养我的那个世界里盛行。

上面这些话也许给人一个印象:我被那些几乎一夜之间涌向我的见证搅得心神不宁。其实不然。这些见证是在一个仍然浮浅的层面上被接收的。它们只是加在我刚得知的另一些事实上,或加在我早已知道、却一直避免去注意的那些事实上。今天,我会这样表述当时学到的教训:“科学家”们,从最显赫的到最无名的,都是些和所有别人一模一样的人!我曾乐于想象“我们”是更好的一类,我们身上多点什么——我足足花了一两年,才摆脱这个顽固透顶的幻觉!

\section{(11) 与克劳德·谢瓦莱相遇——或自由与善意}

在帮助我做到这一点的朋友中,只有一位属于那个我刚离开、不打算再回去的圈子\fnoteP{4}{(6)}{\emph{我在“生存与生活”的朋友们}\par 在这些朋友中,我大概还该算上皮埃尔·萨米埃尔(Pierre Samuel)——我先前主要是在布尔巴基里认识他的,与谢瓦莱一样;而他(和谢瓦莱一样)在“生存与生活”小组里起了重要作用。我不觉得萨米埃尔有那么迷恋那种关于科学家优越性的幻觉。我感到,他带来最多的,是他在共同工作、讨论和待人接物中投入的常识和那份笑呵呵的好脾气,还有他优雅地在一个倾向于激进分析与激进抉择的团体里扮演“讨厌的改良派”的角色。在我退出“生存与生活”之后,他又在那里待了一阵,充当同名公报的主编;当他感到自己的存在对这个团体已不再有用时,便体面地离开了(去了“地球之友”)。\par 萨米埃尔和我出身于同一个狭小的圈子,这并不妨碍他成为那些沸腾年月里的朋友之一——我相信我从他们那里学到了些东西(尽管我是个糟糕的学生……)。他的为人方式,和谢瓦莱的为人方式一样——尽管二人毫不相像——对我那些“功绩主义”的倾向来说,是比最锋利的分析更好的解药!\par 现在我觉得,对那个时期所有让我学到点东西的朋友,我学到的主要不是通过解释、讨论之类,而是通过他们各自的为人方式、他们那与我不同的敏感——而“某种东西”最终就这样传达了过来……说到这里,我特别记得的,除了谢瓦莱和萨米埃尔,还有德尼·盖吉(Denis Guedj,他对“生存与生活”小组有很大的影响力)、达尼埃尔·西博尼(Daniel Sibony,他一直与这个小组保持距离,同时从角落里用一半鄙夷、一半打趣的眼光注视着自己的演变)、戈登·爱德华兹(Gordon Edwards,他是1970年6月在蒙特利尔那场“运动”诞生的共同行动者,并且多年间付出了惊人的精力,维持着《生存与生活》公报的英文“美国版”)、让·德洛尔(Jean Delord,一位与我年纪相仿的物理学家,精细而热情,他喜欢我,也喜欢“生存”这个小圈子)、弗雷德·斯内尔(Fred Snell,另一位在美国布法罗的物理学家, 1972年我在那里逗留了几个月,曾在他乡间的家中做客)。\par 在所有这些朋友中,五位是数学家,两位是物理学家,而全都是科学家——这似乎表明,那些年里与我最接近的圈子仍然是科学家的圈子,尤其是数学家的圈子。} (6)。他就是克劳德·谢瓦莱。他不发议论,对我的议论也不感兴趣,然而我想我可以说:我从他那里学到的东西,比我刚才说的那些更重要、也更隐蔽。

在我相当经常地与他来往的那段日子(“生存”小组的日子,他带着打了折扣的信念加入了这个小组),他常让我不知所措。我说不清是怎么回事,但我感到他握有某种我抓不到的
% ---------------- p.144 ----------------
\origpage{25} 认识——对某些想必极其简单、极其本质之物的理解;这些东西当然可以用简单的话说出来,但理解并不会因此就从一个人“传”到另一个人那里。我现在明白,他和我之间有一种成熟程度上的差别,这使我常常在他面前觉得自己站得不稳,像在进行一场聋子的对话——而这并不是因为彼此缺少好感或敬意。

虽然他没有用这样的词说过(就我记忆所及),但在他看来想必很清楚:我当时——或独自、或凭“生存”小组(后来成了“生存与生活”)内部共同的思考与活动的逻辑——达到的那些“质疑”(关于“科学家的社会角色”、关于科学等等),骨子里仍然是浮浅的。它们当然涉及我生活的世界,甚至涉及我在其中扮演的角色——但并没有真正深刻地牵动我。在那沸腾的年月里,我看待自己的眼光一丝未变。我并不是在那时才开始认识自己的。只是六年之后,我才生平第一次摆脱一个顽固的幻觉——不是关于别人或周围世界的幻觉,而是关于我自己的。那是另一次觉醒,分量比预备了它的第一次更大。它是一连串相继“觉醒”的“瀑布”中最初几次之一;我希望,这瀑布在我余年里还会继续流下去。

我不记得谢瓦莱曾在什么场合提到过认识自己,或者更确切地说,“发现自己”。不过回过头看,很清楚:他准是很早以前就已开始认识自己了。他有时也谈起自己,只在谈这谈那时说上几个字,那份朴素叫人不知所措。他是我没听见说过一句套话的两三个人之一。他话不多,而他说的话所表达的,不是他采纳并据为己有的观念,而是对事物的亲身感知与理解。所以他常叫我不知所措——早在我们还在布尔巴基小组里碰面的时候就是如此。他所说的,往往冲撞了我所珍视的、并因此被我当作“真”的那些看法。他身上有一种我所缺乏的内在自主,我在“生存与生活”的日子里开始隐约察觉它。这种自主不属于智性、不属于话语的层面。它不是一件可以像观念、观点之类那样“采纳”的东西……幸好,我从未想过要把在别人身上察觉到的这份自主“据为己有”。我得找到我自己的自主。这同时意味着:我要学会(或重新学会)做我自己。但在那些年里,我对自己的不成熟、对自己缺乏内在自主,毫无觉察。我若最终发现了这一点,那么与谢瓦莱的相遇,必定是那些在我身上默默发酵的酵母之一——那时我正投身于一些大计划。播下这酵母的,不是议论,也不是言语。要播下它,只需我在路上偶然遇见的某个人不讲议论,只满足于做他自己。

我觉得,在七十年代初,当我们因《生存与生活》公报的出版而定期见面时,谢瓦莱曾不加坚持地试图向我传达一个消息,而我当时太笨手笨脚,抓不住它,或太埋头于我的活动任务。我隐约意识到,他有关于自由——关于内在自由——的东西要教我。我那时倾向于靠一些宏大的道德原则运转,而且从《生存》最初几期起就开始吹那把小号,觉得理所当然;而他对说教式的话语有一种特别的反感。我想,在《生存》创办之初,这是我最摸不透他的地方。在他看来,这种话语只是一种强制的尝试,叠加在窒息人的无数别的外在强制之上。当然,一个人尽可以花一辈子去讨论这样
% ---------------- p.145 ----------------
一种看法,讨论正反两面\origpage{26}。它完全冲撞了我的看法——那看法(可想而知)是由最高尚、最普遍的感情驱动的。我很难过:我无法理解,谢瓦莱——我最敬重的人,我多少把他当作战友的人——竟然会带着一种狡黠的乐趣,不去分享这些感情!

\origpage{27} 我不明白:真相、事物的实际,既不是善良感情的问题,也不是观点或偏好的问题。谢瓦莱看见了一件东西——那东西再简单、再真实不过——而我没有看见。并不是他在什么地方读到过它;看见一件东西,和读到关于它的东西,二者毫无共同之处。人勉强可以用手(盲文)或耳朵(若有人读给你听)去“读”一段文字,但要看见那件东西本身,只能用自己的一双眼睛。我不认为谢瓦莱的眼睛比我的好。但他用它们,我不用。我被我的善良感情之类的东西占得太满,没有闲工夫去看我的善良感情和原则对我自己、对别人——首先是对我自己的孩子们——起了什么效果。

他一定看得很清楚:我常常不用自己的眼睛,甚至一点也不想用。奇怪的是,他从未对我透露过这一点。还是他说了,只是我没听见?还是他忍住了,认为那是白费力气?又或者他压根儿没想到要说——说到底,我用不用自己的眼睛,是我的事,不是他的事!

\section{(12) 功绩与轻蔑}

我想借我自己有限的经验细看一下:轻蔑是何时、如何进入数学家世界的,尤其是在那个由同事、朋友和学生构成的“微观世界”里——它已成了我的第二祖国。同时我也想看看,在这场转变中,我这一份占了多少。

我想我可以毫无保留地说: 1948–49年间,在我前面谈过的那个数学家圈子里(对我来说,它的中心是最初的布尔巴基小组),我没有遇到一丝一毫的轻蔑,甚至没有一丝一毫的鄙夷或居高临下——无论对我自己,还是对任何一个来此学习数学家行当的年轻人,无论法国人还是外国人。在那里起旗舰作用的人物——凭地位或声望,如勒雷、嘉当、韦伊——我并不怕他们,我想我的同学里也没有人怕他们。勒雷和嘉当很有“体面先生”的派头,另当别论;至于其余那些随随便便闯进来、像哥们儿一样跟嘉当“你你我我”、显然又“在行”的家伙,我甚至花了相当一段时间才明白:他们每一个都和嘉当本人一样是大学教授,绝不像我那样挣一天吃一天,而是领着在我看来属于天文数字的薪水,而且还是个有国际影响、国际听众的数学家。

听从韦伊的建议,我在南锡度过了随后的三年。那里当时算是布尔巴基的一个总部:德尔萨特、迪厄多内、施瓦茨、戈德芒(稍后还有塞尔)都在大学里任教。跟我一起在那里的只有四五个年轻人(我记得有利翁斯、马尔格朗热、布吕阿、贝热,若有记错请原谅),所以远不像在巴黎那样“淹在人堆里”。气氛也因此更亲切,大家彼此都认识,我想彼此都以“你”相称。可我翻检记忆时,我亲眼看见一个数学家毫不掩饰地轻蔑对待一个学生的头一遭、也是唯一一遭,恰恰就在那里。那可怜人是当天从另一座城市赶来的,要跟他的老板一起工作。(他当时应该在准备博士论文,后来还体面地通过了,而且此后也
% ---------------- p.146 ----------------
有了相当的名声,我想。)\origpage{28} 那场面让我相当震惊。要是有人敢用这种腔调跟我说哪怕一秒钟的话,我立马就把他摔门轰出去!可这一回,我熟的只是那位“老板”,甚至跟他称兄道弟,而不是那个学生——我只是见过他。我这位年长的同行,除了广博的学养(不只是数学的)和锋利的头脑,还有一种斩钉截铁式的权威感,在当时(而且在之后相当长一段时间里,直到七十年代初)都令我敬畏。他对我有某种影响力。我不记得自己有没有就他的态度问过他什么;我只记得我从那场面里得出的结论:这个可怜学生一定是差劲透了,才该受这样的对待——诸如此类。我当时没有对自己说:这学生若是真的差劲,那该劝他去做别的,该停止跟他合作,但无论如何不该轻蔑地对待他。我认同了像这位显赫前辈那样的“数学强者”,而以那些“废物”为代价——似乎轻蔑他们是正当的。于是我走上了那条早已画好的路:与轻蔑默契,因为这对我方便,它还突出了这一点——我,被接纳入有功者、“数学强者”的兄弟会了!\fnoteP{5}{(7)}{上面这一段,是整篇导论里在我的原稿上涂改得最厉害、增补最多的一段。对那件事的描述、用词的选择,最初都是顶着来的,逆着潮流——显然有一股力量在推着人快点把那件事带过去,好像只为对良心有个交代,好“转到正经事上去”。这些都是一份抵抗的熟悉征兆——这里抵抗的是对这段插曲的澄清,以及对它作为某种内心态度之揭示者的意义的澄清。情形与这篇导论开头(第2段)所描述的十分相似:在数学工作中发现一处矛盾及其意义的那个“关键”时刻——那时起作用的是精神的惰性,是它不肯与一种错误或不充分的理解(但其中根本没有我们这个人牵涉在内)分离的厌恶。后者却是主动的,必要时还能别出心裁,做到无水也能把鱼淹死;而我刚才说的惰性只是一种被动的力量。在现在这个情形里,比在数学工作的情形里更甚:刚刚以其全部的简单、全部的自明显现出来的发现,紧接着就是一阵卸下重负的轻松感,一种解放感。这不止是一种感觉——不如说是对刚刚发生之事的一种敏锐而感激的知觉,那是一次解放。} (7)

当然,我不会比任何人更清楚地对自己说:那些试着搞数学却搞不成的人,活该被轻蔑!如果我在那个时候或任何时候听见有人说这种话,我会漂亮地驳他一顿,真心为如此惊人的精神无知而难过。事实是:我泡在暧昧里,同时在两张互不相通的牌桌上玩:一边是漂亮的原则与感情,另一边是:可怜的家伙,得差劲成什么样才落得这般对待(言下之意:这种事肯定不会落到我头上,那当然!)。

最后我觉得,我讲的那桩事件,尤其是我在其中扮演的那个(看来无伤大雅的)角色,其实是我身上一种暧昧的典型;这暧昧在此后二十年里跟了我整个数学家的生涯,直到1970年“觉醒”之后才消散\fnoteP{6}{(8)}{正如下文将变得清楚的,这种暧昧丝毫没有“在1970年觉醒之后消散”。这里有一个“自我”典型的战略撤退:它把“觉醒之前”那段时期丢给盈亏自负,而那段时期随即成了划分出一个无可指摘的“之后”的分界线!} (8)——而直到今天,写下这些字句的时候,我才清楚地察觉它。也真可惜,我当时竟没有察觉。也许对我来说时机尚未成熟。总之,那些当时传到我耳中的、关于轻蔑之统治的见证——对那统治,我选择了闭上眼睛——并没有把我个人牵涉进去,也没有把我那亲爱的微观世界里离我最近的那部分同行和朋友中的任何一位牵涉进去\fnoteP{7}{(9)}{这并不完全准确:在我最亲近的同行中至少有一个例外,后文会看到。这里有一种记忆特有的“懒”,它往往倾向于把那些与由来已久的、根深蒂固的看法“不合”的事实“随手放过”。} (9)。它们更像是这个调子:唉!竟然不得不%
% 原句在此页末尾中断, 下文接 part1-ch06b.tex (p.147)

% ================================================================
%  第 6 章 梦与做梦者 后半 (原书 p.147-152, 6.9-6.11)
%  源文: _work/src1/task-ch06b.md
%  注意:首段紧接 part1-ch06a.tex 末句 “唉!竟然不得不”, 两者之间不可插入空行
% ================================================================
% ---------------- p.147 ----------------
学到(或者:让人领教)这种事——谁会相信呢——真得是个混账(我差点要说:废物,对不起!),才会那样对待活生生的人!说到底,这跟另一副调子也差不了多少:只消把“废物”换成“混账”,把“让人那样对待”换成“那样对待”,戏法就变完了!而荣誉嘛,当然,毫发无损——对善事业的旗手来说!

\origpage{29} 从这里清楚显出的,是我与种种轻蔑态度的默契。它至少上溯到五十年代初,也就是我在嘉当和他朋友们那里受到善意接待之后的那些年份。后来,当轻蔑在四面八方变得司空见惯时,我之所以“什么也看不见”,是因为我不想看见——并不比那桩孤立而格外扎眼的事里更想看:那一回,真得下足功夫,才能装作既看不见、也觉不出!

这份默契与我新的身份紧密共生——一个受尊敬的群体成员的身份,有功者、数学强者的群体。我记得我特别满意,甚至骄傲:在我为自己选定、又把我增选进去的那个世界里,起作用的不是社会地位,甚至(才不是呢!)也不单是名声——名声还得是挣来的;你尽可以是大学教授、院士或随便什么,可你要只是个平庸的数学家(可怜的家伙!),那你就什么也不是;唯一算数的是功绩,是深刻而独到的想法,是技术上的娴熟,是广阔的视野,诸如此类!

这套功绩主义,我曾毫无保留地认同于它(而它当然始终隐而不宣、不曾明言),可它在我身上到底还是挨了狠狠一击——就在我说过的1970年那场著名觉醒之后的日子里。而且我也不能肯定,它从那一刻起就消失了、没留下痕迹。要做到那一点,大概得靠我清清楚楚地在自身中把它识别出来;可我当时似乎主要是在别人身上谴责它。说起来,最早提醒我注意这套意识形态的,正是谢瓦莱,还有德尼·盖吉(我也是通过“生存”认识他的)——他们管它叫“功绩主义”(méritocratie),或者类似的名字——以及它里面那种暴力、那种轻蔑。谢瓦莱对我说(那大概是我们头一次在他家里见面、谈“生存”的时候),正是因为这个,他才再也受不了布尔巴基的气氛,索性不再踏进去。回想起来,我确信他当时一定早已看出:我实实在在是这套意识形态的一分子,也许甚至看出它在我的某些角落里还留着痕迹。但我不记得他曾有过任何暗示。也许这一回,他仍然宁愿让我自己去把那些他替我勾出的 i 点上点——而我一直等到今天才把它们点上。晚做总比不做强!

\section{(13) 力与厚度}

\origpage{30} 我讲的那桩事,很可能也标志着我内心一次倒向的时刻:转向对功绩兄弟会或多或少无条件的认同,而以那些被视为废物的人为代价——或者,按早几代人的说法,只消说“没有天分”(这个词到我那时候已经不流行了):那些平庸、暗淡的人——顶多算是“共鸣箱”(韦伊在什么地方这样写过),供真正要紧的人的大想法在其中回响……我的记忆常常像个掘墓人,连当时耗掉大量心理能量的片断也一并埋掉;可它偏偏留下了这一段,而这一段跟任何别的记忆都没有直接关联,面目又那么平淡无奇——单凭这一点,那个“倒向”当时确曾发生的感受,就变得可信了。

不到五年前的一次沉思中,我终于意识到:那种“我们,伟大而高贵的精神……”的意识形态,以一种格外极端、格外凶猛的形式,从我母亲童年起就在她身上肆虐,并且主宰着她与他人的关系——她乐于从自己的高处俯视那些人,带着一种往往是鄙夷的、甚至轻蔑的怜悯。
% ---------------- p.148 ----------------
\origpage{31} 再说,我对我的父母怀着毫无保留的敬佩。在那个著名的“数学共同体”之前,我认同的第一个、也是唯一一个群体,就是缩小到母亲、父亲和我三个人的家庭——而我有幸被母亲认可为配得上他们做我的父母。这就是说,轻蔑的种子想必从我的童年起就已播进我这个人里了。也许时机已经成熟,可以追一追这些种子在我的童年和成年生活中经历的种种变故,以及其中一些种子长成的收成——幻觉的收成、孤立的收成、冲突的收成。但这里不是做这件事的地方,此处我追着一个更有限的目标。我想我可以说:这种轻蔑态度在我一生中,从未达到我在母亲一生中见到的那种激烈、那种毁灭性的力量(那是我肯费心去看我父母一生的时候看到的——在母亲去世二十二年之后,在父亲去世三十七年之后)。但现在正是时候,过了这村就没这店:在这里至少仔细考察一下,这种态度在我的数学生涯中占了什么位置。

在这之前,为了把上一段讲到的那桩事放进它的一般背景里,我想强调这样一个事实:在我的五十年代记忆里——甚至更晚的记忆里——它完全是孤立的。即便在今天,虽然我发觉在我曾经所属的那个圈子里,某些最基本的礼貌和对他人的尊重,有时正以一种叫人错愕的方式被侵蚀\fnoteP{8}{(10)}{比方说,我已数不清有多少封信——谈的既有数学问题,也有实际问题或私事——寄给我当作朋友的同事或旧学生,却从未收到回音。看来这并非只是专门留给我一个人的特殊待遇,而确实是一种风气变化的迹象——据一些同方向的反响来看。(诚然,那些反响涉及的情形是:寄数学信件的人,收信的那位知名数学家并不认识他……)} (10),老板对学生那种直接而不加掩饰的轻蔑,也还应当是相当罕见的事。至于五十年代,我很少有记忆指向那种情形:一个有名望的人物周围笼罩着恐惧,或者有谁摆出轻蔑、哪怕只是鄙夷的态度。要往这个方向翻找,我可以说:头一回被请到南锡迪厄多内家做客时——他对我一向有着那份体贴入微的和气——这个文雅亲切的人谈起他的学生来,那副样子真叫我有点发懵:简直全是些蠢货!给他们上课是件苦差,而他们显然什么也听不懂……1970年以后,我听到从阶梯教室那一侧传来的回声,才知道迪厄多内确实让学生们害怕。可是,尽管他以见解果决、并且常常以雷鸣般的直率把它们讲出来而闻名,我却从未见他有过伤人、羞辱人的举动——即使当着他不怎么看重的同事的面,即使在他那些传奇式的雷霆大怒之际;而那怒气,来得快,消得也一样快、一样轻易。

我不认同迪厄多内关于他学生的那些说法,可也没有同他的态度拉开距离——那态度被呈现为世上最理所当然的事,对一个热爱数学的人来说几乎是不言自明的。有这位年长同行满怀善意的权威作后盾,那种态度当时在我眼里,至少是人们在面对学生和教学任务时可以有理由采取的态度之一。

看来,迪厄多内和我一样,都浸透了这套功绩主义;但当我们面对一个血肉之躯时,这种主义的隔离效果在很大程度上就被抵消了:单是那个人的在场,就默默地向我们提醒一些比所谓“功绩”更本质的现实,并重新接上一根被遗忘的纽带。我们的大多数同事或朋友想必也一样——他们浸染那种如此普遍的优越感,并不比迪厄多内或我少。今天,他们中的许多人肯定还是这样。

% ---------------- p.149 ----------------
\origpage{32} 韦伊同样有让学生害怕的名声;而在五十年代我的微观世界里,他是唯一一个让我觉得连那些地位(或者干脆说性情)比较谦卑的同事也怕他的人。他时不时会摆出那种不容回嘴的居高临下,足以叫最牢靠的自信也发慌。我的敏感也在其中起了作用,于是有那么一两次,我们闹了阵短暂的别扭。我在他的做派里没有察觉出一丝轻蔑的意味,也没有察觉出存心伤人、压人的意图;那更像是被惯坏的孩子的作派,以叫人难堪为乐(有时还带点促狭),仿佛以此使自己确信他握有某种权力。再说,他对布尔巴基小组有一种真正惊人的支配力;有时他给我的印象是,他拿指挥棒指挥着它,有点像幼儿园女教师指挥一群乖孩子。

在五十年代,我只记得另外一次,感到轻蔑以一种粗暴、不加掩饰的方式表达出来。它来自一位外国同事兼朋友,年纪跟我差不多。他的数学力量非同寻常。几年前,这股力量已经相当明显了,而我当时却被他对那位大教授的顺从(在我看来近乎谄媚)所触动——那时他还只是那人的一名谦卑助手。他出众的才能很快为他赢得了国际声誉,以及一所格外显赫的大学里的一个关键职位。他在那里统治着一小支助手兼学生的队伍,那副绝对的派头,显然跟他当年的老板统治他和他的同伴们时一样。我问他(如果我记得不错)有没有几个学生(言下之意:跟着他做出好东西的学生),他带着一副假作随意的神气(我把它译成法语)答道:“十二件!”——所以“件”就是他用来指他的学生和助手的词。一个数学家同时有这么多学生在他指导下做研究,确实少见——而我的这位对话者想必暗自引以为傲,又想用那副漫不经心的神气把它藏起来,仿佛在说:“哦,不过十二件,连提都不值一提!”那大概是1959年前后。我当时想必已经有了一层不错的甲壳,可我还是感到一阵恶心!我准是当场就用某种方式对他说了,而且我不认为他为此记恨我。也许他与学生的关系并不像他那说法所能让人设想的那样阴森(我没有得到过他任何一位学生的证词),而他只是落进了自己那孩子气的欲望的陷阱:想在我面前炫耀他全部的荣光。回头看,我看得出这桩事必定是我们关系的一个转折点,而先前的那些关系是友谊的关系——我在他身上感到某种脆弱,还有一种细腻,它们吸引着我亲切的同情。这些品质被磨损了、被腐蚀了,被他那位重要人物的地位——受人敬佩、叫人害怕的地位。这桩事之后,我心里对他留下了一种不适——说到底,我觉得自己不属于他那个世界……

\origpage{33} 然而我们确实属于同一个世界——而且,我并不比他更觉察到:我也在变厚。关于这一点,我留下了一段鲜明的记忆,那是在1958年的爱丁堡国际数学家大会上。从上一年的黎曼–罗赫定理工作起,我就被晋升为大明星了;而且(当时我并不曾用明确的话对自己说过)我也是大会的明星之一。(我在会上作了一个报告,讲同年概形理论的强劲起步。)希尔策布鲁赫(当日的另一位明星,带着他自己的黎曼–罗赫定理)致开幕辞,为的是向那年即将退休的霍奇致敬。说到某处,希尔策布鲁赫暗示:数学主要靠年轻人的工作,而不是靠年长数学家的。这在年轻人占多数的大会厅里激起了一片轰然的赞同。我满心欢喜,当然也完全同意:我刚满三十岁,还满可以算年轻,而世界是属于我的!一兴奋,我准是大声叫喊,还在桌子上猛捶了好几下。碰巧我坐在霍奇夫人(Lady Hodge)旁边——她就是那位
% ---------------- p.150 ----------------
\origpage{34} 本该在这次场合受到致敬的卓越数学家的妻子,而他今年就要退休了。她睁大眼睛转向我,对我说了几句话,我已不记得了——但我准是从她惊讶的眼睛里看到了映出的那份没有分寸的厚度:它刚刚就在这位已近暮年的女士面前毫无节制地发作了一通。那时我感到了一点什么,而“羞耻”这个词给它的形象也许是走了样的——那更像是一条朴素的真相,关于当时的我这个人。那天,我大概再没有在桌子上猛捶过。

\section{(14) 恐惧的诞生}

我猜想,正是在那个时候——当我(并非刻意去求)开始被看作数学界的一个明星的时候——某种恐惧也必定开始围绕着我这个人,对许多我不认识或不太认识的同事来说。我是猜的:我无法用一个确切的记忆把它定位,无法用一幅曾击中我、并留在我记忆里的画面来定位,就像前面讲的那桩事那样(那想必标志着我头一回在接纳我的圈子里遇上轻蔑)。

这事想必是不知不觉地发生的,没有引起我的注意,也没有表现为某个特别而典型的、会被记忆保留下来的事件;而记忆打在它上面的那道光,也许跟打在那另一桩事上的一样,是刻意平淡的。我对那些过渡年月的记忆“整块地”还给我的,是这样一件事:前来接近我的人——无论是在我的讨论班之后,还是在布尔巴基讨论班之类的聚会上、在某个学术会议或大会上——往往得克服一种怯场,而在我们交谈(如果有交谈的话)的过程中,这种怯场多多少少还看得出来。当交谈超过几分钟,这份局促往往就在我们谈话、谈话热起来的过程中渐渐消失。有时——很少——也会发生这种情形:局促一直保持着,以至于成了交流的真正障碍,甚至在最不涉及个人的层面、在一场数学讨论的层面上;而我当时就模糊地感到对面有一种无力的痛苦,一种对自己恼火的痛苦。我说这一切时,并不是真的“记得”,而像隔着一层雾;可这层雾仍然把一些想必被记录过、又多半随即被排遣掉的印象还给了我。要我说出这份局促——这种恐惧的表现——何时出现,除了靠猜,我实在办不到。

我不认为这种恐惧源自我这个人,也不认为它只限于某种态度、某些使我区别于同事的举止。若真是那样,我想我在七十年代初就该听到些回声了——那时我走出一个自己一直扮演的角色,正是明星、“大老板”的角色。我相信,受人畏惧的是这个角色,而不是我这个人。而这个角色,连同它那圈与尊重毫无共同之处的恐惧光晕,在五十年代初并不存在,还没有存在——至少在我1948年一遇上就接纳了我的那个数学圈子里并不存在。

再说,在1970年那场“觉醒”之前,我根本不会想到把这种怯场、这种我有时在不那么亲近的同事身上遇到的局促叫作“恐惧”。它一显露,我自己就感到不自在,于是尽力去消除它。有一件事很值得注意,也典型地说明我亲爱的微观世界对这类事多么不关心:在我属于这个圈子的二十年里,我不记得有过一次,这个问题在我和某位同事之间被谈起,或者被别的人当着我的面谈起!\fnoteP{9}{(11)}{\emph{阿尔多·安德烈奥蒂(Aldo Andreotti)、约内尔·布库尔(Ionel Bucur)}\par{}当然,我这边也未必没有遗漏——何况我那时候那副十足“书呆子”的脾气,大概也不会鼓励别人跟我谈这类事,更不会促使我去记住一场很可能确实有过的、朝这个方向的谈话。可以肯定的是,至少可以说,这个问题被谈起的场合想必极其罕见(甚至不用“恐惧”这个名字去称呼它……),今天想必也一如其旧,尤其是在“上流社会”里。\par 在那个世界里的众多朋友中,除了谢瓦莱(他至少在六十年代期间必定意识到了这种恐惧的气氛),我觉得确实清楚察觉过它的,只有另外一个人,就是阿尔多·安德烈奥蒂。我认识他,以及他的妻子芭芭拉(Barbara)和他们的一对双胞胎孩子(当时还很小),是在1955年(我想是在芝加哥韦伊家的一次晚会上)。我们一直关系很密,直到1970年“大转折”的时候——那时我离开了曾属于我们的那个圈子,同他们也稍稍失去了联系。阿尔多极其敏感,而这份敏感在与数学、与像我这样的“书呆子”的往来中丝毫没有磨钝。他身上有一种对接近他的人自然生发的同情之才。这使他跟我所认识的数学圈内、乃至圈外的所有别的朋友都不一样。在他那里,友谊总是压过共同的数学兴趣(那些兴趣并不缺),而他也是极少数与我稍稍谈过我的一生、他也谈过他自己的一生的数学家之一。他的父亲和我的父亲一样是犹太人,他也因此在墨索里尼的意大利吃过苦,正如我在希特勒的德国吃过苦。我见他总是乐于鼓励和支持年轻研究者——那是在一个越来越难以被当权圈子接纳的气候里。他那种自然的关切,总先落到人身上,而不是落在什么数学“潜力”或名声上。他是我有幸遇到的最叫人亲近的人之一。\par 说到阿尔多,又勾起了对约内尔·布库尔的记忆——他也是出人意料地、未到年纪就被带走了;而且和阿尔多一样,我想,人们更多是把他当作一位乐于重聚的朋友来怀念,而不是当作数学讨论的伙伴。人能从他身上感到一种善良,旁边还有一份不多见的谦逊,一种不断退到后面的倾向。这样一个从不看重自己、也从不给人压力的人,最后竟当上了布加勒斯特大学理学院院长,真是个谜;也许是因为他压根儿没想过要推辞那些他远不渴求、却由同事们或政治当局放到他肩上的职务——那肩膀,说起来,是结实的。他是农民的儿子(在一个“阶级标准”很重要的国家里,这想必起了作用),身上带着农民的通达与朴素。他必定察觉到围绕着名人的那种恐惧,但想必也觉得这事理所当然,是权力地位的自然属性。不过,我不认为他本人曾让任何人畏惧——当然也不曾让他的妻子弗洛丽卡(Florica)、他们的女儿亚历山德拉(Alexandra),或他的同事、学生畏惧;而我能听到的反响,也正是这个方向。} (11)那层替我充作记忆的“雾”,也没有还给我
% ---------------- p.151 ----------------
\origpage{35} 任何印象,像是这类情形曾在我身上引起过有意识或无意识的满足。我不认为在意识层面有过;但我也不敢断言,在最初那些年里,我在无意识层面不曾偶尔被它拂过。若有,那也必定是短暂的,不曾反映到行为里、起一种把局促固定下来的作用。当然,这并不是说我演的那个角色里没有掺进我的自负!但即使我在这角色上不计本钱地投入,当时驱动我这份自我的,并不是想给“普通同事”留下深刻印象的野心,而是不断超越自己,以逼出我的“同侪”不断更新的敬意——而且,也许首先是那些在我还无从施展之前就信任我、把我当作他们中一员的年长者。看来,我面对自己成为对象的恐惧时所持的内心态度——尽力不去理会它,同时在它显露之处勉力把它消除——这态度可以看作整个六十年代我所属那个圈子(“微观世界”)里的典型态度。

在此后过去的十年到十五年里,局面又大大恶化了——至少从那个世界时不时传到我这里的迹象看是如此,从我得以近距离目击、甚至有时参与其中的情形看也是如此。不止一次,在我旧日的朋友或学生中那些与我曾最亲近的人那里,我撞见了轻蔑那种熟悉而无可辩驳的迹象;撞见了那种(表面上“无缘无故”的)使人泄气、羞辱人、压垮人的意志。一阵轻蔑之风起来了——我说不清是什么时候——吹拂着那个曾于我亲切的世界。它吹着,不管什么“功绩”或“过”,用它的气息烧焦卑微的天职,也烧焦最美好的激情。我旧日的伙伴们,各自带着“自己人”,受着坚固城墙的保护,安稳地(像我从前那样)坐在那层裹着各自的、毡子般悄然无声的恐惧里——他们中可有一个人感到这阵风?我确切地认识一个,只有一个,在我的旧友当中:他感到了,还跟我谈起过,只是没有叫出它的名字。还有另外一位,有一天也像是身不由己地察觉到了它,随即赶忙在翌日\fnoteP{10}{(12)}{“翌日”一词在此应按字面理解,不是隐喻。} (12)就把它忘掉。因为感到这阵风并承担它,对我的一位旧友来说,
% ---------------- p.152 ----------------
\origpage{36} 正如对我自己一样,也就是接受把目光落到自己身上。

\section{(15) 收获与播种}

我不打算——我再也不会打算——为一阵吹着的风而愤慨了,既然我已看清:我与这阵风并非无关,尽管我身上的一份自负一心要让我相信相反的事。而即便我与它无关,我的愤慨对受辱的人和施辱的人也都只会是一份太可怜的献礼——而这两种人,我一样爱过。

我与这阵风并非无关:在我为自己选定的那个世界里,我与轻蔑、与恐惧有过默契。对这些污迹,就像对许多别的污迹一样,闭上眼睛对我正方便——在我的职业生活中如此,在我的家庭生活中也一样。在这两者里,我收获了我所播种的——也收获了别人在我之前或与我一同播种的,既包括我的父母(以及我父母的父母……),也包括我那些“旧日的新朋友”。而今天,收获这些发了芽的播种的,还有我之外的别人:既有我的孩子们(以及我孩子的孩子),也有我今天的某个学生——他被我旧日的某个学生轻蔑地对待。

谈到播种与收获,我心里既没有苦涩,也没有认命,更没有自怜。因为我已懂得:即便在苦涩的收成里,也有一种实体的肉,只消我们肯吃,就能拿它来养自己。当这份实体被吃下,成了我们血肉的一部分,苦涩就消失了——它原来只不过是我们面对一份为我们准备的食物时的抵抗的征兆。

我也知道:没有哪一次收获不同时是对另一些收获的播种——而那些收获往往比前一次更苦。当那条看来没有尽头的链条——漫不经心的播种与苦涩的收获,代代相传、代代重演——出现在眼前时,我身上仍有某处会揪紧。但我不再像面对一种残酷而无从逃避的宿命那样被压垮、起而反抗,更不再像从前那样做它心甘情愿而盲目的囚徒。因为我知道,在我遇到的一切里都有一种滋养的实体,无论播种的是我的手还是别人的手——只消我去吃,去看它变成认识。对我的孩子们,对所有我爱过的人、以及此刻我爱着的人,情形也一样:当他们收获我在自负而漫不经心的年月里播下的东西,或者收获我直到今天偶尔还在播下的东西时,也是如此。

% ================================================================
%  第 7 章 双重面孔 (原书 p.153-166, 边码 p.37-p.56)
%  源文: _work/src1/task-ch07.md (= ch07.md)
% ================================================================
\chapter{双重面孔}

\section{(16) 沼泽与头几排}

% ---------------- p.153 ----------------
\origpage{37}可我还没有走到这场反思的尽头:反思的是,在轻蔑的出现与蔓延之中,我占了怎样一份——在这样一个世界里,我至今仍兴高采烈地以“数学共同体”之名来指称它。我现在感觉到,正是这场反思,才是我所能献给那个世界里我所爱过的人的最好的东西;此刻我正准备着——当然不是回到那个世界去,而是要在那里重新开口说话。

我想,我尤其还剩下这样一件事要做:审视一下,在我同他们一样仍是那个世界中一员的那些年月,我与这个那个成员之间,维持的是怎样一种关系。

如今回想起来,有一点让我吃惊:那个世界里有一整块人,我虽说常与他们擦肩而过,却始终没进入我的注意,仿佛他们并不存在。当年我大概把他们感知为某种“沼泽”——在我心里并没有明确的功能,我想,连“共鸣箱”的功能都谈不上;像一团灰色的、无名的人群:在讨论班和学术会议上,他们一成不变地坐在最后几排,仿佛生来就被分派在那儿;报告进行中,他们从不开口冒险提一个问题,因为他们事先就确信无疑:自己的问题只可能问得不着边际。倘若他们向像我这种被公认为“在行”的人提问,那准是在走廊里,在“权威们”显然无意彼此交谈的时候——这时他们便飞快地、像踮着脚尖似的把问题提出来,仿佛为占用我们这种大人物的宝贵时间而羞愧。有时那问题确实显得不着边际,我便(我想)试着用几句话说明为什么;常常它也是切中要害的,我同样尽力作答,我相信。两种情形下都一样:在这样一种心态里(或者我该说,在这样一种气氛里)提出的问题,很少会有第二个问题跟上,把它说得更确切、更深入。也许我们这些坐在头几排的人,在那种场合确实太匆忙了(尽管我们有时无疑也在努力不显出匆忙的样子),以至对面那份畏惧无法消散,交谈也就无从产生。我当然感觉得到,正如我的对话者在他那一侧也感觉得到:我们卷入的这局面有多么虚假、多么造作——只是我那时从未对自己把它说出来,而他也一样,想必从未说出来过。我们两个都像古怪的自动机那样运转,一种古怪的默契把我们连在一起:假装不知道那攫住我们之中一个人的焦虑,而另一个人则隐隐约约地觉察到了它——那是充满焦虑的空气里的一小份焦虑,这空气浸透了那些场所,所有人想必都像我们一样觉察得到,又都选择一致同意地视而不见\fnoteP{1}{(13)}{前面这段描述,显然只求勉强用具体的字眼,把记忆的那团“雾”交给我的东西还原出来;那团雾从未凝结成任何一个稍微确切的具体情形,好让我在这里给出稍微“写实”或“客观”一点的描述。若让这段文字说出这样的意思——那些不愿坐在头几排、或者没有明星与“权威”身份的同行,在同这样一位人物说话时必定满心焦虑——那就是曲解我的话了。在这个圈子里我认识的大多数朋友,显然都不是这样,甚至那些偶尔出入学术会议或讨论班的人也不是。毫无保留地成立的是:“权威”身份造成一道屏障、一道沟壑,把它与没有类似身份的人隔开;而这道沟壑极少消失,哪怕只是一场谈话的工夫。我还要补充:“头几排”与“沼泽”之间那种主观的区分(在我看来却实实在在),绝不能化归为社会学标准(社会地位、职位、头衔等等),甚至也不能化归为“身份”、名声的标准;它还映照出气质上、或者更难把握的禀性上的一些心理特点。我二十岁到巴黎时,就知道自己是个数学家,我做过数学;尽管有我曾有机会谈过的那些水土不服,我心底里觉得自己是“他们中的一员”——只是这一点只有我自己知道,而且我起初甚至并不确信自己还会继续做数学。如今我倒更倾向于坐到最后一排去(在难得出现这个问题的场合)。} (13)。

% ---------------- p.154 ----------------
\origpage{38}这种对焦虑的模糊感知,在我这里要等到1970年第一次“觉醒”之后才变成自觉的;那时,这片“沼泽”从我乐意让它一直待着的那片阴影里走了出来。这件事并不是出于什么有意的决定,我也没有立刻意识到,我就这样离开了一个圈子,进入另一个圈子——从“头几排”人们的圈子进入“沼泽”:突然间,我的新朋友大多正是那些一年前我还会不声不响地归入那片无名无状的疆域的人。所谓沼泽,忽然间动了起来,借着一张张朋友的面孔活了过来——这些朋友与我由一场共同的历险连在一起,另一场历险!

\section{(17) 特里·米尔基尔(Terry Mirkil)}

说真的,早在这场关键的转折之前,我就与几位同学(后来成了“同行”)结下了友谊,假如当年有人问我,我无疑会把他们归入“沼泽”(而且假如他们不是我的朋友……)。非得有这场反思、非得我去翻检记忆,我才想起他们来,那些零散的记忆才聚拢到一处。我认识这三位朋友,是在最初的那些日子里;那时我和他们一样,在南锡学着这门手艺——也就是说,那时我们还在同一条船上,还没有任何东西把我标举为一位“权威”。这多半不是偶然;而此后二十年里再没有过别的这样的友谊,多半也不是偶然。我们四个都是外国人,这肯定是一条不可忽视的纽带——我与那些和我一样被“空投”到南锡的年轻“高师人”的关系,就远没有那么私人化;我们几乎只在系里见面。三位朋友中有一位,一两年后移居南美。他和我一样是法国国家科学研究中心(CNRS)的研究专员,而我总有这么个印象:他自己也不太清楚他在“找”什么;久而久之,他在 CNRS 的处境变得有点尴尬。我们继续偶尔见面或通信,最后失去了联系。我与另外两位朋友的关系持续得更久,也更强烈,远不那么浮浅。何况数学上的共同兴趣在其中只扮演了极黯淡的角色,甚至等于零。

% ---------------- p.155 ----------------
\origpage{39}\origpage{40}与特里·米尔基尔(Terry Mirkil)和他的妻子普雷索西亚(Presocia)——她娇小纤弱,正如他粗壮结实;两个人都带着一股温柔的神气——我们在南锡常常一晚上、有时一整夜地唱歌、弹钢琴(那时是特里弹),谈音乐(那是他们的挚爱),也谈我们各自生活里这样那样重要的事情。老实说,不是最重要的那些——不是那些总被如此小心地闭口不谈的事。然而这份友谊给了我很多。特里有一份我所缺乏的细腻与辨别力,而我大部分精力早已倾注在数学上了。远比我更多地,他保留着对单纯而本质之物的感觉——太阳、雨、土地、风、歌唱、友谊……

特里在达特茅斯学院(Dartmouth College)找到一个合他心意的职位之后——那里离哈佛不算太远,从五十年代末起我常去那儿小住——我们仍继续见面、通信。其间我得知他容易陷入抑郁,这让他一次次长时间住进“疯人院”——在他唯一一次、也是极简略地向我提到此事的信里,他这样称呼那些地方;那是在又一次“可怕的长住”之后。我们见面时,从不谈这件事——只有一两次极偶然地,为了回答我的惊讶:他和普雷索西亚为什么不收养一个孩子。我不认为我脑子里曾冒出过这样的念头:他和我,我们可以谈谈问题的底;或者哪怕只是碰一碰——恐怕连“也许有些问题值得看看”这个念头也不曾有过,无论在我朋友的生活里,还是在我自己的生活里……在这些事情上,有一个不曾说出口、也无法逾越的禁忌。

渐渐地,见面和通信稀疏了。确实,我越来越成了工作和一个角色的囚徒,尤其是成了那份意志的囚徒——它变成了一种固执的念头,也许是对别的东西的一种逃避:在作品的累积中不断地超越自己——而与此同时,我的家庭生活正神秘地、不可挽回地败坏下去……

有一天,我从特里在达特茅斯的一位同行来信中得知,我的朋友自杀了(那已是他死去并被埋葬很久之后……);这个消息像穿过一层雾传到我这里,像来自一个极其遥远、而我不知何时已经离开了的世界的回声。也许那是我心里的一个世界,它早在特里结束自己生命之前就已死去——他的生命被一种焦虑的暴力摧毁,那焦虑他没能、或者不愿去化解,而我也没有能够、或者不愿意去猜出……

\section{(18) 二十年的自命不凡,或:不知疲倦的朋友}

% ---------------- p.156 ----------------
我与特里的关系,我想,在任何时刻都不曾因我们在数学界地位的差异、或因我由此而生的优越感而变味。这份友谊,以及那些年月里生活馈赠给我的另外一两位朋友(全然不管我“配不配”),肯定是当时少有的解药之一,用来对付一种隐秘的自命不凡——它由社会地位滋养,更由我对自己数学力量的自觉以及我自己给予它的价值滋养。我与第三位朋友的关系就不同了。这位朋友,以及后来他的妻子(他是在我们于南锡相识前后认识她的),在这些年里,每一次我们相遇,无论在他们家还是在我家,都对我表现出一种温暖的友谊,满含体贴与质朴。这份友谊里显然没有任何与地位或脑力相关的潜台词。然而,我与他们的关系,二十多年里始终带着我内心那种深刻的暧昧,带着我说过的那道分裂——它给我的数学家生活打上了印记。在他们面前,每一次我都忍不住感到他们亲切的友谊并作出回应,几乎是违心的!与此同时,二十多年里我居然做到了这样一手绝活:从我的伟大高处,轻蔑地俯视我的朋友。这想必从在南锡的头几年就这样开了头;而且很长时间里,我的成见也扩展到了他妻子身上,仿佛\origpage{41}\origpage{42}事先就可以认定:他的妻子只可能同他一样“无足轻重”。母亲和我之间,我们喜欢只用一个嘲弄的绰号来指他;这绰号想必在我母亲1957年去世之后很久还刻在我心里。如今我看出,至少在我这态度背后的诸般力量之一,是我母亲那坚强的个性在她一生中、乃至她死后近二十年里对我的支配——那些年里,我继续浸染在支配了她一生的那些价值之中。我朋友那温顺、和蔼、毫无争强好胜之心的天性,被不声不响地归入“无足轻重”一类,成了嘲讽式轻蔑的对象。直到现在,我第一次费心去审视这段关系究竟是什么样子,才发现这种面对他人温暖善意而执意自我隔绝的状态,范围有多大——它曾如此长久地给这段关系打上印记。我的朋友特里,并不比另一位朋友更好斗、更犀利,却幸好得到了我母亲的认可,不曾成为她嘲笑的对象——我猜,正因如此,我与特里的关系才得以在我心里毫无阻力地绽放。他对数学的投入并不更热忱,“天赋”也并不更出众,可我并没有因此就找到借口,用那层轻蔑与自得的硬壳把自己同他和他的妻子隔开!

在这另一段关系里,我至今仍无法理解的,是我朋友那份亲切的友谊,何以从未因他每次见面都必定感觉到的我的迟疑而灰心。然而今天我很清楚:除了那层硬壳和那份轻蔑,我还是别的东西;除了一块脑肌肉和由它生出虚荣的自命不凡,我还是别的东西。和他们一样,我心里也有那个孩子——那个我假装不知道的孩子,轻蔑的对象。我把自己与他隔断了,可他仍在我心里的某处活着,健康而茁壮,像我出生那天一样。我朋友们的感情,肯定是给那个孩子的——他们不像我那样与自己的根隔断。而肯定也正是他,在大老板背过身去的时候,偷偷地、匆匆地回应了那份感情。

\section{(19) 没有爱的世界}

% ---------------- p.157 ----------------
大老板老了,幸好,他也稍稍风化了一些,那小子从此可以自在些了。至于同这些真正经得起考验的朋友的关系,我想我在这里确实指到了我一生中最触目、最怪诞的一例:某种自命不凡(除了别的之外)在一段私人关系里所起的效用。也许我还在自欺,但我相信这也是唯一的一例:我与数学界(甚至别处)某位同行或朋友的关系,被自命不凡持久地占据,而不是让它满足于偶尔地、隐约而短暂地露一露头。此外我似乎觉得,我当时在数学界乐于来往的众多朋友当中,没有一个我能想象他也经历过类似的事——在一段与同行(朋友也罢,不是也罢)的关系里。在我所有的朋友当中,我也许是最不“酷”的,最“死用功”的,最不肯让一丝幽默透出来的(这本事我晚年才终于有了),最容易把自己看得无比严肃的。肯定也正是因此,我才不会太去寻求与我同类的人作伴(假定真有这样的人)!

奇怪的是,我的朋友们,不管属于“沼泽”与否,都容忍我,甚至喜欢我。这是一件好事,也是一件要紧的事,值得在这里说出来——尽管我们常常一连几小时、几天只见面讨论数学:感情在流动,就像今天它仍在流动一样,在当时的友人(凭一些有时纯属偶然的意气相投)与我之间流动;从1949年在南锡我被亲切接纳的那一刻起——在洛朗·施瓦茨(Laurent Schwartz)与伊莱娜·施瓦茨(Hélène Schwartz)家里(我在那儿有点像家里的一口人),在让·迪厄多内(Jean Dieudonné)家里,在罗杰·戈德芒(Roger Godement)家里(有一阵我也常去\origpage{43})。

这种包裹着我在数学界最初几步的温暖感情,我有点倾向于把它忘掉,它对我整个数学生涯都很重要。肯定是它,让我同那个以长辈们为化身的圈子的关系,也带上了同样温暖的色调;它给了我对这个圈子的认同以全部力量,也给了“数学共同体”这个名号以全部意义。

显然,对今天许多年轻数学家来说,事情是这样:在他们学习的年月、往往还远不止这些年月里,他们被切断了与一切感情、温暖之流的联系;他们的工作只映在一位疏远的老板眼中,映在他吝啬的评语里,有点像在读科研与工业部的一份通令——那通令剪掉了工作的翅膀,剥夺了它比一份阴郁而不牢靠的糊口营生更深的意义。

可我说着说着就提前了——说到七十、八十年代数学界那种失宠,也许是所有失宠中最深的一种;那个数学界里,发号施令的是我昔日的学生,以及我旧日朋友的学生。在那个世界里,老板常常把工作题目分派给学生,像扔给狗一根骨头——就这个,别的没有!像给囚犯分派一间牢房:你就在那儿服你的孤独刑!在那个世界里,某项细致而扎实的工作、多年耐心努力的结果,被那位无所不知、手握权力的人微笑着轻蔑地否掉:“这工作不合我胃口!”事情就这样定了。扔进垃圾桶吧,别再提了……

这样的失宠,我很清楚,在我所认识的那个圈子里、在我五十、六十年代常来往的朋友中间,并不存在。确实, 1970年我得知:在数学之外的科学界,这倒更像是家常便饭——甚至在数学里,显然也不那么罕见:公开露骨的轻蔑,明目张胆的滥用权力(而且无处申诉),甚至在某些有名望、我也曾有机会遇见的同行身上。但在我天真地当作“那个”数学界、或者至少当作这个世界忠实缩影的那个朋友圈子里,我没有见过这样的事。

然而,轻蔑的种子想必已经在那里了,由我的朋友和我播下,在我们的学生身上发了芽。而且不仅在我们的学生身上,也在我某些旧日的伙伴和朋友身上。但我的角色不是揭发,甚至不是斗争:腐败是不能靠斗争解决的。在我爱过的某个学生身上、或在某个旧日伙伴身上看到它,我心里有什么东西会揪紧——而往往,我不是去接受痛苦带给我的认识,而是拒绝这痛苦,挣扎,躲进拒绝和一种斗争姿态里:这种事不该存在!可它偏偏存在——而且我心底知道它的意义何在。从不止一个意义上说,我与此并非无关:我所爱的某个学生或旧日伙伴,乐于不声不响地踩一踩另一个我所爱、并且他在其中认出我的人。

我又跑题了,可以说跑了两重题——仿佛轻蔑的风只绕着我的住处吹!然而,正是它吹在我身上、吹在我亲近和珍爱的人身上的那口气,使我被它触及,认识了它。但现在还不是谈这个的时候,除非只对我自己,在沉默里谈。现在该是我重新拾起这场反思—见证的线索的时候了;它不妨就叫“追踪轻蔑”——追踪我心中和我周围的轻蔑,在我曾属于的那个数学界里,在五十、六十年代。

\section{(20) 一个没有冲突的世界?}

% ---------------- p.158 ----------------
\origpage{44}\origpage{45}我原打算用几行字谈谈“沼泽”,算是尽一份心,只说它在那儿,而我不常去——就像沉思中(数学工作里也如此)常有的事:你注视的那个“什么也没有”,结果却满是生命、奥秘,以及一向被忽略的认识。就像另一个“什么也没有”,它也像碰巧似的就在南锡(看来南锡真是我这个新身份的摇篮!):那个多半有点不中用的学生的“什么也没有”——人家是怎么对待他的,那可真该瞧瞧……刚才我还闪过一下这个念头,就在我写下(也许写得太快了点?)“这些失宠”在我们这儿“还不存在”的时候。这么说吧:我能报告的这类事,只有这独一无二的一桩;它像(这总得承认)我暗指的那种“失宠”,而我不愿过多逗留在细节描述上。受过这种对待的人,用不着我画图,就很清楚我说的是什么。那些没受过、但每次遇上都不急于闭上眼睛的人,也一样清楚。至于其余的人,那些穷开心地轻蔑别人的人,以及那些满足于闭上眼睛的人(就像我自己二十年间成功地做到的那样),哪怕给他们画一本画册也是白费……

我还剩下要做的是:审视这两个十年里我与同行、与学生之间的私人与职业关系,顺便也审视一下我所知道的、我最亲近的同行彼此之间以及他们与学生之间的关系。今天最让我吃惊的是:在所有这些关系里,冲突似乎根本不存在。我得马上补充:这在当时我觉得再自然不过——差不多是最起码的事。冲突,在善良意愿的人之间、在智性与精神上都已成年的人之间、以及诸如此类的人之间(再说一遍:最起码的事!),本就不该有。哪里真有了冲突,我便把它看作一种令人遗憾的误会:只要有应有的善意,把话说清楚,那准能在最短时间内了结,不留痕迹!我年轻时之所以选择数学作我最喜爱的活动,肯定是因为我感到:在这条路上,这种世界观最有可能不至于每每一头撞上令人不安的反证。毕竟,当人们证明了一件事,所有人就一致了——也就是说,善良意愿的人以及诸如此类的人,这自然不在话下。

碰巧,我确实感觉对了。这两个十年,是在我亲爱的“数学共同体”那个“没有冲突”(?)的世界的安宁里度过的;这段历史,也是我内心长期停滞的历史:眼睛和耳朵都塞住了,除了数学几乎什么也没学到——而在我的私人生活里(先是我和母亲之间的关系,然后是我在她死后随即建立的家庭),一场无声的毁灭正猖獗横行,那些年里我任何时候都不敢去看它一眼。可那是另一个故事了……我在这几行里常提到的1970年“觉醒”,不仅是我数学生涯的转折、一次圈子的彻底更换,也是(差一年)我家庭生活的转折。也正是那一年,我第一次在与新朋友的接触中,冒险偶尔——还非常偷偷地——瞥一眼我生活里的冲突。正是那个时候,一个疑问开始在我心里冒出,并在随后几年里成熟:我生活里的冲突,以及我有时在别人生活里感到的那种冲突,并不只是一场误会、一道用海绵一抹就掉的“污痕”。

这种(至少相对的)冲突之缺席,在我选为己有的那个圈子里,事后看来是相当了不起的事——而我后来终于知道:凡有人类生活的地方,冲突都肆虐横行,在家庭里如此,在工作场所也如此,无论是工厂、实验室,还是教授或助教的办公室。简直可以说, 1948年9月或10月,我懵懵懂懂地在巴黎下船,正好一头栽进了宇宙间唯一那座天堂般的小岛,那里的人们彼此相处没有冲突!

% ---------------- p.159 ----------------
\origpage{46}一下子,我觉得这事真了不起——在1970年以来我所知道的一切之后。它当然值得仔细看看——这是神话,还是现实?我看得见流动在我众多朋友与我之间、后来又流动在学生与我之间的那份感情,我用不着去编造它——可在这个仿佛把冲突放逐了的天堂世界里,我几乎不得不去编造冲突!

确实,在这场反思里,我毕竟有机会碰到过两种冲突的情境,每一次都像是我内心某种态度的显影剂:一个是南锡“不中用的学生”那桩事,其中当事人之间的是非曲直我并不清楚。另一个是我自己内心的冲突情境,一道分裂,在我与“不知疲倦的朋友”的关系里——但这一桩从未以人与人之间冲突的形式表现出来,而那本是冲突唯一被普遍承认的形式。值得注意的是,按这个词的通常意义,这些朋友与我之间的关系完全没有任何冲突——它任何时候都不曾有过一丝阴云。分裂在我心里,不在他们那里。

我继续清点。最先想到的之一:布尔巴基团体!在我或多或少定期参加它的那些年里——也就是直到五十年代末——这个团体对我而言体现了一种集体工作的理想:既尊重这项工作中看似极细微的细节,也尊重每个成员的自由。在布尔巴基团体的朋友们中间,我任何时候都不曾感到一丝强制的气味,无论是对我,还是对任何别人,无论资深成员还是受邀而来、试试看他和团体之间是否“来电”的客人。任何时候都不曾有过一丝影响力之争,无论是围绕议程上某个问题上的观点分歧,还是围绕对团体的霸权。这个团体没有首领而运转,而且据我所能察觉,显然没有人内心深处渴望扮演这个角色。当然,如同一切团体,某个成员对团体、或对其他某些成员的影响力大于别人。韦伊在这件事上扮演了一个特殊的角色,我谈过。他在场时,他有点像“主持人”\fnoteP{2}{(14)}{人们会以为这与“没有首领”的说法相矛盾,其实不然。对布尔巴基的老成员来说,我想,韦伊被看作这个团体的灵魂,但从来不是“首领”。他在场而且他乐意的时候,就如我说过的,他成了“主持人”,但他并不发号施令。他心情不好的时候,可以把他厌恶的某个题目上的讨论卡住;大不了等下一次大会韦伊不在时,再平心静气地把那题目捡起来,甚至第二天他不再作梗时就捡起来。决定由在场的成员一致作出,同时考虑到:一个人与所有其他人意见相左而站在正确一边,绝非不可能(甚至也不罕见)。这条原则对集体工作来说,似乎荒谬。了不起的是:它居然行得通!} (14)。有两次,我相信,我的敏感为此不快,我就走掉了——这是我所知道的、我自己造成的唯一“冲突”迹象。渐渐地,塞尔对团体的影响力变得与韦伊相当。在我还是布尔巴基成员的时候,这并没有造成两人之间的对抗局面,我也没有听说他们后来结下了仇怨。又过了二十五年再回头看,我所认识的五十年代的布尔巴基,在我看来始终是一个范例:在一个围绕共同计划形成的团体里,关系质量上取得了不起的成功。这个团体的品质,在我看来,比它产出的那些书的品质还要罕见。我这一生满是特权,其中之一便是遇到了布尔巴基,并在里面待了几年。我之所以没有留下来,绝不是因为冲突,也不是因为我所说的那种品质败坏了,而是因为一些更私人的工作更强烈地吸引着我,我把全部精力都献给了它们。何况这次离开,既没有给我与团体的关系投下阴影,也没有给我与其中任何一位成员的关系投下阴影。

我该逐一回顾一下我所卷入的冲突情境,那些使我与某位同行或某个学生对立起来的情境——1948到1970年之间。
% ---------------- p.160 ----------------
\origpage{47}\origpage{48}唯一稍稍值得一提的,是与韦伊那两次短暂的闹翻,前面已经说过。还有几片短暂的、极其短暂的阴云,落在我与塞尔的关系上,起因是我的敏感:有时一场谈话不再叫他感兴趣,他便有一种令人有时不知所措的洒脱,一下子把它掐断;或者对我投身其中的某项工作、我坚持的某种看法直说他没兴趣,甚至厌恶——也许我坚持得太多、太频繁了!这从未闹到翻脸的地步。除了气质上的差异,我们的数学意气相投格外强,他必定和我一样感到:我们彼此互补。

与我有着相当甚至更强意气相投的另一位数学家,只有德利涅。说到这儿,我想起1969年德利涅被任命到法国高等科学研究所(IHÉS)一事曾引起紧张,当时我并不觉得那是“冲突”(按说冲突会表现为翻脸,或同行关系中的一个转折)。

我想我已经清点完了——就人与人之间的冲突而言,就可由可见迹象察觉的冲突而言,在我常出入的那个圈子里,同行之间、或同行与学生之间的关系中,二十二年间就只有这些,信不信由你。也就是说,在我选中的那个天堂里没有冲突——那么,想必也就没有轻蔑?数学里又一个矛盾?

看来,我非得再仔细看看不可!

\section{(21) 一个守得很好的公开秘密}

昨天我肯定漏掉了一些小事,比如我与某位同行的关系里那些短暂的“冷淡”,尤其因我的敏感而起。我还该补上三四次我的自尊受挫的场合:有时同行和朋友在他们某篇文章里没有提到,我向他们说过的某个想法或结果想必在他们的工作中起过作用(在我看来如此)。我至今还记得,这说明那是个敏感点,也许并没有随年龄完全消失!除了一次,我都没有向当事人提起,他们的诚意当然不容怀疑。相反的情形肯定也发生过,只是我没听到回声。在我的“小宇宙”里,我不曾得知有哪一桩优先权问题成了翻脸或仇怨的由头,甚至不曾有过当事人之间尖酸刻薄的话。不过,唯一一次我进行了这样的讨论(在我看来那是个明摆着的例子),有过一场口角,倒把气氛澄清了,没有留下怨恨的残渣。那是一位格外出色的同行,他的本领之一,是能以惊人的速度吸收他所听到的一切;我似乎觉得,他常常有一种糟糕的倾向:把刚从别人嘴里听来的想法当成自己的。

这里有一个难题,它在所有数学家身上(而且不止在他们身上)都会以或强或弱的形式出现;它也不只是由于那种自我中心的习气——它促使我们大多数人(我也不例外)去归给自己一些“功劳”,无论真的还是想当然的。对一种情境(数学的或别的)的理解,无论我们怎样达到它,有别人帮助与否,本身都是一件本质上属个人的事,一种个人的经验,其果实是一种眼光,那也必然是个人的。一种眼光有时可以传达,但被传达的眼光不同于原初的眼光。话虽如此,要想辨明他人在自己眼光的形成中占的那一份,却需要极大的警觉。我自己肯定并不总有这份警觉——那是我最不关心的事,然而我却期待别人对我有这份警觉\fnoteP{3}{}{(9月30日)不过,事情的另一个侧面,见6月1日的注(晚于本文三个月),“暧昧”(n° 63),那里考察了对自己和他人有所宽纵的种种陷阱。}!
% ---------------- p.161 ----------------
\origpage{49}迈克·阿廷(Mike Artin)是第一个、也是唯一一个让我听到这件事的人。那天他带着一副泄露公开秘密的玩笑神气对我说:当你终于一把抓住一种实质、并在其中弄懂了点什么的时候,还费劲去分辨哪一份“属于自己”、哪一份“属于别人”,这既不可能,也完全徒劳。这话让我有点不知所措,因为这完全不合乎嘉当、迪厄多内、施瓦茨等人以身教给我的职业伦理。然而我隐隐感到,他的话里、以及他那笑盈盈的目光里,有一个我直到那个年纪都没能领会的真理:我的数学(尤其是数学产出)被自我强烈地占据着,而在迈克身上并非如此。他真给人一种印象:做数学就像一个玩得开心的孩子,同时又不忘吃喝。

\section{(22) 布尔巴基,或我的大运气——及其反面}

在往可见表面之下再潜一点之前,有一个判断现在就已摆在我面前:我五十、六十年代常出入的那个数学圈子,确实是一个“没有冲突的世界”,一点不假!这本身就是相当了不起的事,值得我稍作停留。

我得马上说明:这是一个非常狭小的圈子,是我数学小宇宙的核心部分,局限于我近身的环境中——我定期见面、关系最紧密的那二十来位同行和朋友。逐一检视他们时,我吃惊地发现:其中半数以上是布尔巴基的活跃成员。显然,这个小宇宙的核心与灵魂是布尔巴基。它差不多就是布尔巴基,加上与布尔巴基最亲近的那些数学家。六十年代我自己已不再是这个团体的成员,但我与某些成员的关系仍一如既往地紧密,尤其是与迪厄多内、塞尔、泰特(John Tate)、朗(Serge Lang)、卡蒂耶(Pierre Cartier)。此外我仍常去布尔巴基讨论班,或者不如说,正是那时我才成了常客;也正是在那时,我在那里做了我的大部分报告(关于概形理论)。

大概正是在六十年代,布尔巴基团体里的“调子”滑向了越来越明显的精英主义;我当时肯定是其中一分子,因此也就不可能察觉。我还记得1970年的惊讶:我发现布尔巴基这个名字本身,在数学界广大的阶层(此前为我所不知的阶层)里变得多么不受欢迎——它或多或少成了这些词的同义语:精英主义、狭隘的教条主义、对“典范”形式的崇拜而牺牲活生生的理解、晦涩、阉割式的反自发性,等等等等!而且布尔巴基不只是在“沼泽”里名声不好:六十年代,也许更早,我就偶尔从另一些心性不同、对“布尔巴基风格”过敏的数学家那里听到过回声\fnoteP{4}{(15)}{我并没有得到这样的印象:这种对布尔巴基“风格”的“过敏”,曾给这些数学家与我、或与布尔巴基其他成员及同情者之间的交流造成困难——假如这个团体的精神是宗派精神、是精英中的精英的精神,那倒会如此。在风格与时尚之外,这个团体的所有成员对数学实质都有敏锐的感觉,无论那实质来自何处。只是到了六十年代,我才记得我某位朋友把某些数学家称为“讨厌鬼”,因为他对他们的工作不感兴趣。这类事我几乎一无所知,便倾向于把这样的评价当作真话照收,被那份洒脱的自信镇住了——直到有一天我发现,那位“讨厌鬼”是个独到而深刻的心灵,只是没有福气讨我那位才华横溢的朋友喜欢。我似乎觉得,在布尔巴基某些成员身上,面对别人的工作(当自己不懂或懂得不透时)所持的那种谦逊(至少是保留)态度,已经磨损了。起初,那时那种“数学本能”还残留着——它使人感到一种丰富的实质、一项扎实的工作,而无须参照什么名声或名望。从各处传到我这里的回声看,我似乎觉得,谦逊也好,本能也好,今天在我曾属于的那个数学圈子里都已成了稀罕之物。} (15)。作为一个无条件的信徒,我对此感到惊讶, 
% ---------------- p.162 ----------------
\origpage{50}\origpage{51}也有点难过——我原以为数学能使人心一致!然而我本该记得:我起步的时候,硬吞一篇布尔巴基的文字并不总是轻松、也不总是启发人的,尽管它很干脆。那典范的文字,至少可以这么说,几乎不让人感到它写作时的气氛。如今我似乎觉得,这恰恰是布尔巴基文字的主要缺陷——连偶尔的一丝微笑也不能让人猜到:这些文字是人写的,而且是被远比某种对无情严格典范的无条件效忠更重要的东西连在一起的人写的。

不过,滑向精英主义的问题,以及布尔巴基写作风格的问题,在这里都是题外话。在这里让我吃惊的是:我选作职业圈子的这个“布尔巴基小宇宙”,是一个没有冲突的世界。这事在我看来更了不起,因为这个圈子里的主角们各有强大的数学个性,其中不少被视为“大数学家”,他们每一个都当然有分量形成自己的小宇宙,并且在其中当无可争议的中心和首领!\fnoteP{5}{(16)}{说实话,布尔巴基的成员里有好几位,肯定各有自己那个或大或小的“小世界”,在布尔巴基的小世界之外或之上。但在我自己的情形里,这样一个世界只是在我不再是布尔巴基的一员之后,在我全部精力都投入到属于我个人的工作之后,才在我周围形成——这也许并非偶然。} (16)这些强个性的人,两个十年里在同一个微小宇宙、同一个工作团体里和睦甚至亲切地共处——这才是让我觉得如此了不起、也许是独一无二的事。这与昨天谈布尔巴基时已经浮现的“非凡成功”的印象相合。

看来归根到底,我是有了这份非凡的运气:初次接触数学界时,我恰好落在时空里那个优越的地点——那里几年前刚形成了一个品质非凡的数学圈子,也许就品质而言独一无二。这个圈子成了我的圈子,并始终是我心目中理想“数学共同体”的化身;而这样一个共同体,当时(在我这个化身之外)大概并不比数学史上任何别的时刻更存在,也许除了某些同样狭小的团体(比如也许那个围绕毕达哥拉斯形成、精神却完全不同的团体)。

我对这个圈子的认同非常强,与我四十年代末诞生的数学家新身份不可分割。这是我超出家庭之外第一个被温暖接纳、被当作自己人接受的团体。还有另一条纽带,性质不同:我自己对数学的进路,在这个团体的进路里、在我新圈子里成员的进路里得到了确认。它并不等同于“布尔巴基式”的进路,但显然两者是兄弟。

此外,这个圈子对我而言还该代表那个理想的地点(或差不多是!),那个没有冲突的地点——对它的追寻无疑曾把我引向数学,引向一切科学中最不让我感到有一丝冲突意图的科学!而我刚才说到我的“非凡运气”,心里也明白这份运气有它的反面。如果说它让我得以发展出种种手段,在长辈已成同辈的圈子里以数学家的身份充分施展,那么它同样是一种受欢迎的手段,用来逃避我自己生活里的冲突,并造成一场长期的精神停滞。

\section{(23) 自深处}

% ---------------- p.163 ----------------
\origpage{52}这个“布尔巴基式”的圈子,肯定对我的为人、我对世界以及我在世界中位置的看法产生了强烈影响。这里不是试着界定这种影响、以及它如何在我生活中表现的地方。我只想说:我丝毫也不觉得我那些自命不凡的倾向、以及它们那种唯才是举式的合理化,是由我与布尔巴基的接触、由我进入“布尔巴基圈子”所激发的——至少四十年代末和五十年代不是。那些种子早就在我心里播下,在任何别的圈子里都会找到发展的机会。我报告过的“不中用的学生”那桩事,绝不是这个圈子里主导气氛的典型,恰恰相反,我再说一遍;它只是我自己身上一种暧昧态度的典型。布尔巴基的气氛是一种尊重人的气氛,一种自由的气氛——至少我是这样感受的;这种气氛本性上会打消并减弱一切趋向支配或自命不凡态度的倾向,无论是个人的还是集体的。

这个品质非凡的圈子不复存在了。它什么时候死的,我说不上来;想必没有一个人察觉,也没有一个人在内心深处为它敲响丧钟。我猜想,人身上一定发生了不知不觉的败坏——我们想必都“添了些年头”,变得干硬了。我们成了重要人物,有人听,有权势,有人怕,有人求。火花也许还在,但天真在路上丢失了。我们中也许有人会在死前重新找回它,像一次新生——但那个曾接纳我的圈子已不复存在,我若指望它复活,那是徒劳。一切都回到了秩序之中。

而尊重,也许也在路上丢失了。轮到我们有了学生时,也许已经太晚,最好的东西传不下去了——还有一星火花,但没有了天真,也没有了尊重,除了对“同辈”和“自己人”。

风可以刮起来,吹,烧——我们躲在厚厚的城墙后面,各人带着“自己人”。

一切都回到了秩序之中……

\section{(24) 我的告别,或:外国人}

% ---------------- p.164 ----------------
这场对我数学生涯的回顾,走上了与我预想完全不同的路。说实话,我本来根本没想到要回顾,只想用几行、顶多一两页说说:今天我与那个我已离开的世界是什么关系;也许还反过来说说:从我偶尔收到的回声看,我的旧友与我是什么关系。相反,我原本打算稍仔细地考察一下,我在那些紧张的数学工作年月里引入的某些想法和概念的遭际——它们有时颇为奇怪;或者我该说:我有幸瞥见、并把它们从全然未知的黑夜中拉进半明半暗、有时甚至拉到最明亮的日光之下的那些新类型的对象与结构!如今这番话,在这场已经变成对一段过去的沉思、变成一种更好理解并承担某个有时令人不知所措的当下的努力里,似乎显得走调。看来,原计划中对某个几何“学派”的反思,只好等一个更合宜的时机——那个学派在我的推动下形成,又几乎不留痕迹地挥发掉了\fnoteP{6}{}{这个“更合宜的时机”来得比预想的早,那场反思构成了《收获与播种》第二部分“葬礼”的内容。}。所以眼下,我关心的是把这场回顾进行到底:回顾\origpage{53}\origpage{54}我在数学界里的数学生涯,而不是对一部作品及其命运发议论。

在刚过去的五天里,我被反思笔记之外的别的任务占着,一个回忆却相当执拗地回来了。它将作为我停下的那篇《自深处》的尾声。

事情发生在1977年底左右。几周前,我被传唤到蒙彼利埃轻罪法庭,罪名是“无偿收留并供养一名身份不合法的外国人”(也就是说,一名在法国的居留证件不合规定的外国人)。正是这次传唤,让我知道了1945年那部规定外国人在法地位的法令中令人难以置信的一款;这一款禁止任何法国人以任何形式帮助“身份不合法”的外国人。这条法律,甚至在希特勒德国对犹太人也没有类似的条文,显然从未按字面执行过。由于一个非常奇怪的“巧合”,我有幸被选作第一个试验品,用来首次实施这条独一无二的条款。

有几天我一直发懵,像瘫痪了一样,深深地灰心。突然间,我仿佛被抛回三十五年前——那时人命不值几个钱,尤其是外国人的命……后来我反应过来了,我振作起来。几个月里,我把全部精力投进去,试图动员舆论,先是在我所在的大学和蒙彼利埃,然后在全法范围。正是在这段紧张活动的时期——为了一件后来证明注定失败的事业——发生了我今天可以称之为“我的告别”的那个插曲。

为了在全国范围行动,我给科学界五位特别知名的“人物”(其中一位是数学家)写了信,把这条法律告诉他们——它今天在我看来仍像我被告发那天一样难以置信。我在信里提议共同行动,表明我们反对一条邪恶的法律;这条法律等于把几十万住在法国的外国人置于法律之外,并把几百万其他外国人当作麻风病人一样,让民众对他们心存戒备——这些人一下子成了嫌疑犯,可能给那些不提高警惕的法国人招来最糟糕的麻烦。

奇怪的是,完全出乎我意料:这五位“人物”中没有一位给我回音。看来,我是有些东西要学的……

于是我决定去巴黎,趁布尔巴基讨论班的机会——在那里我准能遇到许多旧友——首先在我最熟悉的数学界里动员舆论。我原以为,这个圈子对外国人的事业会格外敏感,因为我所有的数学同行,和我一样,每天都要与外国同行、学生和留学生打交道;他们大多数、或者全部,都曾有过居留证件上的难处,都曾在警察总局的走廊和办公室领教过专横、而且往往是轻蔑。我曾把计划告诉洛朗·施瓦茨,他说讨论班第一天报告结束后会让我发言,把情况提交给在场的同行。

那天我就这样去了,公文包里一大叠传单,是给同行们的。阿兰·拉斯库(Alain Lascoux)帮我分发,在亨利·庞加莱研究所的走廊里,在第一场之前,以及两场报告之间的“幕间”。如果我没记错,他自己还做了一份小传单——他属于那两三位同行之列:听说了这件事之
% ---------------- p.165 ----------------
\origpage{55}\origpage{56}后,也动了心,在我去巴黎之前就联系了我,表示愿意帮忙\fnoteP{7}{(17)}{我投身的那场行动,得到的热烈回响和切实帮助,主要来自科学界之外。除了阿兰·拉斯库与罗杰·戈德芒的友好支持,我在这里还要特别提到让·迪厄多内的支持:他专程到蒙彼利埃出庭,为这件注定失败的事业在别的证词之外添上他热忱的证词。} (17)。罗杰·戈德芒也在其中,他甚至做了一份传单,标题是“诺贝尔奖得主进监狱?”。这在他是一番好意,可我们显然不在同一个频道上:好像丑闻在于动了一位“诺贝尔奖得主”,而不是随便哪个小人物!

布尔巴基讨论班第一天,人确实很多,极多或多或少与我相识的人,包括布尔巴基的旧友和旧伙伴;我想,他们大多数人想必都在场。还有好几位我过去的学生。我差不多有十年没见过这些人了,来的时候我很高兴有这个机会重见他们,哪怕一下子见这么多!不过,总会有机会再小聚的……

然而这次重逢“不是那么回事”,从一开始就相当清楚。许多手伸过来、握上了,这是肯定的;许多“哟,你也在这儿,什么风把你吹来了?”的问题,也是有的——但在快活的语调背后,有一种说不清道不明的尴尬气息。是因为我带来的事业他们并不真正感兴趣,而他们来这里是为了一场一年三次的数学典礼,需要全神贯注?还是与我来做什么无关,而是我这个人本身引起了那种尴尬——有点像一位还俗的神父出现在一群纯正的神学院学生中间所引起的那种尴尬?我说不好——也许两者兼有。在我这一边,我忍不住注意到,某些曾经熟悉、甚至曾是我的朋友的面孔上发生的变化。他们仿佛僵住了,或者塌陷了。我曾在他们脸上认识的那种灵动似乎消失了,仿佛从不曾有过。我觉得自己像是面对一些陌生人,仿佛从来没有什么把我与他们连在一起。我隐约感到,我们并不生活在同一个世界里。我原以为借这个不寻常的机会能找回一些兄弟,结果面对的却是陌生人。教养很好,这得承认:我不记得有过尖酸的话,也不记得有传单被丢在地上。事实上,好奇心帮了忙,分发的传单(几乎)都该被读过了。

然而那条邪恶的法律并没有因此受到威胁!我得到了我的五分钟,也许还占了十分钟,来谈那些在我看来是兄弟、却被叫作“外国人”的人的处境。阶梯教室里坐满了同行,比我做数学报告时还安静。也许对他们讲话的那份确信,已经不在我身上了。不再有从前那种同情与兴趣的流动。人群中想必有些赶时间的人,我一定是这样对自己说的;我缩短了发言,提议立刻与感到关切的同行们留下来,更详细地商议可以做些什么。

宣布散会时,是全体向出口的冲锋——显然,人人都有一班火车或地铁马上要开,绝不能错过!一两分钟之内,埃尔米特阶梯教室就空了,简直不可思议!空荡荡的大阶梯教室里,刺眼的灯光下,只剩下三个人,包括阿兰和我。第三个人我不认识,我敢打赌又是那种见不得人的外国人,交往可疑,而且身份还不合法!我们没有花时间对刚才在眼前上演的那一幕说得太多——它本身已经够雄辩了。也许我还是唯一一个不敢相信自己眼睛的人,而我的两位朋友当时很得体地对此不加评论。显然,我是初来乍到……

晚上在阿兰和他前妻雅克琳家里结束,盘点一下形势,把可以做的事过一遍;也彼此多认识一点。无论是这一天还是后来,我都没有费心把刚经历的这个插曲放进过去的脉络里去看。然而正是那一天,我一定是不言不语地明白了:某个圈子、某个我曾认识并爱过的世界已经不存在了;我原以为能重新找回的一份活生生的温暖,早已消散——多半消散很久了。

% ---------------- p.166 ----------------
这并没有阻止那个世界一年一年传到我这里的回声——那个温暖已逃逸的世界——许多次让我不知所措,让我痛苦地被打动。我怀疑这场反思对将来能否对此有所改变——除非也许是:我再被这样触动时,会少梗着脖子抗拒一点……

% ================================================================
%  第 8 章 师与生 (原书 p.167-174, 边码 p.57-p.66)
%  源文: _work/src1/task-ch08a.md
% ================================================================
\chapter{师与生}

\section{(25) 学生与研究计划}

% ---------------- p.167 ----------------
\origpage{57}我还没有把当年与其他数学家的关系巡视完——那时我觉得自己与他们同属一个世界,同一个“数学共同体”。剩下要看的,主要是我与我那些学生的关系——按我所经历的样子——以及我与另一些人的关系,那些人把我看作前辈。

总的来说,我相信可以毫不保留地说:我与学生的关系,是尊重的关系。至少在这一件事上,我相信,我当年做学生时从前辈们那里得到的东西,并没有随着岁月而变坏。由于我素有做“难”数学的名声(说实话,这是个再主观不过的概念!),而且被人认为比别家老板更苛刻(这一点就主观得少些了),来找我的学生从一开始就相当有劲头:“他们是想干的”!只有过一个学生,开头有点“提不起劲”,说不准他能不能起步——后来还是起了,而且没等我推一把,他自己就发动起来了。

就我所能记起的而言,凡是要求跟我工作的学生,我都收下了。其中有两个,几周或几个月之后发现我的工作方式不合他们的脾胃。说实话,现在我觉得那两次都是“卡住”的局面,而我当时却草率地把它读成了不适于数学工作的迹象。今天,我要作这类预言会谨慎得多。我毫不犹豫地把自己的印象告诉了这两位当事人,劝他们不要在一个——据我看——不合他们天性的行当里继续下去。事实上我后来得知,至少这两个学生中的一个,我是弄错了:这位年轻研究者后来在代数几何与数论交界处的一些难题目上赢得了名声。另一个学生是位年轻女子,她跟我这段不如意之后有没有继续,我就不清楚了。不能排除:我对她能力的看法说得太斩钉截铁,使她泄了气——而她也许和别人一样有能力做出好工作。我觉得我对这两个学生和对别人一样,是给了信任的。我缺的倒是%
% ---------------- p.168 ----------------
辨别力:面对那些肯定是“卡住”、而非“无能”的迹象时,没能分清两者\fnoteP{1}{(18)}{我想,这份缺乏辨别力,在这两回并不是出于我的疏忽,而更是出于不成熟、出于无知。直到大约十年之后,我才开始留意“卡住”的机制——在我自己身上,也在我亲近的人身上,在学生身上——并估量出它们在每个人的生活里扮演的巨大角色,而且不只是在中学或大学里。当然,我后悔这两回没有更成熟一些的辨别力;但我不后悔把自己那些印象——无论有没有根据——明明白白地说出来。当我在某一场合看出工作是敷衍着做的,把这些事照它们本来的样子点明,在我看来是必要而有益的。而如果在另一场合,我由此得出的结论是草率的、没有根据的,那么负有责任的也不只我一个。被这样摇了一下的学生,仍然有得选:要么从中吸取教训(也许头一回就是这样),要么就此泄气,也许还改了行(这也未必是坏事!)} (18)。

从六十年代初起,也就是在大约十年里,有十一个学生在我这里做了国家博士论文\fnoteP{2}{(19)}{\emph{耶稣与十二使徒}\par 从1970年至今,还有一个学生伊夫·拉德加耶里(Yves Ladegaillerie)在我这里准备并通过了论文。前一时期的学生是:P.贝尔特洛(P. Berthelot)、M.德马祖尔(M. Demazure)、J.吉罗(J. Giraud)、M.阿基姆夫人(Mme M. Hakim)、黄春生夫人(Mme Hoang Xuan Sinh)、L.吕西耶(L. Illusie)、P.茹阿诺卢(P. Jouanolou)、M.雷诺(M. Raynaud)、M.雷诺夫人(Mme M. Raynaud)、N.萨维德拉(N. Saavedra)、J.L.维尔迪耶(J.L. Verdier)。(其中六位还是在1970年之后才完成论文工作的,也就是在我很少能抽出时间指导他们的时期。)在这些学生当中,米歇尔·雷诺(Michel Raynaud)占一个特殊的位置:他论文工作所处理的那些本质性的问题与概念,是他自己找到的,而且他完全是独立地把它们发展出来的;所以我这个严格意义上的“论文导师”的角色,只限于读已完成的论文、组建评审委员会,并做其中一员。\par 在我自己出题目的时候,我很当心地只限于那些我足够熟悉、与我有足够深联系的题目,以便在需要时觉得自己能够支撑学生的工作。一个显著的例外是米谢勒·雷诺夫人(Michèle Raynaud)关于基本群的局部与整体Lefschetz定理的工作,这些定理是以适当的平展site上的1-栈的形式表述的。这个问题在我看来是困难的(后来也确实证明如此),而对于我提出的那些猜想(它们其实几乎无可置疑),我并没有证明的思路。这项工作一直延续到七十年代初;雷诺夫人(如同此前她丈夫的情形)发展出一套精细而独到的方法,没有得到我任何帮助,也没有得到别处的帮助。这项出色的工作还引出了一个问题:把雷诺夫人的结果推广到n-栈的情形——在我看来,在概形的框架里,这应当代表“弱Lefschetz定理”一类定理的自然归宿。不过,这里相应的猜想(它同样几乎无可置疑)的表述,本质上要用到n-栈的概念;而n-栈的追寻,据说正是本书的主要对象[其实那是《数学反思》第3卷,不是眼前这部《收获与播种》第1卷——见导论,p.(v)],正如“追寻栈”这个书名所表明的。我们大概会在适当的地方再回到这一点上,\par 另一个相当特殊的情形是黄春生夫人:我最初是1967年12月在河内遇见她的,那是我在被疏散的河内大学讲一个月“课程—讨论班”的时候。第二年,我向她提出了她的论文题目。她是在战争时期格外艰难的条件下工作的,与我的接触只限于零星的通信。1974/75年她得以到法国来(借温哥华国际数学家大会之机),并在巴黎通过了论文(评审委员会由嘉当(Henri Cartan)主持,成员还有施瓦茨(Laurent Schwartz)、德尼(J. Deny)、齐斯曼(M. Zisman)和我)。\par 最后我还得提到皮埃尔·德利涅和卡洛斯·孔图-卡雷尔(Carlos Contou-Carrère):这两位都多少算是学生,前者在1965–68年前后,后者在1974–76年前后。两人显然都有(而且一直有)不寻常的才能,他们运用这些才能的方式大不相同,运气也大不相同。来比雷之前,德利涅多少做过蒂茨(J. Tits)的学生(在比利时)——不过我怀疑,按通常意义上的“学生”来说,他在数学上不曾是谁的学生。孔图-卡雷尔做过圣塔洛(L. Santaló)的学生(在阿根廷),还曾有一段时间——多少算是——做托姆(René Thom)的学生!联系建立起来的时候,两人都已经有一个数学家的身量了;只是孔图-卡雷尔还缺方法、缺手艺。\par 我在数学上对德利涅的作用,只限于零零碎碎地让他知道我关于代数几何懂得的那一点点;而他学它,就像听人讲一个故事——仿佛他向来就知道。此外,还在于沿途提出一些问题,而这些问题他多半当场或在随后几天就找到了答案。这就是我所知道的德利涅最早的工作。1970年以后的(他的,以及我那些“正式学生”的),我仅凭极为零散、遥远的回声才有所知[其实那是《数学反思》第3卷,不是眼前这部《收获与播种》第1卷——见导论,p.(v)]。\par 按孔图-卡雷尔自己在论文开头所说的,我在他身边的作用,只限于把他引入概形的语言。无论如何,他近些年作为国家博士论文所准备的工作,我只是从很远的地方跟着;那题目是最时髦的那一类,超出了我的能力。在广阔世界里经历了几番波折之后,孔图-卡雷尔最近终于在最后一刻——而且(现在在我看来)是违着他的心意的——被逼得来求助于我,请我充当论文导师并组建评审委员会。(这就使他冒着一种风险:显得像是“1970年之后”的格罗滕迪克的学生——在这种情势下,这可能带来严重的不便……)我尽我所能履行了这项任务;这大概是我最后一次行使这个职能(在国家博士论文这一层面上)。在这个有点特殊的情形里,我格外高兴有让·吉罗(Jean Giraud)的友情协助:他也从自己的时间里抽出一两个月,细读那部篇幅庞大的手稿,并写了一份详实而热忱的报告。} (19)。各自选定了合意的题目之后,他们都干得起劲,而且(我是这样感觉的)都与自己选定的题目深深认同。然而有一个例外:有个学生选了一个“该做的”题目——也许并无真正的信念——而这题目也有它不讨人喜欢的一面:它要做的是把一批已经到手的思想作技术上的定稿整理,有时艰难,甚至干枯,再也没有多少意外、多少悬念可言,%
% ---------------- p.169 ----------------
\origpage{58}只是把它们摆出一个全景来\fnoteP{3}{(20)}{这让我想到莫妮克·阿基姆(Monique Hakim)选的题目:说实话,那也不更招人喜欢;我真不知道她是怎么保持住精神的!她固然有时也吃力,但无论如何没到让她变得忧郁或阴沉的地步,而我们之间的工作是在一种亲切、放松的气氛里做的。} (20)。

我当时被一项庞大计划的需要卷着走——我需要人手——所以在出这个题目时,必然是缺乏心理上的辨别力:它肯定不合这位学生特有的性情。他那一头呢,大概也没太看清自己上的是一条什么船!总之,他和我都没能及时看出:这事一开始就起步错了,不如换个题目重新来过。

显然,他干活没有真正的信念,而且始终带着一副多少有些忧郁、阴沉的神气,不曾放下。我想我当时已经到了这样一个地步:不太留意这类事情——而这类事情(我本该记得)在任何研究工作中,而且不只在研究工作中,都能造成天壤之别!我那时的角色,仅限于:工作显出拖沓的样子,我就心烦;工作又动了,我就宽慰地“呼!”一声;等到最后,预定的计划总算“收了尾”。

直到1970年我那次“醒悟”之后好几年,因为要和这位从前的学生通信(他当了教授——再说,在那宽厚的年月里,人人都当上了教授!),我才想到:那一回终究是有什么地方不对头,那也许算不上一次彻底的成功。今天,它在我看来是一次失败——尽管“计划收了尾”(绝非草草了事!),尽管有文凭,还有随之到手的位置。而且我要负很大一部分责任:我把一个计划的需要,放在了一个人的需要之前——而这个人,是怀着信任把自己交到我手上的。我刚才还拿来标榜的那种“尊重”(“毫无保留”!),我说我对学生们做到了——在这里始终是表面的,与尊重的真正灵魂相隔绝:那灵魂就是对一个人的需要怀着情意的关注,至少当这些需要的满足取决于我的时候。这里的需要,是工作中的一份快乐;没有它,工作就失了意义,变成强制。

在这次反思中,我有机会谈到过一个“没有爱的世界”,并且在自己身上寻找那个我拒斥的世界的胚芽。这里就有一粒,而且个头不小——至于它在别人身上怎样长起来,我今天就说不上来了。这种表面的、缺乏关注、缺乏真正爱的尊重,也是我给予自己孩子们的“尊重”。和他们在一起,我有这份特权:看着这粒种子长起来,看着它蔓延。我也多少明白了:面对收成,皱眉头是没有用的……

\section{(26) 严格与严格}

如果把这位学生除外——他肯定并不比别人“天分”差——我可以说,我与学生之间的关系是融洽的,常常甚至是亲热的。出于无奈,他们%
% ---------------- p.170 ----------------
\origpage{60}都学会了耐心对待我作为“老板”的两个主要毛病:一是字迹无法辨认(不过我相信,最后人人都学会了辨认我的字);二是——这一点当然更严重,而我很晚才发现——我从根子上难以跟随别人的思路,除非先把它译成我自己的意象,用我自己的风格重想一遍。我远更倾向于把自己深受浸润的某种对事物的眼光传给学生,而不是鼓励他们心中绽出一种个人的眼光——那眼光也许与我的颇为不同。与学生的关系中这层困难,到今天也还没有消失;但我觉得它的影响减弱了,因为我意识到了自己身上的这种倾向。也许我的气质——无论是先天的还是后天的——使我更适于孤独的工作(我起初十五年的数学活动,约1945至1960年,也正是这样过来的),而不那么适于做“师傅”,去和那些数学上的志业与个性尚未完全成形\fnoteP{4}{(21)}{也许更准确的说法是:就我这种气质而言,我至今缺的是必要的成熟,还不能完全承担起一个教育者的角色。我后天形成的气质,长久以来被一种过分占优势的“男性”(或“阳”)特征所标记;而成熟的一个方面,恰恰是以“女性”(或“阴”)为主的“阴阳”平衡。\par(后来补写。)比我缺一份成熟更甚,我看在我作为教育者的生活里,直到今天,我缺的是一种慷慨——一种比在时间和精力上的投入更为精细地表达出来的慷慨,也是更本质的慷慨。在我教学的第一时期,这份欠缺并没有以可见的方式表现出来(比如说,以一连串失败局面的积累),多半尤其是因为,那些选择来跟我工作的学生身上的强烈动机把它补偿了。相反,在第二个时期,从1970年至今,在我看来,对于我在研究层面(也就是从DEA这一层往上)的教学中所看到的整体失败,这份欠缺至少是原因之一,而且无论如何是最直接牵涉到我本人的那一项原因。关于这一点,见《纲领概要》第8段,以及第9段《一项教学活动的总结》,那里透出这项活动七八年来留在我身上的受挫感[并请对照后来补写的注(23iv)。]。} (21)的学生相处。不过也确实:我从小就喜欢教书;而且从六十年代到今天,我能有的那些学生在我生活里占了重要的位置。这也就是说,我的教学活动、我作为教师的作用,在我的生活里曾经有过、并且仍然保有一个很大的位置\fnoteP{5}{(22)}{不过也许不会太久了,因为我已经决定申请进入法国国家科学研究中心(CNRS),从而结束在大学环境里的教学活动——这几年来,那活动变得越来越成问题。} (22)。

在我教学活动的前一时期,我和任何一个学生之间都没有过表面上的冲突,哪怕只是关系里一阵短暂的“冷淡”也没有过。只有一次,我不得不对一个学生说:他工作不认真;要是照这样下去,我就没兴趣再和他一起干了。那事情究竟怎样,他当然和我一样清楚;他改了过来,这件事就了结了,没有留下一片乌云。另一次,已经是七十年代初,那时我最主要的精力投在“生存与生活”(Survivre et Vivre)小组的活动里;有个学生看了我刚写好的、关于他工作的论文评审报告(我习惯如此),发起火来,认为报告里某些议论是在质疑他工作的质量(这绝不是我的本意)。这一回是我把话收回、把方向拨正了,没有为难。当时我并不觉得这段小插曲会在我们的关系里留下阴影;不过也可能是我弄错了。这位学生和我的关系,比我和别的学生更少个人色彩(我说过的那个“忧伤的学生”除外);那只是一份良好的工作关系,仅此而已,我们之间并没有真正的暖意流过。不过,我并不认为是我身上一种无意识的不厚道,使我把那些他认为于己不利的议论写进了报告——他还说“这事他不会放过”,不像他的一个同伴,那人已经在我这里通过了论文。与另一位学生——他天性敏感而多情——我则有着一份%
% ---------------- p.171 ----------------
\origpage{61}格外亲热的关系;如果说我在关于他论文的报告里也写进了那类曾让他同伴大为不快的议论,那肯定不是出于不厚道!再说,对他们两位,也对我所有的学生,只要我对他们提交的工作不是完全满意,我就不会给答辩开绿灯。而且,那个时期我的学生,论文一通过,没有哪一个不是很快就找到了一个与他相称的位置。

直到1970年:我几乎随时都能抽出时间给学生\fnoteP{6}{(22')}{即便在1970年之后,当我对数学的兴趣在我的生活里变得零散而边缘时,我也不认为有过学生来找我、要我跟他一起工作而我推辞掉的场合。我甚至可以说,除了两三个例外,1970年之后我的学生对他们所做工作的兴趣,远低于我自己对他们题目的兴趣——即便在我除了去学院的日子几乎不操心数学的那些时期。因此,我在1970年之前随时能抽出时间给学生的做法,以及作为它主要标志的那种在工作上的极端严格,对于我后来大多数学生都会毫无意义——他们做数学没有信念,仿佛是靠着一种他们不得不对自己不断施加的努力……} (22')。时机成熟的时候,每逢用得着,我就和这一个或那一个待上一整天、必要时好几天,去啃那些还没弄妥的问题,或者一起把他们文稿的一稿又一稿过一遍。就我对这些工作时段所体验到的而言,我不觉得自己在里面扮演过拍板定案的“导师”角色;每一次都是共同的研究,讨论平起平坐地进行,直到双方都完全满意。学生投入的精力是巨大的,当然无法与我自己应投入的相比——不过我这边有更多的经验,有时还有训练得更敏锐的嗅觉。

然而,在我看来,对任何研究——无论智性的还是别的——的质量最要紧的东西,绝不是经验的问题。那是对自己的严格要求。我要说的这种严格,本质上是精细的;它不属于一丝不苟地遵从无论什么样的规范——严格的规范也好,别的规范也好。它在于对我们自身内部某种细微之物的极度关注,而那东西逃开一切规范、一切尺度。那细微之物,就是:对所考察之物的理解,有,还是没有。更确切地说,我要说的这种关注,是对每一时刻当下那种理解之品质的关注:从一堆杂乱堆叠的概念与陈述(假设的或已知的)发出的那片噪音,直到完全的满足、完美理解那种完成了的和谐。一项研究的深度——无论它的结果是片断的理解还是总体的理解——就在这份关注的品质里。这样一种关注,并不像遵守某条训诫、凭着“当心点”的刻意意图去留神的结果;在我看来,它是从求知的激情里自发地生出来的,是那些把求知的冲动与它那些自我中心的赝品区分开来的标志之一。这种关注有时也被叫作“严格”。这是一种内在的严格,独立于某一时刻在某门(不妨说)特定学科里可能通行的那些严格规范。如果说在这本书里我容许自己对那些严格规范(我曾教过它们,它们自有其存在的理由和用处)放肆一些,我并不认为这本书里这种更本质的严格,比我过去那些规范文体的著作里更少。而如果说我毕竟也许还能传给学生某种比一套语言和一种手艺更贵重的东西,那无疑就是这份严格、这份关注、这份严谨——如果说不是在与他人的关系和对自己的关系里(在那一层,它在我身上和在任何别人身上一样缺),那至少是在数学工作里\fnoteP{7}{(23)}{\emph{孩子与老师}\par“传递”这个词在这里并不真正符合事情的实际,它提醒我要采取一种更谦逊的态度。这种严格不是可以传递的东西,至多只能被唤醒或鼓励——而它从很小的年纪起就被忽视了、被挫伤了,被家庭环境,也被学校和大学。就我所能记起的,这种严格一直在我那些追寻里——至少是智性性质的追寻里;而我不认为它是我的父母传给我的,更不认为是学校的老师们、或我那些数学前辈们传给我的。在我看来,它属于“无辜”的禀性,因而也属于每个人生来就被分派到的东西。这份无辜很早就“青的、没熟的都尝遍了”,于是它不得不或深或浅地沉下去,以致在余生里往往再也见不到它多少痕迹。在我身上,由于一些我还没想到去探究的原因,某种无辜在智性好奇这个相对无害的层面上存活了下来;而在别处一切地方、在所有人身上,它都深深地沉了下去,神不知鬼不觉!也许“教育”——就其完整意义而言——的秘密,或者说谜,就在于重新与这份表面上消失了的无辜取得接触。但是,如果在教育者本人身上这份接触不是已经先在、或已被重新找到,那就谈不上在学生身上重新找到它。而这时教育者“传”给学生的,绝不是这份严格或这份无辜(它们在这两人身上都是天生的),而是一种尊重,一种对那件通常被弃置之物的默默的重新肯定。} (23)。

% ---------------- p.172 ----------------
\origpage{62}这确实是一件很微薄的东西,但也许毕竟聊胜于无。

\section{(27) 失误——或二十年后}

也许除了我谈过的那两个学生——与他们终究没能建立起工作关系——之外,我不记得其他那些来找我、要求跟我工作的学生,是带着“怯场”或畏惧来的。想必他们多少已经认得我,哪怕只在我于法国高等科学研究所(IHÉS)的讨论班上待过一阵子。倘若在我们关系开始时有什么拘束,它最后总会在工作过程中消散,不再留下痕迹。不过这里我要举两个例外。一个就是那位始终没能真正对工作生出兴味的学生,甚至在我们共同工作时,他也一直只是单音节地应声。也许他来的时候,我能留给学生的时间正要减少;和他也没有过那种对着稿子、一坐就是一个下午、一整天的时段。不,确实,我不记得有过那样的时段;我倒觉得我们多半是匆匆见一面,一两个钟头,看看他做到哪儿了。说到底,撞上我,他大概是最倒霉的一个!

我本来要说的另一个学生,则是在我还随时能抽出时间给学生的那个时期跟我工作的。我们的关系从一开始就融洽。他甚至属于那少数几个与我建立起友情的学生之列——我有时去他们家,他们也来我家,有点像两家之间的往来。确实,即便在这些情形里,关系也始终停留在相对表面的层次——至少就我这一方而言。在意识层面上,那时我对自己家里、自己屋顶下发生的事就已经所知无几;至于我那些数学家朋友——学生也好,不是学生也好——的生活,除了妻子和孩子的名字,我最终几乎一无所知(就连这些名字,我有时也会忘掉,而从没有人为此怪我!)。也许我曾是“书呆子”的一个极端例子;但我认为,在我所认识的那个数学圈子里,大多数关系——即便是友好的、亲切的——都停留在这种表面的层次:彼此最终所知极少,除了在“未言明”的层面上被感知到的那一点。这无疑是人与人之间的冲突在这个圈子里如此罕见的原因之一——尽管在我看来很清楚:分裂存在于我的大多数同事和朋友内部,也存在于他们的家庭内部,和我身上、和别处一样多。

我不认为我与这个学生的关系有别于我与别人的关系;当时我也不觉得反过来,他与我的关系与其他学生的关系有什么显著的差别,尤其是与那些同我结下友情的学生相比。只是不久以来,我才得以明白:那一定是一份比我对大多数别的学生更强的关系。一个未曾说出的冲突,其可见的表现像一次意外的启示般到来——那已是他做我学生将近二十年之后。直到那时,我才把它和一件早已被遗忘的“小”事联系起来。在很长一段时间里,也许甚至在整个那一段时期(也就是几年)里——%
% ---------------- p.173 ----------------
\origpage{63}\origpage{64}我们或多或少定期一起工作的那些日子——这个学生一直保留着某种“怯场”。它每次见面都表现出来,带着不会骗人的迹象。随后,在共同工作的过程中,这些迹象又相当快地消失了。这些不安的迹象当然让我不自在,而我觉得他更不自在。我们两个都装作不知道这事,这也理所当然。想必我们谁也不会想到去谈它,甚至不会想到在心里给这样一个古怪、显然值得一顾的处境一点注意!在他那边,也在我这边,这份“怯场”想必被感受为一处单纯的“失误”,一处本不该有的东西。那“失误”定期提醒我们想起它,但每一次,它都识趣地消失,让我们有时间安安稳稳地去忙正经事——数学——同时忘掉“那本不该有的东西”。我不记得自己哪怕一次停下来,对那失误的意义提出什么问题;我也确信,我的学生兼朋友那一头也是如此。想必在我们各自从幼年起在身边所见的一切里,没有任何东西能在他或我心里提示出另一种对待一件令人不快之事的态度,除了尽可能把它挪开、让它不再碍事。在那件事上,这完全做得到,甚至很容易;而我们完全一致地什么都没看见、什么都没感觉、什么都没听见。

然而,借这两三年来回到我这里的许多回声与相互印证,我明白:那被我们当作“本不该有”而挪开的东西,并不因此就不再存在、不再表现出来。有时回到我这里的东西,也同样“不该有”——可它“在”,而且现在已经不能一手背就把它拂开了……

\section{(28) 未完成的收获}

直到1970年第一次“醒悟”之前,我与学生的关系,正如我与自己工作的关系,是满足与快乐的源泉,是我生活中那份和谐感的可触、无可否认的基础之一——那份和谐继续赋予我的生活以意义,而与此同时,一种捉摸不定的毁坏正在我的家庭生活里肆虐。那时候,在我眼里这些关系里没有任何表面冲突的成分;当时它们当中没有哪一个,哪怕在一瞬之间,成为挫折或痛苦的起因。有一件事也许显得悖理:我与某个学生的关系里的冲突,偏偏要等到那次著名的醒悟之后才显出来——也就是在一个转折之后,那转折给了我的生活一种它从前不曾有过的开阔,也许还给了我这个人一点点柔韧的开端——人们会想,这些品质本该有助于化解或避免冲突,而不是挑起或激化它。

可是再仔细看,我就看出这悖理只是表面的;无论从哪个角度看,它都消散了。第一点:一个冲突要有机会化解,首先得它已经表现出来。冲突表现出来的阶段,相对于隐藏着、被忽视的冲突而言,是一种成熟;而后者其实也有表现,而且那些表现恰恰因为借它们表达出来的冲突仍被忽视而更为“有效”。第二点:一个冲突要能以可辨认的方式表现出来,首先得有一段距离缩小了,或者消失了。近十五年来我生活里发生的种种变化——尤其是在一次次“醒悟”过程中——在我看来,都是那种能缩小距离、抹去隔绝的变化。一个冲突面对一位有声望、受人仰慕的老板时难以表达,面对一个被剥去了权力位置的人(在这里是自愿剥去的)、一个从某个握有权威与声望的圈子里自我放逐的人,就自在多了——这个人越来越不被看作某种实体(比如数学)的化身或特权代表,而越来越被看作一个和别人一样的人:%
% ---------------- p.174 ----------------
\origpage{65}\origpage{66}一个不仅可能被伤到,而且越来越不愿把自己的伤或痛藏起来的人。第三点,也是最要紧的:我第一次醒悟以来的演变——尤其是在那个时期以及随后的几年里——本就容易在我那些旧学生井井有条的世界里激起(或者也许是唤醒)疑问、不安、一种“重新审视”。我有充分的机会看出:不仅在他们身上如此,在数学界我昔日的朋友与同伴中间也如此,有时甚至在一些只凭传闻认识我的科学同行中间也如此。

还得说:一个稍有深度的冲突能得到化解,是最罕见不过的事。最常见的情形是:尽管有一切休战与表面的和解,我们那些冲突越排越长的队伍仍一路跟着我们,几乎寸步不离地跟一辈子,最后才在殡葬工那双阴沉的手里把我们松开。我有时有幸看到一个冲突多少解开了,有时甚至看到它在认识中化解——但直到如今,这样的事还不曾在我与某个学生、或与数学界某位昔日朋友的关系中、并借这关系而发生。我也很清楚:这样的事是否终会发生,毫无把握,哪怕我还能再活一百年。

有一件事很值得注意:我与某段过去决裂的那个时刻——我是说离开IHÉS那个插曲(这个机构多少像是我周围形成的那个数学小宇宙的“母模”)——这个决定性的插曲,同时竟是我某个学生对我的深刻敌意第一次表现出来的场合。肯定是这一情形,使那个插曲格外难堪、格外痛苦,像一次分娩或一次出生发生在格外艰难的条件下。当然,那时我还不能在我后来学会的光亮里去看这个插曲——它的意义当时逃开了我。此后很久,这份痛苦的意外仍留在那里。然而,就在那一年的夏天,这场带着苦涩的离开已经显为一次解放——就像一扇门忽然大大敞开(只消我推它一下!),通向一个不曾料想的世界,唤我去发现它。此后每一次新的醒悟,也都是一次新的解放:发现一种从属、一种内在的束缚,并且重新发现一个巨大未知者的在场——它就藏在那本该“已知”之物的熟悉外表后面。但同样在这整整十五年里,直到今天,这份固执、沉默而无懈可击的敌意一直跟着我,成为我在数学生活里所知道的唯一而持久的挫折之源\fnoteP{8}{(23')}{不过七八年来,在我的数学生活里还有另一个慢性的“受挫之源”,只是它在这些年里表达得隐蔽得多。靠着同一种“令人受挫”的局面在我的教学活动中顽固地重复、积累,它最终还是显了出来,并且最后爆发成一种“受够了!”,使我几乎终止了一切所谓“研究指导”的活动。在这次反思的过程中,我有一两次触及这个问题,最后至少在末尾稍作考察。我在那里至少描述了这份受挫感,并考察了它在我“回到数学”这件事上扮演的角色(参看第50段“一段过去的重量”)。} (23')。也许我可以说,它是我为那第一次解放、以及随后的几次解放付出的代价。但我很清楚:解放与内在的成熟,是与“要付的代价”不相干的事;它们不是“赚”与“亏”的问题。或者换个说法:当收获被带到头、当它完成的时候,就没有损失——那看似“损失”的,已经变成了“收获”。而事情就变得清楚:我至今还没能把那场收获带到头;就在我写下这几行的此刻,它仍是未完成的。

% ================================================================
%  第 8 章 师与生 (原书 p.175-186, 后半: 8.5-8.8)
%  源文: _work/src1/task-ch08b.md
%  说明: 本章前半 (8.1-8.4) 见 part1-ch08a.tex
% ================================================================

% ---------------- p.175 ----------------
\section{(29) 敌父(1)}

1970年那个转折之后,在完全不同的环境里——一所外省大学——开始跟我做工作的那类学生,与从前的学生也很不相同。此外,在我这里做到国家博士论文水平的,只有两个人。其余人的工作,都在DEA或第三阶段博士论文的水平上。我还该算上一大批紧紧跟住某些“入门研究”课程的学生:这些课对他们来说,是提出往往出人意料的数学问题的机会,有时也是想象出别出心裁的方法去解决它们的机会。我遇到的最积极的参与,是在给一年级学生开的某些“选修课”上。相反,在那些已经在大学氛围里泡了几年的学生身上,某种清新、某种兴趣的能力、某种个人的眼光,已经或多或少褪了色。选修课的学生里,有几个显然具备成为出色数学家的素质。看当时的情形,我克制住自己,没有鼓励任何一个人走上那条路——尽管那条路本可以吸引他们,他们本可以在那里出类拔萃。

我跟那些为准备硕士文凭而听我这类“课程”的学生,关系大多没有延续到一年以上。每一次,我的印象都是:这些关系很快变得亲切而松弛,总的说来如此。只有一个被无孔不入的“怯场”缠住的学生是例外\fnoteP{9}{(23$''$)}{\emph{不敢玩}\par 这个学生曾和我一起做了一整年的DEA“实习工作”,而在他与我的工作关系里,直到最后都保持着“紧绷”状态。那是一种坦率友好的关系,贯穿着一种无可怀疑的相互好感。然而可就是有这份“怯场”;这份恐惧——它真正的起因肯定不是对我这个人的畏惧,尽管看上去像是如此。要不是这个学生自己跟我谈起,我也许根本不会察觉这件事;他谈,无疑是为了或多或少“解释”一下他在这一年里工作几乎完全卡住的原因。\par 正如在其他学生身上发生过的那样——他们和他一样,起初对某种几何实质抓得很牢——卡住恰恰出现在要“做零件活儿”的时候:也就是把成形的陈述白纸黑字写下来,或者仅仅是把握我提供的那些陈述的意义与含义;我提议把它们当作一种语言的基石、当作“游戏规则”接受下来。“学校式的”反应几乎总是把面对“做研究”局面的学生推向:把导师传下来的那些心照不宣的“游戏规则”当作既定的、又必须服从的东西接受下来,而这些规则恰恰不是要试着说明白的,更不是要弄懂的。这些心照不宣的规则所采取的具体形式,就是语义学或计算上的“配方”,不妨说,像预科班的题解书那样(或者任何一本普通教科书)。此外,学生还期待导师给出“证明……”这种形式的任务——那是他以往经验里遇见的唯一一种数学“思考”形式。(而且我不认为大多数职业数学家、以及其他科学家的禀性有本质不同——只不过“导师”换成了“共识”:共识规定当下的游戏规则,并把它当作不可更动的既定之物。这共识还规定哪些“问题”是要去解决的;在这些问题之间,每个人都觉得自己有随意挑选的自由,甚至允许自己在工作过程中修改它们,乃至发明另一些……)。我注意到,我面对有待探究的数学实质时那种完全不同的态度——因而也是面对学生时那种完全不同的态度——几乎必定会引发一种失措,而焦虑就是它的征兆之一。像一切焦虑一样,它倾向于给自己找一张面孔,投射到某个外部的“理由”上,不管那理由说得通说不通。焦虑最常见的一张面孔,恰恰就是恐惧。\par 这类困难在我教学生涯的第一个时期几乎没有出现,也许只有两例例外:那两次“师生”关系没超过几个星期;也许还有(我说不准)“忧伤的学生”那一例——他也许感到自己被“铆”在一个完全不启发他的题目上,而他本来完全有更换的自由。在那个(我也谈到过的)长期被某种怯场困扰的学生身上,原因显然在别处。他的工作丝毫没有卡住,恰恰相反,他对选定的主题极为自在,并在其上做出了一项有规模的基础性工作。那个时期我的学生大多是巴黎高师的旧生,他们与亨利·嘉当(Henri Cartan)的接触,已经让他们见识了另一种“接近数学”的方式。与之(可以说)相反的另一极,是我在蒙彼利埃大学任教的第二个时期:正是在一年级学生那里,我所说的焦虑最少干扰思考性的工作。在这些学生中的许多人身上,面对一种不同进路时的惊讶既不引起焦虑,也不引起封闭,相反带来的是敞开和兴致——终于可以做点有趣的事了!据我的观察,几年大学对一个学生创造性禀性的作用是彻底而毁灭性的。奇怪的是,在这方面,长年中学的作用似乎相对无害。原因也许是:\par 大学那几年正落在这样一个年纪:我们身上与生俱来的创造力,最迟必须通过个人的工作表达出来,否则就要永远沉船——至少就智力上的创造性工作而言是如此。肯定是出于一种健康的本能,我在做学生的那些年(也是在蒙彼利埃大学),几乎不踏进课堂,把几乎全部精力都投给了个人的数学思考。} (23'),
% ---------------- p.176 ----------------
\origpage{67} 其余则都是如此——那些本应正式在我指导下准备研究工作的学生,无论哪一个层次,也一样。与从前的学生相比,一个不同之处(还有许多别的!)是:我们的关系不那么局限于共同的数学工作。学生与我之间的交流,常常以不那么表面的方式牵涉到我们两个人本身\fnoteP{10}{(23$^{v}$)}{这种差别有一个格外醒目的征兆,出现在“外国人事件”之际——我曾有机会谈到过(第24节)。当时,许多与我完全陌生的人向我表达了同情,而我记不起1970年之前的哪个学生想到过以这种方式表示什么,更不用说在我投身的那场行动中向我提供任何帮助。相反,在我看来,第二个时期的学生或旧学生里,没有一个不曾向我表达同情与声援,还有几个积极参加了我在当地开展的运动。在这个狭小圈子之外,1945年法令那件事也在许多仅仅知道我名字的大学生中间激起了一定的波澜;在我被传唤的那一天,有不少人来到法院,以示声援。这后一个情形还提示:我所察觉到的“之前”与“1970年之后”学生态度上的差别,也许较少表达他们与我之间关系的差别,而更多表达的是一种心态的差别。显然,我“之前”的那些学生已经成了重要人物,而要让重要人物动心,是需要很多条件的……但我1970年离开IHÉS、投身战斗行动那一幕,似乎表明并不只是这一点。那是个他们中间还没有谁太像重要人物的时刻,然而我记不起他们中有谁对我所投身的活动表现出丝毫兴趣。我倒认为,那一定让他们所有人,无一例外地感到不自在。这仍然合乎“心态差别”的方向,但这不能只归因于社会地位的差别。} (23'')。因此,在我教学生涯的第二个时期,与某些学生的关系中的冲突成分表现得更清楚、更直接,甚至更激烈,这就不奇怪了。在第一个时期的旧学生里,有两个人后来表现出系统而毫不含糊的对抗态度(我曾顺带提到过),不过始终停留在未言说的层面,也许甚至停留在无意识的层面。在更长一些的第二个时期,有三个学生让我遭遇了对抗。其中两个,表现得相当尖锐。

这些学生中的一个,对抗是在一段最友好的关系里一夜之间冒出来的——那是在这位朋友不再是我的学生许多年之后。我猜测,冲突的原因与其说是我那不堪形容的举止和人格,不如说是一份压抑已久的不满:他的工作(曾经很出色)没有得到他本有权期待的那份接待。这就是“1970年之后”以我为老板的那份可疑特权的背面;他必定为此怨我,尽管即便在内心深处也不大肯向自己承认。

另一个学生,尖锐的对抗在工作一年半之后就已经出现,而此前的氛围看上去非常亲切。这是第一次也是唯一一次,学生与我之间的关系困难出现在他还是学生身份的时候。它使共同工作的继续成为不可能——而那份工作本来开局吉利,带着最好的兆头般的热情,而且应当说,题目是一个绝好的思考主题。我的感觉是:这位年轻研究者身上有一种暗中作祟的自信不足——不相信自己有能力做出好工作(这能力在我看来毫无疑问);而对抗以尖锐的音调爆发出来,是某种“向前逃窜”,为的是抢在一种他惧怕的失败之前动手,并把责任预先推到一位可憎的老板头上\fnoteP{11}{(23$'''$)}{\emph{两兄弟}\par 这个学生身上的对抗,一上来就采取了“阶级对抗”的形式:我是那个对他的数学前途拥有“生杀之权”的“老板”,可以凭自己的好恶决定……当然,事情本身只能证实这种看法,因为我很快就不再承担指导这个学生的责任(那已变得令人难受)。这把他置于一个微妙的处境——在眼下的世道里,找一个“老板”可不是那么显而易见的事,尤其是题目已经选定的时候。在另一个学生身上,合法期待落了空,对抗采取了类似的形式。我被感受为专横的“学阀”,不能容忍那些被他视为下属的人(学生或级别较低的同事)有任何反驳。\par 这样一种“阶级态度”,在我与第一个时期学生的关系里,从未有过一丝一毫的表现。显而易见的原因是:在1970年以前的形势下,学生通过论文之后毫无疑问会得到一个副教授(maître de conférences)的职位,因而享有与我相同的社会地位,即“大学教授”。数字最能说明问题:1970年以前开始跟我做工作的十一个学生,一完成工作就得到了副教授职位;而在我的指导下或多或少工作过的那二十来个学生里,没有一个获得这样的职位。诚然,他们之中只有两个人有足够的动力去做国家博士论文(而且两人的论文都很出色)。\par 因此,在这第二个时期,某些暧昧(其深层根源始终被遮蔽着)采取了阶级对抗的形式,采取了对“老板”的不信任(这种不信任被呈现、也被感受为“本能的”“发自肺腑的”)——这就不足为奇了。在那些或多或少当过学生样子的人里,有一个与我的友好关系持续了十来年,没有出现过表面上的对抗插曲,却仍带着同一种暧昧,表现为一种不信任的态度,被“储备”在明显的善意之后。说实话,我从未被这种例行公事式的“不信任”骗过;它在我看来,尤其是这位朋友认为应当给自己找的一个理由,好让自己不去冒险越出他选定的、界限分明的领地——无论在职业生活里,还是在整个生活里都是如此。可这件事他本来可以自由地去做,没有人(顶多只有他自己!)会要他交代。\par 这三个例子,也是我整个教书经验中仅有的、学生(或多少像学生的人)与我之间的关系中某种暧昧以“阶级态度”表现出来的情形。当这种态度出现在大学“团体”内部的同事之间时,显得格外暧昧:在这样一个团体里,双方都享有与常人处境相比极为悬殊的特权,而这些特权使级别(以及薪金)的差别显得相对微不足道。我还注意到,一旦当事人自己也被提升到前一天他还在指责别人的那种处境,这些态度就像中了魔法一样消失了(而且理所当然!)。\par 此外,在我于数学界内部(常常也在外部)目睹的多数——如果不是全部——冲突局面里,我都辨认出一种类似的暧昧。那些“有位置的人”,无论其级别是否合乎他们的期待(合理与否),都享有相当闻所未闻的特权,那是任何别的职业或生涯都无法提供的。那些没有位置的人,则向往同样的安全感和同样的特权(这并不必然妨碍他们对数学本身感兴趣,有时还做出漂亮的工作)。在眼下的世道里,抢位置竞争激烈,没有位置的人常被当作流浪汉对待:我不止一次感到那种默契——在乐于羞辱人的人与被羞辱的人之间;后者咽下去,缩起来。他苦涩与敌意的真正对象,并不是那个行使了权力的人,而恰恰是他自己:他自己缩了起来,并且把对方奉为那种他乐意行使的权力的持有者。乐于羞辱人的人,也是那个报复并补偿(却永远抹不掉……)一场长期承受、早已深埋并被遗忘的羞辱的人。而默许自己被羞辱的人,是他的兄弟和效仿者:他暗中羡慕对方,并在苦涩里埋下那场羞辱,以及羞辱带给他的、关于他自己的那条卑微的消息。} (23''')。

% ---------------- p.177 ----------------
\origpage{68} 在学生与我之间所有这些冲突的显现里,有一个共同点——我教人做数学家这一行快二十五年了——就是强烈的暧昧。无一例外,对抗都是在事后才表现出来的,往往来得很阴险,出现在一段就其本身而言无可怀疑的好感关系里。我甚至可以说,在所有这些情形里,正如在许多别的、并未表现出明确对抗成分的情形里一样,我这个人在他们身上曾产生、并且仍在产生强烈的吸引力。肯定正是这吸引力本身的力量,也滋养了对抗的力量,并保证了它的延续。在对抗采取强烈反感、愤然排斥形式的情形里,肯定也仍然如此;同样,在另一个相反的极端情形里也是如此:在友谊般的敬重这件必不可少的外衣之下,(机会合适时)表达的是一种潇洒而精心调配过的轻蔑的做作……

说实话,这样的暧昧局面并不限于我与某些学生或旧学生的关系。事实上,在我整个成年生活中,从三十岁起(也就是说,从我母亲去世起),它们一直大量存在。在我的情感生活或夫妻生活里是这样,在我与男人的关系里也是这样,更确切地说,主要是与那些明显比我年轻的男人。我终于明白:我身上有某种东西——是天生的还是后天得来的,我说不清——似乎使我倾向于充当父亲的形象。看来我必定有理想的体格和相宜的振动频率,正好造就一个完美的养父!应当说,父亲这个角色我穿起来就像手套一样合手——仿佛它生来就是
%
% ---------------- p.178 ----------------
\origpage{69}\origpage{70} 我的。我不会去数自己有多少次对另一个人进入这样的角色,双方默契得完美无间。大多数时候,这种父—子或父—女的角色分配始终留在未言说之中,甚至留在无意识之中;但也有些时候,它被或多或少清楚地说了出来。在某些情形里,我甚至并没有进入什么游戏,却已经成了父亲的样子——我确信如此——对自己意识层面与无意识层面正在酝酿的事情都一无所知。

我第一次察觉自己扮演养父的角色,是在1972年,“生存与生活”时期,当时我突然面对一位年轻朋友的强烈排斥态度。(有趣的巧合:那是个脱离了常规的数学系学生!)我对第三者的某种举止让他失望了。我想,我本来会毫不困难地承认他的失望是有道理的,承认我在那件事上缺少慷慨——但那份反应的激烈,当时简直把我打懵了。那就像一阵突然腾起的猛烈仇恨之火,而它几乎马上又落了下去,因为已经清楚:他并没有真正把我掀下马。(只差一点点,但这一点我没有让他看出来……)。我不知道当时怎么会直觉到:他把我这个人——被他恰如其分地理想化了的——当成了他与父亲之间未解决冲突的投射对象。这个突然的直觉后来被忘掉了,却没能阻止我在此后许多年里仍然带着同样的确信继续进入父亲的角色,丝毫没有提防。当然,后来当我面对冲突的征兆——阴险的或激烈的——时,也总是同样痛苦的惊讶,不敢相信自己的眼睛,也不敢相信别的。

正是经过一场持续六七个月的、关于我父母生平的紧张独处工作——它让我在一种意想不到的光照下看见他们这个人——我才明白:这种养父的角色是虚幻的;它要替代一个确实存在、却被宣布为(哪怕只是通过默契)“失职”的真正父亲(而且要替代得更好,这是预先就说定了的!)。那是在帮别人回避冲突,把冲突从它所在的地方——比如说,在他与父亲的关系里——转移到第三者(此处即我本人)身上,而第三者与它全然无关。自从那场沉思(它发生在1979年8月到1980年3月)之后,我对自己保持警觉,不再闭着眼睛听任自己那份不合时宜的父性天职发作。这并没有阻止那种虚假局面重演(比如我与那个不得不停止合作的学生之间的关系)——但我相信,如今已没有我这方面的默契。

如果把那个合法期待落了空的学生放在一边,那么对我来说毫无疑问:在我遭遇学生或旧学生对抗的所有其他情形里,那都是同一种“与父亲的冲突”原型的重演:那个既受仰慕又受畏惧、既被爱又被恨的父亲——那个要去对抗、去战胜、去取代、也许还要去羞辱的“他”……但也是那个你暗中想成为的人——夺走他的一股力量,据为己有——另一个自我,被畏惧、被憎恨、被逃避……

\section{(30) 敌父(2)}

在某些旧学生与我之间制造出对抗的,并不是1970年那个大转折——当时的背景是一片牧歌式、无云的过去。它只是让一些对抗变得可见,而那些对抗本来很难在一种典型的、更为循规蹈矩的老板—学生(或前老板—前学生)关系框架里表达出来。我猜想,这类冲突在科学界想必并不罕见,只是它们通常表现得比我所卷入的那些关系更迂回、更难辨认。

回想起来,我终究不觉得在与学生的那些关系里,我曾那么倾向于进入父亲的角色——甚至,我抓不住一个哪怕多少朝着那个方向的记忆。

% ---------------- p.179 ----------------
\origpage{71} 就我这个人而言,我似乎觉得,我投入一段师生关系的几乎全部精力,就是我也投入数学、投入实现一个宏大纲领的那份精力。在第一个时期,我只看到一个例子:我身上对某个学生这个人的兴趣——一种亲近或好感的性质——其力量堪与(如果不说是等同)数学兴趣相比。但即便在那个例子里,我也不觉得我对他进入了父亲的角色。至于我对他、或对其他学生,在某个层次上可能施加的影响,那类事情在我与学生的关系里是我从不留意的。(即便今天,我也倾向于不留心,无论对近几年跟我做过工作的学生,还是对别的人。)当然,在所有这些情形里,学生与我之间的关系绝不是“对称”的:至少在师生关系持续的那段时间里(而且多半还要更久),一个学生在我生命里的分量,与我必定在他生命里占据的分量无法相比;关系在我身上和他身上所调动的心理力量也无法相比。除了那五六个以清晰可辨的对抗征兆表现出这些力量的情形,我意识到:在二十多年教书生涯里,我那些不同的学生、后来是旧学生与我之间关系的性质,对我来说始终是个完全的谜!再说,探究这些谜也不是我分内的事,那更多是各人自己那一份的事。不过,既然关心起自己这个人来,也许还有些比“他与旧老板关系的来龙去脉”更灼人的东西可以看。无论如何,尽管我在学生面前没有表现出任何进入父亲角色的倾向,我多半还是常常在他们眼里多少充当了养父——鉴于我前面谈过的自己那份特殊的心理“剖面”,也鉴于一种局面的内在动力:在那种局面里,我起码不能不显出一个长辈的样子。

无论如何,在我提到的几个例子里,学生与我之间关系的这种特殊色调,在我看来毫无疑问。在职业生活之外,还有许多别的例子:不管有没有我的默契,我显然在比我年轻的男人或女人面前充当了养父——他们被这个人吸引,首先与我以相互的好感相连,而绝无亲缘关系。至于我自己的孩子们,我对他们的父性之弦是强烈的;从很小的时候起,他们在我生命里就占据着重要的位置。奇怪的讽刺是:我的五个孩子,没有一个接受“有我做父亲”这件事。在我能够近距离了解的四个孩子的生活里——尤其是近几年——他们与我关系中的这道分裂,是他们自身内在一道深刻分裂的反映;尤其是一种拒绝:拒绝他们身上一切与我这个父亲相像的东西……但这里不是探究这道分裂之根的地方:那些根既扎在一段被撕裂的童年里,也扎在我的童年、我父母的童年里;同样也扎在母亲的童年、她父母的童年里。这里也不是衡量其后果的地方——在他们自己的生活里,或在他们孩子的生命里……

\section{(31) 使人泄气的权力}

为了结束这趟粗略的巡游——穿过1948到1970年间我在数学界里的种种关系——我还得谈谈我与更年轻的数学家的关系:他们或多或少是新手,因而严格说来没有“同事”身份,而我也不曾对他们扮演“老板”的角色。所以,这里说的是些年轻研究者:我在IHÉS的讨论班上与他们相遇一两年,或在哈佛及别处的这类课程、讨论班上,有时也通过
%
% ---------------- p.180 ----------------
\origpage{72}\origpage{73} 通信——比如我收到一位年轻作者寄来的工作,他期待我给出评论,肯定也期待鼓励。

与初入门研究者的关系,属于一个不如“学生们的老板”那么显眼的角色,但同样重要——这是我后来才明白的。那个年代,我并不像这六七年以来那样意识到:对一位有名望的数学家来说,这个角色代表着一种相当大的权力。首先是鼓励、激励的权力:它既存在于明显出色的工作那里(尽管那工作也许因表述笨拙或“行内功夫”不足而吃亏),也存在于仅仅扎实的工作那里;它甚至存在于那种只代表一份极微薄贡献的工作那里——按一位正当壮盛、手段强劲、对主题经验老到、见闻广博的前辈的标准,那贡献甚或微不足道,乃至等于零。只要送到我们手里的工作是认真写成的——这通常从头几页就看得出来——鼓励的权力就在那里。

而使人泄气的权力同样存在,而且无论面对什么样的工作,都可以随意行使。这就是柯西对伽罗瓦、高斯对雅可比行使过的那种权力——它可不是从昨天才有的,杰出而令人畏惧的人物一直都在使用它!历史之所以把这两桩事传给我们,是因为付出了代价的这两个人有足够的信念和自信继续走自己的路,尽管那些当时在数学界呼风唤雨的人怀着的权威毫无善意。雅可比找到了一份杂志发表他的想法,伽罗瓦找到的则是他最后一封信的纸页,权当“杂志”。

今天,对一个无名或不太有名的数学家来说,要让人知道自己,肯定比上个世纪更难。而有名望的数学家的权力,不只在心理层面,也在实践层面。他有权接受或拒绝一篇工作,也就是说:给不给它发表的支持。对也好错也好,我觉得“在我那个时代”,五六十年代,拒绝并非不可上诉——如果工作呈现出“值得注意”的结果,它还有机会得到另一位权威的支持。今天肯定不再如此:如今甚至很难找到一位有影响的数学家,肯(以他乐意采取的心情)去通读一篇属于他专业领域的工作——如果作者尚未获得名声,或者没有一位他认识的同行推荐。

近几年,我曾见有影响、出色的数学家行使他们使人泄气的权力,拒绝这个、拒绝那个:既针对某些显然应当做出来的扎实工作,也针对某些规模可观、清楚显示出作者力量与原创性的工作。有好几次,这样行使自己任意权力的人,竟是我的一个旧学生。这无疑是我作为数学家的生涯里经历过的最苦涩的经验。

但我又跑题了!我本来要考察的是:在我诚心诚意地扮演“知名数学家”角色的那些年月里,我是怎样使用自己手中那份鼓励人和使人泄气的权力的。我还该补充:在1970年以后我的科学活动所处的那个更卑微的层次上——在外省大学里作为普通教师中的一员——这种权力并没有因此就不存在,无论是对我的学生,还是(说实话,很少)对偶尔通信的人。但就我眼前的题目而言,只有我数学生涯的第一个时期才是要紧的。

至于我与学生的关系,从我带过的第一个直到今天,我相信可以毫无保留地说:我做了力所能及的一切去鼓励他们
%
% ---------------- p.181 ----------------
\origpage{74} 在自己选定的工作上往前走\fnoteP{12}{(23$^{iv}$)}{\emph{一次教学的失败(1)}\par 这些字句写下之后,我有机会与1970年之后的两位旧学生谈过,想同他们一起探究我在蒙彼利埃大学研究层次上教学失败的原因。他们对我说:我那种倾向——低估同化那些对我熟悉、对他们却不熟悉的技术对他们可能意味着的困难——在他们身上产生了令人泄气的效果,因为他们感到自己始终达不到我对他们的期望。此外(这件事在我看来意义还要大),他们有时感到受了挫伤:当我把袖子里藏着的一个成形陈述“卖给他们”、向他们泄底的时候,而不是留给他们乐趣,让他们在已经十分接近的时刻靠自己的办法把它发现出来。那之后,他们剩下的就只有做“练习”了(而那并不特别让他们着迷)——也就是证明那个陈述。我先前在一则注(注21)里指出过的自己身上那份“缺乏慷慨”,正落在这一处,当时没有就此多说。主要正是这类扫兴之事,构成了我在他们两人身上——在那么出色的开端之后——研究兴趣消失这件事上的个人贡献。\par 我意识到,我在1970年以前并不比之后更慷慨。如果说那时我没有遇到同样的困难,那无疑是因为:那个时期来找我的那类学生,有足够的动力,甚至能从一场“长长的练习”里品出魅力——那是学习本行、并一路学到一大堆东西的机会;同样,对于一个我“泄了底”的起步陈述,他们也能靠自己的办法从中抽出一大群远远超出第一个的陈述。当我换了教学地点,我在向新学生提议的思考主题上做了必要的调整,办法是选择那些可以被直接直觉把握的数学对象,不依赖任何技术行囊。但这必要的调整本身并不够,因为禀性上的差别(我的新学生与从前的学生相比)比单纯的行囊差别还要重要。这还印证了前面(第25段开头)说过的一点:我在“导师”这个角色上有某种不足;这不足在我教学生涯的第二个时期,比在第一个时期显露得强烈得多。} (23'''')。即便在今天,老板与学生的关系里恐怕也很少是别的样子,尤其当这位老板拥有手段,能够培养出色的学生,并借他们的协助开垦大片待耕的荒地。几乎难以置信、然而千真万确的是:甚至存在这样一种极端情形——一位声望卓著的老板,乐于把曾经激励过他自己年轻时代的数学激情,在才华出众的学生身上熄灭掉。

但我又跑题了!现在要考察的,是我与那些不是学生的年轻研究者的关系。在那样一些关系里,名人身上的自我中心力量,不那么容易把他推向鼓励的方向,因为那个找上他的无名年轻人的成功,对他的荣耀毫无增益,或增益甚微。恰恰相反,我认为:在缺乏真正善意的情况下,仅凭自我中心力量的运作,几乎无一例外地会把人推向相反的方向——去行使使人泄气、拒绝的权力。在我看来,这恰恰就是那条在社会一切部门都能看到的普遍法则:证明自己重要的自我中心欲望,以及满足它时伴随的秘密快乐,通常在这样一种情形下更强、更受享用——手里的权力找到机会去造成别人的失望,甚至羞辱,而不是相反。这条法则在某些特殊处境里表现得格外粗暴,比如战争,比如集中营的世界,监狱或精神病院的世界,或者干脆就是我们这样一个国家里任何人都可能进的医院……但即便在最日常的处境里,我们每个人都遇到过印证这条法则的态度与行为。对这些态度的矫正,首先来自文化的矫正:在某个圈子里,关于什么算“正常”或“可接受”的行为有一种共识;其次则是非自我中心性质的力量,比如对某个特定的人的好感,或者有时,一种自发的善意,甚至与它所指向的人无关。这样的善意无疑是稀罕之物,无论你去哪个圈子里找。至于数学圈里的文化矫正,在我看来,在过去的二十年里已经被大大侵蚀了。至少在我所熟悉的
% ---------------- p.182 ----------------
\origpage{75}\origpage{76} 那些圈子里,情况肯定如此。

我执意跑题,真是没办法!我本来要说的不是一篇关于这个世纪的议论,而是一场对我自己、对我与那些或多或少刚刚入门、又不是我学生的研究者之间关系的沉思。我不认为我提到的那条“法则”在那些关系里找到过表现的机会。由于一些不必在此考察的原因,看来自我中心的力量,虽然在我身上和在任何人身上一样强,在我一生中却没有走那条路去损害别人(除了几桩追溯到童年的情形)。我甚至敢说——既然有机会考察这件事——我面对他人的基本调性是一种善意的调性:能帮的时候就帮,能缓解的时候就缓解,有能力鼓励的时候就鼓励。即便在与那位“不知疲倦的朋友”如此深刻分裂的关系里,我身上的自负也从未把我误导到这样的地步:让我想到(哪怕是无意地)去伤害他。(我本来是有可能那样做的,而且当然会“怀着世上最安宁的良心”。)我相信,在大多数情形里,这种普遍的善意(哪怕只是浮在皮肤表面的一点点)也标记了我在数学界里的关系,包括与那些虽不在学生之列、却可能需要我支持或鼓励的初入门数学家。

我相信,至少在五十年代、直到六十年代初,无一例外都是如此。我觉得至少在那个时期,这份善意并不限于像广中平祐(Heisuke Hironaka)或迈克·阿廷(Mike Artin)那样明显出色的年轻人(当时还没有任何名声证明他们的本领)。但它在六十年代也许或多或少地淡去了,在自我中心力量的作用下。关于这一点,若有任何见证传到我这里,我会格外感激。

我的记忆只交给我一个确切的例子,我这就说;而在这个例子之外,就是那团著名的“雾”——它不凝结成任何别的确切事例或事实,却交给我某种内心的态度。当另一位数学家“踩到我的地盘上”、却装作不问我一声,好像他在自己家里似的——这个毛头小子!——我会感到某种恼火。那多半是年轻人的情形:不大在行,却碰巧重新找到一些我多年来、而且是从高处就知道的东西,有时还偏偏在一些很特别的情形里。我想,这种事不至于常常发生,但也许还是有两三次,也许是四次,我说不准。如我刚才所说,我只记得一个具体的例子,也许是因为那局面与同一位年轻数学家反复出现过好几次,形式各不相同。我可以说:这位年轻研究者——他所属的大学在国外——在各方面都极其得体:他把自己刚做好的工作寄给我,而我理应是消息最灵通的人。每一次,我都反应得极为冷淡,理由就是我说的那个。我甚至说不准自己有没有坦率地告诉他:他做的东西我早就知道了,因此他不至少在前言里向我行个小小的屈膝礼就发表,让我很不痛快。当然,假如他是我的学生,这种作者的虚荣就不会起那么大作用:一方面因为与学生之间已经建立起好感,另一方面因为无论如何不言而喻,学生的工作里也有老板的想法,除非另有说明!我相信这局面与同一位研究者发生过两次,也许甚至三次;每一次我的态度都一样冷淡,一样令人泄气。如果我记得不错,我从未答应过推荐这位研究者的工作到某家杂志发表,也从未答应参加他的论文答辩委员会(我记得这个问题曾被提起过)。几乎可以说,我像是打定主意挑他当出气筒。最妙的是,
%
% ---------------- p.183 ----------------
\origpage{77} 他的工作每一次都完全站得住脚——我相信写得用心,我也没有任何理由假定那些想法不是他自己找到的:那些想法当时还不怎么满街乱跑,只是(或多或少)“众所周知”于一小撮在行的人之间,比如塞尔、卡蒂耶、我,还有一两个人。我不明白的是:这位年轻同事(他当然最终有了论文,也有了一个受之无愧的职位)怎么就没厌烦一次次找上我——那个每次“给他冷脸”的人,而且他显然从未为此怨过我。我还记得有一次,他对我的迟疑表示惊讶;显然他不明白发生了什么。要是他等着我的解释,那可有他难受的!他有一张好看的脸,有点像古典希腊的,非常年轻——线条相当柔和、安详,透出一种内在的平静……现在,我第一次试着去捕捉他这个人和他的相貌所散发的印象,忽然意识到:他真的很像我说过的那位“不知疲倦的朋友”;他们本可以是兄弟——这位与我同龄的朋友,调子是含笑的那种;而这位研究者,比他年轻二十岁,调子略略低沉,却绝不忧伤。这种相似未必没有起过作用:我把一份在另一个人身上找不到机会表达的轻蔑,投射到了这一位身上——因为另一个人以那样忠实的友谊征兆把我解除了武装!而我也确实得长出一副多么厚的甲壳,才没有被这个必定招人喜爱的年轻人的明显诚意和把事情做好的意愿所打动:他不厌其烦地再来,而我连一个微笑都不肯赏他!

\section{(32) 数学家的伦理}

我昨天讲的那个例子,现在我终于肯费心把它白纸黑字记下来,它在我看来具有相当大的意义,在某些方面比我先前讲到的另外三个例子(无疑同样典型)更大——在那三例里,自负的力量深深地扰乱了我身上一种自然的善意与尊重的态度。这一次,我动用了一个实实在在的权力位置(同时像所有人一样,装作不知道这份权力),用它去使一位善意的研究者泄气,拒绝一篇本应发表的工作。这就是人们所说的滥用权力。它并不因为不落入刑法某一条的制裁范围就不那么昭然。幸而当时的形势不如今天严酷,所以这位研究者能够——我相信,没费太大力气——在几位比我更善意的同行的支持下把工作发表出来;他作为数学家的生涯,也没有因为我那滥用权力的行为而受到严重扰乱,更不用说被毁掉。事后我为此高兴,却并不想把它当作“减轻情节”。在更严酷的形势里,我也许会更当心——但这只是一个推测,在这里没有多少关系。不过我仍相信自己可以说:我心里并没有暗藏的恶意,没有由我说的那种恼火引起的害人之心。我对那恼火的反应是“发自肺腑的”,对自己没有丝毫批评的意愿,更没有丝毫意愿去看一眼自己心里发生了什么,哪怕只是看一眼我的反应可能在另一个人的生命里有什么分量。我没有衡量自己手中的权力;与这权力相随的责任这个念头,在那段关系里从未掠过我的心头。这是一桩不负责任行为的典型,就像街头巷尾随处可遇的那样,在科学界和在别处一样。

我记住的这个同类情形中唯一的例子,也许只是几个类似情形中的一个极端例子。引发一种毫无善意的态度的,是虚荣心受到冒犯后的恼火:它
%
% ---------------- p.184 ----------------
\origpage{78}\origpage{79} 容不得“随便哪个初来乍到的人”僭取权利,闯进禁猎区,打走一点只属于这些地盘之主的小猎物……这种恼火有的是现成的合理化,而且派头更高贵,这可想而知。被摆出来的不是我这卑微的个人——不!——而是对艺术、对数学的爱;这个年轻人连“是个天才”这条借口都没有,倒更像笨手笨脚的那一类,他会把一切都弄糟,我们可倒霉了;他要是把事情做得比我好倒也罢了,可我筹划的那些漂亮安排全泡汤了,说真的,总得有点不知羞才行……!而始终如水印般透出的,是那支功绩主义的主导动机:只有最好的那些人(比如我)在我的地盘上才有居留权,或者那些把自己置于这类人中某一位保护之下的人!(至于那种不那么常见的情形——确实是另一位大人物踩到我的地盘上——那又是另一回事了:一天自有一天的烦恼!)在这个具体例子里,还有(对此我如今已不大有疑问)另一股朝同一方向起作用的力量,而它完全是无意识的;这股力量在我与起步时那位不知疲倦的朋友的关系里,已经起过很强的作用:对某一类人的排斥自动机制——那类人不合我从母亲那里接过来的“男子气”标准。但这个情形,虽然对理解我自己有它的意义和趣味,与我眼前的题目却关系不大:我的题目是要在我自己身上,在我还属于某个圈子时曾有过的态度和行为里,找出我今天在那个圈子里所看到的深刻败坏的典型征兆。

我刚考察的这个例子,之所以在我看来比我其他那些缺乏善意与尊重的例子意义更大,是因为正是在它这里,数学家这一行某种基本的伦理遭到了侵犯\fnoteP{13}{(24)}{我想说的这种伦理,同样适用于任何别的围绕研究活动形成的圈子;在那里,把自己的结果公之于众、并从中获得承认的可能性,对每个成员的社会地位、甚至对作为该圈子成员的“存活”,都是“生死攸关”的问题,连带对他和他的家庭意味着一切后果。} (24)。在我起步时接纳我的那个圈子里——也就是布尔巴基及其亲近者的圈子——我想说的这种伦理大体上始终是隐含的,但它毕竟在场、有生命,而且(在我看来)是不可触动的共识的对象。据我记忆所及,唯一把它用清楚明确的话说出来的人是迪厄多内,那无疑是我在南锡做客的最初几次之一。他也许在别的场合又回到这一点。显然他觉得这是一件重要的事;而我当时必定感受到了他对此的看重,因为三十五年后的今天我还能记得。仅仅凭借我那些年长的前辈这个群体的道德权威,以及显然表达了群体共识的迪厄多内,我必定已经默默地把这种伦理认作自己的,却从未给它一刻思考,也从未明白它的重要性何在。说实话,我甚至不会想到:对它加以思考也许有用——因为我早已确信,我的父母和我自己这个人,各自都是伦理态度的完美化身(或者相差无几),负责、全面,而且经得起一切考验\fnoteP{14}{(25)}{\emph{职业伦理的共识——兼及信息控制}\par 除了与迪厄多内的那次谈话,在我作为数学家的生涯里,我记不起有哪次谈话——我参与的或旁听的——谈到这一行的伦理、谈到行业成员之间关系中的“游戏规则”。(这里我排除那些关于科学家与军事机构合作的讨论,它们发生在七十年代初“生存与生活”运动前后。那些讨论其实并不涉及数学家之间的关系。我在“生存与生活”里的许多朋友,包括谢瓦莱和盖吉,都觉得:我当时——尤其起初——把重点放在这个我格外敏感的问题上,使我远离了更本质的日常现实,正是我在眼下这份反思里考察的那类现实。)在学生与我之间,从未谈过这些事。那种心照不宣的共识,我想只限于这一条规则:不把别人的、自己可能已经知道的念头说成是自己的。这共识,在我看来自古就存在,而且至今没有任何科学圈子里的人质疑过。但若缺少另一条补充规则——保障每个研究者有可能让人知道自己的想法和结果——第一条规则就只是一纸空文。在今天科学界,处于声望与权力地位的人,对科学信息握有一种任意处置的控制权。在我所熟悉的那个圈子里,这种控制不再被迪厄多内所说的那种共识所缓和;而那共识,也许从未存在于他充当代言人的那个狭小团体之外。处于权力地位的科学家,几乎能收到他认为值得收到的一切信息(往往还超出),而且对其中很大一部分,他有权阻止发表,同时坐享其利,把那些信息斥为“没有意思”、“多少已知”、“trivial”等等……我在注(27)里还要回到这个局面。} (25)。

% ---------------- p.185 ----------------
\origpage{80} 再说,迪厄多内并没有对我长篇大论——那既不是他的风格,也不是他布尔巴基的任何朋友们的风格。他多半是顺口提起,像提起一件理应不言自明的事。他只是强调一条再简单不过、表面上全然无害的规则,就是这样:任何人发现了一个值得注意的结果,就应当有权利、也有可能把它发表,唯一的条件是这结果尚未成为某份出版物的对象。所以,即便这个结果为一人或数人所知,只要他们不曾费心把它白纸黑字写下来并发表,从而使它可供(哼!)“数学共同体”使用,那么任何别的人(言下之意:包括那位著名的“初来乍到者”)用自己的办法找到这个结果(言下之意:不管他用的是什么办法、什么观点和光照,也不管这些在那些理应比他更在行的人看来是否“褊狭”……),就应当有可能按自己的办法和光照把它发表。我记得迪厄多内还补充说:这条规则若不被遵守,就会给最恶劣的滥用敞开大门——很可能正是在那个场合、经他的口,我才知道了高斯拒绝雅可比工作那桩历史案例——借口是雅可比的想法他早已知道。

这条如此简单的规则,是对迪厄多内身上(以及布尔巴基其他成员身上)那种“功绩主义”态度——正如在我身上一样——的根本矫正。遵守这条规则,是一种廉洁的保证。我很高兴能够说,从至今传到我这里的一切来看,这种根本的廉洁在布尔巴基最初群体的每一位成员身上都保持完好\fnoteP{15}{(26)}{布尔巴基的“创始成员”是亨利·嘉当、克劳德·谢瓦莱、让·德尔萨特(Jean Delsarte)、让·迪厄多内、安德烈·韦伊。他们全都健在,只有德尔萨特例外——他在五十年代英年早逝,而那个时期,这一行的伦理总体上还受到尊重。} (26)。我注意到,对于曾经属于布尔巴基群体或圈子的别的数学家,情况并不都是这样。而在我自己这个人身上,它并没有保持完好。\par{}重读本文时,我曾想删掉这一段:在里头我可能给人这样的印象,像在颁发“廉洁”(或不廉洁)证书,而当事人对此毫不在乎,也不是我该做的事。这段文字可能引起的保留,肯定是有道理的。但我还是把它留着,为的是见证的真实性,也因为这段文字确实再现了我的感受,即便这些感受并不合宜。

迪厄多内用再平实不过的话向我讲的那种伦理,作为一种特定圈子的伦理,已经死了。或者说,那个圈子本身也死了,与那构成它灵魂的廉洁一同死去。这份廉洁在某些孤立的人身上保存了下来;在另一些它曾经败坏的人身上,它重现了,或者将会重现。它在某个人身上出现或消失,是各人精神历险的关键插曲之一。但这历险得以展开的舞台,已经深刻地改变了。一个曾经接纳我、被我认作自己的、我曾暗自引以为傲的圈子,不复存在了。使它有价值的东西,在我自己身上死去了,或者至少被另一种性质的力量侵入并取代——而那远在规范它的默会伦理于使用惯例与信仰表白中被公开否弃之前。如果我此后能够惊讶、能够愤慨,那是出于故意的无知。从这个曾属于我的圈子里回到我身上的东西,本来有一个关于我自己的消息要带给我,而我乐意把它一直回避到今天。

% ---------------- p.186 ----------------
% (本页仅存书眉“8. Maître et élèves”,无正文)


% ================================================================
%  第 9 章 收获 (原书 p.187-196, 边码 p.81-p.94) —— 前半
%  源文: _work/src1/task-ch09a.md
% ================================================================
\chapter{收获}

% ---------------- p.187 ----------------
\section{(33) 笔记——或新伦理}

\origpage{81}当然,一条职业伦理的规则,只有凭一种内在的态度才有意义,那态度就是它的灵魂。它无力创造它所力求表达的那种尊重与公正的态度,最多只能在一个普遍认可这条规则的圈子里,帮助这样的态度维持下去。一旦缺少了内在态度,哪怕这条规则还挂在人们嘴边,它也已失去全部意义、全部价值。

任何注释,无论多么一丝不苟、多么细致入微,也改变不了这一点。

不久前,我昔日的朋友与伙伴中有一位好心地向我解释:如今这年头,唉,数学产出涨得没边儿(这是谁都知道的),“人们”实在不得不——愿意也好,不愿意也罢——在写出来、交上来准备发表的文章里做一番严厉的筛选,只发表其中很小的一部分。他说这话时一脸真诚的遗憾,仿佛他自己也多少是这不可避免的宿命的受害者——同时,他那神情里还带点这样的意思:他自己就属于“法国那六七个人”(唉,是不幸,可事情就是这样!),由他们决定哪些文章能发表、哪些不能。我上了年纪,话少了,便只是默默听着。这个题目上有许多话可说,但我知道那会是白费口舌。一两个月后我还得知,这位同行几年前曾拒绝推荐一篇投给《法国科学院院报》(CR)的笔记,那篇笔记的作者和题目(是我七八年前向他提议的)都让我放在心上。作者为展开这个题目花了两年工夫,这题目确实不时髦(尽管我始终觉得它同样切题)。我认为他做出了出色的工作(作为第三阶段博士论文提交)。我不是这位年轻研究者的“老板”——他碰巧天分极高(我不知道他是否还会在数学里施展这些天分,看看他所受到的接待……),他的工作是在与我毫无接触的情况下做出来的。但也确实,这个题目的来源不可能有任何疑问;这可怜人一开始就处境不妙,而且他肯定还蒙在鼓里!这位同行倒还做得有礼有节,至少还有这点,我也没指望他更多,“真诚地抱歉,可您也明白……”。两%
% ---------------- p.188 ----------------
\origpage{82}年工夫——一个干劲十足的新手研究者的两年工夫,换来的却是一篇三页的《院报》笔记——那篇笔记又能花掉多少公款呢?这里有一种一眼就能看出的荒谬:两者之间那种巨大的不相称。只要肯费心去查一查深层的动机,这种荒谬肯定就消失了。只有这位同行兼旧友才能探测他自己的动机,正如只有我才能探测我的。但不必走很远我就知道,需要照顾的并不是数学产出涨得没边儿,您懂的,也不是公款(或《院报》那位想象出来的“无名读者”的耐心)……

这同一篇预备投给《院报》的笔记,先前还有幸被送到“法国那六七个人……”中的另一位手上,那位把它退给了作者的“老板”,因为这些数学“不叫他觉得好玩”(原话如此!)。(老板虽然恶心,却很谨慎,自己地位本来就不太稳,两次都宁可咽下去,也不愿得罪人……)后来我有机会跟这位同行兼旧日学生谈起此事,才知道他当初还是费心仔细读了送来的笔记,也想过一想(那笔记想必勾起他不少回忆……),并且觉得其中有些陈述本可以用对使用者更友好的方式写出来。然而他并不屑于把宝贵的时间浪费在把意见转达给当事人上:大人物的一刻钟,对一个无名年轻研究者的两年工夫!数学确实叫他“觉得好玩”到足以让他抓住这个机会,重新接触笔记里研究的那种情形(那情形一定在他心里、也在我心里,唤起了各式各样丰富的几何联想),把给出的描述吸收进去,然后——以他的家底与本事,毫不费力地——挑出其中的笨拙与疏漏。他没有浪费时间:他对某个数学情形的知识,由于一个初出茅庐的研究者两年的认真工作,变得更精确、更丰富了;那工作,大师当然几天之内就能做完(只做大略,不做证明)。这些一到手,人就想起来自己是谁了——案子已经判了,无名先生两年的工夫,正好进垃圾桶……

有些人,这种风一吹过来,什么也感觉不到——可我今天仍然被它吹得透不过气。在那件事里,这肯定正是所追求的效果之一(瞧瞧那拒绝所披的漂亮外衣),但肯定不是唯一的效果。就在同一次谈话里,这位旧友带着一种谦逊的得意向我吐露:他只有“所陈述的结果叫他吃惊,或者他不知道该怎么证明”时,才肯把一篇笔记推荐给《院报》\fnoteP{1}{(27)}{\emph{年轻人的“势利眼”,或纯洁的捍卫者}\par Ronnie Brown 向我转述了 J.H.C.怀特黑德(J.H.C. Whitehead)的一句话(布朗是他的学生),说的是“年轻人的势利眼:他们以为,一个定理的证明若是 trivial,定理便是 trivial 的”。我昔日的许多朋友,不妨好好琢磨琢磨这话。这种“势利眼”在今天绝不限于年轻人,我认识不止一位有声望的数学家,就常常如此。我对此格外敏感,因为我在数学里(别处也一样……)做得最好的东西,那些我认为最富成果的概念与结构,以及我凭耐心而顽强的工作从中抽出的本质性质,统统落进了“trivial”这个形容词之下。(这些东西,若作者不是早已成了名人,如今哪一样也难被接受为一篇《院报》笔记!)我做数学家的毕生野心,或者不如说我的激情与快乐,一直是找到那些显而易见的东西;在本书里(包括在眼下这个导论性的章节里……)我唯一的野心也是如此。决定性的东西,往往首先在于看见那个尚未被看见的问题(无论答案如何,也无论答案是否已经找到),或者抽出一个陈述(哪怕是猜想)——它概括并包含一个尚未被看见、或未被理解的局面;至于它是否被证明,证明是否 trivial,完全无关紧要,是彻底的枝节,哪怕一个仓促的、暂定的证明后来被发现是错的。怀特黑德所说的势利眼,是那种饱食终日的浪荡子的势利眼:他非要先弄清一瓶酒贵得离谱,才肯承认它好。这些年里,不止一次,我被旧日的激情一阵阵攫住,捧出我最好的东西,却看着它被那种自命不凡打回来。我感到一种至今犹新的痛,一份快乐落了空——但我还不至于流落街头,而且幸而我并不想把自己写的一篇文章安插出去。\par 怀特黑德所说的势利眼,是一种滥用权力,一种不诚实,而不只是对事物之美的麻木或封闭——当它由一个有权势的人施加于一个任其摆布的研究者身上,那人完全有条件把后者的想法吸收过来、据为己有,同时却以“显而易见”或“trivial”、因而“没有意思”为借口,堵住它们的发表。我这里甚至没有想到通常意义上的抄袭那种极端情形,那在数学界想必还很罕见。然而从实际的角度看,对付出代价的研究者来说,情形是一样的,而使它成为可能的那种内在态度,在我看来也相去不远。只不过它更舒服些:它伴随着一种对他人无限优越的感觉,伴随着那种以数学不可触碰的纯洁之毫不妥协的捍卫者自居者的问心无愧与内心满足。} (27)。这大概正是他发表得很少的一个原因。假如他把自己的标准用在自己身上,%
% ---------------- p.189 ----------------
\origpage{83}他索性就什么也不发表了。(确实,以他眼下的处境,他根本不需要发表。)他什么都晓得,要叫他吃惊,大概就跟找到一条虽可证明、他却不知如何证明的命题一样难。(二十年里,这两样我大约只碰上过两三回,而且最近十年、十五年里一次也没有!)他显然为自己的“质量”标准自豪,这标准使他在数学家的行当里成了把要求推到极致的冠军。我在其中看到的是一种百毒不侵的自我陶醉,而且不止一次看到,在那种笑眯眯、随和厚道的外表后面,是对他人毫无节制的轻蔑。我也看得出,他从中得到了极大的满足。

这位同行的情况,是我在“新伦理”的代表人物中遇到的最极端的一例。但它仍然很典型。在这里,无论是我讲的那件事,还是使之合理化的那篇信念宣言,用常识来衡量都透着一种乌布式的荒诞——其尺度如此之大,以至于这位脑子如此出众的旧友,以及肯定还有许多地位不那么显赫的同行(他们只消不找他推荐笔记就行了),都已经看不见它了。因为要看,至少得去看。只要肯费心去看动机(首先是他自己的动机),荒诞就暴露在光天化日之下,而它们同时也不再是荒诞的——它们交出了自己那份朴素的、一目了然的含义。

如果说这些年里,面对某些态度、尤其是某些行为,我常常痛苦到这般地步,那肯定是因为我隐约在其中辨认出——像是被推到极致、直到可笑或可憎地步的漫画——一些曾经属于我自己的态度与行为,它们又通过我的某位旧学生或旧友回到我身上。不止一次,我心里那套老反射又启动了:去揭发,去同那被清清楚楚指认出来的“恶”作斗争——但即便我偶尔在这里那里顺从了它,那也是带着分裂的确信。说到底,我很清楚,去斗争,仍旧是在事物的表面打滑,仍旧是回避。我的角色不是揭发,甚至也不是“改善”我所身处其中的这个世界,或“改善”我自己这个人。我的天职是学习,是通过我自己去认识这个世界,并通过这个世界认识我自己。如果我的生命能给我自己或他人带来什么好处,那也只是在我懂得忠于这份天职、懂得与自己和谐的限度之内。是时候提醒自己这一点了,好切断我心里那些老机制——它们在这里想推着我去为某个事业辩护(就说是某种已死的伦理吧),或者去说服人(去说服人相信取代它的那种伦理大概“荒诞”吧),而不是去探测、去发现、去认识,或者把描述当作探测的手段。写前面那两三页时,我除了就今日取代昨日之种种的流行态度说几句话之外,并无更确切的目的,却始终感到自己在提防着自己,那种心境,就像一个人随时准备大笔一挥,把刚写完的东西统统划掉、扔进纸篓!不过我还是要把写下的保留下来:它并不假,可它终究造出了一种虚假的局面,因为我在其中牵连别人多于牵连自己。我心底感到,我写着的时候什么也没学到——肯定正是这一点造成了我的不适。总之,是时候回到一种更实在的反思了,一种教我自己的反思,而不是妄图去教别人或说服别人的反思\fnoteP{2}{(28)}{写前面那几页时,起先我夹在“把肚子里的东西倒空”的欲望与一种保留、一种谨慎之间。所以我一直停在差不多上,这肯定是我那份不适、那种“我什么也没学到”之感的主要原因。自从记下这份不适的那几行写出之后,我把这几页改写了两遍——它们此前让我内心不满——写得更清楚地把自己摆进去,也更往事情深处走。一路下来,我当真“学到了点什么”,而且我相信,同时我还成功地用手指触到了一件重要的东西,它既超出这一具体事例,也超出我自己这个人。} (28)。

% ---------------- p.190 ----------------
\section{(34) 淤泥与源泉}

\origpage{85}我觉得,大体上,我已经把自己同其他数学家的关系——各种年纪、各种等级的数学家,那是我还属于他们那个世界、属于数学家的世界的时候——都过了一遍;同时,而且尤其是,也把我以自身态度与行为在其中所占的那一份过了一遍:那是我在今天所看到的某种精神里所占的一份,而这精神肯定不是昨天才有的。在这番反思里,或者更确切地说,在这趟旅行里,我四次遇到一些情形,它们在我看来是我身上某些态度与暧昧的典型:在那里,对他人那种自发的善意与尊重,被一些自我性的力量搅乱了,甚至一扫而光,尤其是(至少在这四例中的三例里)被一种自得搅乱了。这自得主要依仗的是这样一种所谓优越:某种脑力,加上我在数学活动上那种不成比例的投入,据说赋予了我这种优越。它从一种普遍的共识中得到确认与支撑——那共识几乎毫无保留地推崇这种脑力和这种不成比例的投入。

在这几例情形中,最后一例——那个“踩进我的地盘的没规矩的年轻人”——我觉得对眼下的论述最为重要。前三例是我这个人的典型,或我这个人的某些侧面的典型,是在某个时期(也确实是在某种环境里)——但是,如我已经有机会一再说过的,我根本不认为它们是我当时所属那个圈子的典型。我也不认为它们是今天法国数学界的典型,这么说吧——比如说,我与“不知疲倦的朋友”之间那种关系所特有的那种慢性迷误,如今大概很少见,就像当年想必也很少见一样。相反,我在“没规矩的年轻人”这件事上的态度与行为,却是今天数学界里每天都发生的事的典型,无论往哪里看。有影响力的数学家对无名青年怀抱善意、给予尊重,如今倒成了极罕见的例外——如果这青年没福气做他的学生(即便如此也未必……),也不是一位地位相当的同事向他推荐的学生。这大概正是我1970年“觉醒”之后不久就已传到我这里的东西——那场觉醒解开了许多沉默者的舌头——但那时我听到的第一手证词对我仍然是遥远的,因为它们直接涉及的,既不是我本人,也不是我圈子里最亲近的朋友。从某个时候起(大约1976年),我开始受到不止是表面的触动:那时传到我这里的回声,或我亲眼目睹的事实,其主角是这些朋友中的某几位,甚至是一些变得显赫的旧学生;而当那些遭到恶意中伤的人是我很熟悉的人、不止一次是我的学生(1970年以后的学生,那是不用说的!),因而他们的命运牵动我的时候,就更是如此。在某些情况下,而且毫无疑问:那种缺乏善意,甚至那种炫耀式的轻蔑态度,至少因为这样一个事实而变本加厉——某位年轻研究者是我的学生,或者他冒着风险(不一定非是我的学生)去做我昔日的朋友和其他同行也乐于称之为“格罗滕迪克式的玩意儿”的事情……

70年代初,“那个没规矩的年轻人”又给我写过信,非常客气地问我(他本来完全没有义务问我任何事!)是否见怪:他想发表一个定理的证明,那定理人家告诉他是我的,却从未发表过。我记得我回信时还是像过去那样一肚子不快,大概既没说行也没说不行,还让人听出这样的意思——我没看他的证明(他本来%
% ---------------- p.191 ----------------
\origpage{86}\origpage{87}当然愿意给我看,可我全不在乎,我当时正忙着我的那些活动!):那证明对我的证明肯定毫无补益(可它至少有一点好处:白纸黑字写下来,连同陈述本身,可供数学界使用!)。这恰好说明,那场著名的“觉醒”还多么浮浅,对某些根植于自得、根植于“精英至上”态度的行为毫无影响——而同一时期,我肯定正在《生存与生活》(Survivre et Vivre)上那些情真意切的文章里、在公开辩论的发言里,揭发着这些行为……

这就很具体地回答了我先前悬置的一个问题。不如就在这里承认这条朴素的真理:这类自得的态度在我身上绝没有“一劳永逸地”被克服,我怀疑有朝一日会,除非到我死。如果说确有转变,那也不是因为虚荣消失了,而是因为(重新)出现了一种好奇——对我自己这个人、对我身上某些态度与行为的真实本性的好奇。正是凭这种好奇,我才变得多少能察觉自己身上的种种虚荣表现。这深刻地改变了某种内在的动力,也因此改变了“虚荣”的效果;所谓虚荣,就是那股常常推着我去掩盖或伪造我对现实那种健全而细腻的感知的力量,其目的是壮大我自己、把自己置于他人之上,同时却宣称恰恰相反。

也许某位读者会感到困惑,就像我有朝一日也曾困惑:在我作为数学家的生活里,虚荣那种阴险而顽强的存在(他或许有时也在自己的生活里瞥见过它),与我所谓的对数学的爱、或者说激情(它也许同样在他自己对数学、或对某个别人、别的什么事的体验里唤起回声),这两者之间有一种表面的矛盾。如果他确实困惑,那么他自身就具备所需要的一切,可以(像我从前那样)重新接触到事物本身的现实——那是他可以第一手认识的东西——而不是像一只被关在词语与概念的无尽笼子里的松鼠那样打转。

看见一汪浑水的人,难道会说水与泥是一回事?要认识那并非泥的水,只消上到源头,看一看,喝一口。要认识那并非水的泥,只消登上被太阳和风晒干的河岸,在手心里掰开、捻碎一团带沙粒的黏土。野心、虚荣,或多或少能决定一个人在生活中给某种激情——比如数学的激情——留多少地盘,能在回报丰厚时把它变得贪得无厌。但最贪得无厌的野心,本身却无力发现或认识哪怕最微小的一件事,恰恰相反!在工作的时候,当理解一点点开动、成形、加深;当混乱中一点点显出秩序,或者当原本熟悉的东西忽然露出陌生的面貌,继而令人不安,直到终于爆出一个矛盾,把一种看似不可动摇的对事物的看法掀翻——在这样的工作里,没有一丝野心或虚荣的痕迹。那时领着跳舞的,是某种来自比“自我”及其不断自我膨胀的饥饿(哪怕是“知晓”与“知识”的饥饿)远为深远之处的东西——肯定比我们这个人、甚至我们这个物种都远为深远。

那就是源泉,它在我们每个人心里。

\section{(35) 我的激情}

三大激情支配了我成年后的生活,此外还有一些性质不同的力量。我终于在这些激情中认出同一种深层冲动的三种表现;那是我心中认知冲动所走的条条道路中的三条,而在我们这个无限的世界上,呈现在它面前的路有无数条。

% ---------------- p.192 ----------------
\origpage{88}在我生命里最先出现的,是我对数学的激情。十七岁那年,刚出中学,我松开缰绳,任由一种简单的癖好去跑,它就展开成了激情,在随后的二十五年里支配了我生命的进程。我“认识”数学,远在我认识第一个女人之前(除了那个我生下来就认识的女人);如今到了成熟的年纪,我发觉它始终没有燃尽。它不再支配我的生命,我也不再声称要支配它。有时它睡去,有时我甚至以为它熄灭了,它却又悄没声息地回来,一如既往地火热。它不再像从前那样吞噬我的生命——从前我把生命交给它去吞噬。它继续在我生命里留下深深的印痕,就像他爱的女人在一个情人身上留下的印痕。

我生命里的第二种激情,是对女人的追寻。这激情常以追寻伴侣的面貌出现在我面前。直到这追寻快要结束的时候,我才懂得把两者分开:那时我知道,我所追的东西哪里也找不到,或者也知道:我就把它带在自己身上。我对女人的激情,要到我母亲去世之后才真正展开(那是我第一次恋爱五年之后,那次恋爱生了一个儿子)。就在那时,二十九岁,我成了家,又有了三个孩子。起初,对孩子的依恋是对母亲依恋不可分割的一部分,是那股从女人身上发出、把我吸引到她那里去的力量的一部分。那是这份爱的激情的果实之一。

我并没有把这两种激情在我身上的并存体验为冲突,无论起初还是后来。我必定隐隐感到了两者的深层同一,而它明确显现给我,要晚得多,是在第三种激情进入我生命之后。然而,这两种激情对我生命的影响,只能是大不相同。对数学的爱把我吸引到某个世界里,即数学对象的世界;那个世界当然有它自己的“现实”,可它不是人们的生活在其中展开的那个世界。对数学事物的切身认识,可以说没有教给我任何关于我自己的东西,更不用说关于别人的东西——朝向数学的发现之冲动,只能使我远离自己和他人。有时两三个人会在这同一种冲动中相通,但那是一种表面的相通,实际上使每个人既远离自己也远离别人。所以,对数学的激情在我生命里从来不是一种成熟的力量;我也怀疑这样的激情能促进任何人的成熟\fnoteP{3}{(29)}{我这里说的是对数学、或对另一种纯智力活动的紧张而长期的投入。反过来,这样一种激情的展开,可以是与我们身上一种被遗忘的力量重新相识的方式,是同一份不肯就范的质料较量一番的机会,也一路让我们借某种真正属于自己的东西,更新并充实我们的身份感——这样的展开,完全可能是内在旅程中、成熟过程中重要的一步。} (29)。如果说我长期以来给这激情在我的生命里留了如此不成比例的位置,那肯定也正因为:它让我得以逃避对冲突的认识,也逃避对自己的认识。

相反,性的冲动,不管我们愿不愿意,都把我们径直抛向与他人的相遇,径直抛进我们自身、也抛进对方身上那个冲突的结!而我生命里对“伴侣”的追寻,却是对无冲突的幸福的追寻——那不是认知的冲动、性的冲动,并不像我喜欢相信的那样,而是对认识他人身上和我自身之中的冲突的无尽逃避。(为了结束这虚幻的追寻,以及如影随形般伴着它的不安,这是我必须学会的两件事之一。……)幸好,人尽管逃,性总会很快把人带回来!

有一天,我不再拒不接受冲突通过我爱过的、或正爱着的女人,以及通过这些爱情所生的孩子,顽强地带给我的教训。当我终于开始%
% ---------------- p.193 ----------------
\origpage{89}\origpage{90}倾听、开始学习的时候,而且在那之后许多年里,事情一直是这样:凡我学到的东西,都是通过我爱过的或正爱着的女人学到的\fnoteP{4}{(30)}{近几年,接力的是我的孩子们,由他们来给一个有时不大情愿的学生讲授人类存在的奥秘……} (30)。直到1976年,四十八岁,对女人的追寻一直是我生命里唯一巨大的成熟之力。这成熟之所以只是在此后几年、也就是最近七年才完成,是因为我用尽手边的一切手段来保护自己、不去成熟(这手段我是从父母那里、从我认识的那些圈子里学会的)。这些手段中最有效的,就是我在数学激情上的投入。

第三种大激情进入我生命的那一天——1976年10月某个夜晚——学习的大恐惧消失了。那也是对再寻常不过的现实的恐惧,首先是对关乎我这个人的那些朴素真理的恐惧,或对关乎我所珍爱的人的真理的恐惧。说来奇怪,四十八岁之前,我从未在自己身上察觉这恐惧。我是在新激情、认知激情的新表现出现的那一夜发现它的。可以说,后者取代了那终于被认出的恐惧的位置。多年来,我在别人身上把这恐惧看得清清楚楚,却有一种奇怪的盲目,使我对自己身上的它视而不见。对看的恐惧,使我看不见这同一种对看的恐惧!我和所有人一样,牢牢依恋着某种关于自己的形象,那形象大体上从童年起就没有动过。我说的那一夜,也是这旧形象第一次塌下去的一夜。另一些与它相像的形象接替了它,凭着顽固的惯性维持几天、几个月,甚至一两年,然后在审视的目光下也塌下去。懒得去看,常常推迟这样一次新的觉醒——但看的恐惧再也没有回来。凡有好奇之处,恐惧便再无容身之地。当我对自己怀着好奇的时候,我并不更害怕自己会发现什么,就像我想知道一个数学情形的底细时一样:那时有的是一种快活的期待,有时急切,却又执拗,准备好迎接一切愿意来到面前的东西,无论预料的还是意外的——一种满怀激情的专注,守候着那些毫不含糊的迹象,凭它们在假、半真与也许的初始混乱中认出真来。

在对自己好奇之中,有爱;这爱不为任何恐惧所扰——不担心我们所看的东西不合我们所愿。说实话,对自我的爱早已在这一夜之前的几个月里静静地绽开;这一夜也是这爱取得行动形式的一夜,可以说,是变得敢作敢为的一夜,毫不客气地把行头与布景撞翻!如我所说,不久又有别的行头和布景像变戏法一样冒出来,然后轮到它们被撞翻,没有咒骂,也没有咬牙切齿……

这七年里,新激情在我生命中的种种表现,终于在我看来像是一浪接一浪起伏的波涛,像一次辽阔而安详的呼吸的一进一出。这里不是试着去描出它那条蜿蜒变化之线的地方,也不是去描出作为对位的数学激情表现之线的地方。我已放弃去安排这一种或那一种的进程——如今倒不如说是两者的这双重运动在安排我生命的进程——或者更确切地说,它就是那进程本身。

在新激情出现之前的那些月份里——孕育与丰满的月份——对女人的追寻已开始改换面貌。它开始摆脱曾浸润它的那份不安,就像一口气息终于从压迫它的重负下解脱,重新找回它自己的幅度与节律。又像一团火,因为无处出气而半明半灭地闷烧着,一阵凉风吹来,忽然舒展成噼啪作响、轻盈%
% ---------------- p.194 ----------------
\origpage{91}而活泼的火焰!

火已烧足。一种看似无法止息的饥饿得到了饱足。两三年来,那场追寻似乎确已烧尽,不留一丝灰烬,把场地让给两种激情的歌唱与对唱。其一,我少年时代的激情,三十年间一直用来让我与一个被否弃的童年分离。其二,是我成熟年纪的激情,它让我重新找回了孩子,也找回了我的童年。

\section{(36) 欲望与冥想}

我说过的那一夜,一种新激情取代了一种永远消逝的旧恐惧,也是我发现冥想的那一夜。那是我第一次“冥想”之夜,它是在一种专横而急迫的需要压迫下出现的——此前几天,我仿佛被一阵阵焦虑的浪头淹没了。也许像一切焦虑一样,那是一种“起飞”式的焦虑:它执拗地向我指出:关乎我这个人的那份朴素而显眼的现实,已经同那个维持了四十年、我却从未怀疑过的自我形象脱节。肯定有一份巨大的认知渴望,与那些可观的逃跑力量并排着,也与那种逃避焦虑、想跟从前一样安宁的欲望并排着。于是有了一场紧张的工作,持续了几个小时,直到它有了结局,而我当时还不知道正在发生的事意味着什么,更不知道我要往哪里去。在这场工作中,一个个遁词被认了出来;或者更确切地说,正是这场工作让这些遁词一个个现形,每一个都表现为一种内心确信——我终于肯费心把它白纸黑字记下来,仿佛这样才更能浸透它,而它此前一直待在一种有利的模糊里。我满心欢喜地记下它,毫无提防:它必定有本事叫我着迷——那时的心境,就是一个什么都不怀疑的人的心境,对他来说,把一条未说出口的确信白纸黑字写下来,这本身就不可辩驳地证明它真实,证明它有根有据。假如我心里没有那份不知分寸、甚至该说不知羞耻的欲望——我是说认知的欲望——我每次都会停在这个“happy end”上,而每一步确实就是在这happy end 的心境里结束的。然后,我倒霉了!我会心血来潮——天知道怎么回事、为什么——要把我刚满心满意写下的东西凑近再看一看:那不是白纸黑字写着嘛,只要重读一遍就行了!可一认真重读,傻乎乎地重读,我就觉得有点不对劲,觉得不那么清楚,咦咦!然后再费心凑近一点看,就清楚了:根本不是那么回事,可以说全是瞎编,我刚才是拿尿脬当灯笼了!每一次这局部的发现都像个大惊喜,“好家伙!这一条还真不含糊!”——一个快活的惊喜,它带着一股新的劲头,把反思重新发动起来。往前走吧,总会弄到底的,肯定马上就到,只要顺着这股劲儿接着走!小结一下,理一理……于是另一种内心确信已经升上来了,带着“事情底细”的全部外表,我们只愿意相信:这一回总该对了吧;好歹记下来,图个心安,再说记下这么中肯、这么有体会的东西本身就是乐事,不赞成的人准是心术不正,这么明显的诚意,没法做得更好了,就这样,完美!

这就是新的阶段终点,新的happy end,我本来会高高兴兴停在那里的,要不是那个调皮到极点的坏小子又闹了起来——他真是不可救药,竟又把鼻子伸进这最新的“底细”和happy end里。拦不住他,于是又开始了新的一步!

% ---------------- p.195 ----------------
\origpage{92}\origpage{93}就这样,四个小时里,阶段一个接一个过去,像一颗洋葱,我把它的层一层层剥掉(这是那一夜末尾来到我脑中的形象),终于剥到了核心——到达那条再简单不过、再显眼不过的真理;说实话,那真理本就明摆在那里,可我却好几天、好几个星期(说到底,是我一辈子)都成功地把它藏匿在这一堆互相遮掩的“洋葱层”底下。

那条朴素真理终于出现,是巨大的宽慰,是意料之外、彻彻底底的解脱。那一刻我知道,我触到了焦虑的结。这五天的焦虑确确实实被化解了、消融了,转变成刚刚在我心里成形的那种认识。焦虑不只是从我眼前消失——就像整个冥想过程中那样,也像此前五天里好几次那样;而它转变成的认识,根本不属于观念的性质,也不属于那种为了了结、为了安宁而作出的让步(就像同一夜里有几回我做过的那样);那不是一件我从外面拿过来、或取得之后添到自己身上的东西。那是十足意义上的认识,第一手的,朴素而显眼,从此成了我的一部分,就像我的血肉是我的一部分。而且,它是用清楚、毫不含糊的话表达出来的——不是一大段话,而是一句三四个字的、再傻不过的短句。这表达是刚刚进行的那场工作的最后一步;只要这最后一步没跨过去,它就一直是短暂的、可逆的。在这整场工作中,把正在成形的思想、正在出现的念头仔细地、甚至一丝不苟地表达出来,是这工作本质的一部分;而每一次重新出发,都是对刚刚走过那一步的反思——那一步是我凭刚刚写下的见证而知道的(没法把它藏进一个靠不住的记忆的迷雾里!)。

在发现与解脱那一刻之后的几分钟里,我也知道了刚刚发生的事的全部分量。我刚发现的东西,比这几天那条朴素真理还要珍贵。这东西,就是我身上的一种能力——只要我有兴趣——去认识我内心发生之事的底细,认识任何分裂、冲突处境的底细;也因此,就是完全凭自己的手段去解决我心中任何我懂得去觉察的冲突的能力。

这解决不是靠什么恩典的功效完成的——前些年我倾向于那样相信——而是靠一场紧张、执拗、细致的工作,用的是我平常的官能。如果说真有“恩典”,它不在于我们心中某个冲突突然而彻底地消失,也不在于对冲突的一种理解现成地送到我们面前(就像懒人国里烤好的鸡!),而在于这认知欲望的在场或出现\fnoteP{5}{(31)}{我这里想到的是认知欲望的“阳”的形式——那种探测、发现、为出现之物命名的欲望……正是被命名,才使出现的认识变得不可逆、不可磨灭(哪怕它后来被埋没、被遗忘,哪怕它不再起作用……)。认知欲望的“阴”的、“女性”的形式,则在于一种敞开、一种接纳,在于默默迎候一种在我们存在的更深层、思想所不能抵达之处出现的认识。这样一种敞开、以及一种一时抹去一切冲突痕迹的突然认识的出现,仍像一种恩典,它深深地打动我们,尽管其可见的效果也许短暂。不过我猜想,这种如此来到我们身上的、无字的认识,在生命某些罕见的时刻,同样不可磨灭,它的作用甚至延续到我们对它的记忆之外。} (31)。正是这欲望引导我、在几个小时里把我带到冲突的核心——正如爱的欲望让我们万无一失地找到通往所爱女人最深处的那条路。

无论是自我发现还是数学,没有欲望,一切所谓的“工作”都只是装样子,哪里也去不了。最好的情况下,它让乐在其中的人无休无止地“围着罐子转”——罐里的东西是留给饿了要吃的人的!如同世上的%
% ---------------- p.196 ----------------
\origpage{94}所有人一样,我也有欲望与饥饿都不在场的时候。当涉及对自我的认知欲望时,我对自身以及我所卷入的处境的认识便停滞不动,我行事不是凭着对事情的认识,而是听任一些根深蒂固的简单机制摆布,随之而来的种种后果也一并承担——有点像一辆由电脑、而不是由人驾驶的车。但无论是冥想还是数学,没有欲望的时候、没有这份饥饿的时候,我不会想去装模作样地“工作”。所以我不曾有过冥想哪怕几个小时,或者做数学哪怕几个小时\fnoteP{6}{(32)}{\emph{炉中百铁,或:干熬无济于事!}\par 在我还做泛函分析的时候,也就是直到1954年,我有时会没完没了地死抠一个解不出的问题,明明已经没了主意,只是在旧念头的圈子里打转,而那些念头显然已经“咬”不动了。至少有一整年就是这样,特别是关于拓扑向量空间里的“逼近问题”,那问题要到二十年后才被完全不同路数的方法解决——以我当时的地步,那些方法只能与我无缘。那时驱动我的不是欲望,而是一种固执,以及对自己内心所发生之事的无知。那是难受的一年——我一生中唯一一段做数学变得难受的时期!我需要这番经历才明白:“干熬”没有用——一旦工作到了停住的地方,一察觉停住,就得转去做别的——不妨等到更合宜的时候再回到悬置的问题上。那个时刻几乎总是不久就出现——问题自己成熟起来,我根本不必去碰它,只凭在那些看似与此毫无关系的问题上兴致勃勃地工作,就足够了。我确信,当时若一味死抠,哪怕十年我也什么也弄不出来!正是从1954年起,我在数学里养成了同时往炉子里放很多块铁的习惯。我一次只打一块,但靠着一种不断重现的奇迹,我在一块上的工作也惠及所有别的铁,它们各自等着自己的时辰。后来与冥想初次接触时,情形也一样,而且完全不是我事先有意安排的——要审视的灼热问题一天天多起来,随着反思的进行而不断增加。} (32),而没学到一点东西的;而且多半(若非说总是)是某种意外、无法预见的东西。这跟我有、而别人没有的什么官能毫无关系,只在于:我真心不想做的时候,不去装模作样地做。(正是这份“想”的力量,单独就造出了我别处谈到的那种要求:在工作中不满足于差不多,非得把一种理解——无论多么朴素——走到尽头才满意。)在要发现什么的地方,没有欲望的工作就是无意义、就是装样子,跟没有欲望地做爱一样。说实话,我从未受到过这样的诱惑:把精力浪费在假装做一件我根本不想做的事上——何况有那么多让人热血沸腾的事可做,哪怕是在该睡觉的时候睡觉(还有做梦……)。

我相信,也是在同一夜,我明白了:认知的欲望与认知、发现的能力是一回事。只要我们信任它、跟随它,就是欲望把我们一直带到我们想认识的事物的核心。也还是它让我们找到——甚至不必去找——认识这些事物最有效、也最适合我们这个人的方法。至于数学,看来自古以来写作就是一种不可或缺的手段,无论“做数学”的是谁:做数学,首先就是写\fnoteP{7}{(33)}{这并不意味着,工作中没有纸(或黑板——那不过是它的一个变体!)的那些时刻,在数学工作里就不重要。尤其在“敏感时刻”更是如此:那时一种新的直觉刚刚出现,要做的不是“零敲碎打的工作”,而是以一种更整体、更直觉的方式与它“相识”——这种相识由反思的那个非正式阶段预备。在我这里,这类思考多半在床上或散步中进行,我觉得它在整个工作时间里所占的份额相对不大。同样的观察也适用于我迄今所做的冥想工作。} (33)。在一切以智力为主的发现工作里,大概也是如此。但“冥想”——我指的是自我发现的工作——肯定不必如此。然而在我这里,直到目前,写作在冥想中一直是一种有效而不可或缺的手段。就像在数学工作中一样,它是物质载体,把
% (切片止于此句: 源文 p.196 末行 "elle est le support matériel qui fixe le", 下接 p.197)

% ================================================================
%  第 9 章 收获 (后半) 原书 p.197-208, 边码 95-110
%  源文: _work/src1/task-ch09b.md
%  接 part1-ch09a.tex (p.187-196)
% ================================================================
%
% ---------------- p.197 ----------------
\origpage{95} 反思的节奏固定下来,并充当注意力的坐标与集合点——否则我的注意力往往倾向于四散随风;而且,书写还给我们留下一道刚刚做完的工作的有形痕迹,我们可以随时求助于它。在一场旷日持久的冥想里,能够回头翻看那些书面痕迹常常是有用的:它们见证着前几天、甚至几年前某个时刻的冥想。

\origpage{96} 思想,以及对思想的细致表述,因此在我迄今所实践的冥想中扮演着重要角色。但冥想并不因此就限于一种单纯的思想工作。单凭思想,无力把握生命。思想特别擅长的,是侦测矛盾——在我们对自身的看法里、在我们与他人关系的看法里,那些往往大得滑稽的矛盾;但常常,它不足以把握这些矛盾的意义。对于被求知之欲所驱动的人,思想是一件往往有用而有效的工具,甚至不可或缺——只要人始终意识到它的限度,这限度在冥想中显而易见(在数学工作中则更为隐蔽)。要紧的是:在另一样东西出现的那些敏感时刻,思想要懂得退开,踮着脚尖消失——那另一样东西也许以一阵突然而深沉的情绪的形式出现,而与此同时,手也许还在纸上继续奔跑,在同一刻给它一个笨拙而结结巴巴的表达……

\section{(37) 惊叹}

这段关于发现冥想的回顾,完全是意料之外地来到这里的,几乎违背我自己的意志——我起初打算考察的完全不是这个。我本来想谈惊叹。那个夜晚如此丰富,满载着如此多的东西,也满载着面对这些东西的惊叹。其实在工作进行当中,就已经有一种惊叹:面对每一条被揭出的新遁词,一种难以置信的惊叹——那就像一件用粗白线草草缝成的蹩脚戏装,而我竟,简直难以相信!还一本正经地把它当成真得不能再真的东西!此后许多次,在随后的年月里,我又重新遇到这同一种惊叹,如同在冥想的第一夜,面对我所发现的事实的巨大,以及那些让我一直对之视而不见的托辞的粗劣。正是从它那些滑稽可笑的侧面,我开始发现我身上那个未曾被猜想到的世界——一个随着日子、月份与年岁的推移而显露出惊人丰富的世界。然而就在那第一夜,值得我惊叹的就已经不只是些闹剧式的插曲。正是在那一夜,我头一次重新接触到一个被遗忘的、沉睡在我身上的能力;它的性质我还不清楚,只知道它是一种能力,而且随时听我支配。

而此前那几个月,也已经充满了一种无声的惊叹——惊叹于我身上一直带着的某样东西,那东西想必从来就在那里,而我刚刚才重新与它接触。我感受到这东西,不是作为一种能力,而更像是一种隐秘的温柔,一种既极其安宁又令人不安的美。后来,在我发现自己那份被忽视了如此之久的能力的狂喜中,我忘掉了这几个默默孕育的月份——它们只留下几首散落的诗作证:一些情诗,它们夹在我那些冥想笔记中间,大概常常显得格格不入……

只是许多年之后,我才回想起这些惊叹于世界之美、惊叹于我感觉到安歇在自己身上的那份美的时光。那时我才知道:我在自己身上感觉到的这份温柔与美,以及我不久之后发现、并深刻改变了我一生的那份能力,乃是同一件事的两个不可分割的方面。

% ---------------- p.198 ----------------
而我现在也看出:在我们身上那多重的东西——创造性——之中,那温柔的、内敛的、沉默的一面,自发地通过惊叹表达出来。而人认识所爱的女人、她也认识他,同样是在对由所爱者揭示出来的、自身之中那份不可言说的美的惊叹之中。当对所探究之物、对所爱之人的惊叹缺席时,我们与世界的拥抱就被削去了其中最好的部分——被削去了那使它成为对自己、对世界的祝福的东西。不是惊叹的拥抱,是一种无力的拥抱,只是占有姿态的简单复制。它无力生出别的东西,只会生出更多的复制品——也许更大些、更粗些、更厚些,无所谓,永远不是一次更新\fnoteP{8}{(34)}{\emph{无力的拥抱}\par “拥抱”一词在我这里绝不是一个简单的比喻,日常语言在此成了一个深层同一性的反映。人们可以说——而且不无道理——那么“没有惊叹的拥抱是无力的”这话就不对:若按字面理解,大地早该人烟断绝、一片荒芜了。极端的例子是强奸,其中确实没有惊叹,然而有时一个存在就在被强奸的女人身上被孕育出来。由这样的拥抱生出的孩子,想必不能不带这印记,它属于他所分得的那份“包裹”,属于该由他去承担的东西;这并不妨碍一个新存在确实被孕育、并且出生了,确实有过创造,那是力量的标志。也确实,我曾见过某个充满自得的数学家发现并证明出漂亮的定理,那是一次并不缺乏力量的拥抱的标志!但同样真实的是:如果这位数学家的生命被他的自得闷死(在某个时期,我自己的生活在一定程度上就是如此),那么这些与数学的拥抱之果实,对他、对任何人都不是福泽。对一个强奸所生的孩子的父亲和母亲,也可以这么说。我若说“无力的拥抱”,首先指的是无力在那个自以为在创造的人身上生出更新:他其实只造出一个产品,一件外在于他的东西,在他自身之中没有深层的回响;这产品远不能解放他、不能在他身上造出和谐,反而把他更紧地捆在他所囚于其中的那份自得上,那自得不断推着他去生产、再生产。这是深层意义上的一种无力,藏在一种“创造性”的外表背后——那“创造性”说到底只是一种没有刹车的生产力。\par 我也有充分的机会认识到:自得、无力惊叹,其性质是一种真正的盲目,是天然的敏感与嗅觉的阻塞;这阻塞即使不是完全的、永久的,至少在某些特定情境中显而易见。这是一种状态,其中某位声名显赫的数学家有时会在恰恰是他最擅长的东西上,显得比最迟钝的小学生还要愚笨!在另一些场合,他又会做出技术精湛的奇迹。然而我怀疑他是否还有能力去发现那些能够更新一门学科或一门科学的简单而显然的东西。它们在他下面太远,他已不屑于再看它们一眼!要看见谁也不屑于看见的东西,需要一种他已经失去、或被他放逐了的无辜……这肯定不是偶然:近二十年来数学生产的惊人增长,新结果纷至沓来,令那位只想稍稍“跟上形势”的数学家不知所措;然而(就我从各处得到的回声所能判断的),我曾多少熟悉的那些重大思考主题,几乎没有一个是真正更新了的,没有一个是大幅度的转化(而不只是积累)。更新不是量的事,它与投入的量无关——那量可以用某些“水平”的某些数学家花在某个题目上的“数学家人日”数目来衡量。一百万个“数学家人日”也无力生出一件像零那样孩子气的东西,而零更新了我们对数的感知。只有无辜有这份力量,它可见的标志之一,就是惊叹……} (34)。正是在我们还是孩子、准备为世上事物之美和为我们自身之美而惊叹的时候,我们才也准备好了去更新,并准备好作为柔软而顺从的工具,握在工匠手中,好让经祂的手、并借我们,也许有些存在与事物得以更新。

我清楚地记得,在四十年代末及随后的年月里,在这个对我而言代表着数学界的、不拘礼数的小圈子里——它有时喧闹而自信,那种略显断然的口气并不罕见(不过其中从不掺进一丝自得)——在这个圈子里,任何时候都为惊叹留有一席之地。惊叹在谁身上最显眼?是迪厄多内。不论是他作报告,还是他只在下面听,每当那个关键时刻到来、一道豁口突然敞开,就看见迪厄多内喜出望外、满面放光。那是纯粹的惊叹,有传染性、不可抗拒——其中“我”的一切痕迹都消失了。此刻我提起他,才意识到:这惊叹本身就是一种力量,它在他整个人周围立即发生作用,像一种以他为源头的辐射。要说我见过哪位数学家动用一种强大而基本的“鼓励之力”,那就是他!在此之前我从未再想到这一点,但现在我记起来了:也正是带着这样的心情,他当初接纳了
% ---------------- p.199 ----------------
\origpage{97} 我在南锡的最初几项结果,解决了他和施瓦茨一起提出的问题(关于 $(F)$ 空间与 $(LF)$ 空间)。那些结果都很平常,既不天才也不出众,可以说并没有什么值得惊叹的。后来我见过气派完全不同的一些东西,被那些自以为是大人物的同行一句不容分说的轻蔑打掉。迪厄多内身上全然没有这类负担——无论那自命是否站得住。没有任何这类东西妨碍他连小事也为之欣喜。

\origpage{98} 这种欣喜的能力里有种慷慨,它既造福于愿意让它在自己身上舒展的人,也造福于他周围的人。这恩惠的施行,并不怀着讨谁欢心的意图。它单纯得像花香,像阳光的温暖。

在我认识的所有数学家中,这份“馈赠”在我眼里以最耀眼、最有感染力、也许也最起作用的方式出现的,就是迪厄多内;我说不准\fnoteP{9}{(35)}{这份“馈赠”不是任何人的特权,我们生来都有。当它在我身上似乎缺席时,那是我自己把它赶走了,而重新接纳它只取决于我。在我身上或在某某人身上,这份“馈赠”的表达方式与另一个人不同,也许不那么有感染力、不那么不可抗拒,但它并不因此就不在,我也说不出它是否就不那么起作用。} (35)。但在我乐于交往的数学界朋友中间,没有哪一个身上缺这份馈赠。它总能找到机会显露,也许方式更含蓄,但随时都在。每当我走向他们中的一位,想与他分享我刚发现的、让我着迷的东西,它就显露出来。

如果说我作为数学家的生活里也有过挫折与痛苦,那首先就是:在我爱过的某些人身上,再也找不回我曾在他们身上认识的那份慷慨、那份对事物之美——无论“小”或“大”——的敏感;仿佛那曾使他们整个存在颤动不已的东西,已经无声无息地熄灭,被一种自得闷死了——那种自得属于这样一个人:世界对他而言已不够美,不值得他屈尊为之一喜。

当然,也还有另一种痛苦:看见我昔日某位朋友以俯就或轻蔑的态度对待我今天的某位朋友。但说到底,这痛苦也是同一种封闭造成的。对一件事物之美——无论它多么卑微——敞开着的人,一旦感受到这美,就不由自主地也对构想或做成它的人怀有敬意。在一件出于人手的事物之美里,我们感受到做它的人身上那份美的反光,感受到他制作它时所倾注的爱。当我们感受到这美、这爱,我们心里就不可能有俯就或轻蔑——正如在一个我们感受到某个女人的美、以及这美所标志的她身上那份力量的时刻,我们不可能对她有俯就或轻蔑。

\section{(38) 回归的冲动与更新}

迪厄多内身上时或闪耀的那种欣喜,想必触动了我身上某种深沉而有力的东西,这记忆此刻才如此强烈、如此鲜活地回到我心头,仿佛我刚刚还亲眼目睹。(其实我已经将近十五年没怎么有机会见到迪厄多内了,只有一两次匆匆照面。)当然,在意识层面,我并未对此给予特别注意——那不过是我这位年长的同行兼朋友那外放的性格中一个有点动人的特点,有时几乎滑稽。相反,对我重要的是:我在他身上找到了那个完美的、我甚至可以说是梦想中的合作者,来把那些将要为我眼前展开的广阔远景充当基础的东西,以一丝不苟的细心、一种充满爱意的细心,白纸黑字地写下来。直到此刻,当我把这两件事一并提起,
% ---------------- p.200 ----------------
\origpage{99} 那条联系才突然在我眼前显现:使迪厄多内成为一项伟大任务之理想仆人的——无论是在布尔巴基内部,还是在我们为另一项奠基性大工程而进行的合作中——是慷慨,是他的工作中、在他那些重大投入的抉择中,全无一丝虚荣的痕迹。我经常看见他隐没在他甘为仆役的那些任务后面,毫不吝惜地把无穷的精力献给它们,不求任何回报。毫无疑问,他什么也不求,却在工作本身、在他投入其中的那份慷慨本身里,得到一种充实与舒展——所有认识他的人都必定感受过。

我常常感到从他身上闪耀出来的那种发现之欣喜,在我心里立刻与一种相似的欣喜联系在一起——我曾在一个很小的孩子身上目睹过那种欣喜。两个记忆争相涌来——两个都让我重新看见我那时还很小的女儿。第一幅画面里,她大概才几个月,想必刚刚开始会爬。她大概是从人家把她放坐的那块草坪,爬到了旁边一条砾石小径上。她在发现那些小石子,一种无声的——却是起作用的——狂喜:她双手大把抓起石子,往嘴里送!——另一幅画面里,她大概一两岁,有人刚往金鱼缸里撒了鱼食。鱼儿们争先恐后地游过去,大张着嘴,吞下那些悬浮着、慢慢沉向缸底的小小黄色碎屑。小家伙此前从没意识到鱼也像我们一样吃东西。她身上仿佛突然一亮,迸出一声纯粹欣喜的叫喊:“妈妈看,它们在吃饭!”这确实值得惊叹——她在一瞬间发现了一个大奥秘:我们与一切其他生物同属一族的奥秘……

\origpage{100} 一个小孩子的欣喜里有一种无法用言语传达的、有传染性的力量,一种从他身上辐射出来、作用于我们的力量——而我们大多数时候总是竭力躲开它。在内心里安静的时刻,人能感觉到这力量随时都在孩子身上。只是有些时候它的作用更强些。这种围绕孩子的“力场”,在新生儿身上、在生命最初的日子和月份里最强。大多数情况下,它在整个童年期间都还感觉得到,随年岁渐长而一丝丝磨损,直到青春期——那时往往已经看不出痕迹了。然而,人们也能发现它辐射在各个年龄的人周围:在某些人身上,在得天独厚的时刻;在少数另一些人身上,则像一种气息或光晕,时时刻刻环绕着他们这个人。我有幸在童年认识这样一个人,一位男子,如今已经去世。

我也想到另一种力量,或能力:人们有时感到它从某个女人身上辐射出来,尤其是在她舒展于自己的身体、与身体交融的时刻。我常想到的词是“美”,它道出了其中一个侧面。这种美与任何美的规范、与所谓“完美”毫无关系;它不是青春的特权,也不是成熟的特权。它更像是一个人身上某种深层和谐的标志。这种和谐常常仍是残缺的,然而它正是通过这种辐射显现出来——那是力量的标志。这是一种把我们引向它所自出的中心的力量——或者不如说,它在我们身上唤起一种深沉的冲动:回到那位“女人—母亲”的身体里去——我们在生命破晓时,正是从那里出来的。它的作用有时强不可挡,当它来自所爱的女人时,令人翻覆。但对于不刻意对它封闭的人,它在每一个让这种美、这种深层和谐在自己身上舒展的女人身上都感觉得到。

从孩子身上辐射出来的力量,与从那个爱着自己身体的女人身上散发的力量是近亲。后者不断生出前者,正如孩子不断从母亲生出。但童年的力量,其本性不是吸引,也不是排斥。这种力量对不躲开它的人所施的谦卑而谨慎的作用\origpage{101},
% ---------------- p.201 ----------------
是一种更新的作用。

\section{(39) 夜之美、日之美(或:奥吉亚斯的牛圈)}

对我某个孩子的惊叹的记忆,是在五十年代末、六十年代初。如果说后来出生的其他孩子没有给我留下类似的记忆,那也许是因为我自己惊叹的能力已经钝了,我变得太疏远,无法与某个孩子的欣喜相通,甚至无法仅仅做一个见证者。

我还从未想过要追索这种能力在我一生中——从童年到今天——的变迁。那里想必有一条线索,一个极其灵敏的“探测器”。我从未想过追随这条线索,想必是因为这种能力的性质如此谦卑、外观几乎如此微不足道,以至于我几乎不会想到去特别注意它——我当时正全神贯注于发现和探测我所谓的我一生的“伟大力量”(它们今天仍在我身上显现)。然而,这种外观如此谦卑的能力,却提供了一个胜过一切的标志,指示着我们身上那份最稀有、最珍贵的力量之在与不在……

在我整个成年生活中,我从未与这种力量完全隔绝。无论我的生活在别的方面变得多么荒芜,我总在爱中重新找到孩子的惊叹、发现的欣喜。穿过许多片沙漠,爱的激情始终是与某个我已离开的东西之间那条活生生而茁壮的纽带,一根脐带,默默继续以温热而慷慨的血液滋养我。而在很长一段时间里,对所爱女人的惊叹,也与对她所生育的新生存在的惊叹不可分割——那些崭新的存在,无限娇嫩而生机勃勃,见证着她的力量,也承继着它。

但我这里的用意,主要是想稍稍追索这种“无辜之力”在我作为数学家的生活中的变迁,即我属于“数学家的世界”那个时期, 1948到1970年。当然,惊叹从未像在爱的激情中那样渗透我的数学激情。奇怪的是,若要我回想数学工作中某个欣喜或惊叹的特定时刻,我一个也想不起来!我接近数学的方式,从十七岁全力投入时起,就是给自己提出大任务。这些任务从一开始就总是“整理”、大扫除的任务。我看见一片表面的混乱,一堆杂乱异质之物,有时是不可称量的雾气,它们显然应当有一个共同的本质,藏着一个秩序、一种尚被遮蔽的和谐,要靠耐心、细致、往往旷日持久的工作把它揭示出来。这工作常常先是拖把和大刷子的活计——那粗活已经耗去大量精力——然后才轮到用掸子做收尾;收尾我不那么着迷,但也自有其魅力,而且无论如何显然有用。日复一日的工作中有一种强烈的满足:看见那个被猜到的秩序一点点显露出来,它总是比先前瞥见和猜到的更为精致,质地更为丰富。工作不断充满意外插曲,它们多半从考察某个看似极微小、此前一直被忽视的细节中冒出来。常常,对某个“细节”的打磨,给多年前做过的工作投下意想不到的光。有时它也带来新的直觉,而对它的深化又成了另一项“大任务”的对象。

这样,在我的数学工作中(除了1954年前后那个我有机会谈过的“艰难的一年”),有一种持续的悬念,注意力始终被吊着。何况,对我那些“任务”的忠诚不许我跑得太远,我咬着嚼子,急不可耐地想
% ---------------- p.202 ----------------
\origpage{102} 把它们全都做完,终于纵身跃入未知、真正的未知——而这些任务的规模此时已经变得如此之大:要把它们做完,即便有那些终于赶来救援的善意相助,我的余生也不够用!

我的主要向导,是不断寻求一种完美的融贯、一种完整的和谐——我猜到它就在事物动荡的表面之下,我耐心地努力把它揭示出来,从不厌倦。那确实是一种敏锐的“美”感,它是我的嗅觉,也是我唯一的罗盘。我最大的快乐,与其说是在它出现在光天化日之下时凝视它,不如说是看着它一点点从它喜欢不断躲藏其中的阴影与雾气的衣袍下显露出来。当然,不把它带到最明亮的日光下,我是不肯罢休的。那时我有时也体会到凝视的圆满——当一切可闻的声音汇成同一片浩瀚的和谐。但更常见的是:被带到日光下的东西,立刻又成了再一次潜入雾中的动机与手段,去追寻那位永远神秘、永远未知的“她”的又一次化身——她不断呼唤我,为了再认识她……

\origpage{103} 迪厄多内的快乐与欣喜,我觉得,主要在于看见事物之美在光天化日之下显现;而我的快乐,首先在于在雾与夜的皱褶里追寻它。这也许就是迪厄多内与我接近数学的方式之间的深刻差异。对事物之美的感觉,至少在很长一段时间里,在我身上想必不比他弱——尽管它也许在六十年代里、在一种自得的作用下变钝了。但看来,在迪厄多内身上通过惊叹显现出来的对美的感知,在我身上采取了不同的形式:不那么静观,更爱动手,在所感受和表达的情绪层面上也不那么外显。若果真如此,我要追索的,就应当是我身上对数学事物之美的那份敞开之变迁,而不是那份神秘的“惊叹之馈赠”。

\section{(40) 竞技数学}

相当清楚的是,对数学事物之美的敞开,在我身上从未完全消失,即使在六十年代直到1970年——那时自得在我与数学、与其他数学家的关系中逐步占据了越来越大的位置。若没有最低限度对事物之美的敞开,我根本不可能作为数学家“运转”起来,哪怕只开到最低档——而且我怀疑,任何人若身上不多少还活着这种美感,能否在数学里做出有用的工作。依我看,造成这位数学家与那位数学家之间、或同一位数学家的这项工作与那项工作之间差别的,与其说是某种所谓的“脑力”,不如说是这种敞开或敏感的精细度、细腻度——研究者彼此之间不同,同一个研究者不同时刻也不同。最深刻、最丰产的工作,也是那种见证着最敏锐的敏感性的工作——那敏感性用于把握事物隐藏的美\fnoteP{10}{(36)}{在我看来,这种对美的精细敏感,与我曾有机会以“严苛”(对自身)或“严格”(就其完整意义而言)之名谈过的东西密切相连;我曾把它描述为“对我们自身之中某种精细之物的注意”,一种对所受探究之物的理解品质的注意。对一件数学之物的这种理解品质,不能与对该物特有之“美”或深或浅、或完美或粗略的感知分开。} (36)。

若果真如此,那就得相信,这种敏感在我身上一直鲜活到最后,至少某些时刻如此——因为正是在六十年代末\fnoteP{11}{}{(8月8日)核实之后看来,我关于 Motif 的思考之开端是在六十年代初,而不是六十年代末。},我开始瞥见并稍稍揭示出我所发现的最隐蔽、最神秘的数学对象:我称之为“Motif”的东西。
% ---------------- p.203 ----------------
\origpage{104} 它也正是在我作为数学家的生活中对我施以最大魅惑的东西(若不算最近几年的一些思考——它们其实与 Motif 的实在密切相关)。毫无疑问,若不是我的生活突然转向一条完全意外的道路,把我远远带离数学事物那宁静的世界,我最终会追随这强大魅惑的召唤,把那些一直把我囚禁着的“任务”撂在一边!

也许我可以说:在我工作室的孤独里,美感一直保持不变,直到1970年我第一次“觉醒”的时刻,并没有真正受到那种常常标示着我与同类关系的自得的影响?某种“嗅觉”甚至想必随年岁而愈发精细,在与数学事物每日而亲密的接触中。我们对事物所能拥有的亲密认识,有时使我们能把握超出当下所知的东西,并更深入地进入认识——这种认识或这种成熟,以及作为其最显眼标志的这种“嗅觉”,是向事物之美与真之敞开的近亲。它助长、刺激这样的敞开,而且它是此前一切敞开时刻、一切“真理时刻”的总和与果实。

那么,我接下来要考察的,就是:这种对美的自发敏感,在我与这位或那位同行的关系中、在它有机会显现的那些时刻,在多大程度上受到了或深或浅的扰乱。

记忆在这方面交给我的东西,并不能凝成一个有形的、确切的事实,让我在这里多少详尽地摆出来。这里的记忆仍然只是一团雾,但它还是交给我一个总体印象,我得设法把它界定出来。那是某种内心态度留在我身上的印象,它想必最终变成了第二天性,每当我接到一条数学信息、而所涉之物多少“在我的射程内”时,它就显现出来。说实话,从某个相对无害的侧面看,这种态度想必一直属于我,它属于某种气质,我也曾有机会顺带触及。说的就是这种反射:首先只肯了解一个陈述,绝不看它的证明,先试着把它安放到我所已知的东西里,看看用这些已知的东西来说,这陈述是否变得透明、显然。这常常使我以或深或浅的方式重新表述那个陈述,使之更一般或更精确,往往两者兼有。只有当我无法用我的经验和我的意象去“安放”这个陈述时,我才准备好(有时几乎违背我自己的意愿!)去听(或读……)那些来龙去脉——它们有时给出事情的“那个”道理,或者至少给出一个证明,懂不懂另说。

这是我接近数学的一个特点,我觉得它使我有别于当时布尔巴基的所有其他成员,也使我实际上不可能像他们那样插入集体工作。这特点想必也在我教书时构成一种障碍,我的所有学生想必都感受过——直到今天(多亏上了年纪)它才终于松动了一些。

我身上这个特点,想必已经是一种敞开上的欠缺。它只意味着一种部分的敞开,只肯接纳“来得正好”的东西,至少对别的一切都极不情愿接纳。在我的数学投入之抉择中,在我同意把时间花在这样或那样不期而至的信息上的抉择中,这种有意“部分封闭”的打算,今天比任何时候都更强。它甚至是一种必要,如果我想能追随那最令我着迷之物的召唤,而又不再把“我的命拿去吃”献给数学夫人!

% ---------------- p.204 ----------------
\origpage{105} 然而这团“雾”交给我的,不止这个特点——我几年前就已经察觉它了(晚知道总比不知道好!)。在某个时刻,这种反射变成了某种关乎荣誉的事:我要是在比说这句话更短的时间里还“拿不下”这个陈述(假定它对我还不算太熟悉),那才叫见鬼!如果陈述的作者是个有名有姓的陌生人,还多一层微妙:那可不行,我(毕竟人家以为我是在行的人!)岂能手里没有这一套!而常常我确实有,而且绰绰有余——那时我的态度就倾向于这个方向:“行了,您可以回去穿好衣服了——等您做得好一点再来吧!”

这恰恰就是我在“那个在我园子里乱走的毛头小子”一事上的态度。我甚至不敢发誓:他做的东西里没有一些有趣的细节,是我那些“秘密笔记”里所做的东西不曾覆盖的——不过这终归是无关紧要的事\fnoteP{12}{}{(8月8日)后来我发觉,这件事并不那么“无关紧要”:它构成了从“竞技态度”通向某种不诚实之开端的界线,而这条界线,我也许有几次越过了……}。说到底,这段插曲也照亮了我在这里考察的问题:对数学事物之美的那份敞开之深层扰乱。仿佛从我把某样东西“做出来”的那一刻起,它的美对我而言就消失了,剩下的只有一种虚荣,要求为它记功受利。(而我却又不屑于花时间去发表它——诚然,要发表的也太多了。)那是一种典型的占有态度,类似一个男人,认识了一个女人之后,再也感觉不到她的美,又去追逐另外一百个,却仍容不得别人认识她。这种态度我在爱情生活中是斥责的,自以为远在这样一份虚荣之上,却又小心翼翼地不去确认这个显而易见的事实:我对数学的态度恰恰正是如此!

\origpage{106} 我有个印象:我刚才在自己身上指出的这些粗鄙的竞争心性、不妨说“竞技”心性,想必大约就在它们于我自己身上变得普遍的那个时候,开始在我的数学“圈子”里变得普遍起来。要我在时间上定位它们出现的时刻,或它们变成这个圈子里所呼吸的空气之一部分的时刻,或我的学生们在与我接触中所呼吸的空气的时刻,我却无能为力。我唯一敢说的,是这大概在六十年代,也许早在六十年代初,或五十年代末。(若果真如此,我的学生全都摊上了——他们只能要么接受,要么拉倒!)要定位它,我还需要别的确切例子,而此刻它们完全逃出了我的记忆。

当然,这个卑微的实情,与我为自己同数学、同一般年轻研究者的关系所描绘的高贵形象完全相反。用来蒙骗我自己的那个粗劣托辞,灵感来自功绩主义:对于这个形象,我所留住的只是我与学生们的关系(他们为我的声望添砖加瓦,是我声望中最高贵的花朵!),以及我与那些特别出色的年轻数学家的关系——我懂得承认他们的功绩,并且像对待我的学生一样平等待他们,不等他们的头上戴上桂冠(当然,桂冠也没等多久——嗅觉这东西,有就是有,没有就是没有!)。至于那些既没福气做我学生、也没福气做我某位朋友的学生、又不是年轻天才的年轻人,我丝毫不在意我与他们是什么关系。他们不算数。

我相信,当我与年轻研究者有了个人的接触——无论是在讨论班上遇见他,还是他写信给我——上述实情多半会松动、缓和一些。从这一点看,“毛头小子”那件事也许是个稍显特殊、例外的情形。我觉得,对我刚才谈到的那些研究者,我大概会多少把他们看作“受我庇护”的人,这想必在我身上唤醒了一种更和善的态度。在这种情形里,我想出头的欲望也能找到出口:对当事人发表我的评论,建议他也许从更广阔的视角、或者更深地进入事情本身去重做他的工作。在这种情形下,那位在有限时间里多少扮演了学生角色的年轻研究者,很可能也从中有所得,并且对我与他的关系留下好印象。(这方面任何或正或反的回声,我都欢迎。)

% ---------------- p.205 ----------------
\origpage{107} 我在这里主要想到的是较年轻的研究者,而“竞技”态度绝不限于我与他们的关系,这不用说。但确实是在与年轻研究者的关系中,一位有名望的数学家的影响——无论心理上还是实际上的——往往最强,对他们未来的职业生涯也最有后果。

\section{(41) 克里希那穆提,或:成为羁绊的解放}

\origpage{108} 这一夜,我停在一份宽慰之情、一份极大的满足上——那是一个没有白费时间的人的满足!我突然觉得轻松而快乐——一种有时带点促狭的快乐,迸成顽皮的笑——一个捣蛋鬼恶作剧的笑。其实我并没做什么大事,我不过是从一个略微不同的角度,重看了一段已经“知道”的插曲,即那个著名的“毛头小子……”。这个角度显示出的,是我与数学本身的关系——在某些情形下,不只是我与数学家的关系。就这一点,已经足够让一个我曾珍爱的神话化为青烟。

说实话,这并非我第一次审视我与数学的关系。两年半以前,我就曾被引向为此花上几个星期或几个月。那时我(在许多别的事情之外)意识到:在我过去对数学的投入中,自我之力、自我膨胀之力有多么重要。但昨夜我刚刚点到了一个当时逃过我注意的侧面。现在我回过头来谈它,发现这个侧面——即我与数学的关系中那种嫉妒态度——与第一夜“冥想”结束时到来的那个“再平常不过”的发现接上了(当时我并不知道自己在冥想,就像茹尔丹先生不知道自己是在说散文……)。这很可能在我随后那份欢欣的狂喜中占了一份。即便没有被自觉地察觉,那也多少像是对我从前发现的某样东西在新光照下的再次确认——而这时心里的快乐,与数学里的一样:当你并没有去找,却从一条完全不同的路径撞上你认识的、也许多年前就发现的东西。每一次,这都伴随着一种深切的满足感:事物的和谐再一次显露,而我们对它的认识也或多或少被更新。

而且我相信,这一次我确实“转完了”!好几天来我都清楚地感到还有什么东西没被挖出来,却又说不清到底是什么。我没有试着硬来,我觉得只需让它自己来,让所追随的那条线自由展开,穿过一片片既熟悉又意外的风景。说意外,是因为我此前从未费心去看它们。我是以散步的步子走近那个剩下的“热点”的。而我相信,这是我刚刚走完、且已接近尾声的这趟旅程中的最后一个。

而刚一到达这一点,我就有了一种感觉,像一个登上观景台的人:从那里,他看见\origpage{109}自己刚刚走过的那片风景铺展开来,而在途中每一刻,他只能看见其中一角。
% ---------------- p.206 ----------------
\origpage{110} 而现在有了这种辽阔与空间的感觉,那是一种解放。

如果我试着用文字表述眼前这片风景交给我的东西,出来的就是这些:这些年来,在我作为数学家的生活中到来的一切——常常来得不是时候,也不受欢迎——都是我在还属于数学家的世界那些年月里所播种之物的收获与讯息。

当然,这件事,这些年来我已经对自己说过又说,在我刚写下的这些笔记里也说过。我对自己说过,多少类比于另一些执拗地到来的收获——我长久拒斥,最终接纳并化为己有。从我这样接纳的第一份收获起,甚至在我懂得冥想之前,我就明白了:每一份收获都必有其意义,而赌气只是回避一个意义、推迟一个了结的期限。这一认识对我很宝贵,因为它常使我免于自怜,也免于那种义愤——义愤往往是自怜的一种伪装形式。这认识在我身上像是一种半成熟,它还远没有终结那种积习难改的反射:当收获显得苦涩时加以拒绝。当我说“赌气没有用”,收获并没有因此被接纳。我也许既不自怜也不义愤,可我仍在“赌气”!只要那盘菜没吃下去,它就还没被接纳——而不吃,就是赌气。

接纳并吃下去是一件工作:某种能量在“工作”,一种工作在光天化日下或在暗处进行,有东西在转变……而赌气则是一种能量的浪费,那能量散掉了——散在“赌气”上!而吃、消化、吸收这件工作,是省不掉的。仅仅经历一些事件、仅仅“取得”或“获得”一种经验,与一件工作毫无共同之处。它只是一份可供工作的材料,做不做由你。从我遇到数学家的世界起,三十六年了,我一直在使用这份自由:回避一件工作,而可供吃、可供消化的材料却年复一年地增多。我从昨天起感到的这份欢欣的解放感,确实是一个标志:我面前那件工作——我不断地以别的工作或任务为借口推开的——终于做完了。也确实到时候了!

现在还不能确定事情果真如此,还不能确定没有剩下某个阴暗而顽固的角落逃过了我的注意、我还得回到那里去。但同样真实的是:这种解放感不会骗人——我一生中每一次感到它,事后都能证实它确实是一次解放的标志;是某种持久的、已经获得的东西,是一种理解、一种认识的果实,而那认识已成为我的一部分。若我愿意,我有自由无视这认识,把它埋在我想要的地方、按我想要的方式埋掉。但要毁掉它,却不在我、也不在任何人的能力之内——正如人不能毁掉一只果子的成熟,不能把它弄回一种已不属于它的青涩状态。

再一次证实我并不比别人“更好”,这是一大宽慰。当然,这也是我相当经常地对自己重复的话——可是重复和看见,毕竟不是一回事!没有孩子的无辜与灵动——他看见,就像他呼吸——要看见明显之事,往往需要一件工作;好了,做完了,我终于看见了这一件:我并不比几天前还让我“喘不过气”的那些同行或旧学生“更好”!想想我如今卸下的这副担子有多重吧!自认为比别人强,从某种意义说也许挺受用,但那也非常累人。那甚至是精力的极大浪费——正如每一次要维持一个虚构时那样。人很少意识到这一点,但光是顶着风浪维持那个虚构,就已经需要精力了——而明显之事在每一场合都冲着我那仔细塞住的耳朵喊:这是假的,看看吧,傻瓜!看见,有时也许是一件工作,
% ---------------- p.207 ----------------
但一旦做完就做完了。这就一劳永逸地省掉了两件事:一是像那样动不动就堵上眼睛和耳朵地四处走动——这本身也是桩苦差!二是每当有件我自己不小心搁在那里的东西落到我头上时,就把它当作不可容忍的侮辱而自苦。

受够了这套把戏!看清了把戏,就已经出来了。付了钱,没错,我有权一直在里面转到老死,甚至有义务——那也无妨,所有人都会这么对我说:权利、义务——看人下菜碟。所有这些是权利又是义务、是义务又是权利的东西,也累人得很,当我自认为比别人强的时候,它们就黏上我。说到底这很正常:既然是强者,就悄悄收着(这就是“权利”),并且“付账”,为了人类精神与数学的荣誉尽全部义务——真的很美,没错,荣誉、精神、数学,谁还能说得更好,好极了!再来一个!是很美,没错,可也非常累人,到头来会让人落枕。我落过枕了,现在够了——我把位置让给别人去挺直腰板吧。

这也正常(既然我谈到学生):学生超过老师。我曾为此恼怒,我真有精力可浪费!这些都结束了!

多么轻松!

% ---------------- p.208 ----------------

% ================================================================
%  第 10 章 孩子嬉戏 (原书 p.209-218, 边码 111-121)
%  源文: _work/src1/task-ch10.md
% ================================================================
% ---------------- p.209 ----------------
\chapter{孩子嬉戏}

\section{(42) 孩子}

肯定还有些角落,扫帚没扫到。没关系,它们自会引起我的注意,到那时再收拾也不迟。至于我那大名鼎鼎的“数学家过往”,大扫除已经做完,这一点毫无疑问。

既然我又一次看清我并不比旁人更好,那就别再掉进那个老掉牙的陷阱:把自己看成\textbf{比我自己更好}!——把如今的我,退出了那台旋转木马、诸如此类的我,看成比十五年前、或者十五天前的我更好。这十五年里我学到了些东西,这是肯定的;这十五天里也是,甚至从昨天起也是。我学到东西的时候,就成熟了,就不再完全是同一个人了。学到某样东西之后的我,并不比那样东西还摆在我面前时的我“更好”。一个更熟的果子并不比一个不太熟的、或者青的果子“更好”。一个季节并不比它前一个季节“更好”。最熟的果子味道可能更可口,也可能更不可口,这要看各人的口味。一年一年过去,我在自己的皮囊里觉得更自在了,想必我身上发生的变化“合我的口味”——可这些变化并不合我所有朋友或亲近者的口味。每次我重新做起数学,四面八方都来恭维,腔调是这样的:“他当初干吗要去弄别的东西呢!现在一切都归了位,早该如此!”看着一个人变,总叫人不安……

\origpage{111} 我学习,我成熟,我变化——以至于有时我很难在从前那个人身上认出自己,得靠一段回忆、或者旁人一次意想不到的见证,才重新发现他。我在变;同时也有某种东西始终“是同一个”。它一直都在,从永远起就在,肯定从我出生起,也许还更早。近几年,我觉得自己渐渐能把它认准了。我把它叫作“孩子”。凭这东西,此刻的我不比一生中任何别的时刻更好;它就在那里,尽管往往很难猜出它的在场。也凭这东西,\textbf{我不比任何人更好},也没有\textbf{任何人比我更好}。在某些时刻、在某些人身上,孩子更在场一些。而这是件很让人受用的事。这并不意味着某个人比另一个人“更好”,或者比他自己的另一个时刻“更好”。

常常,我做数学的时候,或者做爱的时候,或者冥想的时候,是孩子在玩。并不总是只有他一个在“玩”。可他不在的时候,就没有数学,没有爱,也没有冥想。装样子是没用的——这出戏我很少演。

% ---------------- p.210 ----------------
\origpage{112} 当然不只有孩子。还有“我”,还有“老板”,或者说“大老板”,爱怎么叫就怎么叫。对企业运转来说,老板肯定不可或缺。既然有个老板,总该是为着什么。他照管庶务,而且像所有老板一样,有个讨厌的倾向:越占越宽。他把自己看得无比要紧,非要处处压过对面的老板。占得宽也好,不宽也好,他不过是老板,他不是工人。他组织,他指挥,他也收账——当然!收益他当作自己应得之物收下,亏损他当作侮辱来承受。但他什么也不创造。只有工人有创造的力量,而工人不是别人,正是孩子。

老板和工人合得来的企业,很少见。多数时候,根本看不到工人的影子,他被关在天知道什么地方。是老板做出顶替他站在车间里的样子,结果可想而知。也常常是这样:工人明明就在,老板却对他开战,或是猛烈的战争,或是小股交火——这样的车间里出不了什么像样的东西!有时老板对工人也有一种带戒心的宽容,他一边嘀咕一边由着他干,眼睛却一刻也不离开他。那就像一场从未停歇的战争里,不断续签的停战。工人也就趁着这停战,勉强干上一点活。

我刚才那场冥想的功效,绝不意味着我对数学的那份占有态度就此如中了魔法一般消失了!至少,我得近得多地去看占有欲的种种表现——刚才我不过是叫出了其中一种的名字,轻轻碰到了它一下。这里不是地方:这篇“导论”已经变成了“导论性的一章”,而这一章自己也已经显得长了!不过昨夜有一件事“咔嗒”一声对上了,此刻我想稍稍回过头去说它一说;那是两三年前我带着几分惊讶记下的一件事。

当时我正钻在一个数学问题上,已说不上是什么问题了;某一刻(不知由于什么情况),我发现我盯着的问题也许早有人盯过,它很可能在某本书里白纸黑字地处理过,只消我去图书馆查一查。单是想到这种可能性,效果就如雷击,把我惊得发呆:转眼之间,欲望不见了。忽然之间,那个我也许已经花了好几星期、还打算再花上几星期的问题,对我完全失去了吸引力!那不是赌气,那是兴趣突然而彻底的丧失。就算那本书就在我手里,我也懒得翻开。

\origpage{113} 事实上,那种可能性并未得到证实,于是欲望又回来了,我照旧一路做下去,像什么都没发生过。我却仍然愣着。当然,如果我正在做的东西真真是我另做一件事所\textbf{必需}的,那就不会有那么惊人的兴趣跌落。我常重做已知的东西,明知或猜到它们已知,却一点也不在乎。那时我正处在一种冲劲里:照我自己的方式、按它们向我呈现的样子去做,比到书里、文章里去翻找更省事,尤其有意思得多。那时我是“顺势”奔向别的东西,奔向欲望把我带往的地方。当然,我够“在行”,知道自己最终要到的那个地方,任何书、任何文章里都没有。

这让我想起来:数学工作即使孤零零地做了许多年,也不是纯粹个人的、个体性的工作,不像冥想那样——至少在我身上不是。我在数学里追的那个“未知”,要能那样强有力地吸引我,就必须不只对我未知,而是对所有人都未知。数学书里写着的东西不是未知,即便我自己从未听说过它。读书或读文章从来不吸引我,只要\textbf{能躲我就躲}。它能告诉我的,从来都不是未知;
% ---------------- p.211 ----------------
我对它的兴趣,也没有欲望的那种质地。那是一种“权宜的”兴趣,是对可能对我有用的一条\textbf{信息}的兴趣,是某个欲望的工具,而这条信息根本不是那个欲望的对象。

细想之后,我不觉得我讲的那件事是什么嫉妒、占有的心性的迹象,也不是什么受了挫的虚荣的迹象。我心里没有半点赌气,没有半点失望,只是一个前一瞬还强烈的欲望突然消失了。那是在一个我压根没想过要发表任何东西的时期,也没想过有一天我会忽然又想发表点什么。这欲望不是虚荣的表现,不是急于积累知识、头衔和信誉的饕餮——它确确实实是一个真实的欲望,是沉迷于游戏的孩子的那种欲望。而忽然之间——什么都没有了!谁懂就懂吧,我是不懂……抱歉!

\section{(43) 扫兴的老板——或高压锅}

\origpage{114} \noindent\pmark 我觉得,这场对我数学家一生的回顾,终究是做完了。当然,我没有把题目做尽——那得要好几卷书,假如这样的题目真能“做尽”的话。那不是我的意图。我的意图是弄清楚:在某种我如今只能一阵一阵感觉到的“气息”的出现里,我到底是不是参与者、是不是共同行动者;如果是,又是以什么方式。现在我弄清楚了,这让人舒坦。再往前走,去深挖那些只是瞥见或触及的东西,可能会很来劲。可看、可做、可发现的有意思的东西太多了!至于我作为数学家的过去,我觉得,为了承担这段过去而\textbf{必须}看的,已经看到了。

肯定,若把这番冥想往深里做,我准会学到许多关于我当下的有意思的东西。这项工作几乎每走一步都让我感到的一件事是:我对这段过去还眷恋得多么深,它直到今天在我的自我形象里、以及在我与他人的关系里,分量有多重;尤其是在我与那些——从某种意义上说——被我抛下的人的关系里。这番工作过程中,我与这段过去的关系肯定变了,变得疏离了些,或者说轻了些。往后自有分晓。但很可能还会留有一份眷恋,只要我的数学激情还没烧尽、还没餍足——只要我还“做数学”。至于它会不会先我而熄灭,我一点也不想去猜、去预言……

十多年里,我一直以为这激情已经熄灭。更确切地说,是我径自宣布了它已熄灭。那正是我暂时停下做数学、并且重新发现世界的那一天!此后三四年,我全力投入一项如此紧张的活动,旧日的激情想必找不到一丝缝隙钻出来现身。那是紧张学习的几年,不过在某个层面上还相当浮浅。在那之后的几年里,数学激情以突如其来的、完全无法预料的发作显现出来。每次发作持续几个星期或几个月,而我执意不去理会它们那相当清楚的含义。我早就一次性决定了:做数学的那股馋劲,既然明摆着百无一用,从此就是过时的玩意儿,句号!可那“百无一用”的东西并不这么听——而我这边,照旧充耳不闻。

\origpage{115} \noindent\pmark 有件事也许显得悖谬:恰恰是在发现冥想之后(1976年)、在我的生命里进来一种新的激情之后,旧激情才重新露面,而且格外猛烈,几乎算得上暴烈——仿佛每次都是一只锅盖在过强的压力下崩开。又过了五年,在形势的推动之下——这话说得不假——我才费心去察看究竟发生了什么。那是我做过的最长的一次冥想,针对一个看来范围分明的问题:我花了六个月顽强而紧张的工作,才绕着一座冰山走了一圈;它那可见的尖顶终于碍事到逼着我——几乎是违心——去看个究竟。不得不确认的是一种冲突局面,而从表面看,那是两股力量、两种欲望的冲突:冥想的欲望,和做数学的欲望。

% ---------------- p.212 ----------------
在这场漫长的冥想中,我一步一步学到:做数学的欲望——我一向鄙弃它——和冥想的欲望一样,都是孩子的欲望,而后者我捧到了天上。孩子才不管大老板、大掌柜那份鄙弃,也不管他那份谦逊的傲气!孩子的欲望一个接一个,随着时辰和日子,像一支舞的动作,一个从另一个里生出来。它们的本性就是如此。它们并不彼此对立,正如一首歌的各个诗节并不彼此对立,或者一首康塔塔、一首赋格的各个乐章并不彼此对立。是老板这个蹩脚的乐队指挥,宣布这个乐章“好”、那个乐章“坏”,在本来和谐的地方制造出冲突。

这场冥想之后,老板安分了些,不太做出把鼻子伸到不关他事之处的样子。这一回工作拖得很长,而我原以为几天就能完。工作一做完,“结果”就显得明摆着,几句话就能说出来\fnoteP{1}{(37)}{我想,几乎不必补充:这项旷日持久的工作,一天一天地还显露出别的东西,远不止我刚才以碑铭式的简短形式交出的那个“结果”。冥想的工作如此,为一个打算考察的特定问题所推动的数学工作也莫不如此。所走的那条路(它通向、或不通向对原问题或多或少的澄清)上的种种曲折,往往比原问题本身、或比“最终结果”更有意思。} (37)。可要是有个眼尖的人,在工作之前或进行中就把这几句话告诉我,恐怕也一点都帮不上我。工作之所以拖这么久,是因为阻力又强又深。再说,老板这回被骂得狗血淋头,却一声没吭,因为当时的气氛根本不容他发火。可以肯定的是,这六个月用得其所,而且是省不掉的;正如一个女人省不掉九个月的妊娠,到头来生下一个像娃娃一样“明摆着”的东西。

\section{(44) 再次倒转蒸汽}

到这里,我已经一年半没冥想了,只在十二月里有那么几个小时,为的是看清一个紧迫的问题。而这一年来,我把大部分精力投在数学上。这一“潮”来得和别的潮一样,数学潮也好,冥想潮也好:它们来时不打招呼。就算打了招呼,我也从来听不见!老板对冥想还留着一点小小的偏好,看来不得不信:每一次,我原以为会永远持续的冥想潮,转眼已被数学潮接替;而那数学潮——在我看来不过是几天、最多几星期的事——却拖宕下去,铺展成几个月,也许甚至,谁知道呢,成几年。不过老板终于明白:这些节律不是他造出来的,想去调节它们,他什么也捞不着。

但也许老板那点“小小的偏好”终究翻转过来了,因为近一年来事情已经说定、已经决定:我至少要用几年工夫去“重做数学”,可以说是正式地——我甚至申请了国家科研中心(CNRS)的一个职位!更重要、而且一年前还完全想不到的是:我又开始发表东西了。即便在刚才提到的1981年那场冥想之后,在做数学的欲望不再被当作穷亲戚对待之后,我也不会想到自己还能重新发表数学。顶多是别的东西:一本谈冥想的书,或者谈梦与做梦者的书——再说,我当时忙得根本没闲心去写一本关于它的书!
% ---------------- p.213 ----------------
写它干吗?!

于是就有了一个相当重要的决定,它牵动我今后若干年的生命进程,而它是在一种捎带着的情形下作出的,我甚至说不清是何时、怎样作出的。\origpage{116}\origpage{117} 有一天,当打字的笔记已经攒成厚厚一沓时(瞧瞧!直到那时,我那些数学上的苦思,我都还只是手写了事……\fnoteP{2}{(38)}{这些笔记其实是那封写给……的长信的继续,那封信成了它的第一章。它们被打字机打出来,是为了让这位旧日的朋友、以及另外两三位(尤其是龙尼·布朗(Ronnie Brown))读得清楚,我本以为他们可能会感兴趣。再说,那封信从未收到回音,收信人根本没有读它;将近一年之后(我问他是否收到了),他竟真心感到惊讶:我居然会有一刻以为他会读它,鉴于人们该期待我做的那种数学……} (38),关于栈、同伦模型等等),事情就成了定局:这个,发表!既然要做,索性做足,开一个小系列的数学反思,名字现成,只要把字母大写:“数学反思”!大致就是这样——此刻那团大名鼎鼎的“迷雾”向我复述的,就是这些;这迷雾常常替我充当记忆。这一回的回忆,肯定大大删削过了。无论如何,值得注意的是:这件事做成的时候,我连停下来\textbf{看看}自己往哪儿走、是什么在推我、带我的工夫都没有……这倒是我还想再做一次的:顺着这场不期而至的冥想的势头,好让我能觉得它真正做完了。

立刻浮上心头的问题是:我刚指出的这件“值得注意的事”,是老板那种(所谓的?)“谨慎”的迹象——他无论如何都不肯插手(哪怕只是投去一瞥不检点的目光……)一个如此美妙、根本不需要他的自发运动,如此等等;还是\textbf{相反},这迹象表明他已经断然站了边,而那所谓的“小小的偏好”正把他朝数学的方向全力推?

只消把问题白纸黑字写下来,答案就自己出现了!并不是那小子——他这一回玩起的把戏也许比别的把戏更长久——就此宣布自己要连着 X 年不停手地干下去,老老实实地涂黑纸张,涂够所需的那么多页,凑成漂亮一套大写字头丛书里数目可观的若干卷!是老板把一切都预见到、安排好了,小子只消照办。也许小子巴不得如此,这没法预先知道——但那是次要问题。再说,小子的兴致至少在一定程度上取决于\textbf{境况},而境况主要取决于老板。

老板选了边,再清楚不过。再说他刚显出过某种灵活:一个多月来,一场冥想就在他和蔼的目光下进行。诚然,他的和蔼绝非无私,因为冥想的有形成果,我正在写的这些笔记,将成为他已经看见自己建起的那座塔最美的奠基石,而砌塔的石头,是那位看来很听话的工人-孩子乖乖凿好的。说真的,现在夸他“灵活”还早了点儿!三个月前那几小时冥想,一年半里头总共就这么点,未免太寒碜了!

\origpage{118} \noindent\pmark 然而,我并不觉得这一整段时间里,有过什么被压抑、被挫败的冥想欲望。十二月的那几个小时里,我把事情理清了,看见了该看见的;这就足以改变一个本来不清楚的局面。我重新接上中断的数学工作,而不必把别的什么一刀斩断。我不觉得有冲突蹑手蹑脚地又回来了——我是说,那个两年多前已经化解、这一次会以颠倒的形式重新出现的冲突。老板有偏好,这是他本性如此,也完全是他权利——他要是装作禁止自己有偏好,那才蠢呢(虽然比这更蠢的事也有……)。这不是冲突的迹象,尽管冲突常由它而起。事情到了这一步,实在看不出有什么该怪他不灵活的!

% ---------------- p.214 ----------------
这一点看清了,剩下的是试着圈定老板的“动机”——为了这次蒸汽倒转;它做得再低调不过,可仔细看来,其实相当惊人。

\section{(45) 不是上师的上师——或三条腿的马}

这立刻把我带回1981年7月到12月的那场冥想;在那之前,我刚在一种数学狂热里过了四个月。这段有点疯的日子(从数学上看还颇有收获\fnoteP{3}{(39)}{这尤其是《穿越伽罗瓦理论的长征》(La Longue Marche à travers la théorie de Galois)的时期,《纲领概要》(Esquisse d'un Programme)里谈到过它(第3节:“与儿童画相关联的数域”)。} (39))随着一个梦,一夜之间结束了。那梦以一个具有野性而不可抗拒之力的比喻,描绘了我生命里正在发生的事——一个关于这场狂热的比喻。信息明如闪电,可我还是花了两天紧张的工作,才接受它那明摆着的含义\fnoteP{4}{(40)}{\emph{来访}\par 关于这个梦的工作,是一封用英文写的长信,写给一位朋友兼同行,他前一天曾在我家匆匆过了一趟。做梦者用来从一片表面的虚无里唤起这个逼真得惊人的梦的某些材料,显然取自这位亲爱的朋友那次短暂来访——我已经快十年没见过他了。因此,工作的第一天,与我以往的经验相反,我竟以为可以断定:我做的这个梦关系到我这位朋友,多于关系到我自己——该做这个梦的是他,不是我!这是一种回避梦境信息的方式,而那信息(我本该凭过去的经验一下子就明白)只关乎我一个人。我最终在工作这第一阶段——浮浅的阶段——之后的夜里明白了这一点;第二天我在同一封信里接着写下去。自那封值得记住的信之后,我再没有收到过这位朋友的任何音讯,他是我有过的朋友中最亲近者之一。\par 这次工作也是唯一一次以信件形式(而且还用的是英文)进行的冥想,因此我手上再没有它的书面痕迹。这个插曲让我格外震动——在许多别的插曲之中,它们都表明:一切超出某种门面的工作的迹象,一切把极简单、却通常被大家当作义务去忽视的事实摆到光天化日之下的工作,会在旁人心里引起多大的不适与恐惧。这一点我后面还要回来谈(见第47节,“孤独的历险”)。} (40)。这之后,我知道了自己该做什么。此后六个月的工作里,我再没回到那个梦上去,可我所做的,无非是更深入地进入它的含义、完整地消化它的信息。梦后第三天,这信息只在相当表面、相当粗糙的层面上被理解了。我尤其需要深挖的,是“我的”关系——我是说老板的——与当下这两种欲望各自的关系,而这两种欲望在我看来相互敌对。

\origpage{119} 自那场冥想以来,我生命里发生了这么多事,以至于它显得像是属于一个很遥远的过去。如果我试着把从它那里学到的、关于“老板”动机的东西说出来,就是这样:从“第一次觉醒”(1970年)算起、到那时已经过去的十二年里,老板押的注,明摆着是“\textbf{那匹劣马}”:\textbf{在数学与冥想之间}(他喜欢把两者对立起来)\textbf{他选了冥想}。

 这是一种说法,因为“冥想”这东西连同这个名字,只是五年前的1976年10月才进入我的生命。但在1970年那副刷新过的、我珍爱的自画像里,冥想六年之后来得正是时候,用它的光彩烘托出某种态度或姿态——这姿态早已被察觉,却直到1981年这场冥想才被审视。我把它叫作“大师综合征”,也有人(有道理地)称之为我的“上师做派”。我采用前一个称呼而不是后一个,多半是因为它便于把事情的性质搅混,而我乐意让自己耽在这种混淆里。我身上确实从小就有一个教人的自发乐趣,它与学的自发乐趣毫不相悖;
% ---------------- p.215 ----------------
也毫无做作的成分。在我与学生的关系里起作用的,主要就是这股力量;这关系是浮浅的,却强而纯正——我所谓“纯正”,是指没有做作。是在我称为1970年“觉醒”之后——一个我曾熟悉的天地退缩到几乎消失,连同学生、连同那些让我“教”点什么、让我把我知道的、对我有意义有价值的东西讲出来的机会一起退去——正是在那时,“老板”尽其所能地报了仇:他不再去教数学(那玩意儿只配用来糊口,除此之外配不上我新的伟大),我仿佛看见自己在用我的生活与榜样,教某种“智慧”。当然,我很小心,无论对自己还是对别人都不把这话说出来;而当我听到这个方向的回声时,我当然要矢口否认,并为这些朋友或亲近者如此不理解我而难过。我怎么解释也没用,他们就是不肯懂,真是些再叫人遗憾不过的学生!

\origpage{120} 我读过克里希那穆提(Krishnamurti)的一两本书,印象很深,脑子转眼就把某种信息和某些价值吸收了下来\fnoteP{5}{(41)}{\emph{克里希那穆提,或变成了羁绊的解放}\par 要说我从这次阅读里得到的只是一套词汇,以及一种把它据为己有、最终又理所当然地用它替换现实的倾向,那并不准确。我拿到手的第一本克里希那穆提的书之所以让我如此震动(何况我不过有闲读了几章),是因为他所说的,彻底掀翻了许多对我而言不言自明的东西;而我立刻发觉,那些都是\textbf{老生常谈},从来就是我呼吸的空气的一部分。同时,这次阅读第一次让我注意到一些影响深远的事实,尤其是“对\textbf{现实}的逃避”这一事实——它是心智最强大、最普遍的制约之一。这给了我一把钥匙,去理解此前一直无法理解、并因此(在五六年之后发现冥想之前,我并未察觉这一点)滋生焦虑的种种处境。我立刻在身边各处看到了这种逃避的实情。这解开了一些焦虑,却没有改变任何本质的东西,因为我只在旁人身上看见那实情,同时(理所当然地)设想它在我身上并不存在,总之我是那验证规则的例外(而且对这个确实了不起的例外,我再没多问一句)。事实上,我对别人、对自己都毫无好奇。这把“钥匙”\textbf{只}在怀着深入之欲的人手里才开锁。在我手里,它变成了驱邪术和做派。\par 是在1974年初,我第一次面对这样一个明证:我生命里的破坏一步步跟着我,不可能只来自别人;我身上有某种东西在吸引它、滋养它、让它延续下去。那是一个谦卑而敞开的时刻,宜于更新。可那次更新当时还停在边缘,而且短暂,因为缺少一种\textbf{深入的工作}。那“\textbf{我身上的某种东西}”还很模糊。我分明看出那是爱的匮乏,但连“做一番工作、更切近地查明我身上的爱在何处、如何匮乏,它如何表现出来,造成了哪些具体后果”等等的念头,都不可能来自我至此认识的任何圈子、任何人,也不可能来自克里希那穆提。(恰恰相反,克氏喜欢强调一切工作的徒劳,他自动把它等同于自我的“急于成为”。)于是,我拿一份借来的“智慧”当全部罗盘,除了耐心等待“爱”像圣灵的恩典一样降临到我身上,我看不出还能做什么。\par 然而,我刚刚在一个浪头的最深处学到的那点朴素的真理,却激起了一个强劲的新能量之浪,其势头可比两年半后推动我第一次投入冥想的那个浪。这股能量当时并没有完全闲置。几个月后,就在我因一场及时降临的事故动弹不得的那时,它推动了一次(写下来的)反思:一生中第一次,我审视了那个曾是我与他人关系之未言明基础的世界观,它来自我的父母,尤其来自我母亲。那时我十分清楚地意识到,这个世界观已经破产,它无法说明人与人之间关系的实情,也无助于我这个人以及我与他人的关系之舒展。这次反思仍带着“克里希那穆提风格”的印记,也带着克氏那条禁忌:任何朝向理解的\textbf{真正的工作}都不许做。可它还是使几个月前诞生、起初模糊而难以捉摸的一种认识变得切实而不可逆转。这种认识,当时世上任何书、任何人都无法给我。\par 这次反思要够得上冥想的品质,最缺的是对我的这个人、对我在我自己眼中的\textbf{形象}投去的目光,而不只是对世界观的目光——也就是对一个公理系统、而我并未真正“有血有肉”地出现在其中的目光。它还缺的,是在\textbf{当下}、在反思进行的那一刻对我自己的目光(这次反思仍停留在真正的工作之下);这目光本会让我同时看出一种借来的文风,以及对这些笔记文学外表某种沾沾自喜,也就是缺乏自发性、缺乏真实。无论它多么不足,对我与他人的关系当下影响也相对有限,这次反思在我看来仍是一个阶段——鉴于起点,这大概是必经的阶段——通向两年后应当发生的更深的更新。正是在那时,我终于发现了冥想——在发现这第一个意想不到的事实时:\textbf{关于我自己这个人,竟还有东西可发现}——那些东西几乎完全决定了我一生的走向,以及我与他人的关系的性质……} (41)。无需再多,就足以让我相信一切都已
% ---------------- p.216 ----------------
到手(当然,嘴上还否认)。我不需要再读;无论口头还是笔下,我都能即兴来一套最纯正的克里希那穆提,一套滴水不漏、环环相扣的话语。可话语再漂亮、再严密,也从来一刻都不像管用,无论对我还是对别人。这样过了好几年,我一点也没做出要从中吸取教训的样子。

随着发现冥想,那套行话一夜之间从我身上脱落,不留痕迹。那时我才懂得话语与认识之间的全部差别。

大老板马上校正了准星:克里希那穆提打入冷宫,冥想捧上台面!自然是悄悄地;如今他得换一套全然不同的手法。时势变了:这小子如今在他腿脚间乱窜,而且有时候眼神挺利。想必小子当时正忙着别处。总之,直到五年之后,某口锅炸开了、小子跑来看出了什么事,大老板的把戏才被识破。

说到底这事过去不算久,才两年多一点:那位“不露相的上师”终于被揭穿——又一副伪装进了冷宫!可怜的老板,他几乎要赤条条地见人了。换个说法:那匹“冥想”马顶替了那匹无名的马(那匹马可千万不能叫“克里希那穆提式”!),如今它带来的押注回报实在少得可笑,尤其比起“数学”马当年那些可观的回报——那时老板还押在它身上。他之所以这么久守着那笔烂注,纯粹是惯性——他已经换过一次注,这本来就不多见,而且当时全靠一记猛击的事件才办到\fnoteP{6}{(42)}{\emph{有益的拔除}\par 所说的那桩“猛击”事件,是1969年底发现:我觉得自己属于的那个机构,有一部分经费来自军事部的款项——这与我基本的公理不相容(而且今天仍然不相容)。这一事件是一整串事件(一个比一个更能说明问题!)中的第一桩,其后果是我离开了高等科学研究所(IHES),并一针引一线地彻底换了圈子、换了投入。\par 在IHES的英雄年代,迪厄多内(Dieudonné)和我是它仅有的成员,也是仅有的为它赢得\textbf{信誉}和听众的人:迪厄多内靠编辑《数学出版物》(Publications Mathématiques)(第一卷早在1959年就出了,即莱昂·莫尚(Léon Motchane)创办IHES的次年);我靠“代数几何讨论班”。最初那些年,IHES的存在极其朝不保夕,经费没有着落(靠几家充当赞助人的公司的慷慨),唯一的场所是巴黎梯也尔基金会(Fondation Thiers)借给的一间屋子(借得明显不情愿),只在我开讨论班的日子用〔IHES最近为建所二十五周年印发的一本小册子(尼科·凯珀(Nico Kuiper)好意寄给我一册)对这些\textbf{艰难}的开端只字不提,也许觉得它们不配去年那场大张旗鼓的庆典的庄严〕。我有点觉得自己和迪厄多内一起,是我所属这个机构的“科学”\textbf{共同创办人},而且我打算在此终老!我最终对IHES产生了强烈的认同,而我的离开(由于同事们的漠然)被体验成一种从另一个“家”被拽离,后来才显出是一场解放。\par 回过头看,我意识到我身上想必早就有一种更新的需要,我说不清从何时起。我离开IHES的前一年,我的精力投入忽然翻了个方向:把前一天还烫手的工作、把我最着迷的问题都撂在一边,转而去搞生物学(是在一位生物学家朋友米尔恰·杜米特雷斯库(Mircea Dumitrescu)的影响下)——这肯定不是单纯的\textbf{巧合}。我当时是以在IHES内部长期投入的架势扑进去的(这与该机构多学科的宗旨相符)。那肯定只是对一种深得多的更新之需的宣泄,而这种更新在IHES那“科学蒸笼”的气氛里无法完成,它是在我前面提到过的那场“觉醒的瀑布”中发生的。前后共有七次,最后一次在1982年。“军费”那段插曲是天意,它触发了这些“觉醒”中的第一次。军事部,和我IHES的老同事们一样,最终都得到了我全部的感激!} (42)。老板们不太喜欢换注——而这一次甚至算是往回走,回到上一笔注。

从1973年我退居乡间起,这匹新马带来的回报,比起从前那匹,就开始少得可怜。三年后冥想的意外降临,又让它们略略活络。甚至还有1979年3月到7月那阵令人头晕的高峰——这一段我在这里不多说——那时我又一次显出使徒的模样,这一回是一种既古老又新鲜的智慧的使徒;这智慧唱在我自己作的一部诗作里,而我最终没有把它交到出版商手上\fnoteP{7}{(43)}{\emph{“我自己作的那部诗作”} 里面有许多我第一手就知道的东西,它们今天在我看来,在我的生命里、以及在一般“生命”里,和它写成时一样重要——当时我是打算发表的。我之所以没发表,主要是因为后来我发觉,它的形式患有一种刻意“作诗”的毛病,以致它的整体构思过于搭建,许多段落缺少自然,有些地方甚至僵硬或浮肿得叫人难受。这种时而浮夸的形式,是我当时心性的映照;在那里,说到底常常是“老板”在领着跳舞——当然是笨重地领……} (43)。可是两年之后,上师彻底歇了业,就好像“冥想”这匹马也断了一条腿(就它带给老板的回报而言)——不管有没有手法,都没法再扮上师了!

% ---------------- p.217 ----------------
\origpage{121} 这以后就没怎么再拖——三条腿的马进了冷宫,连同使徒诗人、“不是上师的上师”,以及“不敢说出自己名字的克里希那穆提”。数学万岁!

下文如何,我们饶有兴致地等着……

% ---------------- p.218 ----------------
% (原书 p.218 为章末空白页)


% ================================================================
%  第 11 章 孤独的历险 (原书 p.219-230, 边码 p.122-p.139)
%  源文: _work/src1/task-ch11.md
% ================================================================
% ---------------- p.219 ----------------
\chapter{孤独的历险}

\section{(46) 禁果}

\origpage{122} 记笔记的事,我不得不停了两天。仔细重读之后,我看先前那套戏路子大致上确实是实情的写照——只是这写照现在还得再往深里挖一挖。我尤其该把冥想与数学这两匹“马”各自的长处圈得更准些;也该试着弄明白:究竟是哪些事件、哪些机缘,最终扣动了老板押注的“翻转”——明明有惰性之力拽着他,要他没完没了地守着一注,哪怕那注已经输了。

也许还该探一探孩子自己的偏爱。如今这点已经说定:他时不时就想换一换局,而老板显然也留着最起码的一点伸缩,不至于硬逼着他总玩这个、从不玩那个。这几年,老板学会了把孩子算进去,学着跟他将就,不必等到锅子炸开。这还不是完全的和谐,但已不再是战争,更像是某种诚挚的协约——偶发的紧张往往让它松动,而不是变硬。

只要不被硬顶得太狠,孩子的偏爱天性上是相当柔软的。(这跟老板可不一样——他到头来学到的那一丁点柔软,全是违着自己的意愿、到了晚年才学会的……)不过孩子柔软,并不等于他没有偏爱,不等于他不会对这一样东西的兴致比对那一样更浓。

要把这里头看清,常常一点也不容易:孩子的欲望、老板的偏爱,乃至老板一劳永逸定下的东西,三者之间很难分清。从前我对自己说:冥想比数学更好、更重要、更严肃,如此等等,理由如此如此(不用问,都是最站得住的理由)——那是老板事后给自己找的好理由,好让自己相信他下的那一注确实“押对了”。孩子从不说这样东西比那样“更好”“更重要”。他不爱讲话。他想做什么,只要没人拦,他就去了,从不问这事“重要”还是“更好”。他的兴致,对这件事对那件事、这一刻下一刻,有强有弱。要察觉他的偏爱,听老板那些解释性的话毫无用处——他自称代表孩子说话,其实只能代表他自己。只有观察孩子在玩什么,或许才勘得破他的偏爱。而且就算这样也不那么容易:他兴冲冲地玩着这一样,并不总意味着,要是老板不伸他那根手指头去拨一拨,他就不会欢天喜地地去玩另一样。

% ---------------- p.220 ----------------
\origpage{123} 显然,比其他一切更吸引他的,是未知——是追进夜里那些含混不清的褶子中去,把未知的东西带到光天化日之下,他自己的未知,也是所有人的未知。我感觉到,我说“也是所有人的”的时候,这确实是孩子的欲望,而不是老板的虚荣——那种想叫满堂看客、也叫自己吃一惊的虚荣。还有一点也是说定了的:孩子每回从阁楼、从取之不竭的地窖的幽暗中带回来的,都是些“显然”的、孩子气的东西。它们越是显得显然,他就越高兴。要是它们不显然,那就是他没把活干到底,在幽暗与天光之间半路上停下了。

\origpage{124} 在数学里,“显然”的东西,也是那种迟早总有人会撞上的东西。它们不是“发明”,可以有人做、也可以没人做。它们本来就在那里,从来就在;人人都在它们边上走来走去而不留意,顶多绕着它们兜个大圈子,或者每次都绊一跤、从它们身上跨过去。过了一年,或者过了一千年,难免有人终于留意到那东西,围着它挖,把它刨出来,翻来覆去地看,擦干净,末了给它起个名字。这类工作——我偏爱的工作——每一回换个人也能做;而且,换个人早晚也非做不可\fnoteP{1}{(44)}{不用说,我在这里把一个假设搁在一边了——至少可以说,这个假设绝非不可能:一场原子战争,或诸如此类的另一种“乐子”,猝然闯了进来,一举而永久地终结了那场叫做“数学”的集体游戏,顺带也终结了许多别的东西……} (44)。

在我自己的发现这件事上就完全不一样了——在这场绝非集体的游戏“冥想”里。我所发现的东西,世上没有任何别的人,今天也好、任何别的时刻也好,能替我发现。它只属于我一个人去发现,也就是说,也属于我一个人去承担。这个未知,并不因为事理之必然、几乎不管我肯不肯费心去理会它,就注定要被知道。它在沉默里等着被知道的那一刻;有时时候到了,我听见它在呼唤——而蒙召去认识它的,只有我一个人,只有我里面的那个孩子。这不是一个“缓期执行”的未知。当然,我尽可以听从它的呼唤,也可以躲开,说“明天”,或者说“总有一天”。但那呼唤是冲着我、也只冲着我来的;除我之外没人听得见,除我之外也没人跟得上。

每当我听从了这呼唤,“企业”里就有什么东西变了,多多少少。效果是即时的,当下就被感到是一种恩惠——有时是一种突如其来的解放,一种巨大的宽释:卸下了一副我常常连自己都没察觉在扛着的重担,而它的真实,恰恰由这宽释、这解放显出来。在振幅小一些的音阶上,这类体验在一切发现的工作里都很常见,我也曾有机会谈到过。然而有一件事,使自我发现的工作(不论它做在光天化日之下,还是留在地下)区别于任何别的发现工作:它真的改变了“企业”本身里的什么东西。这不是量上的改变,不是产量的增加,不是产品在尺寸上、甚至在品质上的差别。这是老板与工人—孩子之间关系的改变。也许连老板自己也变了——如果这话除了就他与工人、与孩子的关系而言,还能有别的意思。比如他也许不那么盯着产量了——但这也仍是他与工人关系的一面,因为一种关切、也许一种敬意出现了,而那从前对他是陌生的。凡我冥想过的场合,变化都朝着老板与工人之间关系变得清朗、变得平静的方向。除了某些冥想停留在表面的情形——那是在某种当下而有限的需要压迫之下做的“应景”冥想——那清朗一直延续到今天,平静也是。

% ---------------- p.221 ----------------
\origpage{125} 这就使自我发现的工作带上了别的发现工作所没有的一层意义,尽管许多本质的方面是共通的。在自我的认识里,在自我发现的工作里,有一个维度,使它们区别于一切别的认识、一切别的工作。也许这就是知识之树上那颗“禁果”。也许冥想在我身上施展的那种魔力——更确切地说,是冥想向我揭示其存在的那些奥秘的魔力——正是禁果的魔力。我跨过了一道门槛,恐惧在那里消失了。认识的唯一障碍是一种惰性,一种有时大得可观、但终归有限的惰性,绝非不可逾越。这惰性,我几乎每走一步都感到它,阴魂不散,无处不在。它有时把我惹得冒火,却从没让我泄气。(数学工作里也一样——那里主要障碍也是它,只是分量轻得没法比。)这惰性成了游戏的基本配料之一;更确切地说,成了这游戏里的一位主角——这游戏精巧,而且全不对称,其中有两位——更确切地说三位主角:一边是纵身一跃的孩子,和那(化作惰性的)老板,他竭尽全力地刹车(一面还声称自己不在场);另一边是那美丽未知者隐约可见的身影,满是神秘,既近又远,一面躲闪,一面召唤……

\section{(47) 孤独的历险}

“冥想”在我身上曾经有极大的魔力——跟从前“女人”的吸引力一样大,而它似乎已经取代了后者。我刚写下“曾经”,这并不表示这魔力今天已经熄灭。这一年我把自己投进数学里,它只是退到了背景中。经验告诉我,这局面一天之内就可能翻过来,正如这局面本身也正是一次全然未料的翻转的结果。其实,在我经历过的四个漫长的冥想时期里(其中一个延续了将近一年半),每一次我都理所当然地以为,我会顺着这股势头一直走到最后一口气,尽可能深入地去探生命之谜、人之存在之谜。当笔记堆成吓人的一摞摞、眼看要把我的工作室淹掉时,我甚至终于请人量身定做了一个柜子来安放它们,而且留得很宽(凭一个快速的算术级数推算),好把往后几年里很快会添进去的也一并安放;要是我没记错,我留了大约十五年的余量(这已经相当可观了!)。这一回老板事情办得漂亮,论管家,这是漂亮的管家功夫!这一手,加上把所有与冥想工作或近或远沾边的私人文件来一次大刀阔斧的整理,是他最后着手、并且(几乎)圆满完成的一件差事——就在偏爱与押注翻转之前。这不禁让人琢磨:他是不是心里早有另一层打算,是不是已经看见一册册《数学反思》塞满了那些空着的格子——那格子据说本是留给日后的“笔记”的。

当然,冥想之爱、自我发现之爱,辽阔得足以填满我这辈子直到最后一天。数学之爱也确实没有被耗尽,但也许这份饥渴在往后几年里终会吃饱。我里面有什么在这样盼着,并且觉得数学是一道羁绊,拦着我走那场只有我一人能走的孤独历险。而且我觉得,我里面这个“什么”不是老板,也不是老板的哪一桩心血来潮(按本性,他是分裂的)。我觉得数学之爱还带着老板的印记;无论如何,跟着它走,只是让我的生命在一个封闭的圈子里打转;一个轻省的圈子,一种惰性的运动,绝不是更新。

% ---------------- p.222 ----------------
\origpage{126} 我曾追问:数学之爱在我生命里这样执拗地持续,究竟是什么意思。跟着它走的时候,它并不真的填满我的生命。它给人快乐,也给人满足,但它本身并不足以给人真正的舒展、一种饱满。像一切纯粹智力的活动一样,强烈而长期的数学活动,效果毋宁说是使人变钝。我在别人身上看到这一点,尤其在我自己身上——每当我重新扑进去的时候。这活动太零碎了,它只动用我们直觉、感受诸能力中极小极小的一部分,那些能力由于久不使用而变钝。长期以来我并不觉察,而显然我的大多数同行,当年也不比我更觉察。我想,正是从我单单做冥想以来,我才留意到这件事。稍一留意,它刺眼得很——大剂量的数学使人愚钝。即便在两年半前的那场冥想之后——那场冥想承认了数学之爱确实是一种爱,是我生命里重要的一件事——如今我献身于这爱时,仍留着一份保留,一份迟疑,不是完全的献出。我知道,所谓“完全的献出”其实会是一种退位,会是顺着一种惰性走,会是一种逃避,而不是献出。

对冥想,我身上没有丝毫这样的保留。我投身其中时,是全身心投入,这献出里没有一点分裂的痕迹。我知道,在献出的时候,我与我自身、与世界完全一致——我忠于我的本性,“我就是道”。这份献出对己对人都是好的。它把我向自己敞开,也向他人敞开,用爱去解开我里面还结着的结。

\origpage{127} 冥想把我向他人敞开,它有力量解开我与他的关系,哪怕对方仍旧结着。可是,要找到机会与人多少谈一谈冥想的工作、谈一谈这工作让我认识的这样那样的事,是极难得的。这绝不是因为那些事“太私人”。打个不贴切的比方:对某个时候让我着迷的数学,我只能跟一位具备必备行头的数学家谈,而且他还得同时乐意对之发生兴趣。有时一连几年,我被某些数学东西迷住,却碰不上(甚至不去找)另一个能与之谈论它们的数学家。但我很清楚:只要我去找,就会找到;即便找不到,那也只是运气或时机的问题;让我感兴趣的东西,不可能不使某个人、甚至某些人感兴趣,是十年后还是一百年后,说到底无关紧要。正是这一点,赋予我的工作以意义,哪怕这工作是在孤独中做的。假如世上再没有别的数学家,而且将来也不会再有,我不相信做数学对我还会有意义——我猜,对任何别的数学家、任何别的“研究者”,情形也一样。这就接上了前面那个观察:对我来说,“数学的未知”是还没有任何人知道的东西——它不取决于我一个人,而取决于一个集体的实在。数学是一场集体的历险,已经延续了几千年。

% ---------------- p.223 ----------------
\origpage{128} 至于冥想,要谈论它,“行头”的问题根本不存在;至少到我这个地步不存在,我怀疑它将来也不会存在。唯一的问题,是别人身上有没有一份兴趣,能与我身上这份兴趣相应。所以那是一种好奇心,对着自己和别人身上真正在发生的事——越过那些必不可少的门面;只要人真的想看看门面底下盖着什么,那些门面就藏不住多少东西。但我已经知道,一个人身上出现这种兴趣的时刻,“真相的时刻”,是稀少而短暂的。当然,常能碰到所谓“对心理学感兴趣”的人,读过弗洛伊德、荣格,还读过许多别人,巴不得来几场“有意思的讨论”。他们随身带着那份行头,或轻或重,就是人们所谓的“文化”。它是他们对自己那幅形象的一部分,又加强那幅形象,而那幅形象他们是从来不肯去查看的——跟另一位对数学、对飞碟或者对钓鱼感兴趣的人一模一样。我刚才想说的,不是这种“行头”,也不是这种“兴趣”——尽管同样的字眼在这里指的是性质不同的东西。

\origpage{129} 换一种说法:冥想是一场孤独的历险。它的本性就是孤独的。不仅冥想的工作是一种孤独的工作——我想,一切发现的工作都是如此,哪怕它嵌在一项集体工作里。但由冥想的工作诞生的认识,是一种“孤独”的认识,一种无法分享、更无法“传达”的认识;或者说,它若能被分享,也只在稀少的瞬间。这是一种工作、一种认识,它与最根深蒂固的共识逆向而行,它让每一个人、让所有人不安。这认识当然表达得简单,用些简单清亮的词。当我向自己表达它时,我在表达中学习,因为表达本身就是工作的一部分,由一股强烈的兴趣托着。可是这些同样简单清亮的词,一旦撞上冷漠或恐惧的紧闭的门,就无力向别人传达任何意思。就连梦的语言——气力完全不同,资源无穷,由一位不知疲倦、心怀善意的“做梦者”不断更新——也穿不过那些门……

没有哪一种冥想不是孤独的。只要有一丝想得到谁认可的操心,想得到证实、得到鼓励,就没有冥想的工作,也没有自我的发现。有人会说,任何真正的发现工作,在工作进行的那一刻,也是如此。不错。但在工作本身之外,他人的认可——无论来自亲近的人、同行,还是自己所属的整个圈子——这认可对于这工作在一个献身者生命中的意义,是重要的。这种认可、这种鼓励,是最有力的激励之一,正是它们使“老板”(还是用这个形象)毫无保留地开绿灯,让孩子尽情尽兴。主要正是它们决定了老板的投入。我自己在数学上的投入,情形也一样——那是被嘉当、施瓦茨、迪厄多内、戈德芒这些人,以及他们之后的另一些人的善意、温暖与信任鼓励起来的。相反,对冥想的工作,丝毫没有这样的激励。这是工人—孩子的一种爱,老板心底里还算厚道,多多少少容着它,因为它“不带来”任何东西。它当然结果,但那不是老板巴望的果子。只要他不在这件事上自己骗自己,事情就很清楚:他不会往冥想里投钱——老板天性合群!

只有孩子,天性上是孤独的。

\section{(48) 馈赠与接纳}

昨天谈到冥想的孤独本性时,有个念头掠过:我写了将近六个星期的这些笔记,末了竟成了某种冥想,可它们却是要拿去发表的。这当然必定以许多方式影响了这冥想的形式,尤其是那份求简洁的用心,还有那份求谨慎的用心。冥想的一个本质方面,即对工作当下我身上发生之事保持持续的留意,只偶尔显现,而且浮在表面。这一切肯定影响了工作的进程和品质。但我仍感到它有冥想的品质,首先凭它果实的性质,凭一种对我自己的认识浮现了出来(在此情形下,尤其是对某一段过去的认识),而那是我此前一直回避的。另一个方面是自发性:近五十个“节”或“段”,反思是自发地聚成这些的,而每一个,我在开头时都说不出它的实质会是什么;每一次,实质都只在途中才显露;每一次,工作都带出一些新的事实,或者以新的光照明了此前被忽略的事实。

% ---------------- p.224 ----------------
\origpage{130} 这工作最直接的意义,是与我自己的对话,因而是一场冥想。然而,这场冥想是要发表的,而且还要充当随后那部《数学反思》的“开场”——这一点绝不是可有可无的处境,绝不是在写作过程中不起作用的东西。对我而言,它是这工作意义的一个本质部分。昨天我暗示老板肯定能在其中“捞到他的好处”(他是“捞好处”的老手,什么都少不了他那一份!),这丝毫不意味着它的意义就缩到那一点上——缩成那匹著名的三脚马迟来的、几乎算死后才发生的“回归”!也不止一次,我感到一个行为深层的意思,有时会超出促成它的那些动机(显的或隐的)。而在这“回归数学”里,我猜度着还有另一层意思,不只是某些心理力量相加的结果——那些力量在某个时刻、出于某些缘由,在我这个人身上碰到了一起。

这场我正在进行的“冥想”,是要献给我在数学世界里认识并爱过的人——如果我感到它是那份隐约见到的意义的重要部分,那并不是因为我在盼着这馈赠被人接纳。它被不被接纳,不取决于我,只取决于它所寄予的那人。它被接纳,我当然绝非无所谓。但那不是我的责任。我唯一的责任,是在我所作的馈赠里是真诚的——也就是说,也是做我自己。

冥想让我认识的,是些卑微而显然的事,一些不登大雅之堂的事。也是那些我在任何书里、任何论著里都找不到的事——不论那书多么博学、多么深刻、多么天才——那些没有人能替我找到的事。我盘问过一团“雾”,我费心去听它,我学到了一点卑微的真相,关于一种“竞技态度”及其显然的意义——在我与数学的关系里,也在我与他人的关系里。就算我把《圣经》、古兰经、奥义书都“照原文”读过,再加上柏拉图、尼采、弗洛伊德和荣格,我是个学问广博精深的天才——这一切也只会让我离那真相更远,那是孩子气的、显然的真相。就算我把基督的话重复一百遍——“像小孩子那样的人有福了,因为天国是他们的”——再细细给它作注,那也不过仍旧使我离我里面的孩子更远,离那些让我不舒服、只有孩子才看得见的卑微真相更远。就是这些东西,是我所能献出的最好的。

而我很清楚:这样的话被说出来、被献出去,用的是简单清亮的词,却并不因此就被人接纳。接纳,并不只是收到一条信息,不管是尴尬地收到,还是甚至感兴趣地收到:“嘿,这谁想得到……!”或者:“说到底也没那么稀奇……”。接纳,常常是在献出者身上认出自己。是借着别人的这个人,与自己相识。

\section{(49) 一条裂缝的确认}

这段关于当前工作意义、关于馈赠与接纳的短小思考,像是反思脉络里的一段离题;或者更像是例证,说明某些方面——它们使“冥想”区别于一切别的发现工作,尤其是数学工作。昨天我明白了:这些方面有双重效果,也就是两种方向相反的效果:对“孩子”是一种独一无二的魔力,对“老板”则是全然的无趣。看来这双重效果就在事物的本性里,绝不可能靠什么妥协、什么安排来缓和。不管怎么做,孩子一跟从他真正的偏爱,老板就在里头捞不到他的好处——一点也捞不到!

% ---------------- p.225 ----------------
\origpage{131} 毫无疑问,已经发生的那次翻转,意义就在这里;它很可能在往后几年里把冥想从我生命中一扫而空(除了像三个月前那样的“应景冥想”)。我并不认为,因此那几年就会完全荒芜,正如过去这一年也并不荒芜。但同样确实的是:我在这一年里学到的东西(数学以外),跟此前四年中任何一年学到的东西相比,都微不足道。怪的是,我经历过的四个漫长冥想时期,每一个都是极其饱满的时光,没有任何迹象让人怀疑我心里还有什么落空的东西。然而,锅子既然炸了,就说明某处有压力,而这压力不会是当天才有的;它一定已经存在,在我看不见的某处,存在了几个星期或几个月,而那时我正专注而全然地沉在冥想里。

\origpage{132} 不过这里我是被笔头的冲劲带着走了(更确切地说,是打字机的冲劲)。实情是:(除了最后一个冥想时期——它正冲到半途,就被一连串事件与情势截断了)冥想的强度从某一刻起渐渐低落,就像一道浪,而后面另一道浪正等着接它的位置……说实话,饱满之感也随着这同一运动,只是有个分别:它只在“冥想浪”的时候在场,不在“数学浪”的时候在场。

我试图圈定的这个局面,在我看来已经不再是冲突的局面,但看来它里面仍旧含着冲突的胚芽、冲突的可能。如今它对我来说,也许是凭它对我生命进程的冲击,我内部分裂的最显眼的迹象。这分裂不是别的,就是老板与孩子的分裂。

我无法把它了结。既然它在这一处表现里已经被勘破,我能做的,只是留意它,在往后的月份和年份里追索它的迹象与演变。也许这份对数学的爱——不得不说,有点不合时宜——会在燃烧中烧尽(像我身上另一份爱已经烧尽那样……),让位给那唯一的爱:对自我的发现,对我命运的发现。

这份爱足够辽阔,我说过,足以填满我的生命——而我的整个生命肯定也不够把它耗尽。

\section{(50) 一段过去的重量}

我把《收获与播种》最后收尾完工,已经是几天前的事了——此前一个多月,我一直以为再过几天就要完工。这一回也是,在“收尾”之后,我仍不完全确定自己是否真的完工了——因为有一个问题我确实悬着没解决。那就是“弄明白:究竟是哪些事件、哪些机缘,最终扣动了‘老板’押注的‘翻转’”,使数学取代了冥想的位子,而且是与巨大的惰性之力相抗。这几天,我的思绪不请自来地、相当执拗地回到这个问题上——而这几天我其实已经开始接上另一些全然不同的事,包括一些数学问题(共形几何方面的)。不如再借一借这股冥想的“末势”,多少往深里挖一挖,把地方腾干净。

% ---------------- p.226 ----------------
当我试着“凭感觉”回答为什么“我又弄起数学来了”(指的是那种大规模的、打算长期干的投入,至少以几年计),好几个联想就冒出来。其中最有力的,也许要数我近六七年来在教学活动中终于感到的那种慢性落空之感。就是那种越来越强的感觉:我“大材小用”;而且常常是,我把力气、把我最好的东西投给一些闷闷不乐的学生,而他们对我能给的东西毫不在意。

\origpage{133} 我到处看见值得做的妙事,而它们只等着人去做。往往,一点可怜的行头就足够上手;倒是这些东西自己告诉我们:该发展出什么样的语言来圈住它们,该置办什么样的工具去挖它们。我没法不看见它们——只要由于教学活动而与数学保持着经常的接触(不论那层次多么低),哪怕在我对数学兴趣最边缘的那些时期,也是这样。每一件瞥见的东西背后,只要稍一发掘,又有别的好东西,遮住一些、又露出一些……不论是在数学里还是在别处,只要带着真正的兴趣把目光放上去,就看见一种富饶显露出来,一种深度展开,叫人猜出它们取之不竭。我说的那份落空,就是怎么也没法把这种富饶之感——深之感——哪怕一星半点传给学生:哪怕只是一点想至少把他们手边的东西摸个遍的兴致,在他们无论如何都打算投进所谓“研究”活动的那几个月或几年里,尽情尽兴地干一场——反正也是为了拿这个那个文凭。十年来我教过的学生里,除了两三个,看起来“尽情尽兴”这个念头本身就让他们害怕;他们宁愿几个月、几年地垂着手原地蹭,或者苦哈哈地干一份鼹鼠式的活,那活的来龙去脉他们一概不知,只要末尾有个文凭就行。关于这种创造力的瘫痪,有很多可说的,它与“天赋”或“才能”的有无毫无关系——这又接上了我反思的最开头,那里我曾顺带触及这类阻塞的深层原因。但这不是我在这里要说的话;我这里要说的,毋宁是确认一种慢性的落空状态——这些情形,在最近七年教学活动中不断重演,终于在我身上造成了它。

\origpage{134} “解决”这样一份落空,最现成的办法——至少就它是我身上那位“数学家”的落空、而不是那位教师的落空而言——就是我自己去动手做那些事,至少做其中一部分;那些事我曾绝望地盼着哪个学生终于有一天会接过去做。这也正是我多少在这里那里做过的事:或者借教学活动之余、借教学活动的机会,花几个小时、甚至几天作一次偶然的思考;或者在数学上大发作的时期(那有时来得像真正的爆炸……),一发作就是几个星期、几个月。这样偶发的、一阵一阵的工作,往往只能对一个问题作最初的一遍粗加工,只能给出最零碎的视野——不如说,那是对尚待进行的工作看得更清楚一些;而那工作本身始终搁在那里等着做,而且正因为看得更清,就显得更烫手。两个月前,我给出了一份总体草图,涉及那些我才刚刚开始约略摸到门径的主要题目。那就是《纲领概要》(Esquisse d'un Programme)——我此前已有机会提到过它,它最终将附在这份反思之后,二者合起来构成《数学反思》的第一卷。

相当清楚:单靠这份(不妨说“私下的”)勘察工作,不足以化解我的落空。“大材小用”这个感觉,无疑翻译出一种想有所行动的欲望(我想,它源于自我,也就是说,“老板”的欲望)。这里说的行动,与其说是对他人的行动(比如对我的学生:推动他们,向他们“传达点什么”,或者帮他们拿到某个文凭,好去申请某个职位,等等),不如说是“数学家的”行动:为发现某些无人料到的事实、为某个理论的绽出出一份力,等等……这立刻接上了前面那个观察:数学是一场“集体的历险”。

% ---------------- p.227 ----------------
\origpage{135} 如果我追问自己近十年来做数学时的心境——那是我生命中一段根本不会想到有朝一日还能重新发表东西的时期,而且当时也多少清楚:我眼下或未来的学生,没有一个会用到我的勘察工作——我立刻就看到:那绝非某个只为个人乐趣做事、或者被一种只关乎自己而与他人无关的内在需要所驱赶的人的心境。我做数学的时候,我相信我心里某处是明白的:这些数学是做来传达给别人的,是要成为一件更广大之物的一部分的,我为它出力,而那东西绝不是个体性的。这“东西”,我可以称之为“数学”,或者更好的说法,“我们对数学事物的认识”。这里“我们”一词,首先无疑具体地指我所认识、与我共有某些关切的那些数学家构成的团体;但毫无疑问,它也超出这个狭小的团体,正如它超出我个人。这个“我们”指的是我们的族类——就这一族类通过它在各个时代的某些成员,曾经并且正在关切数学对象世界的诸般实在而言。直到此刻我写下这几行之前,我从未想到过我生命中存在着这“东西”,更没想到要追问它的本性,追问它在我作为数学家与教师的生活里起什么作用。

我提到的那个想有所行动的欲望,在我身上、在我作为数学家的生活里,似乎采取这样的形式:把所有人都不认识的东西从暗处取出来——不只是我不认识的东西(如前所见)——而且,还要把它交给所有人支配,也就是去丰富一份共同的“家产”。换句话说,这是一种欲望:为那超出我个人的“东西”、那份“家产”的扩大与丰厚出一份力。

这份欲望里,当然不缺借自己的作品来扩大自己的欲望。从这一面看,我又碰上了那种对“成长”、对扩大的饥饿——那是自我、是“老板”的特征之一;那是他侵占性的一面,推到极处便是毁灭性的一面(参见注44,第13.1.1节,原书p.260)。不过我同时也明白:想增加那些(或久或暂)多少会挂上我名字的东西的数目,这欲望远不足以穷尽、远不足以覆盖那股更广大的欲望或力量——它推着我去为扩大一份共同家产出力。我觉得这样一份欲望本来可以得到满足(如果不是“在我的企业里”——那里老板还相当有侵占性——至少在一个更成熟的数学家身上),而他自己这个人的角色仍是无名的。这也许就是自我扩大倾向的一种“升华”形式:通过与一件超出自身之物的认同。除非这类力量本身并非源于自我,而是属于一种更细腻、更深的本性;除非它表达的是一种深刻的需要,独立于一切制约,见证着一个人生命与整个族类生命之间那深深的联结——那联结是我们个体存在意义的一部分。我不知道,而追究这样一些如此宽泛的问题,也不是我在这里要说的话。

我要说的,毋宁是(以一种更朴素的角度)考察一个与我这个人有关的具体局面:一个落空的局面,靠零星断续的数学活动得到局部的、暂时的排遣。这样,这局面的逻辑迟早该把我引向把我找到的东西传达出去。既然直到去年,我根本不肯为我的数学之爱付出那种大规模的、长期的投入——那本是必要的:得靠一番因事制宜的“逐件加工”,才能把我掘出的矿藏“开采”出来拿去发表——那么剩下的另一条路,就是把我至少最放在心上的那些东西,传达给几位足够“在行”的数学家朋友。

我想,要是在最近十年里我碰上一位数学家朋友,对我扮演对话者和消息来源的角色(就像五六十年代那许多年里,塞尔在很大程度所扮演的那样),同时又充当一道中转,把我能传给他的“消息”传出去(这角色从前塞尔不必扮演,因为我自己就担了!)——那么我想“在数学上有所行动”的欲望,就会得到足够的满足,足以化解我的落空,而我只消在数学上断断续续、有节制地投些精力,把最大的份额留给我新的爱。我第一次带着这样的期待(至少心里暗含这期待)去找一位数学家朋友,是1975年,最后一次是1982年,一年半以前。有趣的巧合:两次都是为了想办法“安置”(好让它传开去,谁知道呢,说不定最终还会有人发展它!)同一个“纲领”,同调与同伦代数方面的,其最初的胚芽要上溯到五十年代,而早在六十年代末以前就已经完全“成熟”(依照我内心的确信);这纲领的一个初步的、粗线条的展开,恰恰就是《追寻栈》的主题——而我此刻按说正该给它写导论!不管怎样,由于大概各不相同的缘由,我那些想重新找回一种“特权对话者”关系的尝试,都夭折了——那种关系从前(1970年以前)曾与塞尔有过,后来又与德利涅有过。不过有一个共同的情形:我准备给数学的时间相对有限。在我说的那两次(1975年和1982年),这肯定也促成了交流的瘸腿。其实,我主要想的是把什么东西“安置”出去,并不太费心去做那件必要的事——“(重新)跟上行情”,以便我这方面成为一个让通信者满意的对话者;而对方在行内那些通行技术上比我“在行”得多(至少可以这么说!)。

% ---------------- p.228 ----------------
\origpage{136} 我大可以把《致……的信》——它充当《追寻栈》的第一章,去年二月写的信,才一年多以前——看作我最后一次尝试:在我一位旧日朋友那里,为我如今某些想法与关切找到一点回响。从这封信开始(更确切地说,重新开始)的反思,后来竟要成为1970年以来第一篇注定要发表的数学文本——而这一点,我当时好几个星期都没有察觉。差不多一年之后,我才收到对这封内容充实的信的间接反应(参见注\fnoteP{2}{(38)}{这些笔记其实是那封写给……的长信的续篇,后来那封信成了第一章。它们打字打出来,是为了让这位旧日的朋友读得下去,另外还为了两三个人(尤其是龙尼·布朗[Ronnie Brown]),我以为他们会有兴趣。这封信其实从未得到回音,收信人也没读过它:将近一年之后(我问起他是否收到了),他真心感到惊讶——我居然会有一刻以为他可能读它,看看人们该指望从我这儿得到的是什么样的数学……} (38))。这反应比迄今任何一位数学同行写给我的信都更能说明问题,让我感到某些针对我这点微末之身的姿态——自我离开那个曾与他们共处的圈子以来,这些姿态在我的数学家朋友中间已经成了常事。这封信出自一位我曾当作朋友、怀着热忱的同情去信的人,而信里有一种蓄意的嘲弄,格外猛烈地让我想起一件事——那是我最近几年终于越来越清楚地意识到的事。此前,我主要有机会注意到的是对我这个人本身的疏远,在数学“大世界”里,首先是在那些曾与我或远或近算朋友的人中间(45)。那里说的不再是人际层面上的疏远,而是一种共识,其性质如同一种时尚,并且像时尚那样把自己摆成理所当然的东西,流行于一些多少“在行”的人之间:那种论千页一包一包的数学,还有那些我拿它们在一二十年里磨破了人耳朵的概念(46,47),全面权衡起来并不很严肃;那里头有大量装腔作势的轰响,值得的东西却没多少;而除开围绕概形与平展上同调这两个概念的那几大篇“普遍胡话”——它们有时候倒确实自有其用处;唉,这一点人还是乐意承认的——至少剩下的部分,忘掉更厚道;至于那些不顾雅趣、不顾显而易见的严肃准则,仍旧摆出要吹这种格罗滕迪克式小号的人,就该跟他们那位师父(不管认不认)扔进同一条口袋;他们要是受到应得的对待,也只能怪自己……

\origpage{137} 无疑,自1976年以来(50)、尤其是最近两三年来传到我这里的许多这一路的声音(我刚把它们“照原样”抄了下来),终于唤醒了我身上那根好斗的筋——最近十年里它多少睡过去了。它们像反射一样挑起了一个念头:投入混战,让那些什么也没弄懂的毛头小子闭上嘴——总之是个十足愚蠢的反射,就像公牛:只要给它看一块红布,在它鼻子前晃一晃,它立刻就抖擞起来、冲上去,忘了自己原本正安安静静走着的路,那可是它自己的路!不过我仍觉得这反射相当表皮,单靠它不足以让我动摇。再说,幸好,做数学显然比朝一块红布冲过去、四面挨刀要有滋味得多。可是做数学,而且不顾一切地坚持一种工作风格、一种看事情的路数——那些是我的——这也有点“投入混战”的意思;这是在那些轻蔑与拒斥的迹象面前挺立自己——那些迹象落到我身上,毫无疑问是回应我的旧友们在我身上感到、或者以为感到的轻蔑,那轻蔑若不是针对他们,至少也是针对他们毫无保留地认同的那个圈子。所以这多多少少也是在跟着那块红布跑,而不是走自己的路。

% ---------------- p.229 ----------------
\origpage{138} 这个念头,最近几周里已多次在我这里出现过;而今天这场反思,也许主要就是朝考察这一面去的。走着走着,又显出另一面,那里自我的力量无疑也占了大份,但它不像是一种单纯的好斗反射。不如说是一种我身上的欲望——此刻我还看不清它的本性——要给我最近十年或十二年所做的数学工作一个意义,或者要看到它取得它全部的意义;而这意义(我内心确信)不可能归结为私下里的乐趣或个人的历险。可是,即便这欲望的本性仍未被理解,既然我还没闲下来细看它,这场反思已足以向我表明:真正压在我身上、可以说硬掰着我的手让我往数学投入的那股力量,恰恰就在这欲望里——“翻转”之力。红布有没有,它照样起作用。如果它是依恋一段过去的迹象,那也是最近这十年的过去,也就是“1970年以后”的过去,而不是那些已经白纸黑字写下的东西、已经做成的东西的过去,不是1970年以前的。

\origpage{139} 骨子里,对这些东西、对将来“后世”会怎样处置它们(何况是否真有后世还很可疑……),我心里毫无不安。我在这段过去里感兴趣的,绝不是我在其中做过的事(以及它已有或将会有的运道),而恰恰是没有做成的事——在那份当年摆在我眼前的大纲领里,只有极小一部分,靠我的努力,以及靠那些有时愿意加入我的友人和学生的努力,得以实现。这份纲领本身,在我对数学事物的眼光与路数更新的同时,也未经预谋、未经寻求地更新了。年复一年,重心移动了,题目如此,我要说的东西本身也如此:如今我最首要的意图,不再是完成那些精细基础的大工程,而是去探那些最让我着迷的奥秘,比如“Motif”的奥秘,或者 $\mathbb{Q}$ 在 $\mathbb{Q}$ 上的伽罗瓦群的“几何”描述的奥秘。当然,一路走来,我至少免不了在这里那里勾画些基础,就像我已经开始做的那样(除了别的)在《穿越伽罗瓦理论的长征》里,或者像我现在正在《追寻栈》里做的那样。然而要说的东西变了,表达它的风格也变了。

% ---------------- p.230 ----------------
换个说法:最近十年里,在数学事物的世界里,我瞥见了一些神秘的、极美的东西。这些东西不是我私人的,它们生来就是要被传达的——瞥见它们这件事本身的意义,我这样感到,就在于把它们传达出去,好让人接过去、理解、吸收……然而传达它们,哪怕是向自己传达,也就是深入它们,多少发展它们——那是一件工作。我当然很清楚:我绝不会把这工作做到头,即便再给我一百年去投入。但这不该是我今天操心的事——在我余下活着、余下发现世界的时间里,我要拿多少年、多少个月投到那件工作上;而另一件工作正等着我,那是我唯一能做的。我生命的季节,不在我的权柄之内,也不是我的角色去安排的。

% ================================================================
%  第 12 章 第一部分的注 (原书 p.231-241, 边码 p.141-p.152)
%  源文: _work/src1/task-ch12a.md (12.1–12.22, 注 1–22)
% ================================================================
% ---------------- p.231 ----------------
% 原书 p.231-233 的章内目录(Contents)按体例不排印
\chapter{第一部分的注}

% ---------------- p.232 ----------------
% (章内目录续, 不排印)

% ---------------- p.233 ----------------
\section{$\varnothing$}

\notehead{1}{}
\origpage{141}(1984年3月补写)说我的工作“风格”与“方法”没有变,恐怕失于武断——毕竟我在数学上的表达风格已经深刻改变了。这一年来用于《追寻栈》的大部分时间,是坐在打字机前敲下一些思考;它们就是要几乎原样发表的(只除了后来补加的若干相当简短的注,借互见、勘误之类方便阅读)。没有剪刀,没有浆糊,不必费劲去炮制一份“定稿”(那份定稿尤其不该让任何通向它的过程透出影来)——这毕竟也是“风格”与“方法”的变化!除非把数学工作本身与写作、与陈述结果的工作割裂开来,而那是人为的,因为那不合事情的实情:数学工作与写作是拆不开的。

% ---------------- p.234 ----------------
\section{$\varnothing$}

\notehead{2}{}
\origpage{142}(1984年3月补写)重读最后这两段,我有些不适之感,因为写它们的时候,我牵连的是别人,而不是我自己。显然,写的时候,“我自己也可能与此有关”这个念头并没有碰到我。我肯定什么也没学到:我就这样满足于把多年来在别人身上察知、并以许多方式看着被证实的东西,白纸黑字写下来(无疑还带着某种满意)。在这份反思的后续部分,我不得不回想起:我的一生中,对别人的轻蔑姿态并不缺少。我所把握到的他人之轻蔑与对己之轻蔑之间的关联,若在我这个人身上竟然缺席,那才奇怪;健全的理性(以及类似情形下我对自身视而不见的经验——我终究还是察觉到了)告诉我:肯定不该是这样!然而眼下这只是一条纯粹的推论,它唯一可能的用处,是促使我亲眼看一看究竟发生了什么,去看、去审视(如果它确实存在过、或还存在着)这份仍属假设的对己之轻蔑——它埋得如此之深,至今完全逃过了我的目光。真的,可看的东西并不缺!

\pmark 这一件忽然在我看来是最要紧的几件事之一,恰恰因为它藏得如此之深……\fnote{1}{}{(1984年8月)不过关于这一点,参见“屠杀”一注(n\textsuperscript{o}~87)末两段的反思。}

\section{$\varnothing$}

\notehead{3}{}
我这里想到的尤其是那几条“已故”的猜想——莫德尔(Mordell)的、泰特(Tate)的、沙法列维奇(Chafarévitch)的:去年,这三条猜想在法尔廷斯(Faltings)一篇四十页的手稿里同时得到了证明,而就在那个时候,“圈内人”那套早已稳固的共识还裁定这些猜想“够不着”呢!恰好,我所珍爱的“远阿贝尔代数几何”(géométrie algébrique anabélienne)纲领的那条拱顶石式的基本猜想,就与莫德尔猜想相近。(据说莫德尔猜想甚至是这条基本猜想的推论——这就满可以说明:这个纲领并不是给严肃人预备的故事……)

\section{$\varnothing$}

\notehead{4}{}
再说,即便在今天,也还能碰到地位不明的“证明”。格劳尔特(Grauert)那条以他命名的有限性定理的证明,多年来就是如此:谁(而热心人并不缺!)也读不下去。这种困惑,后来由另一些更透亮的证明解决了;其中有些走得更远,接替了最初那个证明。类似而更极端的情形,是所谓“四色问题”的“解答”:其计算部分是用计算机(外加几百万美元)打发掉的。于是这里所说的“证明”,其根基已不再是理解一种数学处境所生的内心确信,而是对一台没有理解能力的机器的信任——而使用它的数学家,对这台机器的结构与运作一无所知。即便假定这计算由别的计算机、按别的计算程序复核过,我也并不认为四色问题就此了结。它只是换了一副面孔:人们不再去寻一个反例,而只去寻一个证明(当然,得是读得懂的证明!)。

% ---------------- p.235 ----------------
\section{$\varnothing$}

\notehead{5}{}
\origpage{143}\origpage{144}这件事之所以格外值得注意,是因为直到1957年,布尔巴基(Bourbaki)小组里还不止一位成员对我有所保留——而我相信,他们终究把我收编进来时,多少是有些勉强的。一句善意的玩笑话,把我归入“危险的专门家”(指泛函分析方面的)。有时我在嘉当(Cartan)身上感到一种未说出口的、更认真的保留——有那么几年,我一定给了他这样的印象:一个倾向于无端而肤浅地推广的人。我见过他大为惊讶:在我为布尔巴基写的头一篇(也是唯一一篇)稍长的稿子(论流形上的微分体系)里,他发现了一点还算有点内容的思考——当初我提议由我来写时,他并不怎么热心。(这份思考几年后对我又有用:我从凝聚对偶的观点发展留数形式体系时,用到了它。)再说,在布尔巴基大会上,我多半是跟不上的,尤其是在共同宣读稿子的时候:那些宣读与讨论的节奏,我实在跟不下来。也许我真的不适合集体工作。不管怎样,我这种难以融入共同工作的困难,或者我因别的缘故在嘉当等人那里引起的保留,从未在任何时刻给我招来讥讽或冷遇,哪怕一丝居高临下的影子也没有——顶多是韦伊(Weil)那里有过一两次(此人确乎另当别论!)。嘉当对我的善意始终如一,不曾稍减,那里面既有热诚,也有他那特有的、在我看来与他这个人分不开的一点幽默。

\section{《生存与生活》的朋友们}

\notehead{6}{}
在这些朋友里,我大概还该算上皮埃尔·萨米埃尔(Pierre Samuel)——从前我主要是在布尔巴基里认识他的,就像谢瓦莱(Chevalley)一样;他(和谢瓦莱一样)在“生存与生活”(Survivre et Vivre)小组里起了重要作用。我看萨米埃尔并不怎么热衷于那种科学家的优越感的幻觉。我感到,他带来的东西很多,主要靠的是他在共同工作、讨论、与人交往中那份健全的常识和笑呵呵的好脾气,也靠他能在这样一个倾向于激进分析与抉择的群体里,优雅地扮演“讨厌的改良派”的角色。我退出“生存与生活”之后,他还在那里待了一阵,担任同名刊物的主编;当他感到自己在这个群体里的存在已不再有用时,他就体面地离开了(去加入“地球之友”)。

萨米埃尔和我同出那个狭小的圈子;这并不妨碍他属于那些沸腾年月里让我觉得学到了一点东西的朋友之列(尽管我是个多坏的学生……)。他那种为人处世的方式,以及谢瓦莱的方式——两者其实毫不相似——对我那些“精英至上”的倾向来说,是比最犀利的分析更好的解药!

如今在我看来,对那个时期所有让我学到了一点东西的朋友来说,我学到的东西,更多是靠他们各自的为人、靠他们与我不同的感受力(其中“某种东西”终于传达了过来),而不是靠解释、讨论之类。说到这一点,我尤其记得——除了谢瓦莱和萨米埃尔——德尼·盖吉(Denis Guedj)(他在“生存与生活”小组里很有威望)、达尼埃尔·西博尼(Daniel Sibony)(他始终与这个小组保持距离,一边从角落里用半是轻蔑、半是打趣的眼光看着它演变)、戈登·爱德华兹(Gordon Edwards)(他是1970年6月蒙特利尔那场“运动”诞生的共同推动者之一,多年来施展出惊人的精力,维持着《生存与生活》刊物的英文“美洲版”)、让·德洛尔(Jean Delord)(一位年纪与我相仿的物理学家,细致而热忱,他对我连同“生存派”这个小圈子都怀有偏爱)、弗雷德·斯内尔(Fred Snell)(另一位在美国布法罗的物理学家,1972年我在那里住了几个月,做客他的乡间住所)。

% ---------------- p.236 ----------------
\origpage{145}这些朋友当中,五位是数学家,两位是物理学家,而全都是科学家——这似乎表明:那些年里离我最近的环境,仍然是科学家的环境,尤其是数学家的环境。

\section{$\varnothing$}

\notehead{7}{}
上一段是整个导论里第一段在我原稿上被狠狠划掉、并添了许多涂改的。对那件事的描述、对词句的选择,起初都是逆着毛捋、顶着潮流来的——显然有一股力量在推着,要把那件事快快带过,仿佛是了却一桩心事,好“去谈正经事”。这些都是一道抵抗的熟面孔;在这里,它抵抗的是把这段插曲弄清楚,以及弄清它作为内心态度之显影的分量。情形与这篇导论开头(第2段)所描述的完全相似:在数学工作中发现一个矛盾及其意义的那个“关键”时刻——那时充当“抵抗”的,是心智的惰性,是它不肯离开一个错误或不充分的视野(而在那里面,我们这个人丝毫未被牵连)。这种抵抗是主动的,必要时还会发明创造,好把一条鱼在没有水的地方也淹死;而我方才说的惰性,只是一种被动的力。在眼前这桩事上,比在数学工作中尤甚:刚刚以其全部单纯、全部显明浮现出来的那份发现,随即带来的是一种卸下重负的轻松感,一种解放感。这不只是一种感觉——不如说,是对刚刚发生之事的敏锐而心怀感激的知觉,而那件事,就是一次解放。

\section{$\varnothing$}

\notehead{8}{}
后文会看清:这份模棱两可,丝毫没有“在1970年觉醒的次日消散”。这里有一招“自我”典型的战略性撤退:把“觉醒之前”那段时间丢给盈亏自负,而“觉醒”立刻成了一道分界线,划出一个无可指摘的“之后”!

\section{$\varnothing$}

\notehead{9}{}
这不完全确切:在我最亲近的同行里至少有一个例外,后面会看出来。这里有一种记忆特有的“懒惰”:它往往倾向于把那些与一套熟悉且根深蒂固的看法“对不上”的事实“蒙混过去”。

\section{$\varnothing$}

\notehead{10}{}
比如,我写给同事或旧日的学生、那些我当作朋友的人的信——数学问题也好,实际事务或私事也好——多得我已数不清,却从未收到过回音。看来这并不只是专门留给我一个人的“优待”,而是一种风气变化的征兆——我听到的回声也是这个方向。(诚然,那些回声涉及的情形是:寄数学信的人,收信的那位知名数学家并不认识……)

% ---------------- p.237 ----------------
\section{阿尔多·安德烈奥蒂、约内尔·布库尔}

\notehead{11}{}
\origpage{146}\origpage{147}当然。我这边也可能有遗忘——何况我那时格外“硬派”(polar)的性情,恐怕也不会让人乐意跟我谈这类事情,也不会让我记住一场很可能发生过的、这个方向的谈话。可以肯定的是:关于恐惧的问题被提起(甚至不用这个名字来称呼它……),至少该是极其罕见的;今天大概同样罕见,尤其是在“上流社会”里。

在那个圈子里我的众多朋友中,除了谢瓦莱(他至少在六十年代该已察觉到这种恐惧的气氛),我觉得真正清楚地感受到了它的,只有阿尔多·安德烈奥蒂(Aldo Andreotti)。我认识他,认识他妻子芭芭拉(Barbara)和他们那对双胞胎孩子(那时还很小),是在1955年(我想是在芝加哥韦伊家的一次晚会上)。我们一直十分亲近,直到1970年那个“大转折”:我离开了从前属于我们的那个圈子,与他们也就有些失去了联系。阿尔多感受力极敏锐,而这份敏锐丝毫没有因为与数学打交道、与我这样的“硬派”打交道而变钝。他身上有一种对接近他的人的、自然而然的同感之才。这使他与我在这数学圈里、乃至圈外认识的所有别的朋友都不同。在他那里,友情总是压过共同的数学利益(那些利益并不缺),他是极少数与我多少谈过我的一生、他也谈过他一生的数学家之一。他父亲和我父亲一样是犹太人,他因此在墨索里尼的意大利吃过苦头,正如我在希特勒的德国。我见他总是愿意鼓励和支持年轻研究者,而那时的气候,已经越来越难让这些人被建制接纳。他那种自发的关切,总先落在人身上,而不是落在什么数学“潜力”或名声上。他是我有幸遇见的最令人眷爱的人之一。

这番对阿尔多的追忆,又勾起了约内尔·布库尔(Ionel Bucur)的记忆——他也是骤然辞世,未及天年;和阿尔多一样,人们(我想)更怀念的是他这个人、这位乐于重逢的朋友,而不是作为数学讨论伙伴的他。在他身上,人感到一种善良,旁边还有一份少见的谦逊,一种不断让自己退到后面的倾向。这样一个人——如此不把自己当回事、如此不想让谁敬畏——最后竟当上了布加勒斯特大学理学院院长,真是个谜;多半是因为他压根没想到要推辞那些他远不觊觎的职务,而他的同事或政治当局偏要把它们放到他肩上——顺带说,那肩膀是结实的。他是农民的儿子(在一个“阶级标准”要紧的国家里,这一定起了作用),也带着农民那份健全的常识与朴素。他肯定察觉得到围绕着名望之人的那份恐惧,但肯定也觉得这事理所当然,是权力位置的天然属性。不过我并不认为他本人曾让谁感到恐惧——无论是他妻子弗洛丽卡(Florica)、他们的女儿亚历山德拉(Alexandra),还是他的同事或学生;我能听到的回声也确实如此。

\section{$\varnothing$}

\notehead{12}{}
这里的“次日”(lendemain)一词要按字面理解,不是隐喻。

% ---------------- p.238 ----------------
\section{$\varnothing$}

\notehead{13}{}
\origpage{148}很清楚,前面的描述没有别的野心,只是想用具体的词语尽力把记忆的那片“迷雾”所交给我的东西还原出来;这片迷雾从未凝结成任何一个稍为具体的个例,好让我在这里给出某种稍微“写实”或“客观”的描述。若要从这段话里读出:那些不愿坐在前排的同行,或者那些没有明星或泰斗身份的同行,在跟后一类人说话时必定被焦虑缠住——那就是曲解我的意思了。在我于这个圈子里认识的多数朋友身上,显然不是这样,即便是那些偶尔出入讨论会或讨论班的人。毫无保留地成立的是:“泰斗”身份造成一道屏障、一道沟,把它与没有同等身份的人隔开;而这道沟极少消失,哪怕只有一场讨论的工夫。我还要补充:“前排”与“沼泽”之间那条主观的界线(我却觉得它十分真实),绝不能归结为社会学标准(社会地位、职位、头衔等等),甚至也不能归结为“身份”或名声的标准;它还反映着气质上、或者某些更难把握的禀性上的心理特点。我二十岁初到巴黎时,我知道自己是个数学家,我做过数学;尽管有人地生疏之感(我有机会谈到过),我在心底里觉得自己“是他们中的一个”——而只有我知道这一点,甚至一开始我都不确定自己会不会继续做数学。今天,我倒更愿意坐到最后一排去(在难得碰上这个问题的场合)。

\section{$\varnothing$}

\notehead{14}{}
人们可能以为这与“没有首领”的说法相矛盾,其实不然。对布尔巴基的老人们来说,我觉得韦伊被看作小组的灵魂,却从来不是“首领”。他在场、并且他乐意的时候,就成了我所说的“串场人”,但他并不发号施令。他心情不好的时候,会把自己厌恶的某个题目卡住,让别人到下一次大会、趁韦伊不在时再安安静静地重拾这个题目,甚至第二天等他不捣乱了再谈。决议由在场的成员一致通过,因为他们认为:一个人逆着所有其他人的一致意见而站在对的一边,这绝非不可能(甚至并不少见)。这条原则对集体工作来说也许显得荒谬。了不起的是:它居然行得通!

\section{$\varnothing$}

\notehead{15}{}
我并不觉得这种对布尔巴基风格的“过敏”,曾在这些数学家与我、或与布尔巴基其他成员和同情者之间造成交流上的困难——倘若小组的精神真是小圈子精神、是精英中的精英的精神,那倒会如此。在风格与时尚之上,小组的所有成员都有一种对数学实质的敏锐感觉,不论那实质来自何处。只是到了六十年代,我才记得我某位朋友把某些他不感兴趣的数学家称作“讨厌鬼”。这类事我另外几乎一无所知,于是往往把这类评价照单全收,被那份满不在乎的自信震住了——直到有一天我发现,某位“讨厌鬼”是个独到而深刻的心灵,只是没有讨得我那位出色的朋友欢心罢了。我觉得在布尔巴基的某些成员身上,面对别人的工作、在自己不懂或懂得不透的时候所持的谦逊态度(或至少是保留态度),是先被磨掉的;而那种“数学本能”却还留着——那种不必参照名声与声望、就能嗅出丰厚实质或扎实工作的本能。

% ---------------- p.239 ----------------
\origpage{149}从各处传来的回声看,我觉得这两样——谦逊与本能——在从前的那个数学圈子里,今天都已成了稀罕之物。

\section{$\varnothing$}

\notehead{16}{}
说实话,布尔巴基的成员里有好几位肯定有“自己的”小宇宙,大小不一,在布尔巴基小宇宙之外或之上。但就我自己的情形而言,这样一个“自己的宇宙”之所以要等到我不再是布尔巴基的一员、全部精力都投到属于我个人的工作上之后,才在我周围形成——这也许并非偶然。

\section{$\varnothing$}

\notehead{17}{}
主要是在科学圈之外,我才遇到对我所投入的那场行动的热烈回响,以及切实的帮助。除了阿兰·拉斯库(Alain Lascoux)与罗歇·戈德芒(Roger Godement)的友情支持,这里我还尤其要提到让·迪厄多内(Jean Dieudonné)的支持:他特地跑到蒙彼利埃的轻罪法庭去出庭,为一件注定失败的事业,把他那份热忱的证词加在别的证词之上。

\section{$\varnothing$}

\notehead{18}{}
我想,这份缺乏辨别力,在这两个场合并非出于我的疏忽,而是出于不成熟、出于无知。直到十来年后,我才开始注意那些阻滞机制——在我自己身上,也在我的亲人或学生身上——并掂量它们在每个人生活里所起的巨大作用,而且不只是在中学或大学里。当然,我后悔这两个场合没有更成熟一些的辨别力;但我并不后悔把自己的印象清楚地说出来,不论那些印象有没有根据。当我在某个场合看到一件做得不认真的工作,把这事照它的本来面目说出来,我觉得是必要而有益的。若在另一个场合,我得出的结论是仓促的、没有根据的,那么负有责任的也并非只有我一个。被这样摇醒的学生,当时还有得选:要么从中吸取教训(第一次也许正是这样),要么就此泄气,也许还改了行(那也未必是坏事!)。

\section{耶稣与十二使徒}

\notehead{19}{}
从1970年至今,还有一个学生伊夫·拉德加耶里(Yves Ladegaillerie)在我这里准备并通过了论文。第一时期的学生是:P.贝尔特洛(P.~Berthelot)、M.德马祖尔(M.~Demazure)、J.吉罗(J.~Giraud)、M.阿基姆夫人(Mme~M.~Hakim)、黄春生夫人(Mme~Hoang Xuan Sinh)、L.伊利西(L.~Illusie)、P.茹阿诺卢(P.~Jouanolou)、M.雷诺(M.~Raynaud)、M.雷诺夫人(Mme~M.~Raynaud)、N.萨阿韦德拉(N.~Saavedra)、J.-L.维尔迪耶(J.-L.~Verdier)。(其中六位是在1970年之后才完成论文工作的,也就是说,在我数学上的可用时间极为有限的时期。)在这些学生当中,米歇尔·雷诺(Michel Raynaud)占一个特殊位置:他论文工作所处理的那些本质性问题与概念,是他自己找到的,而且他是完全独立地发展它们的;我所谓“导师”的角色,因此仅限于读已完成的论文、组成评审委员会,并成为其中一员。

轮到我出题目的时候,我十分小心,只限于那些我有足够强的关系、因而觉得自己在需要时能撑住学生工作的题目。一个显著的例外,是米歇尔·雷诺夫人(Mme Michèle Raynaud)关于基本群局部与整体莱夫谢茨定理的工作,那些定理是在适当的平展\emph{site}上以 1-栈(1-champs)的措辞表述的。这个问题在我看来是困难的(后来也确实证明如此),而对我所提出的那些猜想,我并没有证明的想法(再说那些猜想也几乎不可能有什么疑问)。这项工作一直做到七十年代初;雷诺夫人(如同此前她丈夫的情形)发展出一套精细而独到的方法,没有得到我或任何别人的帮助。这项出色的工作还引出一个问题:把雷诺夫人的结果推广到 $n$-栈($n$-champs)的情形——在我看来,在概形的语境里,这应当就是“弱莱夫谢茨定理”一类定理的自然归宿。不过,这里相关猜想的表述(它同样几乎不可能有什么疑问)本质上用到了 $n$-栈的概念;而对它的追寻,按说正是本书的主要对象\fnote{2}{}{(其实指的是《数学反思》第3卷,而不是眼前这第1卷《收获与播种》——见导论,第(v)页。)},正如书名《追寻栈》所示。我们大概会在相应的地方再回到这一点。

% ---------------- p.240 ----------------
\origpage{151}\origpage{150}另一个相当特殊的情形是黄春生夫人。我最初是在1967年12月在河内遇见她的,那时我在疏散后的河内大学讲了一个月的课程—讨论班。第二年,我向她提出了她的论文题目。她是在战争时期格外艰难的条件下工作的,与我的接触仅限于断续的通信。1974/75年她得以来到法国(趁温哥华国际数学家大会之便),于是在巴黎通过了论文(评审委员会由嘉当主持,另有施瓦茨(Schwartz)、德尼(Deny)、齐斯曼(Zisman)和我)。

最后我还该提到皮埃尔·德利涅(Pierre Deligne)和卡洛斯·孔图-卡雷尔(Carlos Contou-Carrère):这两人都多少算是学生,前者在1965—68年前后,后者在1974—76年前后。两人显然都有(并且一直有)非同寻常的才能,而他们运用这才能的方式极为不同,际遇也极为不同。到比雷(Bures)来之前,德利涅在比利时多少算是蒂茨(Tits)的学生——不过我怀疑,按通常意义说,他在数学上做过谁的学生。孔图-卡雷尔做过桑塔洛(Santalo)的学生(在阿根廷),有一阵还做过托姆(Thom)的学生(多少算吧!)。接触建立的时候,两人都已具备一个数学家的身量,只差孔图-卡雷尔缺少方法与行当上的训练。

我在数学上对德利涅的作用,仅限于零零碎碎地把我所懂的那点代数几何讲给他听;他学起来就像听故事——仿佛他向来就知道。顺带我也提出一些问题,而他多半当场或在随后几天里就找到了答案。这就是我所知道的德利涅最早的工作。1970年以后的工作(他的,以及我那些“正式学生”的),我只从极为零散遥远的回声里略知一二\fnote{3}{}{(我尤其有机会翻阅过贝尔特洛和德利涅寄给我的几份抽印本,承他们的好意。)}。

至于我对孔图-卡雷尔的作用,按他自己在论文开头所说,仅限于把他引进概形的语言。无论如何,他这几年作为国家博士论文所准备的工作,我只是远远地跟着;那是一个再时新不过的题目,超出了我的能力。

% ---------------- p.241 ----------------
\origpage{152}孔图-卡雷尔在大世界里经历了几番波折之后,最近终于在最后关头——而且(现在在我看来)是违着他的心意——不得不求我帮忙,充当论文导师并组成评审委员会。(这有让他显得像“1970年以后”的格罗滕迪克学生的风险,而在一种处境里,这可能会带来严重的不便……)我尽力履行了这项任务;这大概是我最后一次行使这个职能(在国家博士论文这一级上)。在这个有点特殊的场合,我格外高兴有让·吉罗(Jean Giraud)的友情协助:他也拨出一两个月的时间,仔细读了那部厚厚的稿子,并写了一份详尽而热忱的报告。

\section{$\varnothing$}

\notehead{20}{}
这让我想起莫妮克·阿基姆(Monique Hakim)所选的题目,说实话那也不见得更有吸引力;我真不知道她是怎么保持住精神的!如果说她有时也吃力,那至少没有到让她郁郁寡欢、闷闷不乐的地步;我们之间的工作是在亲切而放松的气氛里进行的。

\section{$\varnothing$}

\notehead{21}{}
也许更确切的说法是:对我这样的气质来说,要完全担起一个教师角色,我缺的仍是那份必要的成熟。我后天形成的气质,长期带着“男性”(或“阳”)特征过度占优的印记;而成熟的一个方面,恰恰是一种以“女性”(或“阴”)为主的“阴阳”平衡。

(后来补写)我看出,比成熟更缺的,是一种慷慨——在我至今的教学生涯里,我缺的正是某种慷慨;这种慷慨的表现方式,比在时间与精力上随时可用更为细微,也更为本质。在我第一个教学时期,这份欠缺并没有明显地表现出来(比如说,通过一连串失败的情形积累起来),多半尤其是因为它被选择来跟我工作的学生身上强烈的动机补偿了。相反,在第二个时期,从1970年至今,我觉得这份欠缺至少是原因之一,而且无论如何是最直接牵涉到我自己的那一条原因——我在研究层面(也就是从DEA这一级起)的教学中所看到的整体失败,原因就在这里。关于这一点,见《纲领概要》(Esquisse d'un programme)第8段,以及第9段“一项教学活动的总结”;在那里面透出的是这七八年来这项活动留给我的挫败感\fnote{4}{}{(另可参见后来补写的注(23iv)。)}。

\section{$\varnothing$}

\notehead{22}{}
也许不会再有多久了,因为我已经决定申请进入国家科学研究中心(CNRS),从而结束大学里的教学活动——这几年来,这活动变得越来越成问题。

% ---------------- p.242 ----------------
\origpage{153}\origpage{154}
\section{$\varnothing$}

\notehead{22$'$}{}
即便在1970年以后,当我对数学的兴趣在我生活里变得时有时无、退到了边缘,我也不认为有过哪一回,学生来叫我同他一起工作、而我推辞掉的。我甚至可以说:除了两三个例外,我1970年以后的学生对自己所做工作的兴趣,远远低于我本人对他们题目的兴趣——即便在我除了去系里露面的日子几乎不碰数学的那些时期,也是如此。因此,我1970年以前对学生的那份随时可供支配,以及作为其首要标志的、工作上的极端严格,对我后来多数学生来说恐怕毫无意义:他们做数学并没有什么信念,仿佛那是一种他们不得不对自己持续施加的努力……

\section{孩子与老师}

\notehead{23}{}
这里“传递”一词,与事情的实情并不真正相符——实情把我拉回到一种更谦逊的态度上来。这种严谨并不是一件可以传递的东西,最多只能被唤醒、被鼓励,而它从小就被忽视、被打消,被家庭环境,也被学校和大学。就我所能回想起来的而言,这种严谨一直存在于我的种种追寻之中——至少是那些智力性质的追寻——而我不认为它是由我父母传给我的,更不认为它传自哪位师长,无论是在学校里,还是在数学界的前辈当中。我觉得它似乎属于“天真”的诸属性之一,因而也属于每个人生来便被分得的东西。这种天真很早就“吃够了苦头”,于是不得不或深或浅地潜下去,以致在往后的一生里往往再也见不到它多少痕迹了。在我身上,由于一些我还不曾想去深究的原因,某种天真在智力好奇心这个相对无害的层面上存活了下来;而在其余一切地方,它都深深地潜了下去,神不知、鬼不觉!就像在每个人身上那样。也许,完全意义上的“教学”的秘密——或者不如说它的奥秘——就在于重新与这种表面看来已经消失的天真接上联系。但是,如果这联系尚未首先存在于、或尚未被重新找回到教师本人身上,那就根本谈不上在学生身上把它重新找到。而这时教师“传给”学生的,绝不是这种严谨或这种天真(两者在师生身上都是天生的),而是一种尊重,一种对这件通常遭人弃置的东西的、无声的重新肯定。

\section{$\varnothing$}

\notehead{23$'$}{}
不过,七八年来,在我的数学生涯里还有另一个慢性的“挫败之源”,只是这些年来它表现得隐晦得多。靠着反复出现,靠着同一种“令人挫败”的情形在我的教学活动中顽固地累积,它终究变得显眼起来,末了还爆发为一声“受够了!”,使我实际上终止了一切所谓的“指导研究工作”。在我的这番反思途中,我有一两次顺带触及这个问题,最后总算在结尾处多少把它审视了一番。我在那里至少描述了这份挫败感,并考察了它在我“重返数学”一事中所起的作用(参见第50段“一段过去的重量”)。

% ================================================================
%  第 12 章 第一部分的注 (原书 p.243-254, 边码 p.155-p.171)
%  源文: _work/src1/task-ch12b.md (12.26–12.50, 注 23″–44)
% ================================================================
% ---------------- p.243 ----------------
\section{怕玩}

\notehead{23$''$}{}
这个学生曾和我一起做了一整年的DEA“实习工作”,而在他与我的工作关系里,直到最后都保持着“紧绷”状态。那是一种坦率友好的关系,贯穿着一种无可怀疑的相互好感。然而可就是有这份“怯场”;这份恐惧——它真正的起因肯定不是对我这个人的畏惧,尽管看上去像是如此。要不是这个学生自己跟我谈起,我也许根本不会察觉这件事;他谈,无疑是为了或多或少“解释”一下他在这一年里工作几乎完全卡住的原因。

正如在其他学生身上发生过的那样——他们和他一样,起初对某种几何实质抓得很牢——卡住恰恰出现在要“做零件活儿”的时候:也就是把成形的陈述白纸黑字写下来,或者仅仅是把握我提供的那些陈述的意义与含义;我提议把它们当作一种语言的基石、当作“游戏规则”接受下来。“学校式的”反应几乎总是把面对“做研究”局面的学生推向:把导师传下来的那些心照不宣的“游戏规则”当作既定的、又必须服从的东西接受下来,而这些规则恰恰不是要试着说明白的,更不是要弄懂的。这些心照不宣的规则所采取的具体形式,就是语义学或计算上的“配方”,不妨说,像预科班的题解书那样(或者任何一本普通教科书)。此外,学生还期待导师给出“证明……”这种形式的任务——那是他以往经验里遇见的唯一一种数学“思考”形式。(而且我不认为大多数职业数学家、以及其他科学家的禀性有本质不同——只不过“导师”换成了“共识”:共识规定当下的游戏规则,并把它当作不可更动的既定之物。这共识还规定哪些“问题”是要去解决的;在这些问题之间,每个人都觉得自己有随意挑选的自由,甚至允许自己在工作过程中修改它们,乃至发明另一些……)。\origpage{155}我注意到,我面对有待探究的数学实质时那种完全不同的态度——因而也是面对学生时那种完全不同的态度——几乎必定会引发一种失措,而焦虑就是它的征兆之一。像一切焦虑一样,它倾向于给自己找一张面孔,投射到某个外部的“理由”上,不管那理由说得通说不通。焦虑最常见的一张面孔,恰恰就是恐惧。

这类困难在我教学生涯的第一个时期几乎没有出现,也许只有两例例外:那两次“师生”关系没超过几个星期;也许还有(我说不准)“忧伤的学生”那一例——他也许感到自己被“铆”在一个完全不启发他的题目上,而他本来完全有更换的自由。在那个(我也谈到过的)长期被某种怯场困扰的学生身上,原因显然在别处。他的工作丝毫没有卡住,恰恰相反,他对选定的主题极为自在,并在其上做出了一项有规模的基础性工作。那个时期我的学生大多是巴黎高师的旧生,他们与亨利·嘉当(Henri Cartan)的接触,已经让他们见识了另一种“接近数学”的方式。与之(可以说)相反的另一极,是我在蒙彼利埃大学任教的第二个时期:正是在一年级学生那里,我所说的焦虑最少干扰思考性的工作。在这些学生中的许多人身上,面对一种不同进路时的惊讶既不引起焦虑,也不引起封闭,相反带来的是敞开和兴致——终于可以做点有趣的事了!据我的观察,几年大学对一个学生创造性禀性的作用是彻底而毁灭性的。奇怪的是,在这方面,长年中学的作用似乎相对无害。原因也许是:大学那几年正落在这样一个年纪——我们身上与生俱来的创造力,最迟必须通过个人的工作表达出来,否则就要永远沉船,至少就智力上的创造性工作而言是如此。肯定是出于一种健康的本能,我在做学生的那些年(也是在蒙彼利埃大学),几乎不踏进课堂,把几乎全部精力都投给了个人的数学思考\origpage{156}。

% ---------------- p.244 ----------------
\section{两兄弟}

\notehead{23$'''$}{}
这个学生身上的对抗,一上来就采取了“阶级对抗”的形式:我是那个对他的数学前途拥有“生杀之权”的“老板”,可以凭自己的好恶决定……当然,事情本身只能证实这种看法,因为我很快就终止了对这个学生(已变得令人难受的)责任。这把他置于一个微妙的处境——在眼下的世道里,找一个“老板”可不是那么显而易见的事,尤其是题目已经选定的时候。在另一个学生身上,合法期待落了空,对抗采取了类似的形式。我被感受为专横的“学阀”,不能容忍那些被他视为下属的人(学生或级别较低的同事)有任何反驳。

这样一种“阶级态度”,在我与第一个时期学生的关系里,从未有过一丝一毫的表现。显而易见的原因是:在1970年以前的形势下,学生通过论文之后毫无疑问会得到一个副教授(maître de conférences)的职位,因而享有与我相同的社会地位,即“大学教授”。数字最能说明问题:1970年以前开始跟我做工作的十一个学生,一完成工作就得到了副教授职位;而在我的指导下或多或少工作过的那二十来个学生里,没有一个获得这样的职位。诚然,他们之中只有两个人有足够的动力去做国家博士论文(而且两人的论文都很出色)。

因此,在这第二个时期,某些暧昧(其深层根源始终被遮蔽着)采取了阶级对抗的形式,采取了对“老板”的不信任(这种不信任被呈现、也被感受为“本能的”“发自肺腑的”)——这就不足为奇了。在那些或多或少当过学生样子的人里,有一个与我的友好关系持续了十来年,没有出现过表面上的对抗插曲,却仍带着同一种暧昧,表现为一种不信任的态度,被“储备”在明显的善意之后。说实话,我从未被这种例行公事式的“不信任”骗过;它在我看来,尤其是这位朋友认为应当给自己找的一个理由,好让自己不去冒险越出他选定的、界限分明的领地——无论在职业生活里,还是在整个生活里都是如此。可这件事他本来可以自由地去做,没有人(顶多只有他自己!)会要他交代……

这三个例子,也是我整个教书经验中仅有的、学生(或多少像学生的人)与我之间的关系中某种暧昧以“阶级态度”表现出来的情形。当这种态度出现在大学“团体”内部的同事之间时,显得格外暧昧:在这样一个团体里,双方都享有与常人处境相比极为悬殊的特权,而这些特权使级别(以及薪金)的差别显得相对微不足道。我还注意到,一旦当事人自己也被提升到前一天他还在指责别人的那种处境,这些态度就像中了魔法一样消失了(而且理所当然!)。

此外,在我于数学界内部(常常也在外部)目睹的多数——如果不是全部——冲突局面里,我都辨认出一种类似的暧昧\origpage{157}。那些“有位置的人”,无论其级别是否合乎他们的期待(合理与否),都享有相当闻所未闻的特权,那是任何别的职业或生涯都无法提供的。那些没有位置的人,则向往同样的安全感和同样的特权(这并不必然妨碍他们对数学本身感兴趣,有时还做出漂亮的工作)。在眼下的世道里,抢位置竞争激烈,没有位置的人常被当作流浪汉对待:我不止一次感到那种默契——在乐于羞辱人的人与被羞辱的人之间;后者咽下去,缩起来。他苦涩与敌意的真正对象,并不是那个行使了权力的人,而恰恰是他自己:他自己缩了起来,并且把对方奉为那种他乐意行使的权力的持有者。乐于羞辱人的人,也是那个报复并补偿(却永远抹不掉……)一场长期承受、早已深埋并被遗忘的羞辱的人。而默许自己被羞辱的人,是他的兄弟和效仿者:他暗中羡慕对方,并在苦涩里埋下那场羞辱,以及羞辱带给他的、关于他自己的那条卑微的消息。

% ---------------- p.245 ----------------
\section{一次教学的失败(1)}

\notehead{23$^{iv}$}{}
\origpage{158}这些字句写下之后,我有机会与1970年之后的两位旧学生谈过,想同他们一起探究我在蒙彼利埃大学研究层次上教学失败的原因。他们对我说:我那种倾向——低估同化那些对我熟悉、对他们却不熟悉的技术对他们可能意味着的困难——在他们身上产生了令人泄气的效果,因为他们感到自己始终达不到我对他们的期望。此外(这件事在我看来意义还要大),他们有时感到受了挫伤:当他们已经十分接近的时候,我不是把发现的乐趣留给他们、让他们靠自己的办法去发现,而是把袖子里藏着的一个成形陈述“卖给他们”、向他们泄底。那之后,他们剩下的就只有做“练习”了(而那并不特别让他们着迷)——也就是证明那个陈述。我先前在一则注(注 21)里指出过的自己身上那份“缺乏慷慨”,正落在这一处,当时没有就此多说。主要正是这类扫兴之事,构成了我在他们两人身上——在那么出色的开端之后——研究兴趣消失这件事上的个人贡献。

我意识到,我在1970年以前并不比之后更慷慨。如果说那时我没有遇到同样的困难,那无疑是因为:那个时期来找我的那类学生,有足够的动力,甚至能从一场“长长的练习”里品出魅力——那是学习本行、并一路学到一大堆东西的机会;同样,对于一个我“泄了底”的起步陈述,他们也能靠自己的办法从中抽出一大群远远超出第一个的陈述。当我换了教学地点,我在向新学生提议的思考主题上做了必要的调整,办法是选择那些可以被直接直觉把握的数学对象,不依赖任何技术行囊。但这必要的调整本身并不够,因为禀性上的差别(我的新学生与从前的学生相比)比单纯的行囊差别还要重要。这还印证了前面(第25段开头)说过的一点:我在“导师”这个角色上有某种不足;这不足在我教学生涯的第二个时期,比在第一个时期显露得强烈得多。

\section{$\varnothing$}

% ---------------- p.246 ----------------
\notehead{23$^{v}$}{}
\origpage{159}这种差别有一个格外醒目的征兆,出现在“外国人事件”之际——我曾有机会谈到过(第24节)。当时,许多与我完全陌生的人向我表达了同情,而我记不起1970年之前的哪个学生想到过以这种方式表示什么,更不用说在我投身的那场行动中向我提供任何帮助。相反,在我看来,第二个时期的学生或旧学生里,没有一个不曾向我表达同情与声援,还有几个积极参加了我在当地开展的运动。在这个狭小圈子之外,1945年法令那件事也在许多仅仅知道我名字的大学生中间激起了一定的波澜;在我被传唤的那一天,有不少人来到法院,以示声援。这后一个情形还提示:我所察觉到的“之前”与“1970年之后”学生态度上的差别,也许较少表达他们与我之间关系的差别,而更多表达的是一种心态的差别。显然,我“之前”的那些学生已经成了重要人物,而要让重要人物动心,是需要很多条件的……但我1970年离开IHÉS、投身战斗行动那一幕,似乎表明并不只是这一点。那是个他们中间还没有谁太像重要人物的时刻,然而我记不起他们中有谁对我所投身的活动表现出丝毫兴趣。我倒认为,那一定让他们所有人,无一例外地感到不自在。这仍然合乎“心态差别”的方向,但这不能只归因于社会地位的差别。

\section{$\varnothing$}

\notehead{24}{}
我要谈的这套伦理,同样适用于其他一切围绕研究活动形成的圈子:在那里面,能不能让人知道自己的结果、并从中得到承认,对每个成员的社会地位都是“生死”问题,甚至是他作为该圈子成员的“存亡”问题,而这对本人及其家庭意味着什么,也就可想而知。

\section{职业道德的共识——以及信息的控制}

\notehead{25}{}
除了与迪厄多内的那次谈话,在我作为数学家的生涯里,我想不起还有哪一次我参与或旁听的谈话\origpage{160},谈到了这一行的伦理、谈到同行之间关系里的“游戏规则”。(这里不算七十年代初围绕“生存与生活”运动发生的那些关于科学家与军事机构合作的讨论。那些讨论并不真正涉及数学家彼此之间的关系。再说,我在“生存与生活”里的许多朋友,包括谢瓦莱和盖吉,都觉得我当时——尤其在初期——如此强调这个我格外敏感的问题,使我偏离了更紧要的日常现实,恰恰就是眼下这份反思所考察的那类现实。)在我与学生之间,从未谈到过这些事。我相信,那份默契只限于一条规则:不把别人的想法——只要自己有机会知道——说成自己的。我觉得这条共识自古就有,直到今天在任何科学圈子里都没有被质疑过。可是,若缺少另一条补充规则——保证每个研究者都有可能让人知道自己的想法与结果——第一条规则就只是一纸空文。在今天科学界,处在声望与权力位置上的人,把持着对科学信息的任意支配权。
% ---------------- p.247 ----------------
在我从前熟悉的那个圈子里,这种把持已不再被迪厄多内所说的那种共识所节制——而那种共识,也许从未存在于他充当代言人的那个小圈子之外。处在权力位置上的科学家,实际上能收到他认为值得收到的一切信息(往往还超出);而对于其中很大一部分,他有权阻止其发表,同时自己留着这份信息的好处,把它打发为“没意思”“多少已经知道了”“平凡”之类……这种局面我在注(27)里还要再谈。

\section{$\varnothing$}

\notehead{26}{}
布尔巴基的“创始成员”是亨利·嘉当、克劳德·谢瓦莱、让·德尔萨特(Jean Delsarte)、让·迪厄多内、安德烈·韦伊。他们全都健在,只有德尔萨特例外——五十年代,他未及天年便被带走了;那个年头,这一行的伦理还普遍受到尊重。

重读本文时,我曾想删掉这一段:在这里,我可能给人一种印象,仿佛在颁发“正直”(或不正直)证书,而当事人对此毫不在乎,这事也不该由我来做。这一段可能引起的保留,肯定是有道理的。\origpage{161}可我仍旧留着它,为的是证词的真实,也因为这一段确实如实复现了我的感受,哪怕这些感受并不合宜。

\section{“年轻人的势利”,或纯粹性的捍卫者}

\notehead{27}{}
罗尼·布朗(Ronnie Brown)向我转述了J.H.C.怀特黑德的一句话(他是怀特黑德的学生),说的是“年轻人的势利:他们以为一条定理平凡,只因为它的证明平凡”。我从前的许多朋友,不妨把这句话好好想一想。这种“势利”如今绝不限于年轻人;我认识不止一位有声望的数学家,常这么做。我对此格外敏感,因为我在数学上(别处也一样……)做得最好的那些东西——我引入的、在我看来最丰饶的概念与结构,以及我靠耐心而执拗的工作从中拔出的那些本质性质——全都落进了“平凡”这个评语里。(如今这些东西,若作者还不是个名人,哪一件都很难被《科学院通报》(CR)当作一篇短讯接受!)我这一辈子作为数学家的抱负——不如说是我的激情与快乐——始终是找出那些显然的东西;在这本书里(包括眼前这一章引论性的文字里……),这也同样是我唯一的抱负。决定性的东西,常常首先就是看见那个尚未被看见的问题(不管答案是什么,也不管答案是否已经找到),或者提出一条陈述(哪怕是猜想性的),把一种尚未被看见、或尚未被理解的局面概括并容纳进来;这条陈述若被证明了,那么证明平凡与否无关紧要——那完全是附带的事——甚至一条仓促的、暂用的证明后来被证明是错的,也无关紧要。怀特黑德所说的势利,是那种厌倦了的享乐者的势利:他要先确认一瓶酒贵得惊人,才肯去欣赏它。这些年里,不止一次,旧日的激情一阵阵回来抓住我,我把自己最好的东西捧出去,却看它被那份自得打了回来。我感到一种至今仍活的痛;一份快乐落了空——可我并没有因此流落街头,而且幸而我并不需要替自己写的一篇文章找地方发表。

% ---------------- p.248 ----------------
\origpage{162}怀特黑德所说的势利,当它由一个有权者施于一个任其摆布的研究者身上时,就不只是一种对事物之美的麻木或封闭,而是一种滥用权力、一种不诚实:那位研究者完全在他的掌握之中,他尽可随意吸收、使用其想法,同时以“显然”“平凡”因而“没意思”为借口,堵住它们的发表。我这里甚至没有想到通常意义上的剽窃那种极端情形——在数学界,那想必还极罕见。可是从实际角度看,对为此付出代价的研究者来说,局面是一样的;而使这局面成为可能的内心态度,在我看来也差不太多。它只是更舒服些:因为它伴着一种对他人无限优越的感觉,伴着良心安然与内心满足——那是一个自命为数学之不可触犯的纯粹性的不妥协捍卫者的人。

\section{$\varnothing$}

\notehead{28}{}
写前面那几页时,我起初在“把心里的东西全倒出来”的欲望与一份保留、谨慎的顾虑之间分裂。于是我就停在差不多、大约莫的地方——这肯定是我那份不适、“我什么也没学到”之感的主要原因。自从写下那几行承认不适的文字之后,我把这几页重写了两遍——它们让我心里不满意;我把自己更清楚地放了进去,也把事往更深处推。一路下来,我确实终于“学到了点东西”;我也相信,与此同时我摸到了某件要紧的东西,它超出了这个具体个例,也超出了我这个人。

\section{$\varnothing$}

\notehead{29}{}
我这里说的,是在数学里、或在另一项纯智力活动里那种强烈而长久的投入。反过来,这样一份激情的展开——它可以是重新认识我们身上一种被遗忘的力量的方式,也是与一种不肯就范的实质较量的机会,并且一路下来,用某种真正属于我们自己的东西去更新、丰富我们的身份感——这样一种展开,很可能是内心路程上的一个重要阶段,是成熟过程中的一个重要阶段。

\section{$\varnothing$}

\notehead{30}{}
这几年,接力的是我的孩子们:由他们去教一个有时不大情愿的学生,人生的种种奥秘……

\section{$\varnothing$}

\notehead{31}{}
\origpage{163}我这里想到的是求知欲的“阳”的形式——那探问、发现、为显现者命名的形式……正是由于被命名,出现的知识才变得不可逆转、不可抹去(哪怕它后来被埋掉、被遗忘,不再活跃……)。求知欲的“阴”的、“女性”的形式,在于一种敞开、一种接纳,在于无声地迎候一种出现在我们存在更深层——思想到不了的地方——的知识。
% ---------------- p.249 ----------------
这样一种敞开、一种突然到来、一时间抹去一切冲突痕迹的知识的出现,仍像一份恩典:它触到深处,而它可见的效果也许转瞬即逝。不过我猜,这种在我们一生中某些稀有时刻如此到来的、无言语的知识,同样不可抹去;它的作用甚至延续到我们对它的记忆之外。

\section{炉中百铁,或:干熬无用!}

\notehead{32}{}
我还在做泛函分析的时候,也就是直到1954年,我有时会没完没了地死咬一个问题,解不出来;那时我已经没有新想法,只是绕着旧想法的圈子打转,而那些想法显然已经“咬不住”了。整整一年都是这样,尤其是拓扑向量空间里的那个“逼近问题”;它要再过二十年才被解决,用的完全是另一路的方法,以我当时所处的境地,那些方法只能与我擦肩而过。当时驱动我的不是欲望,而是一种固执,以及对自己内心所发生之事的不了解。那是难受的一年——我一生中唯一一段做数学变得难受的时期!我需要这段经历才明白:“干熬”没有用——一件工作一旦到了停顿点,而且只要停顿被察觉,就该转到别的事上去;至于悬置的那个问题,等到更合宜的时候再回来。这个时刻几乎总是很快就出现——问题自己成熟起来,我并没有做出碰它的样子;靠的只是兴致勃勃地做那些看来与它毫无关系的问题的功效。我确信,当时若我硬撑下去,哪怕十年也一事无成!从1954年起,我在数学上养成了一个习惯:炉子里总同时烧着许多块铁。我一次只打一块,可由于一种不断更新的奇迹\origpage{164},我在一块上做的工作也惠及所有其余的,它们各自等着自己的时辰。与我第一次接触冥想时也一样,而且完全不是我刻意为之:随着反思继续,有待审视的灼热问题,数目一天天增多。

\section{$\varnothing$}

\notehead{33}{}
这并不意味着:工作中纸(或黑板,那不过是它的一种变体!)不在场的那些时刻,在数学工作里就不重要。尤其如此的是那些“敏感时刻”:一种新的直觉刚刚出现,这时要做的是以一种更整体、更直觉的方式去“认识”它,而不是做“零件活儿”——而反思的这个不定形阶段,正是为那件事作准备。在我身上,这类思考主要在床上或散步时进行;我觉得它只占用于工作的全部时间中一个相对不大的部分。同样的观察也适用于我迄今所从事的冥想工作。

\section{无力的拥抱}

\notehead{34}{}
“拥抱”这个词在我这里绝不是单纯的比喻;日常语言在此映出的是一种深层的同一。人们不无道理地会说:这么说,“没有惊叹的拥抱无能为力”就不对了——若按字面如此,大地就该人口凋零,或者干脆荒无人烟。极端的例子是强奸:惊叹当然不在场,可有时被强奸的女人身上仍孕育出一个生命。由这样的拥抱生下的孩子,肯定免不了带着那个印记;那印记会是他分得的那一份“包裹”的一部分,该由他自己去承担。可这并不妨碍:一个新生命确实被孕育、被生了下来;确实有过创造,那是力量的标志。也确实如此:我见过某个满身自得的数学家找到并证明了漂亮的定理,那正是拥抱不乏力气的标志!但同时也不假:若某位数学家的生活被他的自得闷死(在某种程度上,我自己生活中某个时期也曾如此),那么他与数学的这些拥抱所结的果子,对他、对别人都不是什么福分。从强奸中生下的孩子的父亲与母亲,也可以这样说。我说“无力的拥抱”,首先指的是这样一件事:在那个自以为在创造的人身上,无力生出更新——他只造出一个产品、一件他之外的东西,在他自身里没有深的回响;
% ---------------- p.250 ----------------
\origpage{165}这产品不但不解放他、不在他里面造出和谐,反而把他更紧地捆在他身内那份他作囚徒的自得上,那自得不断催他去生产、再生产。这是深层上的一种无力,藏在“创造力”的表象背后——而那归根到底不过是无缰的生产力。

我也有充分机会看出:自得、无惊叹之能,其性质是一种真正的盲,是自然感受力与嗅觉的堵塞;这堵塞即使不是全面而永久的,至少在某些具体情形里是明显的。在这种状态里,某位有声望的数学家有时会在恰恰他擅长的那些事情上,显得比最死心眼的小学生还蠢!在另一些场合,他会做出技术精湛的奇迹。可我怀疑,他是否还有能力去发现那些单纯而显然、有力量更新一门学科或一门科学的东西。那些东西离他太低、太远,他已经不屑于去看!要看别人不屑于看的东西,需要一份他早已失去、或被他放逐的天真……近二十年间数学产出惊人地增长,新结果的纷繁令人无所适从,连只想稍微“跟上”的数学家都被淹没;而与此同时(就我从四处传来的回声所能判断的而言),我所多少熟悉的那些大思考主题,却几乎没有真正的更新、没有大幅度的转化(不只是积累式的转化)——这肯定不是偶然。更新不是量的事情;它与投入的量无关——那种量可以用若干“数学家人日”来计量,即某些“水平”的某些数学家花在某个题目上的日数。一百万个数学家人日,也无力生出像零那样孩子气的东西——而零更新了我们对数的感知。只有天真拥有这份力量;它的一个可见标志,就是惊叹……

\section{$\varnothing$}

\notehead{35}{}
这份“天赋”不是谁的特权,我们全都生来带着它。当它在我身上显得不在时,那是我自己把它赶走的;要不要重新接纳它,只在于我。在我身上,或在某某身上,这份“天赋”的表达方式与另一个人不同,不那么有感染力,也许不那么不可抗拒,但它照样在场;我也说不出它是否不那么有力。

\section{$\varnothing$}

\notehead{36}{}
\origpage{166}这种对美的细腻感受力,在我看来与一样东西紧密相连;
% ---------------- p.251 ----------------
我曾以“严要求”(对自己的)或“严格”(取其完整意义)之名谈到它,把它描述为“对我们自身中某种细腻之物的注意”,对所要探问之物的理解品质的注意。对一件数学东西的理解品质,离不开对该物特有之“美”的一种或深或浅、或全或欠的感知。

\section{$\varnothing$}

\notehead{37}{}
我想,几乎不必补充:这项长久的工作日复一日地显出许多别的东西,远不止我刚才以碑铭般的简省形式交出的那个“结果”。冥想的工作,与由某个打算考察的具体问题所驱动的数学工作,并无不同。所走那条路(它通向、或不通向对最初问题或多或少的澄清)上的种种曲折,往往比最初的问题或“最终结果”更有意思。

\section{$\varnothing$}

\notehead{38}{}
这些注其实是写给……的那封长信的继续,那封信成了本书第一章。它们用打字机打出来,是为了让这位旧日的朋友能读,也为了另外两三位(尤其是罗尼·布朗),我以为他们可能会感兴趣。再说,那封信从未得到回音,收信人也没有读它:将近一年之后(我问过他是否确实收到),他由衷地惊讶——我怎么会有一刻以为他会读它,鉴于人们该期待我做的那类数学……

\section{$\varnothing$}

\notehead{39}{}
这一时期,除别的之外,也正是《穿越伽罗瓦理论的长征》的时期;在《纲领概要》里(第3段:“与童画相联系的数域”)谈到过它。

\section{来访}

\notehead{40}{}
对这个梦所做的工作,是一封英文长信的内容;收信人是一位朋友兼同行,他前一天在我家匆匆过了一趟。做梦者用来从一片表面的虚无中唤出这个逼真得惊人的梦的那些材料,有些显然取自那位亲爱的朋友来访的短暂片段——我已有近十年没再见过他。\origpage{167}于是,在工作头一天,而且与我过去的经验相违,我竟以为可以由此断定:我做的这个梦关系到我的这位朋友,多过关系到我——该做这个梦的是他,不是我!这是一种回避梦境消息的办法;而那消息(凭我过去的经验,我本该一开始就知道)只关系到我,不关别人。在这第一阶段——浮浅的——工作之后的夜里,我终于意识到了这一点;第二天,我在同一封信里接着做了下去。自那封值得记住的信之后,我再没有收到这位朋友的任何音讯——他是我有过的最亲近的朋友之一。

% ---------------- p.252 ----------------
这项工作,是唯一一次以信件形式(而且用的还是英文)成形的冥想,因此我手里不再有它的书面痕迹。这段插曲让我格外震动——在众多别的例子当中,它们都表明:任何越出某层门面的工作迹象,任何把极简单的事实——而人们通常认为有义务对之视而不见——带到光天化日之下的工作,是多么在别人身上引起不适与恐惧。这一点我在后面还要谈(见第47段《孤独的历险》)。

\section{克里希那穆提,或变成羁绊的解放}

\notehead{41}{}
若说我从这次阅读中得到的只有一套词汇,以及一种把它据为己有、最终又顺理成章地用它取代现实的倾向,那是不准确的。我拿到手的第一本克里希那穆提(Krishnamurti)的书之所以那么震动我(而且我还只来得及读其中几章),是因为他说的东西彻底掀翻了我许多自以为当然的事;我立刻看出,那些都是老生常谈,一直是我所呼吸的空气的一部分。同时,这次阅读第一次把我的注意力引向一些分量很重的事实,尤其是对现实的逃避——它是心智最强大、最普遍的制约之一。这给了我一把关键的钥匙,去理解那些此前一直不可理解、并因此(在我五六年之后发现冥想之前,我并未察觉这一点)滋生焦虑的处境。我立刻就能在我周围各处看到这种逃避的现实。这解开了某些焦虑,却没有改变任何本质的东西,\origpage{168}因为我只在他人身上看到那份现实,同时(理所当然地)以为它在我身上并不存在,以为我总之是那条证实规则的例外(而且对这么一个确实了不起的例外,我没有再问任何别的问题)。事实上,我对别人、对自己都毫无好奇。这把“钥匙”只有在被深入之欲驱动的人手里才能开锁。在我手里,它成了驱邪和做姿态。

直到1974年初,我才第一次认清:我生活里的那种毁灭一步步跟着我,不可能只来自别人;我身上有某种东西在吸引它、喂养它、让它延续。那是一个谦卑与敞开的时刻,适宜于一次更新。可那次更新当时仍是边缘的、短暂的,因为缺少深入的工作。那个“我身上的某种东西”还很模糊。我看得出那是爱的匮乏;可是,要有一项工作去更贴近地界定:爱的匮乏在我身上何处、如何发生,它如何显现,它具体造成了什么后果,等等——这样的念头,既不可能来自我至今认识的任何圈子或任何人,也不可能来自克里希那穆提。(恰恰相反,K.喜欢强调一切工作的徒劳,他自动把工作等同于“我”的“成为之饥”。)就这样,我拿一份借来的“智慧”当唯一的罗盘,看不见还有别的事可做,只好耐心等着“爱”像圣灵的恩典一样降到我身上。

然而,我刚在一个浪的深处学到的那点卑微真相,已经激起了一股强劲的新能量之浪,可与两年半之后承载我第一次投入冥想的那股相比。这份能量当时并没有完全闲置。几个月后,我因一场凑巧的事故动弹不得,它承载了一份(写下来的)反思:我一生中第一次审视那套世界观——它一直是我与他人关系未曾言说的基础,来自我的父母,尤其来自我母亲。当时我非常清楚地看出:这套看法已经破产,它无法说明人与人之间关系的实际,也无法促成我这个人以及我与他人关系的舒展。这份反思仍带着
% ---------------- p.253 ----------------
“克里希那穆提风格”的印记,也带着克里希那穆提式的禁忌——禁止一切朝向理解的真正工作。可它毕竟把几个月前生出的一份认识变得可触而不可逆转——\origpage{169}那认识起初一直模糊、游移。那份认识,当时世上任何一本书、任何别的人都不可能带给我。

要够得上冥想的品质,这份反思最缺的,是对我自己这个人、对我关于自己的看法的注视——而不只是对世界观的注视,对一个公理系统的注视,而在那系统里我并不真的“有血有肉”地出场。它还缺另一样:在当下、在反思进行的那一刻对我自己的注视(而那份反思还停留在真正工作的此岸);那样一种注视,本来会让我既察觉一种借来的风格,也察觉我对这些笔记的文学面貌有某种自得,也就是说,缺少自发与真实。尽管它如此不足,尽管它对我与他人关系的直接影响相对有限,这份反思在我看来仍是一个阶段——考虑到起点,这大概是必需的——通向两年后该发生的那次更深的更新。正是在那时,我终于发现了冥想——同时发现这第一个未曾料到的事实:关于我自己这个人,是有东西可发现的;而那些东西几乎完全决定了我一生的走向,以及我与他人关系的性质。

\section{有益的剥离}

\notehead{42}{}
那个“有冲击力”的事件,是1969年底的发现:我觉得自己属于的那个机构,部分经费来自国防部的款项——这与我基本的公理不相容(今天也仍然不相容)。这个事件是一整串别的(一个比一个更有揭示力!)事件中的第一个;它们的结果是我离开高等科学研究所(IHÉS),并由此一环扣一环地彻底换了圈子与投入。

在IHÉS的英雄年代,迪厄多内和我是它仅有的成员,也是仅有的为它在科学界赢得信誉与听众的人:迪厄多内靠主编《数学出版物》(头一卷早在1959年就出了,即莱昂·莫尚(Léon Motchane)创办IHÉS的次年),我靠《代数几何讨论班》。最初那些年,IHÉS的存在极为飘摇:经费没有着落(靠几家\origpage{170}充当赞助人的公司的慷慨),唯一的场所是巴黎蒂耶尔基金会(Fondation Thiers)借出的一间厅(它显然很不情愿),专供我讨论班开讲的日子\fnote{5}{}{IHÉS最近为纪念其成立二十五周年出的一本小册子(承尼科·凯珀(Nico Kuiper)好意寄我一本),对这样艰难的开端只字不提——也许认为它们配不上这个场合的庄重吧;去年,这个周年被大张旗鼓地庆祝了一番。}。我多少觉得自己,连同迪厄多内,是我所属这个机构的“科学”共同创办人;我还真打算在那儿过完这辈子!我终于对IHÉS有了很强的认同;而我的离开(由于同事们的漠然),当时被体验为从另一个“家”被剥离,后来才显出是一次解放。

回过头看,我意识到自己身上想必早有更新的需要,我说不清是从什么时候起。我离开IHÉS的前一年,我的精力投入突然翻了个面:丢下那些昨天还烫手的工作、那些最让我着迷的问题,转而去投身生物学(是在一位生物学家朋友米尔恰·杜米特雷斯库(Mircea Dumitrescu)的影响下)——这肯定不是单纯的巧合。我投身其中,是打算在IHÉS内部做
% ---------------- p.254 ----------------
\origpage{171}长期投入的(这也符合该机构多学科的宗旨)。这肯定只是给一种深得多的更新需要找了个出口;那种更新在IHÉS那种“科学蒸笼”的气氛里无从完成,而是在我已提及的那场“觉醒的瀑布”中发生的。共有七次,最后一次在1982年。“军费”那段插曲是天赐的,因为它触发了这些“觉醒”中的第一次。国防部,以及我在IHÉS的旧同事们,最终都配得上我全部的感激!

\section{$\varnothing$}

\notehead{43}{}
“我写的那部诗作”里包含许多我第一手知道的东西;它们今天在我生活中、在一般的“生活”中,显得与它写成之时同样重要——那时我是打算把它发表的。我之所以没有发表,主要是因为后来意识到:那形式被一种刻意“作诗”的意图所累,以至于它整体构思过于搭建,许多段落缺少自发,有些地方甚至僵硬或浮肿得难受。这种有时浮夸的形式,是我当时禀性的映照:在那里面,常常确实是“老板”在领舞——笨重得很,那是不用说的……

\section{$\varnothing$}

\notehead{44}{}
不用说,我在这里是把这样一个假设撇开的——至少可以说,它并非不可能:一场原子战争,或另一场同类的“乐事”,突然闯来,把那个叫做“数学”的集体游戏一举粗暴地终结,顺带还终结别的好些东西……


\end{document}
